\chapter{Brave-New Lie Algebroids and Differentiation}
\label{ch:brave-new-lie-algebroids}

\section{Purpose, continuity, and the differentiation problem}
\label{sec:bnla-purpose}

The preceding chapter constructed brave-new formal Lie groupoids as the
unitwise relative de Rham-complete fixed points of the centred completion
coreflection.  Thus, for a morphism
\[
    X\longrightarrow S
\]
and a formal Lie groupoid
\[
    \mathfrak G_\bullet\rightrightarrows X,
\]
the nonlinear infinitesimal anchor is not additional structure.  It is the
contractibly unique fixed-object groupoid morphism
\[
    \phi_{\mathfrak G}:
    \mathfrak G_\bullet
    \longrightarrow
    R_{X/S,\bullet}
\]
to the terminal formal relation.  On effective quotients it is represented by
\[
\begin{tikzcd}
X \arrow[r,two heads,"q_{\mathfrak G}"]
  \arrow[dr,two heads,"\epsilon_{X/S}"'] &
Q_{\mathfrak G}
  \arrow[d,two heads,"r_{\mathfrak G}"]\\
&Q_{X/S},
\end{tikzcd}
\]
and both displayed quotient maps are relative de Rham equivalences.

The same chapter constructed formal orbit disks, formal isotropy, the
anchor-kernel equivalence
\[
    I^{\mathrm{for}}_{\mathfrak G}
    \simeq
    \mathfrak G_1\times_{R_{X/S}}X,
\]
the formal-isotropy bitorsor, descended formal adjoint transport, and the
effective-image disk--jet calculus.  In the represented range it also
constructed
\[
    \mathfrak a_{\mathfrak G}^{\vee}
    :=u^*L_{\mathfrak G_1/X,s},
    \qquad
    a_{\mathfrak G}^{\vee}:
    L_{X/S}\longrightarrow\mathfrak a_{\mathfrak G}^{\vee},
\]
and proved that this cotangent anchor is the first structured coefficient of
the canonical nonlinear anchor.

The present chapter differentiates that nonlinear package at three compatible
levels and proves the comparison theorems needed to pass between them.

First, the pointed arrow geometry is stabilized over $X$.  This gives an
intrinsic stable shadow of the complete arrow object and a stable anchor
without representability.  Formal isotropy is differentiated separately, and
its comparison with the homotopy fibre of the stable anchor is retained as a
morphism rather than silently replaced by an equivalence.

Second, actual formal isotropy is a group object.  Fibrewise suspension gives
a cocommutative Hopf object, and the spectral Lie--Hopf adjunction gives an
integral spectral Lie algebra of derived primitives
\cite{Ching2005,Heine2024}.  This Lie structure belongs to the isotropy group
object.  No binary bracket is declared on the reduced stable coefficient of
the complete arrow family.

Third, the effective formal quotient determines finite infinitesimal lifting
functors.  In the coherent spectral coefficient sector we construct a
spectral Artinian deformation category intrinsically, prove that its objects
admit finite square-zero decompositions without retaining such decompositions,
construct the derived adic filtration, prove the relative
cotangent--square-zero comonadicity theorem, dualize the resulting relative
coLie comonad by pro-coherent spectral Koszul duality, prove spectral Artinian
excision, and obtain the many-object spectral partition-Lie recognition
theorem.  The one-object spectral Koszul-dual input is the coherent
$E_\infty$ theory of \cite{BrantnerCamposNuiten2025}; the relative
organization follows the cotangent--square-zero and cellular-deformation
mechanism of \cite{BrantnerMagidsonNuiten2025}, but every relative spectral
statement used below is constructed and proved in the spectral category.

\paragraph{Coefficient-domain convention.}
The intrinsic stable and Hopf constructions are defined for every brave-new
formal Lie groupoid in the native big spectral Nisnevich $\infty$-topos.
The spectral partition-Lie construction uses coherent spectral coefficient
rings, because that is the coefficient sector in which pro-coherent spectral
Koszul duality is available.  On an affine chart
\[
    \operatorname{Spec}(B)\longrightarrow\operatorname{Spec}(R),
\]
coherence is a property of the coefficient chart on which the algebraic
functor is evaluated.  It is not a field of structure attached to
$\mathfrak G$.  Local coherence and qcqs hypotheses in globalization are
likewise theorem-scope conditions.

\paragraph{Non-conflation convention.}
The following constructions remain distinct unless a theorem identifies them:
\begin{enumerate}
    \item the intrinsic stable arrow coefficient
    $\mathbb A^{\mathrm{st}}_{\mathfrak G}$;
    \item the reduced stable coefficient
    $\mathbb I^{\mathrm{st}}_{\mathfrak G}$ of actual formal isotropy;
    \item the homotopy fibre
    $\mathbb F^{\mathrm{st}}_{\mathfrak G}$ of the stable anchor;
    \item the integral spectral Lie algebra
    $\mathfrak l^{\mathrm{sp}}_{\mathfrak G}$ of derived primitives of the
    unreduced isotropy suspension Hopf object;
    \item the represented first-order coefficient
    $\mathfrak a^\vee_{\mathfrak G}$;
    \item the tangent-fibre functor of the spectral Artinian formal-moduli
    envelope;
    \item the spectral partition-Lie algebroid
    $\mathbb A^{\pi,E_\infty}_{\mathfrak G}$;
    \item the characteristic-zero dg comparison object
    $\mathbb A^{\mathrm{Lie}}_{\mathfrak G}$.
\end{enumerate}

The distinction between actual isotropy and a homotopy anchor fibre is
essential.  In any stable category
\[
    \operatorname{fib}(0\to T)\simeq\Omega T,
\]
although the ordinary kernel of the zero map is zero.  The universal pattern
retained throughout the chapter is therefore
\[
    \boxed{
    \text{differentiated actual isotropy}
    \longrightarrow
    \text{homotopy anchor fibre}
    \longrightarrow
    \text{isotropy discrepancy}.}
\]

\paragraph{No hidden package conditions.}
No class of admissible algebroids, formal-completion presentation, chosen
finite infinitesimal factorization, chosen Artinian square-zero chain, chosen
adic filtration, chosen anchor, connection, splitting, transitivity or
surjectivity condition, primitive-exactness condition, or formal-moduli
recognition datum is part of the input.  The adic filtration is a left-adjoint
construction and is completed functorially.  Completeness is a property of
the resulting arrow.  Spectral Artinian objects are defined by intrinsic
nilpotence, finiteness, and boundedness conditions; finite square-zero
factorizations are proved to exist and are used only inside proofs.  The
affine Schlessinger equations are imposed by the explicit Artinian
localization.  Globalization is then constructed in the nonlinear target:
an accessible spectral deformation localization, imposing Postnikov
convergence and nilpotent excision, is applied once to the effective formal
quotient.  Affine formal leaves are extracted from that single target, and
their spectral etale cocartesianity is proved by the coherent pre-tangent
construction, the cartesian anchored right-left coefficient family,
almost-finitely-presented approximation, and conservative affine recognition.  Pro-coherent etale transport is
obtained from the relative right-left completion of the coherent spectral
coefficient family; finite-Tor Beck--Chevalley and spectral fppf universal
descent then prove etale descent for that family.  The spectral partition-Lie
algebraic structure
is produced by a monad obtained from the cotangent--square-zero comonad by
pro-coherent duality and a proved completion theorem.  Canonical cobar
resolutions are constructions, not presentations retained on objects.

\paragraph{Comparison scope.}
The characteristic-zero dg comparison is constructed directly by rational
symmetric monoidal rectification of the spectral cotangent--square-zero and
pro-coherent Koszul-dual constructions.  Rational rectification is proved to
identify the intrinsic spectral and dg Artinian categories and their
Schlessinger generators, so the rectified formal-moduli object is defined
literally on dg Artinian tests.  On the coefficient range of
\cite{Nuiten2019} it agrees with dg formal-moduli differentiation of that same
rectified formal-moduli object.  No integral spectral-to-animated Artinian
comparison is assumed.

\paragraph{Scope.}
This chapter does not yet construct the coherent Spencer object, the
co-Kleisli operator algebra, universal enveloping actions on jets, PBW
comparisons with the jet operator calculus, explicit rooted-tree formulas, or
operator-valued Maurer--Cartan transport.  Those belong to the
Atiyah--Spencer chapter.

\paragraph{Chapter-level output.}
The chapter constructs the intrinsic stable differential and its isotropy
comparison, the isotropy Hopf object and derived primitives, the represented
first-order interface, including retractive invariance under unitwise
formalization, the spectral adic and cotangent--square-zero theory,
the relative spectral partition-Lie algebroid monad, including the complete
anchor-weight filtration and the zero-anchor subcomonad, intrinsic spectral
Artinian tests and their Schlessinger localization, the anchored cellular
attachment calculation, affine and global spectral formal-moduli/partition-Lie
recognition, the relative pro-coherent etale coefficient family, the global
spectral deformation formalization of effective quotients, the coherent
pre-tangent and almost-finitely-presented tangent construction, the cartesian
anchored right-left transport, the etale-cocartesian formal-leaf construction
and its spectral controller,
fixed-object and spectral-etale varying-object functoriality, rational Artinian and dg
rectification, spectral and dg
isotropy-to-anchor-fibre comparisons, gauge effective-image corrections, and
the classical smooth Lie-groupoid and Atiyah calculations.  The final package
theorem records the precise theorem scopes.

The logical progression is
\[
\boxed{
\begin{aligned}
\text{formal Lie groupoid}
&\longrightarrow \text{intrinsic stable differential}\\
&\longrightarrow \text{integral spectral-Lie isotropy}\\
&\longrightarrow \text{spectral Artinian formal moduli}\\
&\longrightarrow \text{spectral adic and cotangent--square-zero duality}\\
&\longrightarrow \text{spectral partition-Lie algebroid}\\
&\longrightarrow \text{rational dg comparison}\\
&\longrightarrow \text{Atiyah--Spencer operation calculus}.
\end{aligned}}
\]
\section{Parameterized stable coefficients and reduced stabilization}
\label{sec:bnla-stable-ambient}

Retain from Chapter~10 the parameterized stable coefficient category
\[
    \mathcal E_X
    :=
    \operatorname{Sp}(\mathcal X_{/X}),
\]
the stabilization of the slice $\infty$-topos over $X$.  Its stable,
presentable, closed symmetric monoidal structure and its pullback
functoriality were constructed there.  Write
\[
    \mathbb S_X\in\mathcal E_X
\]
for its tensor unit and
\[
    \Sigma^\infty_{+,X}:
    \mathcal X_{/X}
    \longrightarrow
    \mathcal E_X
\]
for fibrewise suspension with a disjoint basepoint.

\subsection{Reduced fibrewise stabilization}

\begin{definition}[Pointed object over $X$]
\label{def:bnla-pointed-object}
A pointed object over $X$ is a diagram
\[
    X\xrightarrow{e}Y\xrightarrow{p}X,
    \qquad
    pe\simeq\operatorname{id}_X.
\]
We regard it as an object of $(\mathcal X_{/X})_*$.  A morphism of pointed
objects is a morphism over $X$ preserving the displayed section.
\end{definition}

The section $e$ induces
\[
    \mathbb S_X
    \simeq
    \Sigma^\infty_{+,X}X
    \longrightarrow
    \Sigma^\infty_{+,X}Y.
\]

\begin{definition}[Reduced fibrewise stabilization]
\label{def:bnla-reduced-stabilization}
For a pointed object $(Y,e)$ over $X$, define
\[
    \widetilde\Sigma^\infty_X(Y,e)
    :=
    \operatorname{cofib}
    \left(
        \mathbb S_X
        \longrightarrow
        \Sigma^\infty_{+,X}Y
    \right)
    \in\mathcal E_X.
\]
When the section is understood we write
$\widetilde\Sigma^\infty_XY$.
\end{definition}

\begin{proposition}[Functoriality of reduced stabilization]
\label{prop:bnla-reduced-stabilization-functoriality}
Reduced fibrewise stabilization defines a functor
\[
    \widetilde\Sigma^\infty_X:
    (\mathcal X_{/X})_*
    \longrightarrow
    \mathcal E_X.
\]
In particular, a commutative diagram
\[
\begin{tikzcd}
X \arrow[r,"e"] \arrow[d,equal] &
Y \arrow[r,"p"] \arrow[d,"f"] &
X \arrow[d,equal]\\
X \arrow[r,"e'"'] &
Y' \arrow[r,"p'"'] &
X
\end{tikzcd}
\]
induces canonically
\[
    \widetilde\Sigma^\infty_X(f):
    \widetilde\Sigma^\infty_X(Y,e)
    \longrightarrow
    \widetilde\Sigma^\infty_X(Y',e').
\]
The construction preserves identities, composition, equivalences, and all
higher functorial coherences.
\end{proposition}

\begin{proof}
The pointed morphism determines a morphism of the defining cofiber diagrams
\[
\begin{tikzcd}
\mathbb S_X \arrow[r] \arrow[d,equal] &
\Sigma^\infty_{+,X}Y \arrow[d]\\
\mathbb S_X \arrow[r] &
\Sigma^\infty_{+,X}Y'.
\end{tikzcd}
\]
Cofiber is a functor on the arrow category of the stable $\infty$-category
$\mathcal E_X$.  Applying it gives the displayed morphism.  All coherence is
inherited from the suspension functor and functoriality of finite colimits.
\end{proof}

\begin{proposition}[Base change of reduced stabilization]
\label{prop:bnla-reduced-stabilization-base-change}
Let $f:X'\to X$ be a morphism in $\mathcal X$.  For every pointed object
$(Y,e)$ over $X$, there is a canonical equivalence
\[
    f^*\widetilde\Sigma^\infty_X(Y,e)
    \xrightarrow{\sim}
    \widetilde\Sigma^\infty_{X'}
    \left(
        X'\times_XY,
        f^*e
    \right).
\]
These equivalences preserve identities, composition, and all higher
pullback-pasting coherences.
\end{proposition}

\begin{proof}
Pullback between slices of an $\infty$-topos preserves all small colimits and
finite limits.  Stabilization therefore gives an exact pullback functor
\[
    f^*:\mathcal E_X\longrightarrow\mathcal E_{X'}
\]
which sends $\mathbb S_X$ to $\mathbb S_{X'}$.  Fibrewise suspension is the
stabilization of free pointing and hence satisfies
\[
    f^*\Sigma^\infty_{+,X}Y
    \simeq
    \Sigma^\infty_{+,X'}(X'\times_XY).
\]
Apply the exact functor $f^*$ to the defining cofiber sequence.
\end{proof}

\begin{lemma}[The identity-pointed object stabilizes to zero]
\label{lem:bnla-reduced-identity-zero}
There is a canonical equivalence
\[
    \widetilde\Sigma^\infty_X(X,\operatorname{id}_X)
    \simeq 0.
\]
\end{lemma}

\begin{proof}
The defining map is the identity of $\mathbb S_X$.
\end{proof}

\begin{remark}[Stable coefficient, not tangent complex]
\label{rem:bnla-stable-not-tangent}
The object $\widetilde\Sigma^\infty_X(Y,e)$ is the stable shadow of the full
pointed nonlinear object.  It is not, by definition, the first cotangent or
tangent coefficient along $e$.  No HKR collapse, exterior presentation,
principal-parts model, or first-order truncation is implicit in the notation.
\end{remark}

\section{The intrinsic stable differential of a formal Lie groupoid}
\label{sec:bnla-stable-differential}

Fix
\[
    \mathfrak G_\bullet
    \in
    \operatorname{FormLieGpd}(X/S).
\]
Write
\[
    s,t:\mathfrak G_1\rightrightarrows X,
    \qquad
    u:X\longrightarrow\mathfrak G_1
\]
for source, target, and unit.

\begin{definition}[Intrinsic stable arrow coefficient]
\label{def:bnla-stable-arrow-coefficient}
Regard $\mathfrak G_1$ over $X$ by $s$.  Define
\[
    \mathbb A^{\mathrm{st}}_{\mathfrak G}
    :=
    \widetilde\Sigma^\infty_X
    (\mathfrak G_1,u).
\]
\end{definition}

For the maximal relative relation write
\[
    R_{X/S}
    =
    X\times_{Q_{X/S}}X
\]
and regard it over $X$ through its source projection, pointed by the diagonal
unit $u_R:X\to R_{X/S}$.

\begin{definition}[Stable relative relation coefficient]
\label{def:bnla-stable-relative-tangent}
Define
\[
    \mathbb T^{\mathrm{st}}_{X/S}
    :=
    \widetilde\Sigma^\infty_X
    (R_{X/S},u_R).
\]
The superscript $\mathrm{st}$ is essential: no identification with
$L_{X/S}^{\vee}$ is asserted.
\end{definition}

The canonical nonlinear anchor in degree one
\[
    \phi_{\mathfrak G,1}:
    \mathfrak G_1\longrightarrow R_{X/S}
\]
is a morphism of pointed objects over $X$.

\begin{definition}[Intrinsic stable anchor]
\label{def:bnla-stable-anchor}
Define
\[
    a^{\mathrm{st}}_{\mathfrak G}:
    \mathbb A^{\mathrm{st}}_{\mathfrak G}
    \longrightarrow
    \mathbb T^{\mathrm{st}}_{X/S}
\]
to be $\widetilde\Sigma^\infty_X(\phi_{\mathfrak G,1})$.
\end{definition}

\begin{theorem}[Intrinsic stable differentiation]
\label{thm:bnla-intrinsic-stable-differentiation}
For fixed $X/S$, the assignment
\[
    \mathfrak G_\bullet
    \longmapsto
    \left(
        \mathbb A^{\mathrm{st}}_{\mathfrak G}
        \xrightarrow{a^{\mathrm{st}}_{\mathfrak G}}
        \mathbb T^{\mathrm{st}}_{X/S}
    \right)
\]
is functorial on $\operatorname{FormLieGpd}(X/S)$.

More generally, let a formal-groupoid functor
\[
    F:(\mathfrak G,X)\longrightarrow(\mathfrak H,Y)
\]
cover a morphism $f:X\to Y$ over $S$.  Then there is a canonical commutative
square in $\mathcal E_X$
\[
\begin{tikzcd}
\mathbb A^{\mathrm{st}}_{\mathfrak G}
  \arrow[r]
  \arrow[d,"a^{\mathrm{st}}_{\mathfrak G}"'] &
f^*\mathbb A^{\mathrm{st}}_{\mathfrak H}
  \arrow[d,"f^*a^{\mathrm{st}}_{\mathfrak H}"]\\
\mathbb T^{\mathrm{st}}_{X/S}
  \arrow[r] &
f^*\mathbb T^{\mathrm{st}}_{Y/S}.
\end{tikzcd}
\]
For an arbitrary spectral base change $S'\to S$, these comparison morphisms
specialize to equivalences for the base-changed groupoid.

No represented, smoothness, dualizability, or formal-moduli condition enters
the construction.
\end{theorem}

\begin{proof}
For fixed $X$, automatic functoriality of canonical anchors from the preceding
chapter gives a homotopy-commutative square of pointed objects
\[
\begin{tikzcd}
\mathfrak G_1 \arrow[r,"F_1"]
  \arrow[d,"\phi_{\mathfrak G,1}"'] &
\mathfrak H_1 \arrow[d,"\phi_{\mathfrak H,1}"]\\
R_{X/S} \arrow[r,equal] & R_{X/S}.
\end{tikzcd}
\]
Apply Proposition~\ref{prop:bnla-reduced-stabilization-functoriality}.

If $F$ covers $f:X\to Y$, form the pulled-back pointed objects
\[
    X\times_Y\mathfrak H_1,
    \qquad
    X\times_YR_{Y/S}.
\]
The degree-one functor and the functoriality of the maximal formal relation
give pointed maps
\[
    \mathfrak G_1\longrightarrow X\times_Y\mathfrak H_1,
    \qquad
    R_{X/S}\longrightarrow X\times_YR_{Y/S}
\]
compatible with canonical anchors.  Reduced stabilization and the base-change
equivalence of Proposition~\ref{prop:bnla-reduced-stabilization-base-change}
produce the displayed square.

For a base change $S'\to S$, the preceding chapter identifies the
base-changed formal groupoid and maximal relation by equivalences.  Apply the
same proposition.
\end{proof}

\begin{remark}[Stable shadow of the complete arrow object]
The construction uses the whole completed arrow object $\mathfrak G_1$ rather
than its first structured neighbourhood.  It is therefore available before
represented cotangent theory and can retain stable information invisible to
first-order differentiation.  The phrase ``stable shadow'' makes no claim to
a Goodwillie or Taylor-tower expansion.
\end{remark}

\section{The stable anchor fibre, cofiber, and isotropy comparison}
\label{sec:bnla-anchor-triangle}

\begin{definition}[Stable anchor fibre and cofiber]
\label{def:bnla-stable-kernel-normal}
Define
\[
    \mathbb F^{\mathrm{st}}_{\mathfrak G}
    :=
    \operatorname{fib}
    \left(a^{\mathrm{st}}_{\mathfrak G}\right)
\]
and
\[
    \mathbb N^{\mathrm{st}}_{\mathfrak G}
    :=
    \operatorname{cofib}
    \left(a^{\mathrm{st}}_{\mathfrak G}\right).
\]
We use the word \emph{fibre} rather than \emph{kernel} to retain the stable
homotopy information.
\end{definition}

\begin{theorem}[Stable anchor triangle]
\label{thm:bnla-stable-anchor-triangle}
There is a canonical exact triangle
\[
    \mathbb F^{\mathrm{st}}_{\mathfrak G}
    \longrightarrow
    \mathbb A^{\mathrm{st}}_{\mathfrak G}
    \xrightarrow{a^{\mathrm{st}}_{\mathfrak G}}
    \mathbb T^{\mathrm{st}}_{X/S}
    \longrightarrow
    \mathbb N^{\mathrm{st}}_{\mathfrak G}.
\]
Moreover there is a canonical equivalence
\[
    \boxed{
    \mathbb N^{\mathrm{st}}_{\mathfrak G}
    \simeq
    \Sigma\mathbb F^{\mathrm{st}}_{\mathfrak G}.}
\]
The triangle and the suspension comparison are natural in $\mathfrak G$ and
compatible with arbitrary spectral base change.
\end{theorem}

\begin{proof}
The first sequence is the canonical fibre/cofiber triangle of a morphism in a
stable $\infty$-category.  In a stable category, the cofiber of a map is the
suspension of its fibre.  Exact pullback between parameterized stable
categories preserves both constructions and their canonical comparison.
\end{proof}

\begin{remark}[The zero-anchor warning]
For a zero anchor $0\to T$ one has
\[
    \operatorname{fib}(0\to T)\simeq\Omega T,
    \qquad
    \operatorname{cofib}(0\to T)\simeq T.
\]
Thus the homotopy anchor fibre cannot be identified with ordinary isotropy by
formal terminology alone.
\end{remark}

\subsection{Differentiated actual isotropy}

The preceding chapter constructed the formal isotropy group
\[
    I^{\mathrm{for}}_{\mathfrak G}\longrightarrow X
\]
and the canonical Cartesian square
\[
\begin{tikzcd}
I^{\mathrm{for}}_{\mathfrak G}
  \arrow[r]
  \arrow[d] &
\mathfrak G_1
  \arrow[d,"\phi_{\mathfrak G,1}"]\\
X \arrow[r,"u_R"'] & R_{X/S}.
\end{tikzcd}
\]

\begin{definition}[Stable formal-isotropy coefficient]
\label{def:bnla-stable-formal-isotropy}
Define
\[
    \mathbb I^{\mathrm{st}}_{\mathfrak G}
    :=
    \widetilde\Sigma^\infty_X
    \left(I^{\mathrm{for}}_{\mathfrak G},e_I\right).
\]
\end{definition}

The inclusion into the arrow object induces
\[
    i^{\mathrm{st}}_{\mathfrak G}:
    \mathbb I^{\mathrm{st}}_{\mathfrak G}
    \longrightarrow
    \mathbb A^{\mathrm{st}}_{\mathfrak G}.
\]

\begin{proposition}[Canonical isotropy nullhomotopy]
\label{prop:bnla-stable-isotropy-null}
The composite
\[
    \mathbb I^{\mathrm{st}}_{\mathfrak G}
    \xrightarrow{i^{\mathrm{st}}_{\mathfrak G}}
    \mathbb A^{\mathrm{st}}_{\mathfrak G}
    \xrightarrow{a^{\mathrm{st}}_{\mathfrak G}}
    \mathbb T^{\mathrm{st}}_{X/S}
\]
has a canonical nullhomotopy.
\end{proposition}

\begin{proof}
The nonlinear composite factors, with its specified pullback homotopy, through
\[
    I^{\mathrm{for}}_{\mathfrak G}\longrightarrow
    X\xrightarrow{u_R}R_{X/S}.
\]
After reduced stabilization, the middle pointed object becomes zero by
Lemma~\ref{lem:bnla-reduced-identity-zero}.  The specified factorization
therefore supplies the canonical nullhomotopy.
\end{proof}

\begin{definition}[Stable isotropy comparison and discrepancy]
\label{def:bnla-stable-isotropy-comparison}
The preceding nullhomotopy determines, by the universal property of the fibre,
a canonical morphism
\[
    \boxed{
    \gamma^{\mathrm{st}}_{\mathfrak G}:
    \mathbb I^{\mathrm{st}}_{\mathfrak G}
    \longrightarrow
    \mathbb F^{\mathrm{st}}_{\mathfrak G}.}
\]
Define the stable isotropy discrepancy by
\[
    \mathbb C^{\mathrm{iso,st}}_{\mathfrak G}
    :=
    \operatorname{cofib}
    \left(\gamma^{\mathrm{st}}_{\mathfrak G}\right).
\]
\end{definition}

\begin{proposition}[Functoriality and base change of isotropy comparison]
\label{prop:bnla-isotropy-comparison-functoriality}
The morphisms $\gamma^{\mathrm{st}}_{\mathfrak G}$ and the discrepancy objects
$\mathbb C^{\mathrm{iso,st}}_{\mathfrak G}$ are functorial for formal-groupoid
morphisms with the same variance as
Theorem~\ref{thm:bnla-intrinsic-stable-differentiation}.  Under arbitrary
spectral base change they are carried to the corresponding constructions by
canonical equivalences.
\end{proposition}

\begin{proof}
Formal isotropy, its inclusion, the canonical anchor, and the Cartesian
anchor-fibre square are functorial by the preceding chapter.  Reduced
stabilization, fibres, and cofibres are functorial.  The nullhomotopies are the
images of the same functorial nonlinear pullback squares.  Exact base change
preserves every step.
\end{proof}

\begin{remark}[No collapse of isotropy into the anchor fibre]
Reduced stabilization need not preserve the nonlinear pullback defining
formal isotropy.  Thus no unconditional equivalence
\[
    \mathbb I^{\mathrm{st}}_{\mathfrak G}
    \simeq
    \mathbb F^{\mathrm{st}}_{\mathfrak G}
\]
is asserted.  The discrepancy is not auxiliary structure: it is the canonical
cofiber of the canonical comparison.
\end{remark}

\section{Formal isotropy as a Hopf object and its integral spectral Lie algebra}
\label{sec:bnla-spectral-lie-isotropy}

\subsection{The suspension Hopf object and its augmentation ideal}

\begin{definition}[Isotropy suspension Hopf object]
\label{def:bnla-isotropy-hopf}
Define
\[
    H^{\mathrm{iso}}_{\mathfrak G}
    :=
    \Sigma^\infty_{+,X}
    I^{\mathrm{for}}_{\mathfrak G}
    \in\mathcal E_X.
\]
\end{definition}

\begin{proposition}[Cocommutative Hopf structure]
\label{prop:bnla-isotropy-hopf}
The group structure on $I^{\mathrm{for}}_{\mathfrak G}\to X$ canonically
equips $H^{\mathrm{iso}}_{\mathfrak G}$ with a cocommutative Hopf-object
structure in $\mathcal E_X$.  The construction is functorial and compatible
with pullback.
\end{proposition}

\begin{proof}
Fibrewise suspension with disjoint basepoint is strong symmetric monoidal from
the cartesian monoidal slice to parameterized smash products.  Multiplication,
identity, inversion, and diagonal therefore become multiplication, unit,
antipode, and cocomultiplication.  The cartesian diagonal is cocommutative.
All group identities and higher coherences are preserved functorially.
\end{proof}

The identity section and structural morphism give a split augmentation
\[
    \mathbb S_X
    \xrightarrow{\eta}
    H^{\mathrm{iso}}_{\mathfrak G}
    \xrightarrow{\varepsilon}
    \mathbb S_X,
    \qquad
    \varepsilon\eta\simeq\operatorname{id}_{\mathbb S_X}.
\]

\begin{theorem}[Reduced isotropy is the augmentation ideal]
\label{thm:bnla-isotropy-augmentation-ideal}
There is a canonical equivalence
\[
    \boxed{
    \mathbb I^{\mathrm{st}}_{\mathfrak G}
    \simeq
    \operatorname{fib}
    \left(
        H^{\mathrm{iso}}_{\mathfrak G}
        \xrightarrow{\varepsilon}
        \mathbb S_X
    \right).}
\]
Equivalently,
\[
    H^{\mathrm{iso}}_{\mathfrak G}
    \simeq
    \mathbb S_X\oplus
    \mathbb I^{\mathrm{st}}_{\mathfrak G}
\]
in $\mathcal E_X$, naturally in $\mathfrak G$ and under base change.
\end{theorem}

\begin{proof}
By definition,
\[
    \mathbb I^{\mathrm{st}}_{\mathfrak G}
    =
    \operatorname{cofib}
    \left(
        \mathbb S_X\xrightarrow{\eta}
        H^{\mathrm{iso}}_{\mathfrak G}
    \right).
\]
The retraction $\varepsilon$ splits $\eta$.  In a stable $\infty$-category, a
split morphism exhibits its middle object as the direct sum of the source and
the complementary fibre of the retraction.  Therefore
\[
    H^{\mathrm{iso}}_{\mathfrak G}
    \simeq
    \mathbb S_X\oplus\operatorname{fib}(\varepsilon),
\]
and the cofiber of $\eta$ is canonically the second summand.  Naturality
follows from functoriality of unit, augmentation, and the stable splitting.
\end{proof}

\subsection{Derived primitives}

We use the spectral Lie operad and the enveloping-Hopf/derived-primitives
adjunction in the normalization of the spectral Lie--cocommutative Koszul
duality of \cite{Ching2005,Heine2024}.  This convention is introduced here;
it was not part of the definition of a formal Lie groupoid.

\begin{definition}[Integral spectral Lie algebra of formal isotropy]
\label{def:bnla-spectral-lie-isotropy}
Define
\[
    \mathfrak l^{\mathrm{sp}}_{\mathfrak G}
    :=
    \operatorname{Prim}^{\mathrm{der}}_X
    \left(
        H^{\mathrm{iso}}_{\mathfrak G}
    \right)
    \in
    \operatorname{Alg}_{\mathrm{Lie}^{\mathrm{sp}}}(\mathcal E_X).
\]
\end{definition}

\begin{theorem}[Integral spectral-isotropy construction]
\label{thm:bnla-integral-spectral-isotropy}
For every brave-new formal Lie groupoid, formal isotropy canonically determines
$\mathfrak l^{\mathrm{sp}}_{\mathfrak G}$.  There is a canonical counit of
cocommutative Hopf objects
\[
    U^{\mathrm{Hopf}}_X
    \left(
        \mathfrak l^{\mathrm{sp}}_{\mathfrak G}
    \right)
    \longrightarrow
    H^{\mathrm{iso}}_{\mathfrak G}.
\]
No assertion that this counit is an equivalence is made integrally.
\end{theorem}

\begin{proof}
The category $\mathcal E_X$ is stable, presentable, and presentably symmetric
monoidal.  The spectral enveloping-Hopf/derived-primitives adjunction therefore
applies to the Hopf object of
Proposition~\ref{prop:bnla-isotropy-hopf}.  The displayed morphism is its
counit.  Functoriality is functoriality of the right adjoint
$\operatorname{Prim}^{\mathrm{der}}_X$.
\end{proof}

\begin{remark}[Three intrinsic stable objects]
The three objects
\[
    \mathbb I^{\mathrm{st}}_{\mathfrak G},
    \qquad
    \mathbb F^{\mathrm{st}}_{\mathfrak G},
    \qquad
    \mathfrak l^{\mathrm{sp}}_{\mathfrak G}
\]
come from three different constructions: the augmentation ideal of the
suspension Hopf object, the fibre of the stabilized nonlinear anchor, and
derived primitives of the unreduced Hopf object.  No pairwise equivalence is
built into the notation.
\end{remark}

\section{Differentiated formal adjoint transport and primitive base change}
\label{sec:bnla-adjoint-transport}

The preceding chapter constructed coherent conjugation
\[
    \operatorname{Ad}_{\mathfrak G}:
    s^*I^{\mathrm{for}}_{\mathfrak G}
    \xrightarrow{\sim}
    t^*I^{\mathrm{for}}_{\mathfrak G}
\]
over $\mathfrak G_1$.

\begin{proposition}[Adjoint transport on isotropy Hopf objects]
\label{prop:bnla-hopf-adjoint-transport}
Fibrewise suspension carries formal conjugation to an equivalence of
cocommutative Hopf objects over $\mathfrak G_1$
\[
    \operatorname{Ad}^{H}_{\mathfrak G}:
    s^*H^{\mathrm{iso}}_{\mathfrak G}
    \xrightarrow{\sim}
    t^*H^{\mathrm{iso}}_{\mathfrak G}.
\]
It satisfies the unit, composition, and all higher Cech transport coherences
inherited from nonlinear conjugation.
\end{proposition}

\begin{proof}
Conjugation is a homomorphism of group objects over $\mathfrak G_1$.  Apply the
strong symmetric monoidal fibrewise suspension functor.  All coherence diagrams
are carried functorially to coherence diagrams of Hopf morphisms.
\end{proof}

For an arbitrary morphism $f:Y\to X$, a comparison between pullback of
primitives and primitives after pullback can be constructed without asserting
that it is an equivalence.

\begin{definition}[Primitive base-change comparison]
\label{con:bnla-primitive-base-change}
Let $H$ be a cocommutative Hopf object in $\mathcal E_X$.  Pull back the counit
\[
    U_X^{\mathrm{Hopf}}
    \operatorname{Prim}^{\mathrm{der}}_X(H)
    \longrightarrow H.
\]
Strong symmetric monoidality of $f^*$ and functoriality of the enveloping
construction give
\[
    U_Y^{\mathrm{Hopf}}
    \left(
        f^*\operatorname{Prim}^{\mathrm{der}}_X(H)
    \right)
    \longrightarrow
    f^*H.
\]
Adjunction defines the canonical comparison
\[
    \boxed{
    \beta^{\mathrm{Prim}}_{f,H}:
    f^*\operatorname{Prim}^{\mathrm{der}}_X(H)
    \longrightarrow
    \operatorname{Prim}^{\mathrm{der}}_Y(f^*H).}
\]
The construction is natural in $f$, $H$, identities, composition, and higher
base-change pasting.
\end{definition}

\begin{proof}
Only the counit, strong symmetric monoidality of pullback, and the target
adjunction are used.  Naturality and composition follow from uniqueness of
adjoint transposes.
\end{proof}

\begin{theorem}[Differentiated adjoint transport]
\label{thm:bnla-differentiated-adjoint}
There is a canonical equivalence of spectral Lie algebras over $\mathfrak G_1$
\[
    \operatorname{Prim}^{\mathrm{der}}_{\mathfrak G_1}
    \left(s^*H^{\mathrm{iso}}_{\mathfrak G}\right)
    \xrightarrow{\sim}
    \operatorname{Prim}^{\mathrm{der}}_{\mathfrak G_1}
    \left(t^*H^{\mathrm{iso}}_{\mathfrak G}\right).
\]
Together with Definition~\ref{con:bnla-primitive-base-change}, this produces
a canonical coherent diagram
\[
\begin{tikzcd}[column sep=large]
s^*\mathfrak l^{\mathrm{sp}}_{\mathfrak G}
  \arrow[r,"\beta^{\mathrm{Prim}}_{s}"] &
\operatorname{Prim}(s^*H^{\mathrm{iso}}_{\mathfrak G})
  \arrow[r,"\sim"] &
\operatorname{Prim}(t^*H^{\mathrm{iso}}_{\mathfrak G}) &
 t^*\mathfrak l^{\mathrm{sp}}_{\mathfrak G}
  \arrow[l,"\beta^{\mathrm{Prim}}_{t}"'].
\end{tikzcd}
\]
No base-change equivalence for the right adjoint $\operatorname{Prim}$ is
postulated.
\end{theorem}

\begin{proof}
Apply the functor of derived primitives over $\mathfrak G_1$ to the Hopf
equivalence of Proposition~\ref{prop:bnla-hopf-adjoint-transport}.  The outer
maps are the always-defined comparisons of
Definition~\ref{con:bnla-primitive-base-change}.  Higher coherence follows
by applying the same functors to the higher conjugation diagrams.
\end{proof}

\begin{corollary}[Literal pullback transport when the comparisons are equivalences]
\label{cor:bnla-literal-primitive-transport}
Whenever a separate theorem identifies both
$\beta^{\mathrm{Prim}}_{s}$ and $\beta^{\mathrm{Prim}}_{t}$ as equivalences,
the preceding diagram canonically contracts to an equivalence
\[
    s^*\mathfrak l^{\mathrm{sp}}_{\mathfrak G}
    \xrightarrow{\sim}
    t^*\mathfrak l^{\mathrm{sp}}_{\mathfrak G}.
\]
\end{corollary}

\begin{proof}
Invert the two outer equivalences and compose.  The coherence follows from the
coherent diagram of Theorem~\ref{thm:bnla-differentiated-adjoint}.
\end{proof}

\section{The intrinsic differentiated package}
\label{sec:bnla-intrinsic-package}

\begin{definition}[Intrinsic differentiated package]
\label{def:bnla-intrinsic-package}
Write
\[
    \operatorname{Diff}^{\mathrm{st}}(\mathfrak G)
\]
for the canonically constructed tuple
\[
\left(
\mathbb A^{\mathrm{st}}_{\mathfrak G},
 a^{\mathrm{st}}_{\mathfrak G},
\mathbb F^{\mathrm{st}}_{\mathfrak G},
\mathbb N^{\mathrm{st}}_{\mathfrak G},
\mathbb I^{\mathrm{st}}_{\mathfrak G},
\gamma^{\mathrm{st}}_{\mathfrak G},
\mathbb C^{\mathrm{iso,st}}_{\mathfrak G},
H^{\mathrm{iso}}_{\mathfrak G},
\mathfrak l^{\mathrm{sp}}_{\mathfrak G}
\right),
\]
together with the Hopf adjoint transport, primitive comparison maps, and all
constructed coherence.
\end{definition}

\begin{remark}[Package name, not an axiomatic category]
The notation is only an abbreviation.  No abstract object is defined by
listing fields or imposing compatibility axioms.  Every component is the
output of a functorial construction from $\mathfrak G$.
\end{remark}

\begin{theorem}[Intrinsic differentiated package theorem]
\label{thm:bnla-intrinsic-package}
The package $\operatorname{Diff}^{\mathrm{st}}(\mathfrak G)$ is functorial
with the variance of
Theorem~\ref{thm:bnla-intrinsic-stable-differentiation}.  Under arbitrary
spectral base change, the stable arrow coefficient, relation coefficient,
anchor triangle, actual-isotropy coefficient, isotropy comparison,
discrepancy, and suspension Hopf object commute with pullback by canonical
equivalences.  The spectral Lie object has the canonical comparison
\[
    c_X^*\mathfrak l^{\mathrm{sp}}_{\mathfrak G}
    \longrightarrow
    \mathfrak l^{\mathrm{sp}}_{\mathfrak G'}
\]
constructed as $\beta^{\mathrm{Prim}}_{c_X,H^{\mathrm{iso}}_{\mathfrak G}}$.
\end{theorem}

\begin{proof}
Combine Theorems~\ref{thm:bnla-intrinsic-stable-differentiation} and
\ref{thm:bnla-stable-anchor-triangle},
Proposition~\ref{prop:bnla-isotropy-comparison-functoriality},
Theorem~\ref{thm:bnla-isotropy-augmentation-ideal}, and
Definition~\ref{con:bnla-primitive-base-change}.
\end{proof}

\section{The represented first-order interface without quotient presentation data}
\label{sec:bnla-represented-first-order}

The intrinsic stable package is complete without representability.  We now
return to the represented range already constructed at the end of the preceding
chapter in order to identify the first structured coefficient and compare it
with effective quotient geometry.

Let $S$, $X$, and $\mathfrak G_1$ be spectral schemes representing the
corresponding native objects, with
\[
    \mathfrak G_1
    \mathop{\rightrightarrows}^{s}_{t}
    X,
    \qquad
    u:X\to\mathfrak G_1.
\]
No representation of $Q_{\mathfrak G}$ is used initially.

\subsection{Vertical split deformations of the unit}

For $M\in\operatorname{QCoh}(X)_{\geq0}$, let
\[
    X\xrightarrow{i_M}X[M]\xrightarrow{p_M}X
\]
be the split structured square-zero object constructed from $M$ by the global
square-zero theory of Chapter~2.

\begin{definition}[Relation-vertical deformation space]
\label{def:bnla-relation-vertical-deformations}
Define
\[
    \operatorname{Der}^{\mathrm{vert}}_{\mathfrak G}(M)
\]
to be the mapping space of pointed objects over $X$
\[
    \operatorname{Map}_{(\operatorname{Sch}^{\mathrm{sp}}_{/X})_*}
    \left(
        (X[M],i_M),
        (\mathfrak G_1,u)
    \right).
\]
Equivalently, it is the space of morphisms
\[
    h:X[M]\to\mathfrak G_1
\]
with specified compatibilities
\[
    sh\simeq p_M,
    \qquad
    hi_M\simeq u.
\]
\end{definition}

\begin{theorem}[The source coefficient corepresents vertical deformations]
\label{thm:bnla-source-coefficient-corepresents}
There is a natural equivalence
\[
    \boxed{
    \operatorname{Der}^{\mathrm{vert}}_{\mathfrak G}(M)
    \simeq
    \operatorname{Map}_{\operatorname{QCoh}(X)}
    \left(
        \mathfrak a^\vee_{\mathfrak G},M
    \right),}
\]
where
\[
    \mathfrak a^\vee_{\mathfrak G}
    =
    u^*L_{\mathfrak G_1/X,s}.
\]
The equivalence is natural in $M$, formal-groupoid morphisms in the represented
range, and represented base change.
\end{theorem}

\begin{proof}
The relative cotangent universal property for
$s:\mathfrak G_1\to X$ identifies maps from a split square-zero object over
$X$ centred at the unit with relative derivations at $u$.  Explicitly, after
pullback along $u$, the universal relative derivation gives
\[
    \operatorname{Der}^{\mathrm{vert}}_{\mathfrak G}(M)
    \simeq
    \operatorname{Map}_{\operatorname{QCoh}(X)}
    \left(
        u^*L_{\mathfrak G_1/X,s},M
    \right).
\]
This is the displayed formula.  Naturality is the functoriality of the
cotangent--derivation equivalence and of the constructed split square-zero
objects.
\end{proof}

\begin{remark}[No quotient representation is needed]
The coefficient $\mathfrak a^\vee_{\mathfrak G}$ is therefore intrinsically
characterized by the represented arrow relation itself.  Representability of
the effective quotient is only needed for the comparison theorem below, not
for construction of the coefficient or anchor.
\end{remark}

\subsection{Comparison with a represented effective quotient}

Suppose now that the already constructed effective quotient
\[
    q_{\mathfrak G}:X\twoheadrightarrow Q_{\mathfrak G}
\]
is represented by a spectral scheme.  Groupoid effectivity gives
\[
    \mathfrak G_1
    \simeq
    X\times_{Q_{\mathfrak G}}X,
\]
with $s=\operatorname{pr}_1$ and $t=\operatorname{pr}_2$.

\begin{theorem}[Relative cotangent identification]
\label{thm:bnla-relative-cotangent-identification}
There are canonical equivalences
\[
    L_{\mathfrak G_1/X,s}
    \simeq
    t^*L_{X/Q_{\mathfrak G}}
\]
and
\[
    \boxed{
    \mathfrak a^\vee_{\mathfrak G}
    \simeq
    L_{X/Q_{\mathfrak G}}.}
\]
Under the boxed equivalence, the cotangent anchor of the preceding chapter
\[
    a^\vee_{\mathfrak G}:
    L_{X/S}\longrightarrow\mathfrak a^\vee_{\mathfrak G}
\]
is the transitivity morphism
\[
    L_{X/S}\longrightarrow L_{X/Q_{\mathfrak G}}
\]
for $X\to Q_{\mathfrak G}\to S$.
\end{theorem}

\begin{proof}
The Cartesian square
\[
\begin{tikzcd}
\mathfrak G_1 \arrow[r,"t"] \arrow[d,"s"'] &
X \arrow[d,"q_{\mathfrak G}"]\\
X \arrow[r,"q_{\mathfrak G}"'] &
Q_{\mathfrak G}
\end{tikzcd}
\]
gives the first equivalence by cotangent base change.  Pulling back along
$u$ and using $tu\simeq\operatorname{id}_X$ gives the second.

The first-order anchor in the preceding chapter is induced by
$(s,t):\mathfrak G_1\to X\times_SX$ and the target component.  Naturality of
cotangent transitivity for the displayed Cartesian square identifies it with
the transitivity map for $X\to Q_{\mathfrak G}\to S$.
\end{proof}

\begin{corollary}[Internal-dual tangent comparison]
\label{cor:bnla-relative-tangent-identification}
Internal duality gives
\[
    \mathfrak a_{\mathfrak G}
    \simeq
    \left(L_{X/Q_{\mathfrak G}}\right)^\vee
\]
and the tangent anchor
\[
    \left(L_{X/Q_{\mathfrak G}}\right)^\vee
    \longrightarrow
    L_{X/S}^\vee.
\]
When the two cotangent complexes are finite locally free this is the usual
tangent map
\[
    T_{X/Q_{\mathfrak G}}
    \longrightarrow
    T_{X/S}.
\]
\end{corollary}

\begin{proof}
Apply the internal-Hom functor
$\underline{\operatorname{Hom}}_X(-,\mathcal O_X)$ to
Theorem~\ref{thm:bnla-relative-cotangent-identification}.  In the finite
locally free range internal duality agrees with dualizability.
\end{proof}

\section{First structured neighbourhoods and the complete stable anchor}
\label{sec:bnla-first-order-vs-stable}

Retain the constructed first neighbourhoods
\[
    X\xrightarrow{i_u}D^{(1)}_{\mathfrak G}
    \xrightarrow{\nu_u}\mathfrak G_1
\]
and
\[
    X\xrightarrow{i_1}D^{(1)}_{X/S}
    \xrightarrow{\nu_1}X\times_SX,
\]
together with
\[
    d^{(1)}(s,t):
    D^{(1)}_{\mathfrak G}
    \longrightarrow
    D^{(1)}_{X/S}
\]
and
\[
    \ell_1:
    D^{(1)}_{X/S}
    \longrightarrow
    R_{X/S}.
\]

\begin{definition}[First-order stable coefficients]
\label{def:bnla-first-order-stable-coefficients}
Define
\[
    \mathbb A^{(1),\mathrm{st}}_{\mathfrak G}
    :=
    \widetilde\Sigma^\infty_X
    \left(D^{(1)}_{\mathfrak G},i_u\right),
\]
\[
    \mathbb T^{(1),\mathrm{st}}_{X/S}
    :=
    \widetilde\Sigma^\infty_X
    \left(D^{(1)}_{X/S},i_1\right).
\]
\end{definition}

\begin{theorem}[First-order stable comparison]
\label{thm:bnla-first-order-stable-comparison}
There is a canonical homotopy-commutative square
\[
\begin{tikzcd}
\mathbb A^{(1),\mathrm{st}}_{\mathfrak G}
  \arrow[r]
  \arrow[d,"a^{(1),\mathrm{st}}_{\mathfrak G}"'] &
\mathbb A^{\mathrm{st}}_{\mathfrak G}
  \arrow[d,"a^{\mathrm{st}}_{\mathfrak G}"]\\
\mathbb T^{(1),\mathrm{st}}_{X/S}
  \arrow[r] &
\mathbb T^{\mathrm{st}}_{X/S}.
\end{tikzcd}
\]
The left vertical map is the reduced stabilization of
$d^{(1)}(s,t)$; the horizontal maps are induced by $\nu_u$ and $\ell_1$.
The coefficient morphism of the underlying structured first-neighbourhood
comparison is exactly
\[
    a^\vee_{\mathfrak G}:
    L_{X/S}\longrightarrow\mathfrak a^\vee_{\mathfrak G}.
\]
\end{theorem}

\begin{proof}
The preceding chapter proves that the two composites in
\[
\begin{tikzcd}
D^{(1)}_{\mathfrak G}
  \arrow[r,"\nu_u"]
  \arrow[d,"{d^{(1)}(s,t)}"'] &
\mathfrak G_1
  \arrow[d,"\phi_{\mathfrak G,1}"]\\
D^{(1)}_{X/S}
  \arrow[r,"\ell_1"'] &
R_{X/S}
\end{tikzcd}
\]
are canonically homotopic through a contractible space of compatible lifts.
Apply reduced fibrewise stabilization.  The coefficient statement is the
construction of $d^{(1)}(s,t)$.
\end{proof}

\begin{remark}[No first-order collapse]
Neither horizontal map is asserted to be an equivalence.  The first-order
coefficient is a constructed shadow of the stable anchor of the complete arrow
object, not a presentation of it.
\end{remark}

\subsection{First-order invariance under unitwise formalization}

The preceding chapter formalizes every fixed-object internal groupoid
$\mathcal H_\bullet$ by unitwise centred completion.  The completed arrow
object need not be represented, so its first-order invariance is formulated by
the split square-zero probe functor rather than by assigning it an
unconstructed cotangent complex.

\begin{definition}[Split first-order probe functor]
\label{def:bnla-split-probe-functor}
Let
\[
    X\xrightarrow{e}Y\xrightarrow{p}X,
    \qquad pe\simeq\operatorname{id}_X,
\]
be a pointed native object over a represented spectral scheme $X$.  Define
\[
    \operatorname{Probe}^{(1)}_{(Y,e)}:
    \operatorname{QCoh}(X)_{\geq0}
    \longrightarrow\mathcal S
\]
by
\[
    \operatorname{Probe}^{(1)}_{(Y,e)}(M)
    :=
    \operatorname{Map}_{(\mathcal X_{/X})_*}
    \left(
        (X[M],i_M),(Y,e)
    \right),
\]
where $X[M]$ is the constructed split structured square-zero object over $X$.
When this functor is corepresented by a quasi-coherent module $K$, we call $K$
the split first-order coefficient of $(Y,e)$.
\end{definition}

\begin{lemma}[Retractive centred coreflection]
\label{lem:bnla-retractive-centred-coreflection}
Let
\[
    X\xrightarrow{e}Y\xrightarrow{p}X,
    \qquad pe\simeq\operatorname{id}_X,
\]
be a pointed object and give the centred completion
$C_X(Y,e)$ the retraction
\[
    p^c:
    C_X(Y,e)\xrightarrow{\kappa_{Y,e}}Y\xrightarrow{p}X.
\]
Then
\[
    X\xrightarrow{\widetilde e}C_X(Y,e)\xrightarrow{p^c}X
\]
is again pointed.  If
\[
    X\xrightarrow{d}Z\xrightarrow{r}X,
    \qquad rd\simeq\operatorname{id}_X,
\]
is pointed and the centred object $(Z,d)$ is relative de Rham-complete, then
composition with the centred-completion counit induces a canonical equivalence
\[
\boxed{
    \operatorname{Map}_{(\mathcal X_{/X})_*}
    \left(
        (Z,d),(C_X(Y,e),\widetilde e)
    \right)
    \xrightarrow{\sim}
    \operatorname{Map}_{(\mathcal X_{/X})_*}
    \left(
        (Z,d),(Y,e)
    \right).}
\]
The equivalence is natural in maps of pointed targets and complete pointed
sources and is coherently compatible with identities and composition.
\end{lemma}

\begin{proof}
The preceding chapter gives, in the undercategory $X/$, the centred-coreflection
equivalence
\[
    \kappa_{Y,e}\circ(-):
    \operatorname{Map}_{X/}
    \left(
        Z,C_X(Y,e)
    \right)
    \xrightarrow{\sim}
    \operatorname{Map}_{X/}(Z,Y)
\]
for every centred-complete $(Z,d)$.  Postcomposition with the retractions gives
a homotopy-commutative square
\[
\begin{tikzcd}
\operatorname{Map}_{X/}(Z,C_X(Y,e))
  \arrow[r,"\sim"]
  \arrow[d,"{(p^c)_*}"'] &
\operatorname{Map}_{X/}(Z,Y)
  \arrow[d,"p_*"]\\
\operatorname{Map}_{X/}(Z,X)
  \arrow[r,equal] &
\operatorname{Map}_{X/}(Z,X).
\end{tikzcd}
\]
The mapping space in the pointed category is the homotopy fibre of the
corresponding vertical map over the specified retraction $r:Z\to X$.  Taking
homotopy fibres therefore gives the displayed equivalence.  Naturality and all
higher coherence are inherited from the centred coreflection and functoriality
of homotopy fibres.
\end{proof}

\begin{theorem}[First-order probe invariance under centred completion]
\label{thm:bnla-first-probe-centred-completion}
Let $X$ be represented and let $(X\xrightarrow{e}Y\to X)$ be a pointed object
whose uncompleted target $Y$ is represented.  Then, for every connective
$M\in\operatorname{QCoh}(X)$, composition with the centred-completion counit
induces a canonical equivalence
\[
    \operatorname{Probe}^{(1)}_{(C_X(Y,e),\widetilde e)}(M)
    \xrightarrow{\sim}
    \operatorname{Probe}^{(1)}_{(Y,e)}(M).
\]
Consequently the completed pointed object has split first-order coefficient
\[
    e^*L_{Y/X,p}.
\]
If $C_X(Y,e)$ is itself represented, then there is a canonical equivalence
\[
    \widetilde e^*L_{C_X(Y,e)/X}
    \xrightarrow{\sim}
    e^*L_{Y/X,p}.
\]
\end{theorem}

\begin{proof}
For a connective module $M$, the centre
\[
    i_M:X\longrightarrow X[M]
\]
is a structured square-zero thickening.  The relative de Rham localization of
Chapters~9--10 inverts every such thickening, so $(X[M],i_M)$ is complete as a
centred object.  Lemma~\ref{lem:bnla-retractive-centred-coreflection} therefore
gives
\[
\begin{aligned}
&\operatorname{Map}_{(\mathcal X_{/X})_*}
\left(
    (X[M],i_M),(C_X(Y,e),\widetilde e)
\right)
\\
&\qquad\simeq
\operatorname{Map}_{(\mathcal X_{/X})_*}
\left(
    (X[M],i_M),(Y,e)
\right).
\end{aligned}
\]
This is the first assertion.

Because $Y\to X$ is represented, the relative cotangent universal property
identifies the right-hand mapping space with
\[
    \operatorname{Map}_{\operatorname{QCoh}(X)}
    \left(
        e^*L_{Y/X,p},M
    \right).
\]
Thus the split first-order probe functor of the completion is corepresented by
$e^*L_{Y/X,p}$.

If the completion is represented, its own relative cotangent complex
corepresents the same functor.  Uniqueness of corepresenting objects gives the
final equivalence.  Every comparison is natural in pointed represented maps,
and represented base change preserves the split square-zero probes, centred
completion, and the cotangent--derivation equivalence.
\end{proof}

\begin{theorem}[The first structured unit neighbourhood survives formalization]
\label{thm:bnla-first-neighbourhood-formalization}
Let $\mathcal H_\bullet\rightrightarrows X$ be a fixed-object groupoid such
that $X$ and $\mathcal H_1$ are represented spectral schemes.  Let
\[
    X\xrightarrow{i_u}D^{(1)}_{\mathcal H}
    \xrightarrow{\nu_u}\mathcal H_1
\]
be the constructed first structured neighbourhood of the unit.  Then
$\nu_u$ has a contractibly unique pointed lift
\[
    \boxed{
    \widehat\nu_u:
    D^{(1)}_{\mathcal H}
    \longrightarrow
    C_X(\mathcal H_1,u)
    =
    \widehat{\mathcal H}^{\,\mathrm{dR}}_1}
\]
through the centred-completion counit.  Moreover, for every connective
$M\in\operatorname{QCoh}(X)$,
\[
\boxed{
\operatorname{Map}_{(\mathcal X_{/X})_*}
\left(
    (X[M],i_M),
    (\widehat{\mathcal H}^{\,\mathrm{dR}}_1,u)
\right)
\simeq
\operatorname{Map}_{\operatorname{QCoh}(X)}
\left(
    u^*L_{\mathcal H_1/X,s},M
\right).}
\]
Thus the first-order probe functor of the formalized groupoid is canonically
corepresented by the already constructed coefficient
$u^*L_{\mathcal H_1/X,s}$, whether or not the completed arrow object is
represented.

Let
\[
    C_X(s,t):
    C_X(\mathcal H_1,u)
    \longrightarrow
    C_X(X\times_SX,\Delta)
    \simeq
    R_{X/S}
\]
be the completed source--target map.  There is a contractibly unique compatible
homotopy
\[
\boxed{
    C_X(s,t)\widehat\nu_u
    \simeq
    \ell_1d^{(1)}(s,t)}
\]
whose image under the completion counit
$R_{X/S}\to X\times_SX$ is the specified source--target comparison.  Hence the
coefficient morphism
\[
    L_{X/S}\longrightarrow u^*L_{\mathcal H_1/X,s}
\]
of $d^{(1)}(s,t)$ is the first structured coefficient of the canonical
nonlinear anchor of the unitwise formalization.
\end{theorem}

\begin{proof}
The centre
\[
    X\longrightarrow D^{(1)}_{\mathcal H}
\]
is a structured square-zero thickening and is therefore a relative de Rham
equivalence.  Hence the pointed source
$(D^{(1)}_{\mathcal H},i_u)$ is centred-complete.  Apply
Lemma~\ref{lem:bnla-retractive-centred-coreflection} to the pointed target
$(\mathcal H_1,u)$: the point $\nu_u$ has a contractibly unique pointed
preimage, which is $\widehat\nu_u$.

The split-probe formula is
Theorem~\ref{thm:bnla-first-probe-centred-completion} applied to
$(\mathcal H_1,u)$.

For the anchor comparison, put
\[
    Y:=X\times_SX
\]
with first projection to $X$.  Functoriality of centred completion applied to
$(s,t):(\mathcal H_1,u)\to(Y,\Delta)$ gives
\[
    C_X(s,t):C_X(\mathcal H_1,u)\to C_X(Y,\Delta)\simeq R_{X/S}.
\]
Both composites
\[
    C_X(s,t)\widehat\nu_u,
    \qquad
    \ell_1d^{(1)}(s,t):
    D^{(1)}_{\mathcal H}\rightrightarrows R_{X/S}
\]
map under the centred-completion counit
$R_{X/S}\to X\times_SX$ to the same specified map
\[
    (s,t)\nu_u:
    D^{(1)}_{\mathcal H}\longrightarrow X\times_SX.
\]
Since $D^{(1)}_{\mathcal H}$ is centred-complete, composition with that counit
induces an equivalence
\[
    \operatorname{Map}_{X/}
    \left(D^{(1)}_{\mathcal H},R_{X/S}\right)
    \xrightarrow{\sim}
    \operatorname{Map}_{X/}
    \left(D^{(1)}_{\mathcal H},X\times_SX\right).
\]
The homotopy fibre over $(s,t)\nu_u$ is therefore contractible.  The two
composites are points of this fibre, so the space of compatible paths between
them is contractible.  Finally $d^{(1)}(s,t)$ was constructed from the
coefficient morphism
$L_{X/S}\to u^*L_{\mathcal H_1/X,s}$, which is therefore the first structured
coefficient of the completed nonlinear anchor.  Naturality and represented
base change follow from the same coreflection, cotangent, and structured
square-zero base-change comparisons.
\end{proof}

\section{Finite spectral quotient-lifting functors}
\label{sec:bnla-formal-moduli-envelope}

The nonlinear quotient
\[
    q_{\mathfrak G}:X\twoheadrightarrow Q_{\mathfrak G}
\]
already determines its finite infinitesimal lifting spaces.  We construct
those spaces here on the full finite spectral infinitesimal test category of
Chapter~9.  The coherent Artinian subcategory and the Schlessinger reflection
used for algebraic recognition are constructed later from intrinsic
finiteness conditions; no finite factorization is retained at either stage.

\subsection{Centred finite spectral tests}

Fix a represented affine centre
\[
    x:U\longrightarrow X
\]
over $S$.

\begin{definition}[Centred finite infinitesimal tests]
\label{def:bnla-centred-inf-tests}
Let
\[
    \operatorname{Inf}^{\mathrm{fin}}_{U/S}
\]
be the $\infty$-category whose objects are finite spectral infinitesimal
thickenings
\[
    j:U\longrightarrow T
\]
over $S$ and whose morphisms are commutative triangles over $S$ restricting
to $\operatorname{id}_U$ on the centre.  The phrase ``finite infinitesimal
thickening'' has the meaning fixed in Chapter~9: it is an arrow in the
equivalence-closed class generated by structured spectral square-zero
immersions under finite composition.  A factorization, coefficient module,
and classifying derivation are not part of an object.
\end{definition}

Write
\[
    q_x:=q_{\mathfrak G}x:U\longrightarrow Q_{\mathfrak G}.
\]

\begin{definition}[Raw spectral quotient-leaf functor]
\label{def:bnla-raw-formal-leaf}
Define
\[
    \mathcal L^{\mathrm{raw}}_{\mathfrak G,x}:
    \left(\operatorname{Inf}^{\mathrm{fin}}_{U/S}\right)^{\mathrm{op}}
    \longrightarrow\mathcal S
\]
by
\[
\boxed{
    \mathcal L^{\mathrm{raw}}_{\mathfrak G,x}(U\hookrightarrow T)
    :=
    \operatorname{fib}_{q_x}
    \left(
        \operatorname{Map}_S(T,Q_{\mathfrak G})
        \longrightarrow
        \operatorname{Map}_S(U,Q_{\mathfrak G})
    \right).}
\]
A point is an extension of the quotient map from the centre to the finite
infinitesimal thickening, together with the specified homotopy on the centre.
\end{definition}

\begin{proposition}[Functoriality of the raw leaf functor]
\label{prop:bnla-raw-leaf-functoriality}
The raw leaf construction is contravariantly functorial in finite
infinitesimal tests and functorial in the centre and in formal-groupoid
morphisms.  For an arbitrary spectral base change $c:S'\to S$, write
\[
    X':=X\times_SS',
    \qquad
    \mathfrak G':=\mathfrak G\times_SS'.
\]
If $x':U'\to X'$ is an affine centre and $U'\hookrightarrow T'$ is a finite
infinitesimal test over $S'$, then restriction of scalars along $c$ gives a
canonical equivalence
\[
\boxed{
    \mathcal L^{\mathrm{raw}}_{\mathfrak G',x'}(U'\hookrightarrow T')
    \simeq
    \mathcal L^{\mathrm{raw}}_{\mathfrak G,cx'}
    (U'\hookrightarrow T'),}
\]
where the right-hand side regards $U'$ and $T'$ as objects over $S$.  Starting
instead with a test over $S$ and base changing the test to $S'$ gives the
canonical mate comparison; no equivalence for that mate is asserted in
general.  Moreover
\[
    \mathcal L^{\mathrm{raw}}_{\mathfrak G,x}(\operatorname{id}_U)
    \simeq *.
\]
\end{proposition}

\begin{proof}
A morphism of thickenings restricts an extension to the smaller target.
Functoriality in $x$ and $\mathfrak G$ is composition with the corresponding
quotient maps.

For base change, the preceding chapter gives
\[
    Q_{\mathfrak G'}
    \simeq
    Q_{\mathfrak G}\times_SS'.
\]
For every $S'$-object $T'$ the pullback--restriction adjunction in the slice
therefore gives
\[
    \operatorname{Map}_{S'}(T',Q_{\mathfrak G'})
    \simeq
    \operatorname{Map}_{S}(T',Q_{\mathfrak G}),
\]
where on the right $T'$ is regarded over $S$; the same formula holds for
$U'$.  These equivalences carry the base point $q_{x'}$ to the restricted
base point $q_{cx'}$, and hence induce the displayed equivalence on homotopy
fibres.  If a test over $S$ is first pulled back to $S'$, the unit of the
pullback--restriction adjunction instead supplies the stated mate comparison.

At the identity thickening the map whose fibre is taken is the identity of
$\operatorname{Map}_S(U,Q_{\mathfrak G})$.
\end{proof}

\begin{lemma}[Connective pullbacks along finite infinitesimal extensions]
\label{lem:bnla-connective-pullback-infinitesimal}
Consider a pullback square in commutative ring spectra
\[
\begin{tikzcd}
A \arrow[r] \arrow[d] & A_0 \arrow[d]\\
A_1 \arrow[r] & A_{01}
\end{tikzcd}
\]
in which the three displayed input rings are connective and one of the maps
$A_i\to A_{01}$ is a finite spectral infinitesimal extension.  Then $A$ is
connective.  Consequently the square is also a pullback in connective
commutative ring spectra.
\end{lemma}

\begin{proof}
It is enough to treat one structured square-zero extension and then iterate.
For a structured square-zero map $A_0\to A_{01}$, Chapter~2 gives
surjectivity on $\pi_0$.  On underlying spectra there is a fibre sequence
\[
    A\longrightarrow A_0\oplus A_1\longrightarrow A_{01}.
\]
The only possible negative group is
\[
    \pi_{-1}A
    \simeq
    \operatorname{coker}
    \bigl(\pi_0A_0\oplus\pi_0A_1\to\pi_0A_{01}\bigr),
\]
which vanishes because one component is surjective.  All lower negative
homotopy groups vanish by connectivity of the inputs.  Finite composition
gives the result.
\end{proof}

\subsection{Reduced-classical consequence of de Rham formality}

\begin{proposition}[Formal quotients are reduced-classical equivalences]
\label{prop:bnla-formal-quotient-reduced-equivalence}
For every formal Lie groupoid,
\[
    D_{\mathrm{red}}(q_{\mathfrak G}):
    D_{\mathrm{red}}X
    \xrightarrow{\sim}
    D_{\mathrm{red}}Q_{\mathfrak G}
\]
and
\[
    D_{\mathrm{red}}(r_{\mathfrak G}):
    D_{\mathrm{red}}Q_{\mathfrak G}
    \xrightarrow{\sim}
    D_{\mathrm{red}}Q_{X/S}
\]
are equivalences.
\end{proposition}

\begin{proof}
The preceding chapter proves that both quotient maps are inverted by relative
de Rham localization.  Chapter~9 constructs $D_{\mathrm{red}}$ as a further
left-exact localization receiving the de Rham localization.  Hence every
relative de Rham equivalence is sent to a reduced-classical equivalence.
\end{proof}
\section{Coherent spectral coefficients and pro-coherent duality}
\label{sec:bnla-spectral-procoherent}

We now construct the many-object spectral partition-Lie coefficient theory
needed to differentiate the formal-leaf functor.  The construction is local on
coherent spectral coefficient charts and will later be descended along
spectral etale maps.

\subsection{Coherent affine spectral coefficients}

Fix a morphism of connective commutative ring spectra
\[
    R\longrightarrow B.
\]

\begin{definition}[Coherent spectral coefficient ring]
\label{def:bnla-coherent-spectral-ring}
A connective $E_\infty$-ring $B$ is \emph{coherent} when $\pi_0(B)$ is a
coherent ordinary ring and every $\pi_n(B)$ is finitely presented over
$\pi_0(B)$.  We write
\[
    \operatorname{APerf}(B)\subset\operatorname{Mod}_B
\]
for the standard stable $\infty$-category of almost perfect $B$-modules and
\[
    \operatorname{Coh}(B)\subset\operatorname{APerf}(B)
\]
for its eventually coconnective coherent objects, in the normalization of
\cite{BrantnerCamposNuiten2025}.
\end{definition}

\begin{definition}[Almost finite presentation]
\label{def:bnla-spectral-almost-finite-presentation}
Let $A\to C$ be a morphism of connective $E_\infty$-rings.  We say that
$C$ is \emph{almost finitely presented} over $A$ if, for every integer $n$,
the functor represented by $C$ on $n$-truncated connective commutative
$A$-algebras preserves filtered colimits.  Equivalently, $C$ is almost
compact in the connective commutative $A$-algebra category in the sense of
Postnikov truncations.  \end{definition}

\begin{definition}[Pro-coherent spectral modules and continuous dual]
\label{def:bnla-procoherent-spectral}
For coherent $B$ define
\[
    \operatorname{QC}^{\vee}_B
    :=
    \operatorname{Ind}
    \left(
        \operatorname{Coh}(B)^{\mathrm{op}}
    \right).
\]
For an arbitrary $B$-module $M$, define its continuous pro-coherent dual
\[
    M^\vee\in\operatorname{QC}^{\vee}_B
\]
by the exact functor
\[
    K\longmapsto
    \operatorname{map}_{B}(M,K),
    \qquad
    K\in\operatorname{Coh}(B).
\]
When $M$ is almost perfect this agrees with the continuous duality functor on
the almost-perfect sector.  Put
\[
    \mathbb T^{\mathrm{pc}}_{B/R}
    :=
    \left(L_{B/R}\right)^\vee.
\]
\end{definition}

\begin{proposition}[Pro-coherent coefficient properties]
\label{prop:bnla-procoherent-properties}
For coherent $B$, the category $\operatorname{QC}^{\vee}_B$ is stable,
presentable, and compactly generated by continuous duals of coherent
$B$-modules.  Continuous duality restricts to a fully faithful functor
\[
    \operatorname{APerf}(B)^{\mathrm{op}}
    \longrightarrow
    \operatorname{QC}^{\vee}_B.
\]
The pro-coherent composition product used for spectral partition-Lie
operations preserves sifted colimits in the algebra variable.
\end{proposition}

\begin{proof}
These are the pro-coherent spectral coefficient theorems of
\cite{BrantnerCamposNuiten2025}.  The formula of
Definition~\ref{def:bnla-procoherent-spectral} is the exact-functor
description of the continuous dual.  On almost-perfect modules the resulting
functor is fully faithful, while continuous duals of coherent modules are
compact generators of the Ind-completion.
\end{proof}

\begin{lemma}[Finite-Tor pro-coherent Beck--Chevalley]
\label{lem:bnla-procoherent-finite-tor-bc}
Consider a homotopy-pushout square of coherent connective $E_\infty$-rings
\[
\begin{tikzcd}
B \arrow[r,"f"] \arrow[d,"g"'] &
B' \arrow[d,"g'"]\\
C \arrow[r,"f'"'] &
C'=C\otimes_BB'.
\end{tikzcd}
\]
Suppose that $f$ has finite Tor-amplitude; in particular this applies when
$f$ is spectral etale.  Then extension of scalars on coherent modules induces
a colimit-preserving functor
\[
    f^*:\operatorname{QC}^{\vee}_B
    \longrightarrow
    \operatorname{QC}^{\vee}_{B'}
\]
with a pro-coherent right adjoint $f_*^{\mathrm{pc}}$.  The canonical
Beck--Chevalley transformation
\[
\boxed{
    g^*f_*^{\mathrm{pc}}
    \longrightarrow
    (f')_*^{\mathrm{pc}}(g')^*
}
\]
is an equivalence.  Moreover, for every almost-perfect $B$-module $M$,
\[
\boxed{
    f^*(M^\vee)
    \simeq
    (B'\otimes_BM)^\vee.}
\]
All these equivalences are coherent for identity squares, horizontal and
vertical composition, and iterated pullback.
\end{lemma}

\begin{proof}
Use the right-left exact-functor model of pro-coherent modules of
\cite{BrantnerCamposNuiten2025}, which in the coherent case agrees with the
exact-functor and Ind-descriptions of
Proposition~\ref{prop:bnla-procoherent-properties}.  In this model pullback
along a morphism of connective coefficient rings is induced by precomposition
with restriction of scalars on the eventually coconnective module sector;
write this sector as $\operatorname{Mod}^{+}$.  If $f$ has finite
Tor-amplitude, then its pro-coherent right adjoint is induced by precomposition with extension of scalars
\[
    f^*_{\mathrm{mod}}:
    \operatorname{Mod}^{+}_B
    \longrightarrow
    \operatorname{Mod}^{+}_{B'},
    \qquad
    N\longmapsto B'\otimes_BN.
\]
If $f$ has Tor-amplitude contained in a finite interval, extension of scalars
is exact, preserves filtered colimits, and carries an $n$-coconnective module
to an $(n+d)$-coconnective module for a bound $d$ independent of the module.
Consequently precomposition preserves the almost-finite-presentation
conditions defining the right-left pro-coherent subcategory.  This constructs
$f_*^{\mathrm{pc}}$ and, because it is pointwise precomposition, shows that
it preserves colimits.

Let $F\in\operatorname{QC}^{\vee}_{B'}$ and let
$N\in\operatorname{Mod}^{+}_C$.  Evaluate the canonical
Beck--Chevalley transformation on $N$.  By the preceding pointwise
presentations its source is
\[
    F\left(B'\otimes_B\operatorname{Res}^{C}_{B}N\right),
\]
while its target is
\[
    F\left(
        \operatorname{Res}^{C'}_{B'}
        (C'\otimes_CN)
    \right).
\]
Since $C'=C\otimes_BB'$, associativity of relative tensor products gives a
canonical equivalence of $B'$-modules
\[
    B'\otimes_B\operatorname{Res}^{C}_{B}N
    \xrightarrow{\sim}
    \operatorname{Res}^{C'}_{B'}(C'\otimes_CN).
\]
The evaluated Beck--Chevalley map is exactly the image of this module
base-change equivalence under $F$.  It is therefore an equivalence for every
$F$ and $N$, which proves
\[
    g^*f_*^{\mathrm{pc}}
    \xrightarrow{\sim}
    (f')_*^{\mathrm{pc}}(g')^*.
\]

For the continuous-dual formula, let $K'\in\operatorname{Coh}(B')$.
Then
\[
\begin{aligned}
    (f^*M^\vee)(K')
    &\simeq
    \operatorname{map}_B
    \left(M,\operatorname{Res}^{B'}_BK'\right)\\
    &\simeq
    \operatorname{map}_{B'}
    \left(B'\otimes_BM,K'\right)\\
    &=
    (B'\otimes_BM)^\vee(K').
\end{aligned}
\]
Yoneda in the exact-functor model gives
\[
    f^*(M^\vee)\simeq(B'\otimes_BM)^\vee.
\]
Every comparison above is induced by restriction of scalars and the canonical
associativity and unit equivalences for tensor products.  Their standard
pasting coherences therefore give the identity, horizontal, vertical, and
iterated-pullback coherences claimed in the statement.
\end{proof}

\begin{remark}[Conservativity is obtained from etale descent]
\label{rem:bnla-procoherent-conservativity-from-descent}
No conservativity assertion for an arbitrary faithfully flat finite-Tor map is
used here.  For the faithfully spectral etale maps needed in globalization,
conservativity is proved inside the descent argument of
Proposition~\ref{prop:bnla-procoherent-etale-descent} from spectral fppf
universal descent and is then recorded separately in
Corollary~\ref{cor:bnla-procoherent-etale-conservative}.
\end{remark}


\begin{definition}[Anchored pro-coherent spectral modules]
\label{def:bnla-spectral-anchor-category}
Define
\[
    \operatorname{Anch}^{\mathrm{pc}}_{B/R}
    :=
    \left(
        \operatorname{QC}^{\vee}_B
    \right)_{/\mathbb T^{\mathrm{pc}}_{B/R}[1]}.
\]
Let
\[
    \operatorname{dAPerfAnch}_{B/R}
    \subset
    \operatorname{Anch}^{\mathrm{pc}}_{B/R}
\]
be the full subcategory whose source is the continuous dual of a connective
almost perfect $B$-module.
\end{definition}

No anchor in this category is attached to a formal groupoid by hand.  The
category is the target in which continuous duality transports relative
derivation theory.
\section{The spectral cotangent--square-zero comonad}
\label{sec:bnla-spectral-cot-sqz}

The relative spectral partition-Lie construction begins with the cotangent
fibre of an augmented spectral algebra.  The point of this section is to prove,
inside connective commutative ring spectra, the finiteness, completion, and
comonadicity results needed later.  No notion of an admissible augmentation is
introduced.

\subsection{Relative spectral derivations}

\begin{definition}[Relative derivation category]
\label{def:bnla-spectral-derivation-category}
Define
\[
    \operatorname{Der}^{\mathrm{sp}}_{B/R}
    :=
    \left(\operatorname{Mod}_B\right)_{L_{B/R}[-1]/}.
\]
Thus an object is a pair $(M,\alpha)$ with
\[
    \alpha:L_{B/R}[-1]\longrightarrow M.
\]
Write
\[
    \operatorname{Der}^{\mathrm{aperf}}_{B/R,\geq0}
\]
for the full subcategory in which $M$ is connective and almost perfect.
\end{definition}

The shifted morphism
\[
    L_{B/R}\longrightarrow M[1]
\]
classifies, by Chapter~2, a structured square-zero extension
\[
    B_\alpha\longrightarrow B
\]
over $R$.

\begin{definition}[Cotangent and square-zero functors]
\label{def:bnla-cot-sqz-functors}
For an $R$-$E_\infty$-algebra map $A\to B$, define
\[
    \operatorname{cot}_{B/R}(A)
    :=
    \left(
        L_{B/R}[-1]\longrightarrow L_{B/A}[-1]
    \right),
\]
where the arrow is induced by cotangent transitivity.  For an arbitrary
relative derivation $(M,\alpha)\in\operatorname{Der}^{\mathrm{sp}}_{B/R}$,
let
\[
    B\xrightarrow{d_\alpha}B\oplus M[1]
\]
be the section classified by the shifted derivation
$L_{B/R}\to M[1]$, let $d_0:B\to B\oplus M[1]$ be the zero section, and
define in unrestricted commutative ring spectra
\[
    B_\alpha
    :=
    B\times_{B\oplus M[1]}B,
\]
where the two maps are $d_\alpha$ and $d_0$.  Define
\[
    \operatorname{sqz}_{B/R}(M,\alpha)
    :=
    \left(B_\alpha\longrightarrow B\right).
\]
Thus $\operatorname{sqz}_{B/R}$ is defined on the whole relative derivation
category; connectivity is imposed only when restricting the adjunction below.
\end{definition}

\begin{proposition}[Spectral cotangent--square-zero adjunction]
\label{prop:bnla-cot-sqz-adjunction}
There is a canonical adjunction
\[
    \operatorname{cot}_{B/R}
    \dashv
    \operatorname{sqz}_{B/R}.
\]
More explicitly,
\[
\begin{aligned}
&\operatorname{Map}_{\operatorname{Der}^{\mathrm{sp}}_{B/R}}
\left(
    \operatorname{cot}_{B/R}(A),(M,\alpha)
\right)
\\
&\qquad\simeq
\operatorname{Map}_{\operatorname{CAlg}(\operatorname{Sp})_{R//B}}
\left(A,B_\alpha\right).
\end{aligned}
\]
\end{proposition}

\begin{proof}
A point of the right-hand mapping space is a lift of $A\to B$ through the
structured square-zero extension classified by
$L_{B/R}\to M[1]$.  The relative cotangent lifting theorem identifies such
lifts with maps
\[
    L_{B/A}\longrightarrow M[1]
\]
whose restriction along $L_{B/R}\to L_{B/A}$ is the specified classifying
morphism.  After shifting by $[-1]$, this is precisely a morphism from
$\operatorname{cot}_{B/R}(A)$ to $(M,\alpha)$ in the undercategory of
$L_{B/R}[-1]$.
\end{proof}

\begin{lemma}[Connective restriction of structured square-zero formation]
\label{lem:bnla-sqz-connective-restriction}
Assume that $B$ is connective and that
$(M,\alpha)\in\operatorname{Der}^{\mathrm{sp}}_{B/R}$ has connective
underlying $B$-module $M$.  Then $B_\alpha$ is connective, the augmentation
\[
    B_\alpha\longrightarrow B
\]
is surjective on $\pi_0$, and its homotopy fibre is canonically $M$ as a
$B_\alpha$-module through the augmentation.  In particular, the unrestricted
right adjoint $\operatorname{sqz}_{B/R}$ restricts to connective commutative
ring spectra on connective derivations.
\end{lemma}

\begin{proof}
Chapter~2's structured square-zero classification identifies the pullback
\[
    B_\alpha
    =B\times_{B\oplus M[1]}B
\]
with the structured extension classified by
$L_{B/R}\to M[1]$.  Its underlying fibre sequence is
\[
    M\longrightarrow B_\alpha\longrightarrow B.
\]
Since both $M$ and $B$ are connective, the long exact homotopy sequence makes
$B_\alpha$ connective.  The same sequence gives surjectivity on $\pi_0$, and
the displayed fibre identification is part of the structured square-zero
construction.  No connectivity assertion is used in the unrestricted
adjunction itself.
\end{proof}

\subsection{Derived adic filtration}

The completion used below is constructed before any finiteness condition is
imposed.  We work in the unbounded category while forming filtered limits and
recover connectivity from the associated graded.

\begin{definition}[Filtered spectral algebras]
\label{def:bnla-filtered-spectral-algebras}
Let
\[
    \operatorname{Fil}(\operatorname{Sp})
    :=
    \operatorname{Fun}(\mathbb N^{\mathrm{op}},\operatorname{Sp})
\]
with Day convolution.  Write
\[
    \operatorname{Gr}(\operatorname{Sp})
    :=
    \operatorname{Fun}(\mathbb N,\operatorname{Sp})
\]
and let
\[
    \operatorname{Gr}:
    \operatorname{Fil}(\operatorname{Sp})
    \longrightarrow
    \operatorname{Gr}(\operatorname{Sp})
\]
be associated graded.
\end{definition}

For a filtered commutative $R$-algebra $F^\bullet C$, weight zero gives the
augmented arrow
\[
    F^0C\longrightarrow\operatorname{Gr}^0C.
\]
The resulting functor
\[
    \operatorname{aug}:
    \operatorname{CAlg}(\operatorname{Fil}(\operatorname{Sp}))_R
    \longrightarrow
    \operatorname{CAlg}(\operatorname{Sp})^{\Delta^1}_R
\]
is accessible and limit-preserving.  Hence it admits a left adjoint.

\begin{definition}[Derived adic filtration and completion]
\label{def:bnla-derived-adic-filtration}
For an augmented arrow $A\to B$ of commutative $R$-ring spectra define
\[
    \operatorname{adic}(A\to B)
\]
to be its image under the left adjoint
$\operatorname{adic}\dashv\operatorname{aug}$.  Define
\[
    \widehat{\operatorname{adic}}(A\to B)
    :=
    \lim_n
    \operatorname{adic}(A\to B)/F^{n+1},
\]
and put
\[
    \widehat A_B
    :=F^0\widehat{\operatorname{adic}}(A\to B).
\]
The unit followed by completion gives a canonical morphism
\[
    A\longrightarrow\widehat A_B.
\]
\end{definition}

No ideal, topology, filtration, or generators are part of an augmented arrow.
The filtration above is functorially produced by an adjunction.

\begin{definition}[Filtered and graded cotangent fibres]
\label{def:bnla-filtered-graded-cotangent}
Let $C$ be a filtered or nonnegatively graded commutative $R$-algebra and
write
\[
    C_0:=\operatorname{Gr}^0C.
\]
Regard $C_0$ as a filtered or graded algebra concentrated in weight zero and
define
\[
    \operatorname{cot}(C)
    :=
    L_{C_0/C}[-1].
\]
For an augmented arrow $A\to B$, put
\[
    \operatorname{cot}(A\to B)
    :=
    L_{B/A}[-1].
\]
\end{definition}

\begin{lemma}[Cotangent compatibility with grading, augmentation, and adic filtration]
\label{lem:bnla-adic-cotangent-compatibility}
There are canonical natural equivalences
\[
    \operatorname{Gr}\operatorname{cot}(C)
    \simeq
    \operatorname{cot}(\operatorname{Gr}C),
\]
\[
    \operatorname{aug}\operatorname{cot}(C)
    \simeq
    \operatorname{cot}(\operatorname{aug}C),
\]
and, for every augmented arrow $A\to B$,
\[
    F^{\leq1}\operatorname{cot}(A\to B)
    \simeq
    \operatorname{cot}\operatorname{adic}(A\to B).
\]
Here $F^{\leq1}M$ denotes the filtered $B$-module obtained by placing $M$ in
filtration weight one, so that its only nonzero associated-graded piece is
$M$ in weight one.
\end{lemma}

\begin{proof}
Associated graded and augmentation are symmetric monoidal functors which
commute with derived symmetric algebras.  The universal derivation
characterization of the cotangent complex therefore gives the first two
equivalences.

For the third, test the filtered cotangent fibre against a filtered
$B$-module concentrated in positive filtration.  A map from
$\operatorname{cot}\operatorname{adic}(A\to B)$ to such a module classifies a
filtered structured square-zero extension of the weight-zero augmentation.
Applying $\operatorname{aug}$ gives the corresponding structured square-zero
extension of the arrow $A\to B$.  Since
$\operatorname{adic}\dashv\operatorname{aug}$ and augmentation preserves
structured square-zero pullbacks, the resulting derivation space is
represented by $L_{B/A}[-1]$ placed in filtration weight one.  Yoneda gives
the displayed equivalence.
\end{proof}

\begin{lemma}[Cotangent detection for positive graded spectral algebras]
\label{lem:bnla-graded-cotangent-detection}
Let
\[
    f:C\longrightarrow D
\]
be a morphism of nonnegatively graded commutative ring spectra.
Then $f$ is an equivalence if and only if
\[
    C_0\xrightarrow{\sim}D_0
\]
and
\[
    \operatorname{cot}(C)
    \xrightarrow{\sim}
    \operatorname{cot}(D).
\]
\end{lemma}

\begin{proof}
Only the converse requires proof.  Identify $C_0\simeq D_0$ and suppose that
$f$ is not an equivalence.  Let $n\geq1$ be the least weight for which
\[
    \operatorname{cofib}(C_n\to D_n)\not\simeq0.
\]
Let $C_{\geq n}$ and $D_{\geq n}$ denote the weight-$\geq n$ parts, regarded
as graded $C$-modules, and form
\[
    D'
    :=
    C
    \otimes_{\operatorname{Sym}_C(C_{\geq n})}
    \operatorname{Sym}_C(D_{\geq n})
    \longrightarrow
    D.
\]
The terms of symmetric degree at least two in
$\operatorname{Sym}_C(D_{\geq n})$ have weight at least $2n$, hence at least
$n+1$.  Since $C\to D$ is an equivalence in weights below $n$, the map
$D'\to D$ is an equivalence in weights at most $n$.  Resolving $D$ by free
$D'$-algebras therefore shows that
\[
    (L_{D/D'})_k\simeq0
    \qquad (k\leq n).
\]
Cotangent transitivity then gives
\[
    (L_{D/C})_n
    \simeq
    (L_{D'/C})_n
    \simeq
    \operatorname{cofib}(C_n\to D_n),
\]
which is nonzero.

Now apply cotangent transitivity to
\[
    C\longrightarrow D\longrightarrow D_0.
\]
After shifting by $[-1]$ it identifies
\[
    D_0\otimes_D L_{D/C}[-1]
\]
with the fibre of
\[
    \operatorname{cot}(C)\longrightarrow\operatorname{cot}(D).
\]
Because $C\to D$ is an equivalence in weights below $n$, the same free-algebra
resolution argument shows that $L_{D/C}$ vanishes in weights below $n$.
Compute $D_0\otimes_D L_{D/C}$ by the two-sided bar construction.  In total
weight $n$, every bar summand containing a positive-weight factor from $D$
would have to be paired with a component of $L_{D/C}$ of weight strictly
below $n$, and is therefore zero.  Thus the weight-$n$ term of the bar
construction is canonically the weight-$n$ term of $L_{D/C}$.  Consequently
\[
    \left(D_0\otimes_D L_{D/C}\right)_n
    \simeq
    (L_{D/C})_n
    \not\simeq0.
\]
The fibre of the cotangent comparison is therefore nonzero, contradicting
the assumed equivalence.  Hence $f$ is an equivalence.
\end{proof}

\begin{theorem}[Associated graded of the spectral adic filtration]
\label{thm:bnla-adic-associated-graded}
For every augmented arrow $A\to B$ of commutative ring spectra there is a
natural equivalence of graded commutative $B$-algebras
\[
\boxed{
    \operatorname{Gr}
    \widehat{\operatorname{adic}}(A\to B)
    \simeq
    \operatorname{Sym}^{\mathrm{gr}}_B
    \left(L_{B/A}[-1](1)\right).}
\]
Completion does not change the associated graded.
\end{theorem}

\begin{proof}
Put
\[
    M:=L_{B/A}[-1],
    \qquad
    G:=\operatorname{Gr}\operatorname{adic}(A\to B).
\]
By Lemma~\ref{lem:bnla-adic-cotangent-compatibility},
\[
    \operatorname{cot}(G)
    \simeq
    M(1),
\]
so the graded cotangent fibre is concentrated in weight one.

For every nonnegatively graded commutative algebra $C$, multiplication
induces a canonical map
\[
    \operatorname{Sym}^{\mathrm{gr}}_{C_0}
    \left(C_1(1)\right)
    \longrightarrow C,
\]
and the induced map
\[
    C_1\longrightarrow\operatorname{cot}(C)_1
\]
is an equivalence because there are no positive decomposable terms in weight
one.  Applying this to $G$ identifies $G_1\simeq M$ and produces
\[
    \operatorname{Sym}^{\mathrm{gr}}_B(M(1))
    \longrightarrow G.
\]
This map is an equivalence in weight zero and induces an equivalence on
cotangent fibres.  Lemma~\ref{lem:bnla-graded-cotangent-detection} therefore
makes it an equivalence.

Finally, for fixed weight $n$, the associated-graded piece is already
detected on every quotient by $F^m$ with $m>n$.  The inverse system is
therefore eventually constant in each weight, so passing to the completed
inverse limit leaves the associated graded unchanged.
\end{proof}

\begin{proposition}[Zeroth homotopy of the spectral adic filtration]
\label{prop:bnla-adic-pi-zero}
Let $A\to B$ be a connective augmented arrow inducing a surjection on
$\pi_0$, and put
\[
    I:=\ker\bigl(\pi_0(A)\to\pi_0(B)\bigr).
\]
Then the filtered ordinary commutative ring
$\pi_0\operatorname{adic}(A\to B)$ has terms
\[
    F^n\pi_0\operatorname{adic}(A\to B)
    \simeq
    \operatorname{Sym}^n_{\pi_0(A)}(I)
    \qquad(n\geq1),
\]
with structural maps induced by multiplication and the inclusion
$I\hookrightarrow\pi_0(A)$.  In particular, the image of
\[
    \pi_0F^n\operatorname{adic}(A\to B)
    \longrightarrow
    \pi_0(A)
\]
is $I^n$, and every finite filtration quotient satisfies
\[
\boxed{
    \pi_0\left(
        F^0\operatorname{adic}(A\to B)/
        F^n\operatorname{adic}(A\to B)
    \right)
    \simeq
    \pi_0(A)/I^n.}
\]
\end{proposition}

\begin{proof}
To determine $\pi_0$ it is enough to test the left-adjoint universal property
against discrete filtered commutative rings.  A map of augmented arrows
\[
    (A\to B)
    \longrightarrow
    \left(F^0C\to\operatorname{Gr}^0C\right)
\]
into a discrete filtered ring is determined on components by a ring map
\[
    \pi_0(A)\longrightarrow F^0C
\]
together with a map
\[
    I\longrightarrow F^1C
\]
compatible with the inclusion into $F^0C$.  The latter extends uniquely,
multiplicatively, to maps
\[
    \operatorname{Sym}^n_{\pi_0(A)}(I)
    \longrightarrow
    F^nC.
\]
Thus the displayed symmetric-power filtration represents the same
left-adjoint functor as
$\pi_0\operatorname{adic}(A\to B)$.

The composite
\[
    \operatorname{Sym}^n_{\pi_0(A)}(I)
    \longrightarrow
    \pi_0(A)
\]
is multiplication, so its image is $I^n$.  Since the adic filtration of a
connective arrow is connective, the cofiber sequence
\[
    F^n\operatorname{adic}(A\to B)
    \longrightarrow
    F^0\operatorname{adic}(A\to B)
    \longrightarrow
    F^0\operatorname{adic}(A\to B)/
    F^n\operatorname{adic}(A\to B)
\]
gives the displayed finite-stage formula on $\pi_0$.  No interchange of
$\pi_0$ with the completed inverse limit is used here.  The completed
$\pi_0$-calculation in the finitely generated sector is constructed in
Lemma~\ref{lem:bnla-filtered-derived-completion}.
\end{proof}

\begin{lemma}[Complete filtered detection]
\label{lem:bnla-complete-filtered-detection}
Let $F^\bullet M\to F^\bullet N$ be a morphism of complete descending
filtered objects in a stable $\infty$-category.  If
\[
    \operatorname{Gr}M\xrightarrow{\sim}\operatorname{Gr}N,
\]
then
\[
    F^0M\xrightarrow{\sim}F^0N.
\]
\end{lemma}

\begin{proof}
Inductively on $n$, the cofiber sequences comparing successive finite
filtration quotients show that
\[
    F^0M/F^{n}M\longrightarrow F^0N/F^{n}N
\]
is an equivalence for every $n$.  Taking inverse limits and using completeness
gives the result.
\end{proof}

\begin{lemma}[Connectivity of a complete positive filtration]
\label{lem:bnla-complete-filtered-connectivity}
Let $F^\bullet C$ be a complete descending filtered spectrum such that every
associated-graded piece is connective.  Then $F^0C$ is connective.
\end{lemma}

\begin{proof}
Put $C_n=F^0C/F^{n+1}C$.  Induction on $n$ shows that every $C_n$ is
connective and that
$\pi_0(C_{n+1})\to\pi_0(C_n)$ is surjective.  Completeness gives
$F^0C\simeq\lim_nC_n$.  The Milnor exact sequence then kills all negative
homotopy groups: for degree $-1$ the relevant $\lim^1$ vanishes by the
surjectivity just established, and below degree $-1$ both outer terms vanish.
\end{proof}

\subsection{Spectral almost-finite surjections}

The next results supply the finiteness theorem needed for the spectral
version of the relative cotangent--square-zero argument.

\begin{lemma}[Relative spectral Hurewicz comparison]
\label{lem:bnla-relative-spectral-hurewicz}
Let $A\to C$ be a morphism of connective $E_\infty$-rings inducing an
isomorphism on $\pi_0$, and suppose
\[
    K:=\operatorname{cofib}(A\to C)
\]
is $n$-connective for some $n\geq1$.  Then the canonical relative-cofiber
morphism
\[
    C\otimes_A K
    \longrightarrow
    L_{C/A}
\]
induces an isomorphism
\[
    \boxed{
    \pi_n(K)\xrightarrow{\sim}\pi_nL_{C/A}.}
\]
The comparison is natural in commutative diagrams of such morphisms.
\end{lemma}

\begin{proof}
The relative cotangent connectivity theorem used in Chapter~2 says that, for
an $n$-connective relative cofiber, the canonical morphism
\[
    C\otimes_AK\longrightarrow L_{C/A}
\]
has $(n+1)$-connective fibre.  It therefore induces an isomorphism on
$\pi_n$.

Since $K$ is $n$-connective and
$\pi_0(A)\xrightarrow{\sim}\pi_0(C)$, the connective tensor-product spectral
sequence gives
\[
    \pi_n(C\otimes_AK)
    \simeq
    \pi_0(C)\otimes_{\pi_0(A)}\pi_n(K)
    \simeq
    \pi_n(K).
\]
Combining the two identifications gives the displayed isomorphism.
Naturality is the naturality of the relative-cofiber-to-cotangent morphism,
already used in the derived conormal comparison of Chapter~2.
\end{proof}

\begin{lemma}[Positive spectral relation cells]
\label{lem:bnla-positive-relation-cell}
Let $C$ be a connective $E_\infty$-ring, let $n\geq0$, let
\[
    P\simeq C[n]^{\oplus r}
\]
be a finite free connective cell module, and let
\[
    \chi:P\longrightarrow C
\]
be a $C$-linear morphism.  Write
\[
    S:=\operatorname{Sym}_C(P)
\]
and let
\[
    \epsilon_\chi,\epsilon_0:S\rightrightarrows C
\]
be the commutative $C$-algebra maps classified respectively by $\chi$ and the
zero map.  Define the corresponding positive relation-cell attachment by
\[
\boxed{
    C//\chi
    :=
    C_{\epsilon_\chi}\otimes_S C_{\epsilon_0}.}
\]
Then:
\begin{enumerate}
    \item $C//\chi$ is a finite-cell commutative $C$-algebra;
    \item the two-sided bar construction carries a natural exhaustive
    multiplicative Koszul filtration whose associated graded is
    \[
    \boxed{
        \operatorname{Gr}^k(C//\chi)
        \simeq
        \operatorname{Sym}^k_C(P[1])
        \simeq
        \left(P[1]^{\otimes_C k}\right)_{h\Sigma_k},
        \qquad k\geq0;}
    \]
    in particular the weight-zero term is $C$ and the first positive term is
    $P[1]$;
    \item $C//\chi$ is almost perfect as a $C$-module.
\end{enumerate}

Now let $C\to D$ be a morphism of connective $E_\infty$-rings inducing an
isomorphism on $\pi_0$, and assume
\[
    K:=\operatorname{cofib}(C\to D)
\]
is $(n+1)$-connective.  Suppose that a morphism
\[
    u:P\longrightarrow K[-1]
\]
is surjective on $\pi_n$.  Let
\[
    \chi:
    P\xrightarrow{u}K[-1]\longrightarrow C
\]
be the composite with the boundary of the cofiber sequence.  The canonical
nullhomotopy of $P\to C\to D$ induces a commutative $C$-algebra map
\[
    C//\chi\longrightarrow D,
\]
and its cofiber is $(n+2)$-connective.
\end{lemma}

\begin{proof}
The pushout defining $C//\chi$ is the attachment of finitely many
commutative-algebra relation cells in homological degree $n+1$, so it is a
finite-cell commutative $C$-algebra.

Compute the pushout by the two-sided bar construction
\[
    \left|
    \operatorname{Bar}_\bullet
    (C_{\epsilon_\chi},S,C_{\epsilon_0})
    \right|.
\]
Give the normalized bar object the increasing Koszul filtration by the number
of unresolved copies of the free generator $P$.  The endpoint differential
coming from $\epsilon_\chi$ contracts one generator against $\chi$ and
therefore lowers this filtration.  Hence it disappears on associated graded.
The associated graded is consequently the bar object with both endpoint
actions equal to the zero augmentation, namely the zero--zero derived
self-intersection
\[
    C_{\epsilon_0}\otimes_{\operatorname{Sym}_C(P)}C_{\epsilon_0}.
\]
The pushout square of $C$-modules
\[
\begin{tikzcd}
P \arrow[r] \arrow[d] & 0 \arrow[d]\\
0 \arrow[r] & P[1]
\end{tikzcd}
\]
and the free-commutative-algebra pushout formula identify this zero--zero
self-intersection with
\[
    \operatorname{Sym}_C(P[1]).
\]
The Koszul filtration is multiplicative, exhaustive, and has symmetric degree
as its filtration weight.  Thus
\[
    \operatorname{Gr}^k(C//\chi)
    \simeq
    \operatorname{Sym}^k_C(P[1])
    \simeq
    \left(P[1]^{\otimes_C k}\right)_{h\Sigma_k}.
\]
This is the spectral specialization of the standard filtered cell attachment:
no passage to animated rings is involved.

Because $P[1]$ is $(n+1)$-connective, the weight-$k$ term is
$k(n+1)$-connective.  For each fixed $k$, the homotopy-orbit object is almost
perfect.  Indeed, the skeletal filtration of the free contractible
$\Sigma_k$-space gives
\[
    E\Sigma_k^{(m)}{}_+
    \otimes_{\Sigma_k}
    P[1]^{\otimes_C k}
    \longrightarrow
    \left(P[1]^{\otimes_C k}\right)_{h\Sigma_k},
\]
where the source is perfect and, as $m\to\infty$, the cofiber becomes
arbitrarily connective.  Hence every associated-graded term is almost
perfect.  Since the connectivity tends to $+\infty$ with $k$, every fixed
Postnikov range of the bar filtration sees only finitely many weights.
Finite extensions of the corresponding perfect approximations therefore give
perfect $C$-module approximations to $C//\chi$ in every Postnikov range.  Thus
$C//\chi$ is almost perfect over $C$.

For the final assertion, the cofiber sequence
\[
    K[-1]\longrightarrow C\longrightarrow D
\]
provides a specified nullhomotopy of the composite
$P\xrightarrow{\chi}C\to D$.  By the pushout universal property this is
exactly the datum of a map $C//\chi\to D$.  Put
\[
    C':=C//\chi.
\]
The preceding filtration shows that
\[
    \operatorname{cofib}(C\to C')
\]
is $(n+1)$-connective, with first nonzero associated-graded term $P[1]$;
all nonlinear weights are at least $2(n+1)$-connective.  Functoriality of the
Koszul filtration for the map $C'\to D$ identifies its weight-one map with
the suspension
\[
    P[1]\xrightarrow{u[1]}K.
\]
Consequently the induced map
\[
    \pi_{n+1}\operatorname{cofib}(C\to C')
    \longrightarrow
    \pi_{n+1}K
\]
is exactly the map on $\pi_n$ induced by $u$, and is surjective.  Octahedrality for
\[
    C\longrightarrow C'\longrightarrow D
\]
gives a cofiber sequence
\[
    \operatorname{cofib}(C\to C')
    \longrightarrow
    K
    \longrightarrow
    \operatorname{cofib}(C'\to D).
\]
The first two terms are $(n+1)$-connective and the map on their first possible
homotopy groups is surjective.  Hence the last term is $(n+2)$-connective.
\end{proof}

\begin{lemma}[Spectral cell criterion for almost finite presentation]
\label{lem:bnla-spectral-cell-aft}
Let $A\to B$ be a morphism of connective $E_\infty$-rings inducing a
surjection on $\pi_0$.  The following conditions are equivalent:
\begin{enumerate}
    \item $B$ is almost finitely presented as a connective commutative
    $A$-algebra;
    \item $L_{B/A}$ is an almost perfect $B$-module and $\pi_0(B)$ is finitely
    presented as a $\pi_0(A)$-algebra;
    \item $B$ is almost perfect as an $A$-module.
\end{enumerate}
\end{lemma}

\begin{proof}
The equivalence of (1) and (2) is the Postnikov form of the spectral
cotangent finiteness criterion: almost compactness on every truncated
connective algebra category is detected by finite presentation on $\pi_0$ and
almost-perfectness of the relative cotangent complex.  We compare these
conditions with (3) using the positive relation cells of
Lemma~\ref{lem:bnla-positive-relation-cell}.

Assume first (2).  Since
\[
    \pi_0(A)\twoheadrightarrow\pi_0(B)
\]
is a finitely presented quotient, its kernel is finitely generated.  Choose
finitely many generators, lift them to
$\pi_0\operatorname{fib}(A\to B)$, and choose representatives.  Their
boundary maps define
\[
    \chi_0:A^{\oplus r_0}\longrightarrow A.
\]
Set
\[
    B_0:=A//\chi_0.
\]
The chosen lifts supply the required nullhomotopies in $B$, hence a map
$B_0\to B$.  Since the defining pushout is a pushout of connective
commutative ring spectra, applying $\pi_0$ gives the ordinary quotient by the
chosen generators.  Therefore
\[
    \pi_0(B_0)
    \simeq
    \pi_0(A)/(x_1,\ldots,x_{r_0})
    \xrightarrow{\sim}
    \pi_0(B).
\]
By Lemma~\ref{lem:bnla-positive-relation-cell}, $B_0$ is finite-cell as an
$A$-algebra and almost perfect as an $A$-module.

Suppose inductively that a finite-cell $A$-algebra $B_n$ has been constructed
together with a map $B_n\to B$ whose cofiber
\[
    K_n:=\operatorname{cofib}(B_n\to B)
\]
is $(n+1)$-connective.  Because $B_n$ is finite-cell over $A$,
$L_{B_n/A}$ is perfect.  Cotangent transitivity
\[
    B\otimes_{B_n}L_{B_n/A}
    \longrightarrow
    L_{B/A}
    \longrightarrow
    L_{B/B_n}
\]
therefore makes $L_{B/B_n}$ almost perfect.  The relative spectral Hurewicz
comparison gives
\[
    \pi_{n+1}(K_n)
    \xrightarrow{\sim}
    \pi_{n+1}L_{B/B_n},
\]
so this first unresolved group is finitely generated over
$\pi_0(B)\simeq\pi_0(B_n)$.  Choose finitely many generators and represent
them by a morphism
\[
    u_n:B_n[n]^{\oplus r_n}\longrightarrow K_n[-1]
\]
surjective on $\pi_n$.  Lemma~\ref{lem:bnla-positive-relation-cell} attaches
the corresponding relation cell
\[
    B_{n+1}:=B_n//\chi_n
\]
and produces a map $B_{n+1}\to B$ whose cofiber is $(n+2)$-connective.
Moreover $B_{n+1}$ remains finite-cell over $A$ and, by transitivity of
almost-perfectness, almost perfect as an $A$-module.

Fix a Postnikov degree $d$ and choose $n>d$.  Since $B_n$ is almost perfect
over $A$, choose a perfect $A$-module $P_d$ and a map
\[
    P_d\longrightarrow B_n
\]
whose cofiber is $(d+1)$-connective.  The composite
\[
    P_d\longrightarrow B_n\longrightarrow B
\]
has the same property because $B_n\to B$ is an equivalence through degree
$d$.  Thus $B$ admits perfect $A$-module approximations through every
Postnikov degree, proving (3).

Conversely assume (3).  Then $\pi_0(B)$ is finitely presented as a
$\pi_0(A)$-module.  Since it is a quotient of $\pi_0(A)$, the kernel of
\[
    \pi_0(A)\twoheadrightarrow\pi_0(B)
\]
is finitely generated.  Construct the same finite relation-cell algebra
$B_0\to B$ as above.  It induces an isomorphism on $\pi_0$ and is almost
perfect as an $A$-module.

Suppose inductively that $B_n\to B$ has $(n+1)$-connective cofiber $K_n$, that
$B_n$ is finite-cell as an $A$-algebra, and that $B_n$ is almost perfect as an
$A$-module.  Since $B$ and $B_n$ are almost perfect over $A$, the cofiber
$K_n$ is almost perfect over $A$.  Hence its first possible nonzero homotopy
group
\[
    \pi_{n+1}(K_n)
\]
is finitely generated over $\pi_0(A)$, and therefore over the quotient
$\pi_0(B_n)\simeq\pi_0(B)$.  Choose finitely many generators and a morphism
\[
    u_n:B_n[n]^{\oplus r_n}\longrightarrow K_n[-1]
\]
surjective on $\pi_n$.  The positive relation-cell lemma produces
$B_{n+1}\to B$ with $(n+2)$-connective cofiber.  It also makes
$B_{n+1}$ finite-cell over $B_n$ and almost perfect over $B_n$; hence it is
finite-cell over $A$ and almost perfect over $A$.

For every fixed $d$, choose $n\geq d$.  Then
\[
    \tau_{\leq d}B_n
    \xrightarrow{\sim}
    \tau_{\leq d}B.
\]
A map from a connective $A$-algebra to a $d$-truncated connective
$A$-algebra depends only on its $d$-truncation.  Therefore mapping out of $B$
into any $d$-truncated target is computed by the finite-cell, hence compact,
$A$-algebra $B_n$.  It follows that mapping out of $B$ on $d$-truncated
connective $A$-algebras preserves filtered colimits.  Since $d$ is arbitrary,
$B$ is almost finitely presented over $A$.  This proves (3)$\Rightarrow$(1)
and completes the equivalence.
\end{proof}

\begin{lemma}[Finite presentation of the ordinary augmentation ideal]
\label{lem:bnla-augmentation-ideal-finite-presentation}
Let $A\to B$ be a morphism of connective $E_\infty$-rings which is
surjective on $\pi_0$.  Assume that $B$ is coherent and that $B$ is almost
perfect as an $A$-module.  Put
\[
    I:=\ker\bigl(\pi_0(A)\to\pi_0(B)\bigr).
\]
Then $I$ is finitely presented as a $\pi_0(A)$-module.
\end{lemma}

\begin{proof}
Put
\[
    K:=\operatorname{fib}(A\to B).
\]
Because $A$ and $B$ are connective and $\pi_0(A)\twoheadrightarrow\pi_0(B)$,
the long exact homotopy sequence shows that $K$ is connective.  The
subcategory of almost-perfect $A$-modules is stable under finite limits and
colimits.  Since $A$ is perfect over itself and $B$ is almost perfect over
$A$, the fibre $K$ is almost perfect over $A$.

For a connective almost-perfect $A$-module, degree-zero homotopy is finitely
presented over $\pi_0(A)$.  Indeed, almost-perfectness implies that mapping
out of $K$ commutes with filtered colimits on every fixed
Postnikov-truncated module category.  Restricting to $0$-truncated
$A$-modules and using
\[
    \operatorname{Mod}_{A,[0,0]}
    \simeq
    \operatorname{Mod}_{\pi_0(A)}
\]
shows that
\[
    \operatorname{Hom}_{\pi_0(A)}\bigl(\pi_0K,-\bigr)
\]
preserves filtered colimits.  Hence $\pi_0K$ is finitely presented.

The long exact homotopy sequence gives an exact sequence
\[
    \pi_1(B)
    \longrightarrow
    \pi_0(K)
    \longrightarrow
    \pi_0(A)
    \longrightarrow
    \pi_0(B)
    \longrightarrow0.
\]
Consequently
\[
    I
    \simeq
    \pi_0(K)\big/
    \operatorname{im}\!\left(\pi_1(B)\to\pi_0(K)\right).
\]
Since $B$ is coherent, $\pi_1(B)$ is finitely presented, hence finitely
generated, over $\pi_0(B)$ and therefore also finitely generated over
$\pi_0(A)$.  Its image in $\pi_0(K)$ is thus finitely generated.

Finally, the quotient of a finitely presented module by a finitely generated
submodule is finitely presented.  Choose a finite presentation of $\pi_0K$
and lift finitely many generators of the displayed image to the finite free
module of generators; adjoining those lifts to the finite set of relations
gives a finite presentation of the quotient.  Therefore $I$ is finitely
presented.
\end{proof}

\begin{lemma}[Positive graded symmetric algebras are almost finitely augmented]
\label{lem:bnla-positive-graded-aft}
Let $B$ be connective and let $M$ be a connective almost perfect $B$-module.
Then the augmentation
\[
    \operatorname{Sym}^{\mathrm{gr}}_B(M(1))\longrightarrow B
\]
is almost finitely presented in connective nonnegatively graded commutative
$B$-algebras.
\end{lemma}

\begin{proof}
Put
\[
    C:=\operatorname{Sym}^{\mathrm{gr}}_B(M(1)).
\]
Let $D$ be an $m$-truncated connective nonnegatively graded commutative
$C$-algebra.  By the free graded-algebra adjunction, a morphism of augmented
graded $B$-algebras
\[
    B\longrightarrow D
\]
over $C$ is equivalently a nullhomotopy, compatible with the augmentation, of
the image of the universal weight-one map
\[
    M\longrightarrow D_1.
\]
Consequently its mapping space is a finite-limit expression built from
$\operatorname{Map}_B(M,D_1)$.  Evaluation in weight one preserves filtered
colimits of graded algebras.  Since $M$ is connective and almost perfect,
mapping out of $M$ into $m$-truncated connective $B$-modules preserves
filtered colimits.  Filtered colimits of spaces commute with finite limits.
Thus mapping out of the augmentation $C\to B$ into $m$-truncated connective
graded algebras preserves filtered colimits.  Since $m$ was arbitrary, the
augmentation is almost finitely presented.
\end{proof}

\begin{lemma}[Almost-perfectness from a complete positive filtration]
\label{lem:bnla-complete-filtered-aperf}
Let $A$ be a connective complete filtered $E_\infty$-ring and let $N$ be a
connective complete filtered $A$-module.  Suppose that
$\operatorname{Gr}N$ is almost perfect over $\operatorname{Gr}A$ and that the
connectivity of the positive filtration pieces of $\operatorname{Gr}N$ tends
to infinity with the weight.  Then $F^0N$ is almost perfect as an
$F^0A$-module.
\end{lemma}

\begin{proof}
Fix a Postnikov degree $d$.  By almost-perfectness of $\operatorname{Gr}N$,
there is a perfect graded module $Q_d$ mapping to $\operatorname{Gr}N$ with
$d$-connective cofiber.  Only finitely many weights affect degrees below $d$
by the increasing-connectivity hypothesis.  Lift the finite cell construction
of those weights to a perfect filtered $A$-module $P_d\to N$.  Its cofiber
has $d$-connective associated graded.  Every finite filtration quotient of the
cofiber is therefore $d$-connective; completeness and the Milnor argument of
Lemma~\ref{lem:bnla-complete-filtered-connectivity} imply that its filtration
zero term is $d$-connective.  Thus $F^0N$ admits perfect approximations through
every Postnikov degree and is almost perfect.
\end{proof}

\begin{lemma}[Completeness of almost-perfect filtered modules]
\label{lem:bnla-aperf-filtered-complete}
Let $F^\bullet A$ be a complete connective filtered $E_\infty$-ring.  Every
connective almost-perfect filtered $A$-module is complete.
\end{lemma}

\begin{proof}
A perfect filtered $A$-module is complete: it is obtained from finite sums of
shifts of $A$ by finite limits, finite colimits, and retracts, all of which
preserve completeness.

Let $M$ be connective and almost perfect.  For every $d\geq0$, choose a
perfect filtered module $P_d$ and a morphism
\[
    P_d\longrightarrow M
\]
with $d$-connective cofiber $C_d$.  Since $P_d$ is complete,
\[
    \lim_nF^nM
    \simeq
    \lim_nF^nC_d.
\]
The filtration t-structure is pointwise, so every $F^nC_d$ is
$d$-connective.  The Milnor exact sequence shows that the inverse limit is
$(d-1)$-connective.  Since $d$ is arbitrary,
\[
    \lim_nF^nM\simeq0.
\]
Thus $M$ is complete.
\end{proof}

\begin{lemma}[Cotangent completeness for completed spectral adic filtrations]
\label{lem:bnla-completed-adic-cotangent}
Let $A\to B$ be a connective augmented arrow inducing a surjection on
$\pi_0$, and suppose
\[
    M:=L_{B/A}[-1]
\]
is connective and almost perfect.  Then
\[
    \operatorname{cot}
    \widehat{\operatorname{adic}}(A\to B)
\]
is a connective almost-perfect complete filtered $B$-module, and the
completion morphism induces an equivalence
\[
\boxed{
    \operatorname{cot}\operatorname{adic}(A\to B)
    \xrightarrow{\sim}
    \operatorname{cot}
    \widehat{\operatorname{adic}}(A\to B).}
\]
\end{lemma}

\begin{proof}
Associated graded commutes with cotangent fibres, so
Theorem~\ref{thm:bnla-adic-associated-graded} gives
\[
\begin{aligned}
    \operatorname{Gr}
    \operatorname{cot}
    \widehat{\operatorname{adic}}(A\to B)
    &\simeq
    \operatorname{cot}
    \operatorname{Sym}^{\mathrm{gr}}_B(M(1))
    \\
    &\simeq
    M(1).
\end{aligned}
\]
Thus the associated graded of the filtered cotangent fibre is connective,
almost perfect, and concentrated in the single positive weight one.

Choose, in every Postnikov range, a finite perfect graded approximation to
$M(1)$.  The finite cells of such an approximation lift along the filtered
free-module adjunction to a perfect filtered module mapping to the filtered
cotangent fibre.  The cofiber has arbitrarily high connective associated
graded; complete filtered connectivity then makes its filtration pieces
arbitrarily connective.  Hence the filtered cotangent fibre is almost
perfect.  Lemma~\ref{lem:bnla-aperf-filtered-complete} makes it complete.

By Lemma~\ref{lem:bnla-adic-cotangent-compatibility},
\[
    \operatorname{cot}\operatorname{adic}(A\to B)
    \simeq
    F^{\leq1}M,
\]
which is already complete.  The natural comparison to the cotangent fibre of
the completed adic algebra is an equivalence on associated graded, since both
associated gradeds are $M(1)$.  Complete filtered detection gives the
displayed equivalence.
\end{proof}

\subsection{Derived completeness and the complete almost-finite sector}

Let $A\to B$ be a connective augmented arrow surjective on $\pi_0$ and put
\[
    I:=\ker\bigl(\pi_0(A)\to\pi_0(B)\bigr).
\]
Derived $I$-completeness has its usual module-theoretic meaning.

\begin{lemma}[Filtered completeness and the augmentation ideal]
\label{lem:bnla-filtered-derived-completion}
Assume $I$ is finitely generated.
\begin{enumerate}
    \item If $F^\bullet M$ is a complete module over a filtered algebra whose
    weight-zero augmentation is $A\to B$, then $F^0M$ is derived
    $I$-complete.
    \item For the constructed adic filtration,
    \[
    \boxed{
        \pi_0(\widehat A_B)
        \simeq
        \pi_0(A)^\wedge_I
        :=
        \lim_n\pi_0(A)/I^n.}
    \]
    \item If $A$ is derived $I$-complete and
    \[
        B\otimes_A\widehat A_B\xrightarrow{\sim}B,
    \]
    then $A\to\widehat A_B$ is an equivalence.
\end{enumerate}
\end{lemma}

\begin{proof}
For (1), write
\[
    F^0M\simeq\lim_n F^0M/F^nM.
\]
Every $x\in I$ has filtration weight at least one.  Multiplicativity therefore
makes
\[
    x^n:
    F^0M/F^nM
    \longrightarrow
    F^0M/F^nM
\]
null.  Thus every finite filtration quotient is derived $I$-complete.
Derived $I$-complete modules are closed under limits, so $F^0M$ is derived
$I$-complete.

For (2), choose generators $x_1,\ldots,x_r$ of $I$ and lifts of them to
$\pi_0(A)$.  The finite-generator model for derived $I$-completion computes
completion by the Koszul tower obtained from powers of these generators.
Applied to a complete positive filtration in which every $x_i$ has filtration
weight at least one, this model identifies the zeroth Postnikov truncation of
the completed filtration with the inverse limit of the zeroth Postnikov
truncations of its finite quotients.  Proposition~\ref{prop:bnla-adic-pi-zero}
identifies those truncations by
\[
    \pi_0\left(
        F^0\operatorname{adic}(A\to B)/
        F^n\operatorname{adic}(A\to B)
    \right)
    \simeq
    \pi_0(A)/I^n.
\]
The resulting tower of discrete rings has surjective transition maps, so its
ordinary inverse limit computes its zeroth derived inverse limit.  Therefore
\[
    \pi_0(\widehat A_B)
    \simeq
    \lim_n\pi_0(A)/I^n
    =
    \pi_0(A)^\wedge_I.
\]

For (3), part (1) makes $\widehat A_B$ derived $I$-complete.  We first verify
the connectivity needed for derived Nakayama.  Since $A$ is connective and
derived $I$-complete, the finite-generator Koszul model gives
\[
    A
    \simeq
    \lim_n A/(x_1^n,\ldots,x_r^n).
\]
The quotients are connective and their maps on $\pi_0$ are surjective.  The
Milnor exact sequence therefore gives
\[
\begin{aligned}
    \pi_0(A)
    &\simeq
    \lim_n
    \pi_0\bigl(A/(x_1^n,\ldots,x_r^n)\bigr)\\
    &\simeq
    \lim_n
    \pi_0(A)/(x_1^n,\ldots,x_r^n).
\end{aligned}
\]
The filtrations by $(x_1^n,\ldots,x_r^n)$ and by $I^n$ are cofinal: one has
\[
    (x_1^n,\ldots,x_r^n)\subseteq I^n,
    \qquad
    I^{r(n-1)+1}\subseteq(x_1^n,\ldots,x_r^n).
\]
Hence
\[
    \pi_0(A)
    \xrightarrow{\sim}
    \pi_0(A)^\wedge_I.
\]
Combining this with (2), the canonical map
\[
    \pi_0(A)\longrightarrow\pi_0(\widehat A_B)
\]
is an isomorphism.  The associated-graded formula of
Theorem~\ref{thm:bnla-adic-associated-graded} and
Lemma~\ref{lem:bnla-complete-filtered-connectivity} show independently that
$\widehat A_B$ is connective.  Therefore
\[
    K:=\operatorname{fib}(A\to\widehat A_B)
\]
is connective.

The full subcategory of derived $I$-complete $A$-modules is stable.  Since
both $A$ and $\widehat A_B$ are derived $I$-complete, $K$ is derived
$I$-complete.  The assumed base-change equivalence gives
\[
    B\otimes_AK\simeq0.
\]
The canonical truncation morphism
\[
    B\longrightarrow H\pi_0(B)
    \simeq H\!\left(\pi_0(A)/I\right)
\]
therefore gives
\[
    K\otimes_AH\!\left(\pi_0(A)/I\right)
    \simeq
    (K\otimes_AB)\otimes_BH\pi_0(B)
    \simeq0.
\]
Derived Nakayama for the connective derived $I$-complete module $K$ now gives
$K\simeq0$.  Thus $A\to\widehat A_B$ is an equivalence.
\end{proof}

\begin{proposition}[Intrinsic characterization of complete almost-finite augmentations]
\label{prop:bnla-spectral-aft-criterion}
Let $A\to B$ be a morphism of connective $E_\infty$-rings inducing a
surjection on $\pi_0$, and let $I$ be the kernel.  The following are
equivalent:
\begin{enumerate}
    \item $A\to B$ is almost finitely presented and $A$ is derived
    $I$-complete;
    \item $B$ is almost perfect as an $A$-module and $A$ is derived
    $I$-complete.
\end{enumerate}
Under these conditions the canonical map
\[
    A\xrightarrow{\sim}\widehat A_B
\]
is an equivalence, and the completed adic filtration has associated graded
\[
    \operatorname{Sym}^{\mathrm{gr}}_B(L_{B/A}[-1](1)).
\]
\end{proposition}

\begin{proof}
The equivalence of the two conditions is
Lemma~\ref{lem:bnla-spectral-cell-aft}.  Under either condition,
$\pi_0(B)$ is finitely presented as a $\pi_0(A)$-module.  Since
$\pi_0(A)\twoheadrightarrow\pi_0(B)$, the ideal
\[
    I=\ker\bigl(\pi_0(A)\to\pi_0(B)\bigr)
\]
is therefore finitely generated.  Put
\[
    F^\bullet A_{\mathrm{ad}}
    :=
    \operatorname{adic}(A\to B)
\]
and let
\[
    \widehat{F^\bullet A_{\mathrm{ad}}}
\]
be its filtered completion.  Theorem~\ref{thm:bnla-adic-associated-graded}
identifies its associated graded with
\[
    \operatorname{Sym}^{\mathrm{gr}}_B(L_{B/A}[-1](1)).
\]
This graded algebra is connective and almost finitely augmented by
Lemma~\ref{lem:bnla-positive-graded-aft}.  Hence the completed filtered
algebra is connective by
Lemma~\ref{lem:bnla-complete-filtered-connectivity}, and its filtration-zero
term is $\widehat A_B$.

It remains to prove
\[
    A\xrightarrow{\sim}\widehat A_B.
\]
By Lemma~\ref{lem:bnla-filtered-derived-completion}, it is enough to prove
that
\[
    f:
    B\otimes_A\widehat A_B
    \longrightarrow
    B
\]
is an equivalence.  Lemma~\ref{lem:bnla-filtered-derived-completion} gives
\[
    \pi_0(\widehat A_B)
    \simeq
    \pi_0(A)^\wedge_I.
\]
For a finitely generated ideal, reduction of the classical completion modulo
$I$ is the original quotient:
\[
    \pi_0(A)^\wedge_I/I\pi_0(A)^\wedge_I
    \simeq
    \pi_0(A)/I
    \simeq
    \pi_0(B).
\]
The connective tensor-product spectral sequence therefore gives
\[
\begin{aligned}
    \pi_0(B\otimes_A\widehat A_B)
    &\simeq
    \pi_0(B)\otimes_{\pi_0(A)}\pi_0(A)^\wedge_I\\
    &\simeq
    \pi_0(B),
\end{aligned}
\]
so $\pi_0(f)$ is an isomorphism.

We next show $L_f\simeq0$.
Lemma~\ref{lem:bnla-completed-adic-cotangent} identifies the filtered
cotangent fibres before and after completion:
\[
    \operatorname{cot}(F^\bullet A_{\mathrm{ad}})
    \xrightarrow{\sim}
    \operatorname{cot}
    \left(
        \widehat{F^\bullet A_{\mathrm{ad}}}
    \right).
\]

Applying augmentation gives an equivalence
\[
    L_{B/A}[-1]
    \xrightarrow{\sim}
    L_{B/\widehat A_B}[-1].
\]
Let
\[
    C:=B\otimes_A\widehat A_B,
    \qquad
    f:C\longrightarrow B.
\]
Cotangent transitivity for $A\to\widehat A_B\to B$, together with cotangent
base change for $B\to C$, identifies the cofiber of the preceding shifted
cotangent comparison with
\[
    L_{B/C}[-1]=L_f[-1].
\]
Hence $L_f[-1]\simeq0$, equivalently $L_f\simeq0$.  Chapter~2's spectral
equivalence criterion
($\pi_0$-isomorphism together with vanishing relative cotangent complex)
gives $f\simeq\operatorname{id}_B$.  Now
Lemma~\ref{lem:bnla-filtered-derived-completion} gives
$A\simeq\widehat A_B$.

The associated-graded statement is
Theorem~\ref{thm:bnla-adic-associated-graded}.
\end{proof}

\begin{definition}[Complete almost-finite augmented arrows]
\label{def:bnla-complete-aft-augmented}
Let
\[
    \operatorname{CAlg}^{\mathrm{sp},\wedge\mathrm{aft}}_{R/B}
    \subset
    \operatorname{CAlg}(\operatorname{Sp})^{\mathrm{cn}}_{R//B}
\]
be the full subcategory of arrows $A\to B$ for which
\begin{enumerate}
    \item $\pi_0(A)\twoheadrightarrow\pi_0(B)$;
    \item $B$ is almost perfect as an $A$-module;
    \item $A$ is derived-complete for the kernel of the preceding surjection.
\end{enumerate}
By Proposition~\ref{prop:bnla-spectral-aft-criterion}, these conditions are
equivalent to almost finite presentation together with derived completeness,
and the constructed adic completion is automatically $A$ itself.  No
filtration or completion presentation is retained.
\end{definition}

\subsection{Square-zero finiteness}

\begin{lemma}[Spectral square-zero finiteness]
\label{lem:bnla-spectral-square-zero-finiteness}
Let $B$ be coherent, let $M$ be a connective almost perfect $B$-module, and
let
\[
    B_\alpha\longrightarrow B
\]
be the structured square-zero extension classified by
$\alpha:L_{B/R}[-1]\to M$.  Then $B$ is almost perfect as a
$B_\alpha$-module and $B_\alpha\to B$ is almost finitely presented.
\end{lemma}

\begin{proof}
Let $\mathcal K$ be the thick subcategory of
$\operatorname{Mod}_{B_\alpha}$ generated by restrictions of almost-perfect
$B$-modules.  Tensoring the square-zero fibre sequence
\[
    M\longrightarrow B_\alpha\longrightarrow B
\]
with an almost-perfect $B_\alpha$-module $N$ gives
\[
    M\otimes_B(B\otimes_{B_\alpha}N)
    \longrightarrow N
    \longrightarrow B\otimes_{B_\alpha}N.
\]
The two outer terms are restrictions of almost-perfect $B$-modules, hence
$N\in\mathcal K$.

Conversely $B_\alpha\in\mathcal K$ by the original fibre sequence.  If a
connective $N\in\mathcal K$ has first nonzero homotopy in degree $k$, that
homotopy group is finitely generated.  Choose finitely many generators and a
map $B_\alpha[k]^{\oplus r}\to N$.  Its cofiber is $(k+1)$-connective and
remains in $\mathcal K$.  Iteration produces perfect approximations through
every Postnikov degree, so every connective object of $\mathcal K$ is almost
perfect.  In particular $B$ is almost perfect over $B_\alpha$.  The
surjective spectral cell criterion
Lemma~\ref{lem:bnla-spectral-cell-aft} then makes the augmentation almost
finitely presented.
\end{proof}

\begin{lemma}[Square-zero objects lie in the complete almost-finite sector]
\label{lem:bnla-sqz-complete-aft}
If $(M,\alpha)\in
\operatorname{Der}^{\mathrm{aperf}}_{B/R,\geq0}$, then
\[
    B_\alpha\longrightarrow B
\]
belongs to
$\operatorname{CAlg}^{\mathrm{sp},\wedge\mathrm{aft}}_{R/B}$.
\end{lemma}

\begin{proof}
Lemma~\ref{lem:bnla-spectral-square-zero-finiteness} makes $B$ almost perfect
as a $B_\alpha$-module and makes the augmentation almost finitely presented.
Since $B$ is coherent, Lemma~\ref{lem:bnla-augmentation-ideal-finite-presentation}
applied to $B_\alpha\to B$ shows that
\[
    I:=\ker\bigl(\pi_0(B_\alpha)\to\pi_0(B)\bigr)
\]
is finitely presented, hence finitely generated.  The structured square-zero
multiplication gives $I^2=0$, so $I$ is nilpotent.  Every module on which a
power of $I$ acts trivially is derived $I$-complete; in particular
$B_\alpha$ is derived $I$-complete.  Apply
Proposition~\ref{prop:bnla-spectral-aft-criterion}.
\end{proof}

\subsection{Comonadicity}

\begin{theorem}[Spectral cotangent--square-zero comonadicity]
\label{thm:bnla-spectral-cot-sqz-comonadic}
For coherent $B$, the adjunction of
Proposition~\ref{prop:bnla-cot-sqz-adjunction} restricts to
\[
    \operatorname{cot}_{B/R}:
    \operatorname{CAlg}^{\mathrm{sp},\wedge\mathrm{aft}}_{R/B}
    \rightleftarrows
    \operatorname{Der}^{\mathrm{aperf}}_{B/R,\geq0}:
    \operatorname{sqz}_{B/R},
\]
and the left adjoint is comonadic.
\end{theorem}

\begin{proof}
The right adjoint lands in the displayed source category by
Lemma~\ref{lem:bnla-sqz-complete-aft}.  For an arrow in the source,
$L_{B/A}[-1]$ is connective because $\pi_0(A)\to\pi_0(B)$ is surjective, and
it is almost perfect by the almost-finite cotangent criterion.  Thus the left
adjoint lands in the displayed derivation category.

We verify Barr--Beck directly.  First, $\operatorname{cot}_{B/R}$ is
conservative.  If $A\to A'$ over $B$ induces
$L_{B/A}\simeq L_{B/A'}$, Theorem~\ref{thm:bnla-adic-associated-graded}
identifies the associated gradeds of the two completed adic filtrations.
Both augmentations are fixed by the constructed completion, so complete
filtered detection gives $A\simeq A'$.

Now let $A^\bullet$ be a cosimplicial object of the complete almost-finite
source category such that
$\operatorname{cot}(A^\bullet)$ admits a split totalization in
$\operatorname{Der}^{\mathrm{aperf}}_{B/R,\geq0}$.  Form
\[
    A:=\operatorname*{Tot}A^\bullet
\]
first in the unbounded category
$\operatorname{CAlg}(\operatorname{Sp})_{R//B}$.
Apply completed adic filtration termwise.  In the unbounded filtered category,
associated graded commutes with the totalization, and
\[
\begin{aligned}
\operatorname{Gr}
\operatorname*{Tot}
\widehat{\operatorname{adic}}(A^\bullet\to B)
&\simeq
\operatorname*{Tot}
\operatorname{Sym}^{\mathrm{gr}}_B
\bigl(\operatorname{cot}(A^\bullet)(1)\bigr)
\\
&\simeq
\operatorname{Sym}^{\mathrm{gr}}_B
\left(
\operatorname*{Tot}\operatorname{cot}(A^\bullet)(1)
\right).
\end{aligned}
\]
The last equivalence uses only that the augmented cosimplicial cotangent
object is split: a split totalization is absolute and hence is preserved by
every functor, including the graded free-commutative-algebra functor.

The total filtered algebra is complete as a limit of complete filtered
algebras.  Its associated graded is connective; hence its filtration-zero
term is connective by
Lemma~\ref{lem:bnla-complete-filtered-connectivity}.  The total cotangent
fibre
\[
    M:=\operatorname*{Tot}\operatorname{cot}(A^\bullet)
\]
is connective and almost perfect by the Barr--Beck hypothesis.  The positive
graded symmetric algebra on $M$ is almost finitely augmented by
Lemma~\ref{lem:bnla-positive-graded-aft}.  The same filtered approximation
argument as in Proposition~\ref{prop:bnla-spectral-aft-criterion} therefore
shows that the filtration-zero augmentation belongs to
$\operatorname{CAlg}^{\mathrm{sp},\wedge\mathrm{aft}}_{R/B}$.

Finally compare cotangent fibres.  In the unbounded filtered category there is
a canonical morphism
\[
    \operatorname{cot}
    \left(
      \operatorname*{Tot}
      \widehat{\operatorname{adic}}(A^\bullet\to B)
    \right)
    \longrightarrow
    \operatorname*{Tot}
    \operatorname{cot}
    \left(
      \widehat{\operatorname{adic}}(A^\bullet\to B)
    \right).
\]
On associated graded this is the cotangent comparison for the split
cosimplicial graded free algebra and is therefore an equivalence.  The source
is an almost-perfect filtered module over a complete almost-finite filtered
algebra and is complete by
Lemma~\ref{lem:bnla-aperf-filtered-complete}; the target is a
limit of complete filtered modules and is complete.  Complete filtered
detection makes the comparison an equivalence.  Weight one gives
\[
    \operatorname{cot}_{B/R}(A)
    \xrightarrow{\sim}
    \operatorname*{Tot}
    \operatorname{cot}_{B/R}(A^\bullet).
\]
Thus the totalization exists in the source and is preserved by cot.  Together
with conservativity, Barr--Beck--Lurie comonadicity applies.
\end{proof}

\begin{definition}[Spectral relative coLie comonad]
\label{def:bnla-spectral-colie}
Define
\[
    \operatorname{coLie}^{\pi,E_\infty}_{B/R}
    :=
    \operatorname{cot}_{B/R}\operatorname{sqz}_{B/R}
\]
on $\operatorname{Der}^{\mathrm{aperf}}_{B/R,\geq0}$, with the comonad
structure induced by Theorem~\ref{thm:bnla-spectral-cot-sqz-comonadic}.
\end{definition}

\begin{corollary}[Canonical square-zero resolution]
\label{cor:bnla-spectral-square-zero-resolution}
Every
$A\to B\in\operatorname{CAlg}^{\mathrm{sp},\wedge\mathrm{aft}}_{R/B}$
admits a canonical augmented cosimplicial cobar resolution
\[
    A\longrightarrow A^\bullet
\]
whose positive terms are constructed structured square-zero extensions of
$B$, such that
\[
    A\xrightarrow{\sim}\operatorname*{Tot}A^\bullet
\]
and $\operatorname{cot}(A^\bullet)$ is split cosimplicial.  No square-zero
presentation is chosen.
\end{corollary}

\begin{proof}
This is the canonical cobar resolution of the comonadic adjunction.  Every
positive term lies in the image of $\operatorname{sqz}$; the standard extra
codegeneracies split the cotangent image; comonadicity identifies the
augmentation with the totalization.
\end{proof}

\begin{corollary}[Relative cobar resolution over a square-zero target]
\label{cor:bnla-spectral-relative-square-zero-resolution}
Let
\[
    A\to B
    \in
    \operatorname{CAlg}^{\mathrm{sp},\wedge\mathrm{aft}}_{R/B},
    \qquad
    M\in\operatorname{APerf}(B)_{\geq0},
\]
and let
\[
    \eta:A\longrightarrow B\oplus M[1]
\]
be a morphism over $B$.  There is a canonical coaugmented cosimplicial diagram
\[
    A\longrightarrow A^\bullet\longrightarrow B\oplus M[1]
\]
such that every positive term is obtained by applying
$\operatorname{sqz}$ to a morphism of relative derivations into the
zero-class derivation associated to $M[1]$,
$\operatorname{cot}(A)\to\operatorname{cot}(A^\bullet)$ is split, and
\[
    A\xrightarrow{\sim}\operatorname*{Tot}A^\bullet.
\]
\end{corollary}

\begin{proof}
A comonadic adjunction remains comonadic after passage to an overcategory over
an object in the image of its right adjoint.  Apply this to
$\operatorname{cot}\dashv\operatorname{sqz}$ over the split square-zero target
$B\oplus M[1]$.
\end{proof}
\section{The relative spectral partition-Lie algebroid monad}
\label{sec:bnla-spectral-lie-algebroid-monad}

We now dualize the relative coLie comonad.  The only external algebraic input
at this stage is the pro-coherent spectral Koszul-dual theory of
\cite{BrantnerCamposNuiten2025}; the relative anchor dependence is constructed
from the comonad of the preceding section.

\subsection{Dualization on dually almost-perfect anchors}

Continuous duality gives a fully faithful functor
\[
    \left(
        \operatorname{Der}^{\mathrm{aperf}}_{B/R,\geq0}
    \right)^{\mathrm{op}}
    \longrightarrow
    \operatorname{Anch}^{\mathrm{pc}}_{B/R},
\]
sending
\[
    L_{B/R}[-1]\xrightarrow{\alpha}M
\]
to
\[
    M^\vee\xrightarrow{\alpha^\vee}
    \mathbb T^{\mathrm{pc}}_{B/R}[1].
\]
Its essential image is
$\operatorname{dAPerfAnch}_{B/R}$.

\begin{definition}[Relative spectral partition-Lie functor on dually almost-perfect anchors]
\label{def:bnla-relative-spectral-lie-compact}
On $\operatorname{dAPerfAnch}_{B/R}$ define
\[
    \operatorname{Lie}^{\pi,E_\infty,\mathrm{dap}}_{B/R}
\]
by transporting the opposite comonad
\[
    \left(
        \operatorname{coLie}^{\pi,E_\infty}_{B/R}
    \right)^{\mathrm{op}}
\]
through continuous duality.
\end{definition}

\begin{proposition}[One-object spectral identification]
\label{prop:bnla-one-object-spectral-partition-monad}
For $R=B$, the functor
$\operatorname{Lie}^{\pi,E_\infty,\mathrm{dap}}_{B/B}$ agrees on dually
almost-perfect generators with the coherent spectral partition-Lie monad
\[
    \operatorname{Lie}^{\pi}_{B,E_\infty}
\]
constructed by pro-coherent PD Koszul duality in
\cite{BrantnerCamposNuiten2025}.
\end{proposition}

\begin{proof}
When $R=B$ the anchor target vanishes and a relative derivation is simply a
$B$-module.  The square-zero functor is the nonunital $E_\infty$ square-zero
algebra construction, and
$\operatorname{cot}\operatorname{sqz}$ is its Koszul-dual coalgebraic
comonad.  The refined continuous duality of
\cite{BrantnerCamposNuiten2025} identifies the opposite comonad on dually
almost-perfect objects with the PD Koszul dual of the nonunital
$E_\infty$-operad, which is their spectral partition-Lie monad.
\end{proof}

\subsection{Anchor-weight filtration and continuity}

The filtration used to extend the relative coLie comonad is a filtration of
the already constructed coLie object.  It is not completed by definition.
The next lemma proves that, in the almost-perfect range in which it is used,
the filtered cotangent object is automatically complete.

\begin{lemma}[Filtered spectral square-zero cotangent finiteness]
\label{lem:bnla-filtered-square-zero-cotangent-finiteness}
Let $B$ be coherent and connective.  Let $F^\bullet M$ be a connective
almost-perfect filtered $B$-module which is generated, under finite extensions
and retracts, by finitely many filtration shifts of connective almost-perfect
$B$-modules.  Let
\[
    L_{B/R}[-1]\longrightarrow F^\bullet M
\]
be a filtered relative derivation, and form the corresponding structured
square-zero extension in filtered connective commutative ring spectra
\[
    C^\bullet\longrightarrow B,
\]
where $B$ is concentrated in filtration weight zero.  Then
\[
    L_{B/C^\bullet}[-1]
\]
is a connective almost-perfect filtered $B$-module and is complete.
\end{lemma}

\begin{proof}
The proof of spectral square-zero finiteness is internal to the filtered
stable category.  We spell out the points needed here.

The category $\operatorname{Fil}(\operatorname{Sp})$ is stable and
presentable, its $t$-structure is pointwise, and Day convolution is exact and
preserves colimits separately.  A finite filtration shift of a perfect
$B$-module is perfect as a filtered $B$-module.  Consequently $F^\bullet M$
is almost perfect: choose perfect approximations to its finitely many
almost-perfect generators in every Postnikov range and take the same finite
extensions and retracts in the filtered category.

Let $\mathcal K$ be the thick subcategory of filtered
$C^\bullet$-modules generated by restrictions of connective almost-perfect
filtered $B$-modules.  Tensoring the square-zero fibre sequence
\[
    F^\bullet M\longrightarrow C^\bullet\longrightarrow B
\]
with a connective filtered $C^\bullet$-module $N$ gives
\[
    F^\bullet M\otimes_B
    \bigl(B\otimes_{C^\bullet}N\bigr)
    \longrightarrow
    N
    \longrightarrow
    B\otimes_{C^\bullet}N.
\]
Thus $C^\bullet$ itself belongs to $\mathcal K$, and every connective object
of $\mathcal K$ has finitely generated first nonzero filtered homotopy: this
property holds on the generators and is preserved by finite fibres, cofibres,
and retracts.  If $N\in\mathcal K$ is $k$-connective, choose finitely many
filtered generators of $\pi_kN$.  They give a morphism from a finite direct
sum of filtration shifts of $C^\bullet[k]$ to $N$ whose cofiber is
$(k+1)$-connective and remains in $\mathcal K$.  Iterating produces perfect
filtered approximations to $N$ through every Postnikov degree.  Hence every
connective object of $\mathcal K$ is almost perfect.  In particular $B$ is
almost perfect as a filtered $C^\bullet$-module.

We now repeat the positive relation-cell construction of
Lemma~\ref{lem:bnla-positive-relation-cell} internally in
$\operatorname{Fil}(\operatorname{Sp})$.  Its proof uses only stability, Day
convolution, the two-sided bar construction, and finite skeleta of
$E\Sigma_k$, so it applies verbatim to filtered connective modules.  In
particular, attachment of a finite positive filtered relation cell produces a
finite-cell filtered algebra which is almost perfect over the previous stage,
and a surjection onto the first unresolved relative homotopy group raises the
connectivity of the comparison map by one.

Starting from the filtered augmentation $C^\bullet\to B$, choose finite
filtered generators for its degree-zero kernel and attach the corresponding
relation cells.  Inductively suppose that a finite-cell filtered
$C^\bullet$-algebra $B_n\to B$ is an equivalence through Postnikov degree $n$
and is almost perfect as a filtered $C^\bullet$-module.  Since $B$ is almost
perfect over $C^\bullet$, its relative cofiber is almost perfect, so its first
unresolved homotopy group is finitely generated.  The natural
relative-cofiber-to-cotangent morphism is compatible with evaluation in every
filtration degree, and the relative spectral Hurewicz comparison identifies
this group with the first unresolved homotopy group of $L_{B/B_n}$.  Finitely
many filtered positive relation cells therefore produce $B_{n+1}\to B$, an
equivalence one degree further, and the filtered relation-cell lemma keeps
$B_{n+1}$ almost perfect over $C^\bullet$.

For every fixed Postnikov degree only finitely many cells and symmetric word
lengths occur.  The associated graded of each relation cell has homogeneous
pieces
\[
    \left(P[1]^{\otimes k}\right)_{h\Sigma_k}
\]
for finite perfect filtered cell modules $P$.  A sufficiently large finite
skeleton of $E\Sigma_k$ computes the required fixed Postnikov range, so each
visible homogeneous piece is approximated there by a finite colimit of finite
tensor products of perfect filtered modules.  Hence, for every $d$, there is
a finite-cell filtered commutative $C^\bullet$-algebra $B^{(d)}\to B$ inducing
an equivalence after $d$-truncation.  Mapping out of $B$ into $d$-truncated
filtered connective $C^\bullet$-algebras is therefore computed by the compact
object $B^{(d)}$ and commutes with filtered colimits.  Thus
$C^\bullet\to B$ is almost finitely presented in filtered connective
commutative ring spectra.

The filtered cotangent universal property now shows directly that
$L_{B/C^\bullet}$ is almost perfect: for every filtered system of
$d$-coconnective filtered $B$-modules $N_\lambda$, the comparison
\[
\operatorname*{colim}_\lambda
\operatorname{Map}_B(L_{B/C^\bullet},N_\lambda)
\longrightarrow
\operatorname{Map}_B
\left(L_{B/C^\bullet},\operatorname*{colim}_\lambda N_\lambda\right)
\]
is the fibre, over the fixed augmentation, of the corresponding filtered
mapping-space comparison for the almost finitely presented algebra $B$ and is
therefore an equivalence.  Surjectivity on $\pi_0$ gives connectivity after
shifting by $[-1]$.

Finally $B$, concentrated in filtration weight zero, is a complete connective
filtered algebra.  Lemma~\ref{lem:bnla-aperf-filtered-complete} therefore
applies to the connective almost-perfect filtered $B$-module
$L_{B/C^\bullet}[-1]$ and proves its completeness.
\end{proof}

\begin{construction}[Anchor-weight filtration]
\label{con:bnla-colie-anchor-filtration}
Let
\[
    \alpha:L_{B/R}[-1]\longrightarrow M
\]
be a relative derivation with $M$ connective and almost perfect.  Put the
source $L_{B/R}[-1]$ in filtration weight zero and put $M$ in filtration
weight one.  Thus
\[
    F^0M=M,
    \qquad
    F^1M=M,
    \qquad
    F^2M=0,
\]
with the evident filtered $B$-module structure, and the classifying map becomes
a filtered derivation
\[
    L_{B/R}[-1]\longrightarrow F^\bullet M.
\]
Form the structured filtered square-zero extension
\[
    B^\bullet_\alpha\longrightarrow B
\]
and define
\[
\boxed{
    F^\bullet
    \operatorname{coLie}^{\pi,E_\infty}_{B/R}(M,\alpha)
    :=
    L_{B/B^\bullet_\alpha}[-1].}
\]
Evaluation at filtration weight zero commutes with the square-zero pullback and
with relative cotangent formation, so
\[
\boxed{
    F^0
    \operatorname{coLie}^{\pi,E_\infty}_{B/R}(M,\alpha)
    \simeq
    \operatorname{coLie}^{\pi,E_\infty}_{B/R}(M,\alpha).}
\]
Lemma~\ref{lem:bnla-filtered-square-zero-cotangent-finiteness} shows that the
displayed filtered cotangent object is already complete.  Thus no completion
changes the coLie object.

Associated graded is symmetric monoidal and preserves the finite pullback
which defines the structured square-zero extension.  The source of the
classifying map lies in weight zero and its target in weight one, so the
associated-graded classifying map vanishes.  Consequently there is a canonical
equivalence of complete graded cotangent objects
\[
\boxed{
    \operatorname{Gr}
    F^\bullet\operatorname{coLie}^{\pi,E_\infty}_{B/R}(M,\alpha)
    \simeq
    \operatorname{Gr}
    F^\bullet\operatorname{coLie}^{\pi,E_\infty}_{B/R}(M,0).}
\]
All constructions are natural in morphisms of relative derivations.
\end{construction}

\begin{definition}[Almost eventually constant tower]
\label{def:bnla-almost-eventually-constant}
A tower
\[
    \cdots\longrightarrow M_2\longrightarrow M_1
\]
of connective $B$-modules is \emph{almost eventually constant} if, for every
Postnikov degree $d$, the induced tower after $d$-truncation is eventually
constant.
\end{definition}

\begin{definition}[Finite cosimplicial diagram]
\label{def:bnla-finite-cosimplicial}
A cosimplicial diagram is \emph{finite} if it is $n$-coskeletal for some
$n\geq0$.  Equivalently, under the stable Dold--Kan description its normalized
cochain object is supported in a finite range.  These are the finite
cosimplicial totalizations used in the right-left extension formalism below.
\end{definition}

\begin{lemma}[Continuity properties of relative spectral coLie]
\label{lem:bnla-spectral-colie-continuity}
Let $R\to B$ be a morphism with $B$ coherent.  The functor
\[
    \operatorname{coLie}^{\pi,E_\infty}_{B/R}
    =\operatorname{cot}_{B/R}\operatorname{sqz}_{B/R}
\]
on connective almost-perfect derivations has the following properties.
\begin{enumerate}
    \item It preserves the connective almost-perfect sector.
    \item It preserves sifted colimits whenever the colimit remains in that
    sector.
    \item It preserves finite cosimplicial totalizations whose terms and
    totalization have connective almost-perfect underlying modules.
    \item It preserves inverse limits of almost eventually constant towers in
    that sector.
\end{enumerate}
\end{lemma}

\begin{proof}
For (1), Lemma~\ref{lem:bnla-sqz-complete-aft} puts the square-zero extension
in the complete almost-finite sector.  Cotangent transitivity then makes the
new cotangent fibre connective and almost perfect.

For (2), structured square-zero formation is computed from the derivation by
a finite pullback in commutative ring spectra.  The forgetful functor from
commutative ring spectra creates limits and preserves sifted colimits, while
sifted colimits in spectra commute with finite limits because the category is
stable.  Thus square-zero formation preserves the relevant sifted colimits.
The cotangent functor is a left adjoint and preserves them as well.

For (3), let
\[
    (M^\bullet,\alpha^\bullet)
\]
be a finite cosimplicial diagram as in the statement and put
\[
    (M,\alpha)
    :=
    \operatorname*{Tot}
    (M^\bullet,\alpha^\bullet).
\]
Structured square-zero formation is the right adjoint of
Proposition~\ref{prop:bnla-cot-sqz-adjunction}, so it preserves this
totalization.  Applying the anchor-weight construction gives a canonical map
of filtered cotangent objects
\[
\begin{aligned}
&F^\bullet
\operatorname{coLie}^{\pi,E_\infty}_{B/R}(M,\alpha)
\\
&\qquad\longrightarrow
\operatorname*{Tot}
F^\bullet
\operatorname{coLie}^{\pi,E_\infty}_{B/R}
(M^\bullet,\alpha^\bullet).
\end{aligned}
\]
The source is connective almost perfect by
Lemma~\ref{lem:bnla-filtered-square-zero-cotangent-finiteness}; the target is
a finite limit of connective almost-perfect filtered modules and is therefore
connective almost perfect.  Both are complete by
Lemma~\ref{lem:bnla-aperf-filtered-complete}.

Pass to associated graded.  Construction~\ref{con:bnla-colie-anchor-filtration}
reduces the comparison to the zero-anchor functor
\[
    N\longmapsto
    \operatorname{cot}(B\oplus N(1))
\]
on strictly positively graded $B$-modules.  Using the bar resolution, this
functor is obtained by colimits and compositions from the homogeneous derived
symmetric-power functors
\[
    \operatorname{Sym}^n_B.
\]
These functors are $n$-homogeneous, hence $n$-excisive, and preserve
filtered colimits.  We use the following finite-totalization consequence in
the stable setting.  If
\[
    F:\mathcal C\longrightarrow\mathcal D
\]
is an $n$-excisive filtered-colimit-preserving functor between cocomplete
stable $\infty$-categories and $X^\bullet$ is right Kan extended from
$\Delta_{\leq m}$ for some finite $m$, then the canonical map
\[
    F(\operatorname*{Tot}X^\bullet)
    \longrightarrow
    \operatorname*{Tot}F(X^\bullet)
\]
is an equivalence.  Indeed, induction on the Goodwillie degree reduces to an
$n$-homogeneous functor
\[
    X\longmapsto
    G(X,\ldots,X)_{h\Sigma_n}
\]
for a symmetric functor $G:\mathcal C^n\to\mathcal D$ which preserves
colimits, hence is exact, separately in every variable.  Right Kan extension
from $\Delta_{\leq m}$ is a finite-limit condition, so separate exactness
preserves it in each variable.  The diagonal
$\Delta\to\Delta^n$ is right cofinal, and therefore
\[
\begin{aligned}
G(\operatorname*{Tot}X^\bullet,\ldots,
  \operatorname*{Tot}X^\bullet)
&\simeq
\operatorname*{Tot}
G(X^\bullet,\ldots,X^\bullet).
\end{aligned}
\]
Homotopy orbits are exact in the stable target, so the same holds after
$(-)_{h\Sigma_n}$.  The induction through the Goodwillie fibre sequences
proves the claim.

Applying this to every derived symmetric power gives the finite-totalization
mechanism used in \cite{BrantnerMagidsonNuiten2025}.  Since finite limits are
finite colimits in a stable category, preservation of the fixed finite
totalization is stable under the colimits appearing in the bar resolution;
it is also stable under composition.  Hence the associated-graded comparison
is an equivalence.
Complete filtered detection then gives an equivalence before passing to
associated graded.  Taking filtration weight zero yields
\[
    \operatorname{coLie}^{\pi,E_\infty}_{B/R}(M,\alpha)
    \xrightarrow{\sim}
    \operatorname*{Tot}
    \operatorname{coLie}^{\pi,E_\infty}_{B/R}
    (M^\bullet,\alpha^\bullet).
\]

For (4), let
\[
    \cdots\longrightarrow(M_2,\alpha_2)
    \longrightarrow(M_1,\alpha_1)
\]
be an almost eventually constant tower and put
\[
    (M,\alpha):=\lim_n(M_n,\alpha_n).
\]
Fix a Postnikov degree $d$.  For $n$ sufficiently large the fibre
\[
    K_n:=\operatorname{fib}(M\to M_n)
\]
is $(d+1)$-connective.  Since structured square-zero formation is the right
adjoint in Proposition~\ref{prop:bnla-cot-sqz-adjunction}, it preserves the
inverse limit, and the induced morphism
\[
    B_\alpha\longrightarrow B_{\alpha_n}
\]
has $(d+1)$-connective underlying fibre.  The relative cotangent connectivity
theorem then makes its relative cotangent complex $(d+2)$-connective.
Cotangent transitivity, followed by the shift $[-1]$, shows that
\[
    \operatorname{cofib}
    \left(
        \operatorname{coLie}^{\pi,E_\infty}_{B/R}(M,\alpha)
        \longrightarrow
        \operatorname{coLie}^{\pi,E_\infty}_{B/R}(M_n,\alpha_n)
    \right)
\]
is $(d+1)$-connective.  Hence the two coLie objects have the same
$d$-truncation for all sufficiently large $n$.  This holds for every $d$.
Postnikov completeness therefore identifies
\[
    \operatorname{coLie}^{\pi,E_\infty}_{B/R}(M,\alpha)
    \xrightarrow{\sim}
    \lim_n
    \operatorname{coLie}^{\pi,E_\infty}_{B/R}(M_n,\alpha_n),
\]
which proves (4).
\end{proof}

\begin{theorem}[Spectral partition-Lie algebroid monad]
\label{thm:bnla-spectral-lie-algebroid-monad}
Let $R\to B$ be a morphism of connective $E_\infty$-rings with $B$ coherent.
There is a unique sifted-colimit-preserving monad
\[
\boxed{
    \operatorname{Lie}^{\pi,E_\infty}_{B/R}:
    \operatorname{Anch}^{\mathrm{pc}}_{B/R}
    \longrightarrow
    \operatorname{Anch}^{\mathrm{pc}}_{B/R}}
\]
whose restriction to $\operatorname{dAPerfAnch}_{B/R}$ is
$\operatorname{Lie}^{\pi,E_\infty,\mathrm{dap}}_{B/R}$.
\end{theorem}

\begin{proof}
The pro-coherent completion theory of
\cite{BrantnerCamposNuiten2025} identifies
$\operatorname{QC}^{\vee}_B$ as the sifted-colimit completion of continuous
duals of connective almost-perfect modules subject to the relations generated
by the finite totalizations and almost-eventually-constant towers appearing
in their right-left extension formalism.  Passing to the slice over the fixed
object $\mathbb T^{\mathrm{pc}}_{B/R}[1]$ applies the same completion to the
source and transports the anchor morphism.

Under continuous duality, preservation of the defining completion relations
is exactly Lemma~\ref{lem:bnla-spectral-colie-continuity}.  Hence the opposite
dual of the coLie comonad extends uniquely to a sifted-colimit-preserving
endofunctor on all anchored pro-coherent objects.  The coLie counit and
comultiplication dualize on the dense dually almost-perfect subcategory to a
unit and multiplication; the same continuity lemma lets those natural
transformations extend uniquely.  The monad identities hold on the dense
subcategory and therefore on the unique extensions.
\end{proof}

\begin{definition}[Spectral partition-Lie algebroids]
\label{def:bnla-spectral-partition-lie-algebroids}
Define
\[
\boxed{
    \operatorname{LieAlgd}^{\pi,E_\infty}_{B/R}
    :=
    \operatorname{Alg}_{\operatorname{Lie}^{\pi,E_\infty}_{B/R}}
    \left(\operatorname{Anch}^{\mathrm{pc}}_{B/R}\right).}
\]
An object is called a spectral partition-Lie algebroid over $B/R$.
\end{definition}

No partition-Lie structure is supplied as a presentation of a formal
groupoid.  The category has been constructed from spectral cotangent theory,
structured square-zero extensions, comonadicity, and pro-coherent Koszul
duality.

\begin{corollary}[One-object recovery]
\label{cor:bnla-spectral-one-object-recovery}
There is a canonical equivalence
\[
    \operatorname{LieAlgd}^{\pi,E_\infty}_{B/B}
    \simeq
    \operatorname{Alg}_{\operatorname{Lie}^{\pi}_{B,E_\infty}}
    \left(\operatorname{QC}^{\vee}_B\right).
\]
\end{corollary}

\begin{proof}
For $R=B$ the anchor target vanishes.  The monads agree on the dense
dually almost-perfect subcategory by
Proposition~\ref{prop:bnla-one-object-spectral-partition-monad}, and the
uniqueness clause in
Theorem~\ref{thm:bnla-spectral-lie-algebroid-monad} extends the identification.
\end{proof}

\begin{lemma}[Zero derivations form the one-object subcomonad]
\label{lem:bnla-zero-derivation-subcomonad}
Let
\[
    z:
    \operatorname{APerf}(B)_{\geq0}
    \longrightarrow
    \operatorname{Der}^{\mathrm{aperf}}_{B/R,\geq0},
    \qquad
    M\longmapsto
    \left(L_{B/R}[-1]\xrightarrow{0}M\right).
\]
Then there is a canonical natural equivalence of comonads on the indicated
connective almost-perfect sector
\[
\boxed{
    \operatorname{coLie}^{\pi,E_\infty}_{B/R}\,z
    \xrightarrow{\sim}
    z\,\operatorname{coLie}^{\pi,E_\infty}_{B/B}.}
\]
Under this equivalence the underlying module on both sides is
\[
    L_{B/(B\oplus M)}[-1],
\]
and the relative anchor to $L_{B/R}[-1]$ is canonically nullhomotopic.
\end{lemma}

\begin{proof}
For the zero derivation, structured square-zero formation is the split
extension
\[
    B\oplus M\longrightarrow B.
\]
Its $R$-algebra structure factors through the canonical section
\[
    s:B\longrightarrow B\oplus M,
\]
and the augmentation
$q:B\oplus M\to B$ satisfies $qs=\operatorname{id}_B$.  Cotangent
transitivity for
\[
    R\longrightarrow B\oplus M\xrightarrow{q}B
\]
gives
\[
    B\otimes_{B\oplus M}L_{(B\oplus M)/R}
    \longrightarrow
    L_{B/R}
    \xrightarrow{\delta}
    L_{B/(B\oplus M)}.
\]
Functoriality of the cotangent complex along the $R$-algebra section $s$
produces a section
\[
    L_{B/R}
    \longrightarrow
    B\otimes_{B\oplus M}L_{(B\oplus M)/R}
\]
of the first morphism.  The transitivity triangle supplies a canonical
nullhomotopy of $\delta$ after precomposition with that first morphism.
Precomposing further with its section gives a canonical nullhomotopy
\[
    \delta\simeq0.
\]
After shifting, the object
$\operatorname{cot}_{B/R}(B\oplus M\to B)$ therefore lies in the zero-anchor
subcategory of relative derivations.

The underlying augmented $B$-algebra
$B\oplus M\to B$ is independent of the map $R\to B$.  Consequently the
underlying cotangent module is exactly the one occurring for base $B/B$:
\[
    L_{B/(B\oplus M)}[-1]
    =
    \operatorname{coLie}^{\pi,E_\infty}_{B/B}(M).
\]
This identification is natural in $M$.

More precisely, restriction of scalars along $R\to B$ and the zero-section
functor fit into a commutative square of the two cotangent--square-zero
adjunctions on the split square-zero image.  The preceding nullhomotopy is the
mate of that square.  Hence the induced identifications commute with the
adjunction counits and with the iterated square-zero maps which define
comultiplication.  They therefore identify the induced comonad counit and
comultiplication, proving the displayed equivalence of comonads.
\end{proof}

\begin{theorem}[Zero-anchor spectral partition-Lie algebroids]
\label{thm:bnla-zero-anchor-spectral-algebroids}
The zero morphism
\[
    0\longrightarrow\mathbb T^{\mathrm{pc}}_{B/R}[1]
\]
carries a canonical spectral partition-Lie algebroid structure.  The full
subcategory of zero-anchor algebroids is canonically equivalent to spectral
partition-Lie algebras:
\[
\boxed{
    \operatorname{LieAlgd}^{\pi,E_\infty,0}_{B/R}
    \simeq
    \operatorname{Alg}_{\operatorname{Lie}^{\pi}_{B,E_\infty}}
    \left(\operatorname{QC}^{\vee}_B\right).}
\]
Taking the fibre of the anchor is right adjoint to the inclusion of this
zero-anchor subcategory.  On underlying pro-coherent objects it is the
ordinary homotopy fibre.  Moreover this anchor-fibre functor preserves sifted
colimits.
\end{theorem}

\begin{proof}
Continuous duality carries the zero derivation
\[
    L_{B/R}[-1]\xrightarrow{0}M
\]
to the zero anchored pro-coherent object
\[
    M^\vee\xrightarrow{0}\mathbb T^{\mathrm{pc}}_{B/R}[1].
\]
Lemma~\ref{lem:bnla-zero-derivation-subcomonad} and
Proposition~\ref{prop:bnla-one-object-spectral-partition-monad} identify, on
the dense dually almost-perfect zero-anchor subcategory, the restriction of
the relative monad with the one-object spectral partition-Lie monad
$\operatorname{Lie}^{\pi}_{B,E_\infty}$.  Both functors preserve the
completion relations used in
Theorem~\ref{thm:bnla-spectral-lie-algebroid-monad}; uniqueness of the
sifted-colimit extension therefore gives an equivalence of monads on the whole
zero-anchor pro-coherent category.  Passing to algebras gives
\[
    \operatorname{LieAlgd}^{\pi,E_\infty,0}_{B/R}
    \simeq
    \operatorname{Alg}_{\operatorname{Lie}^{\pi}_{B,E_\infty}}
    (\operatorname{QC}^{\vee}_B),
\]
and the zero source carries the canonical zero algebra structure.  Write
\[
    \mathbf 0_{B/R}
    :=
    \left(0\longrightarrow\mathbb T^{\mathrm{pc}}_{B/R}[1]\right)
\]
for the resulting zero-anchor algebroid.

Put
\[
    T:=\mathbb T^{\mathrm{pc}}_{B/R}[1].
\]
The terminal object $(T\xrightarrow{\operatorname{id}}T)$ of the anchored
category carries a canonical algebra structure for the relative monad: there
is a unique anchored morphism from the image of the terminal object under the
monad back to the terminal object, and the unit and associativity identities
are automatic by uniqueness.  Denote this terminal spectral partition-Lie
algebroid by
\[
    \mathbf T_{B/R}.
\]
There is a unique algebroid morphism
\[
    \mathbf 0_{B/R}\longrightarrow\mathbf T_{B/R}.
\]

For a spectral partition-Lie algebroid $\mathfrak a$, let
$\mathfrak a\to\mathbf T_{B/R}$ be its unique map to the terminal algebroid and
define
\[
\boxed{
    \operatorname{Fib}_{B/R}(\mathfrak a)
    :=
    \mathfrak a
    \times_{\mathbf T_{B/R}}
    \mathbf 0_{B/R}.}
\]
The forgetful functor from algebras over a monad creates limits.  Hence, if
$\rho:E\to T$ is the underlying anchor of $\mathfrak a$, then
\[
    U\operatorname{Fib}_{B/R}(\mathfrak a)
    \simeq
    \operatorname{fib}(\rho)
\]
with zero anchor.

Let $j$ denote the zero-anchor inclusion.  For a zero-anchor partition-Lie
algebra $\mathfrak g$, the unique morphism
$j\mathfrak g\to\mathbf T_{B/R}$ factors through
$\mathbf 0_{B/R}$ by the unique morphism $j\mathfrak g\to\mathbf 0_{B/R}$ in
the zero-anchor subcategory.  The pullback universal property therefore gives
natural equivalences
\[
\begin{aligned}
\operatorname{Map}_{\operatorname{LieAlgd}^{\pi,E_\infty}_{B/R}}
\left(j\mathfrak g,\mathfrak a\right)
&\simeq
\operatorname{Map}_{\operatorname{LieAlgd}^{\pi,E_\infty,0}_{B/R}}
\left(
    j\mathfrak g,
    \operatorname{Fib}_{B/R}(\mathfrak a)
\right)\\
&\simeq
\operatorname{Map}_{\operatorname{Alg}_{\operatorname{Lie}^{\pi}_{B,E_\infty}}}
\left(
    \mathfrak g,
    \operatorname{Fib}_{B/R}(\mathfrak a)
\right).
\end{aligned}
\]
Thus anchor fibre is right adjoint to the zero-anchor inclusion.

Finally, both the relative and zero-anchor monadic forgetful functors create
sifted colimits because their monads preserve them.  For a sifted diagram of
underlying anchors $E_i\to T$, stability gives
\[
\begin{aligned}
\operatorname*{colim}_i\operatorname{fib}(E_i\to T)
&\simeq
\operatorname*{colim}_i\operatorname{cofib}(E_i\to T)[-1]\\
&\simeq
\operatorname{cofib}
\left(
    \operatorname*{colim}_iE_i\longrightarrow T
\right)[-1]\\
&\simeq
\operatorname{fib}
\left(
    \operatorname*{colim}_iE_i\longrightarrow T
\right),
\end{aligned}
\]
where the sifted colimit of the constant diagram at $T$ is $T$.  Hence the
canonical comparison from the sifted colimit of the anchor fibres to the
anchor fibre of the sifted colimit is an equivalence on underlying
pro-coherent objects, and therefore an equivalence of zero-anchor algebroids.
\end{proof}

\section{Spectral Artinian extensions and affine formal-moduli recognition}
\label{sec:bnla-spectral-formal-recognition}

The affine recognition theorem is built from an intrinsic category of
spectral Artinian extensions.  Finite square-zero factorizations will be
proved to exist; they are not included in the definition.

\subsection{The intrinsic spectral Artinian category}

\begin{definition}[Spectral Artinian extension]
\label{def:bnla-spectral-artinian}
Let $R\to B$ be a morphism of connective $E_\infty$-rings with $B$ coherent.
A connective augmented $R$-algebra
\[
    A\longrightarrow B
\]
is a \emph{spectral Artinian extension} if:
\begin{enumerate}
    \item $\pi_0(A)\to\pi_0(B)$ is surjective with nilpotent kernel;
    \item $A\to B$ is almost finitely presented;
    \item $\operatorname{fib}(A\to B)$ is eventually coconnective.
\end{enumerate}
Write
\[
    \operatorname{Art}^{\mathrm{sp}}_{R/B}
    \subset
    \operatorname{CAlg}(\operatorname{Sp})^{\mathrm{cn}}_{R//B}
\]
for the full subcategory on these arrows.
\end{definition}

Thus an Artinian object is specified by properties of the arrow itself.  No
square-zero chain, coefficient module, derivation, or nilpotent filtration is
retained.

\begin{lemma}[Coherence across a finitely presented nilpotent extension]
\label{lem:bnla-coherent-nilpotent-ring}
Let $A_0\twoheadrightarrow B_0$ be a surjection of ordinary commutative
rings with nilpotent kernel $I$.  Assume that $B_0$ is coherent and that $I$
is finitely presented as an $A_0$-module.  Then:
\begin{enumerate}
    \item $A_0$ is coherent;
    \item every power $I^k$ is finitely presented as an $A_0$-module;
    \item if $M$ is a finitely presented $A_0$-module, then
    $I^kM/I^{k+1}M$ is finitely presented over $B_0$.
\end{enumerate}
\end{lemma}

\begin{proof}
We first record two elementary finite-presentation facts and then run a
noncircular induction through the nilpotent tower.

Suppose $R\twoheadrightarrow S=R/J$ with $J$ finitely generated.  If an
$S$-module $N$ is finitely presented as an $R$-module, then it is finitely
presented as an $S$-module.  Indeed, choose a presentation
\[
    R^a\longrightarrow R^b\longrightarrow N\longrightarrow0.
\]
If $K\subset R^b$ is the kernel, then $K$ is finitely generated and
$JR^b\subset K$ because $J$ annihilates $N$.  Reduction modulo $J$ gives
\[
    S^b\longrightarrow N
\]
with kernel $K/JR^b$, which is finitely generated over $S$.  Conversely, if
$S$ is finitely presented as an $R$-module and $N$ is finitely presented over
$S$, then $N$ is finitely presented over $R$ by composing finite
presentations.

We next prove the square-zero coherence step.  Let
\[
    C\twoheadrightarrow D=C/J,
    \qquad J^2=0,
\]
with $D$ coherent and $J$ finitely presented as a $C$-module.  Since $J$ is
finitely generated, $D$ is finitely presented as a $C$-module; since $J$ is
annihilated by itself, the preceding quotient argument also makes $J$ a
finitely presented $D$-module.

Let $K=(x_1,\ldots,x_r)\subset C$ be a finitely generated ideal and let
\[
    \varphi:C^r\twoheadrightarrow K
\]
be the evident surjection.  Reduction modulo $J$ gives
\[
    \overline\varphi:D^r\twoheadrightarrow\overline K.
\]
The ideal $\overline K\subset D$ is finitely generated and hence finitely
presented because $D$ is coherent.  Thus
\[
    L:=\ker(\overline\varphi)
\]
is finitely generated.  Choose generators
$\ell_1,\ldots,\ell_a$ of $L$, lifts
$\widetilde\ell_j\in C^r$, and generators
$j_1,\ldots,j_b$ of the $D$-module $J$.

We claim that $K\cap J$ is finitely generated over $D$.  Let $y\in K\cap J$
and choose $v\in C^r$ with $\varphi(v)=y$.  Since $y$ reduces to zero,
$\overline v\in L$, so
\[
    \overline v
    =
    \sum_j\overline c_j\ell_j.
\]
After lifting the coefficients,
\[
    v-\sum_jc_j\widetilde\ell_j\in J^r.
\]
Applying $\varphi$ shows that $y$ is a $D$-linear combination of the finite
collection
\[
    \varphi(\widetilde\ell_1),\ldots,
    \varphi(\widetilde\ell_a),
    \qquad
    j_qx_i
    \quad
    (1\leq q\leq b,\ 1\leq i\leq r).
\]
Every displayed element lies in $K\cap J$.  Hence $K\cap J$ is finitely
generated over $D$.  It is a finitely generated submodule of the finitely
presented $D$-module $J$, so coherence of $D$ makes $K\cap J$ finitely
presented over $D$.  The ideal $\overline K$ is finitely presented over $D$
as well.  Since $D$ is finitely presented over $C$, both modules are finitely
presented over $C$.  The exact sequence
\[
    0\longrightarrow K\cap J
    \longrightarrow K
    \longrightarrow\overline K
    \longrightarrow0
\]
therefore makes $K$ finitely presented over $C$.  Thus $C$ is coherent.

Return to $A_0\twoheadrightarrow B_0=A_0/I$ and choose $m$ with $I^m=0$.
For $1\leq n\leq m$ put
\[
    R_n:=A_0/I^n.
\]
Thus $R_1=B_0$ and $R_m=A_0$.  We prove inductively that every $R_n$ is
coherent.  Suppose $R_{n-1}$ is coherent.  The kernel of
\[
    R_n\twoheadrightarrow R_{n-1}
\]
is
\[
    J_n:=I^{n-1}/I^n,
\]
and $J_n^2=0$.  Because $I$ is finitely presented over $A_0$, base change
gives a finitely presented $R_{n-1}$-module
\[
    I\otimes_{A_0}R_{n-1}
    \simeq
    I/I^n.
\]
The submodule
\[
    J_n=I^{n-1}/I^n
    \subset
    I/I^n
\]
is finitely generated: $I$ is finitely generated, hence so is $I^{n-1}$.
Since $R_{n-1}$ is coherent, $J_n$ is finitely presented over $R_{n-1}$.
The square-zero coherence step therefore implies that $R_n$ is coherent.
Induction gives coherence of $R_m=A_0$.

Now every $I^k$ is a finitely generated ideal of the coherent ring $A_0$ and
is therefore finitely presented.  If $M$ is finitely presented over $A_0$,
then $M$ is a coherent $A_0$-module, so the finitely generated submodules
$I^kM$ and $I^{k+1}M$ are finitely presented.  Their quotient
\[
    I^kM/I^{k+1}M
\]
is finitely presented over $A_0$ and is annihilated by $I$.  Applying the
quotient finite-presentation argument from the first paragraph makes it
finitely presented over $B_0=A_0/I$.
\end{proof}

\begin{lemma}[Refining a square-zero coefficient filtration]
\label{lem:bnla-refine-square-zero-filtration}
Let $C\to D$ be a structured square-zero extension classified by a
$D$-module $M$, and suppose $M$ admits a finite filtration by $D$-submodules
whose successive cofibres are $M_1,\ldots,M_r$.  Then $C\to D$ admits a
finite factorization by structured square-zero extensions with coefficient
modules $M_i$.  The factorization is used only as a proof construction and is
functorial with respect to a fixed filtered coefficient.
\end{lemma}

\begin{proof}
A structured square-zero extension is the pullback of its classifying
derivation against the zero derivation.  For a cofiber sequence
$M'\to M\to M''$, compose the classifying derivation with
$M[1]\to M''[1]$ and form the corresponding square-zero extension $C''\to D$.
The original pullback then factors through $C''$, and the relative fibre
$C\to C''$ is the square-zero extension classified by the induced
$M'$-valued derivation.  Iterate along the finite filtration.
\end{proof}

\begin{proposition}[Finite square-zero generation of spectral Artinian extensions]
\label{prop:bnla-spectral-art-square-zero-generation}
An augmented arrow $A\to B$ is spectral Artinian if and only if it admits a
finite factorization
\[
    A=A_m\longrightarrow A_{m-1}\longrightarrow\cdots
    \longrightarrow A_0=B
\]
in which every $A_i\to A_{i-1}$ is a structured square-zero extension whose
coefficient is the restriction of a connective coherent $B$-module.
Moreover every spectral Artinian source $A$ is coherent and
\[
    A\to B
    \in
    \operatorname{CAlg}^{\mathrm{sp},\wedge\mathrm{aft}}_{R/B}.
\]
\end{proposition}

\begin{proof}
Suppose first that a finite square-zero factorization is given.  A structured
square-zero extension by a connective coherent module is almost finitely
presented by Lemma~\ref{lem:bnla-spectral-square-zero-finiteness}; it is
nilpotent on $\pi_0$ and has eventually coconnective fibre.  These properties
are preserved under finite composition.  Coherence of the sources follows
inductively from the square-zero fibre sequences: $\pi_0$ changes by a
finitely presented square-zero ideal and every higher homotopy group is a
finite extension of finitely presented modules.  Finally $B$ is almost
perfect over every source by finite transitivity of almost-perfect modules,
and the nilpotent augmentation ideal makes every source derived-complete.
Definition~\ref{def:bnla-complete-aft-augmented} therefore applies.

Conversely let $A\to B$ be spectral Artinian and put
\[
    I=\ker\bigl(\pi_0(A)\to\pi_0(B)\bigr),
    \qquad I^m=0.
\]
The almost-finite criterion makes $B$ almost perfect as an $A$-module.  Since
$B$ is coherent, Lemma~\ref{lem:bnla-augmentation-ideal-finite-presentation}
shows that
\[
    I=\ker\bigl(\pi_0(A)\to\pi_0(B)\bigr)
\]
is finitely presented as a $\pi_0(A)$-module.  Notice that the proof does not
identify $I$ with $\pi_0\operatorname{fib}(A\to B)$: the long exact homotopy
sequence gives only the canonical quotient
\[
    \pi_0\operatorname{fib}(A\to B)\twoheadrightarrow I,
\]
whose kernel is the image of $\pi_1(B)$.

Lemma~\ref{lem:bnla-coherent-nilpotent-ring} now implies that $\pi_0(A)$ is
coherent and that the nilpotent filtration and every finitely presented
$\pi_0(A)$-module have finitely presented $\pi_0(B)$-linear
associated-graded pieces.  The fibre
\[
    K:=\operatorname{fib}(A\to B)
\]
is nevertheless almost perfect as an $A$-module, since $A$ is perfect and
$B$ is almost perfect over $A$.  Since $K$ is eventually coconnective by the
Artinian hypothesis, choose $t$ so that it is $t$-coconnective.

We first separate the discrete nilpotent part.  Form the pullback
\[
    A^{(1)}:=H\pi_0(A)\times_{H\pi_0(B)}B.
\]
The map $A\to A^{(1)}$ is an isomorphism on $\pi_0$.  To construct the
factorization from $A^{(1)}$ to $B$ with the correct variance, put
\[
    D_k:=H\!\left(\pi_0(A)/I^k\right),
    \qquad
    1\leq k\leq m.
\]
Then
\[
    D_1\simeq H\pi_0(B),
    \qquad
    D_m\simeq H\pi_0(A),
\]
and every morphism
\[
    D_{k+1}\longrightarrow D_k
\]
is a discrete structured square-zero extension with coefficient
$H(I^k/I^{k+1})$.  The connective ring spectrum $B$ has its canonical
truncation morphism
\[
    B\longrightarrow H\pi_0(B)=D_1.
\]
Define
\[
    \widetilde D_k
    :=
    D_k\times_{D_1}B.
\]
Thus
\[
    \widetilde D_1\simeq B,
    \qquad
    \widetilde D_m
    \simeq
    H\pi_0(A)\times_{H\pi_0(B)}B
    =A^{(1)}.
\]
Moreover pullback pasting gives
\[
    \widetilde D_{k+1}
    \simeq
    D_{k+1}\times_{D_k}\widetilde D_k,
\]
so stability of structured square-zero extensions under affine base change
makes
\[
    \widetilde D_{k+1}\longrightarrow\widetilde D_k
\]
a structured square-zero extension.  Its coefficient is the restriction of
$H(I^k/I^{k+1})$.  Since
\[
    I(I^k/I^{k+1})=0,
\]
the coefficient action factors through
$\pi_0(A)/I\simeq\pi_0(B)$; hence this coefficient is the restriction of
the connective $B$-module $H(I^k/I^{k+1})$.  By
Lemma~\ref{lem:bnla-coherent-nilpotent-ring} these modules are coherent.
Consequently
\[
    A^{(1)}=\widetilde D_m
    \longrightarrow\cdots\longrightarrow
    \widetilde D_1=B
\]
is the required finite chain of structured square-zero extensions by
coefficients restricted from coherent $B$-modules.

It remains to factor $A\to A^{(1)}$.  Inductively construct objects
$A^{(i)}$ for $i\geq1$ together with maps
\[
    A\longrightarrow A^{(i+1)}
    \longrightarrow A^{(i)}
\]
such that
\[
    C_i:=\operatorname{cofib}(A\to A^{(i)})
\]
is $i$-connective and $t$-coconnective.  The case $i=1$ is the construction
above.  Suppose $A^{(i)}$ is constructed.  Since
$\pi_0(A)\simeq\pi_0(A^{(i)})$,
Lemma~\ref{lem:bnla-relative-spectral-hurewicz} gives
\[
    \pi_iL_{A^{(i)}/A}
    \simeq
    \pi_i(C_i).
\]
The module $C_i$ is almost perfect: it is obtained by finite extensions from
the almost-perfect fibre of $A\to B$ and the coherent coefficients already
used.  Hence
\[
    N_i:=\pi_i(C_i)
\]
is finitely presented.

Regard
\[
    M_i:=HN_i[i-1]
\]
as an $A^{(i)}$-module through
$A^{(i)}\to H\pi_0(A^{(i)})\simeq H\pi_0(A)$.  The Hurewicz identification
and the canonical projection to the first nonzero Postnikov group give a
classifying morphism
\[
    \eta_i:
    L_{A^{(i)}/A}
    \longrightarrow
    HN_i[i]
    =
    M_i[1].
\]
Let
\[
    A^{(i+1)}\longrightarrow A^{(i)}
\]
be the structured square-zero extension over $A$ classified by $\eta_i$.
Its coefficient is therefore $M_i=HN_i[i-1]$, in the square-zero convention
of Chapter~2.

The underlying square-zero fibre sequence is
\[
    M_i
    \longrightarrow
    A^{(i+1)}
    \longrightarrow
    A^{(i)}.
\]
Apply the octahedral axiom to
\[
    A\longrightarrow A^{(i+1)}
    \longrightarrow A^{(i)}.
\]
It gives a cofiber sequence
\[
    \operatorname{cofib}(A\to A^{(i+1)})
    \longrightarrow
    C_i
    \longrightarrow
    M_i[1].
\]
By construction of $\eta_i$ and
Lemma~\ref{lem:bnla-relative-spectral-hurewicz}, the last arrow is the
canonical Postnikov projection
\[
    C_i\longrightarrow H\pi_i(C_i)[i].
\]
Consequently
\[
    \boxed{
    \operatorname{cofib}(A\to A^{(i+1)})
    \simeq
    \tau_{\geq i+1}C_i.}
\]
Thus the connectivity increases by one while the fixed coconnective bound
$t$ remains.

It remains to refine the coefficient to coefficients descending from $B$.
The finitely presented $\pi_0(A)$-module $N_i$ has the finite $I$-adic
filtration
\[
    N_i\supset IN_i\supset\cdots\supset I^mN_i=0.
\]
Its successive quotients
\[
    I^kN_i/I^{k+1}N_i
\]
are finitely presented $\pi_0(B)$-modules by
Lemma~\ref{lem:bnla-coherent-nilpotent-ring}.  Hence the actual square-zero
coefficient
\[
    M_i=HN_i[i-1]
\]
has a finite filtration by $A^{(i)}$-submodules whose successive cofibres are
\[
    H\!\left(I^kN_i/I^{k+1}N_i\right)[i-1].
\]
These are connective coherent $B$-modules because $i\geq1$.
Lemma~\ref{lem:bnla-refine-square-zero-filtration} therefore refines
$A^{(i+1)}\to A^{(i)}$ into finitely many structured square-zero extensions
with coefficients obtained by restriction from coherent $B$-modules.

After $i>t$, the cofiber vanishes, so $A\simeq A^{(i)}$.  Concatenating the
finite refinements with the discrete nilpotent factorization gives the
desired finite square-zero chain.  No step of that chain is retained as
structure on $A\to B$.
\end{proof}

\begin{corollary}[Artinian arrows are complete almost-finite augmented]
\label{cor:bnla-art-complete-aft}
There is a fully faithful inclusion
\[
    \operatorname{Art}^{\mathrm{sp}}_{R/B}
    \hookrightarrow
    \operatorname{CAlg}^{\mathrm{sp},\wedge\mathrm{aft}}_{R/B}.
\]
\end{corollary}

\begin{proof}
This is the final assertion of
Proposition~\ref{prop:bnla-spectral-art-square-zero-generation}.
\end{proof}

\subsection{Schlessinger localization on spectral Artinian tests}

Put
\[
    \mathcal P^{\mathrm{Art,sp}}_{B/R}
    :=
    \operatorname{Fun}
    \left(\operatorname{Art}^{\mathrm{sp}}_{R/B},\mathcal S\right).
\]

\begin{definition}[Artinian Schlessinger generators]
\label{def:bnla-schlessinger-generators}
Let $\mathcal W^{\mathrm{Sch,Art}}_{B/R}$ be the set of Yoneda morphisms
consisting of:
\begin{enumerate}
    \item the morphism whose locality forces the value at $B\to B$ to be
    contractible;
    \item for every structured square-zero pullback
    \[
    \begin{tikzcd}
    A_\eta \arrow[r] \arrow[d] & B \arrow[d]\\
    A \arrow[r,"\eta"'] & B\oplus M[1]
    \end{tikzcd}
    \]
    with $A\to B$ Artinian and $M$ connective coherent, the Yoneda morphism
    whose locality condition is
    \[
        F(A_\eta)
        \xrightarrow{\sim}
        F(A)\times_{F(B\oplus M[1])}F(B).
    \]
\end{enumerate}
\end{definition}

\begin{definition}[Spectral formal moduli problems]
\label{def:bnla-spectral-fmp}
Define
\[
    \operatorname{FMP}^{\mathrm{sp}}_{B/R}
    \subset
    \mathcal P^{\mathrm{Art,sp}}_{B/R}
\]
to be the full subcategory of
$\mathcal W^{\mathrm{Sch,Art}}_{B/R}$-local objects.
\end{definition}

Thus the Schlessinger equations are fixed-point properties of a localization,
not auxiliary structure carried by a formal-moduli object.

\begin{proposition}[Artinian Schlessinger reflector]
\label{prop:bnla-artinian-schlessinger-reflector}
The inclusion
\[
    \operatorname{FMP}^{\mathrm{sp}}_{B/R}
    \hookrightarrow
    \mathcal P^{\mathrm{Art,sp}}_{B/R}
\]
admits an accessible left adjoint
\[
    L^{\mathrm{Sch,Art}}_{B/R}.
\]
Its local objects are exactly the functors with contractible value at $B$ and
square-zero excision.
\end{proposition}

\begin{proof}
The Artinian category is essentially small in the fixed universe, so the
Yoneda generators form a set.  The functor category is presentable.
Accessible localization at the strongly saturated class generated by that set
therefore exists.  Yoneda identifies locality with the two displayed
conditions.
\end{proof}

\subsection{The spectral Koszul-dual functor}

Comonadicity equips every cotangent fibre with its canonical
$\operatorname{coLie}^{\pi,E_\infty}_{B/R}$-coalgebra structure.

\begin{definition}[Spectral Koszul-dual functor]
\label{def:bnla-spectral-D-functor}
Define
\[
    D^{\mathrm{sp}}_{B/R}:
    \operatorname{CAlg}^{\mathrm{sp},\wedge\mathrm{aft}}_{R/B}
    \longrightarrow
    \left(
        \operatorname{LieAlgd}^{\pi,E_\infty}_{B/R}
    \right)^{\mathrm{op}}
\]
by
\[
\boxed{
    D^{\mathrm{sp}}_{B/R}(A)
    :=
    \operatorname{cot}_{B/R}(A)^\vee,}
\]
with spectral partition-Lie algebroid structure obtained by dualizing the
canonical coLie-coalgebra structure.
\end{definition}

\begin{theorem}[Full faithfulness of spectral Koszul duality on complete arrows]
\label{thm:bnla-spectral-D-fully-faithful}
The functor $D^{\mathrm{sp}}_{B/R}$ is fully faithful.  Its essential image has
dually almost-perfect underlying anchored source.
\end{theorem}

\begin{proof}
Theorem~\ref{thm:bnla-spectral-cot-sqz-comonadic} identifies the source with
the category of coalgebras for the relative coLie comonad on
$\operatorname{Der}^{\mathrm{aperf}}_{B/R,\geq0}$.  Continuous duality is fully
faithful on connective almost-perfect modules and, by construction of the
relative monad, identifies the opposite coLie comonad with the restriction of
$\operatorname{Lie}^{\pi,E_\infty}_{B/R}$.  Passing from coalgebras to the
opposite dual monad algebras gives the result.
\end{proof}

\begin{proposition}[Square-zero extensions give free spectral algebroids]
\label{prop:bnla-square-zero-free-spectral-algebroid}
For a structured square-zero extension $B_\alpha\to B$ classified by
\[
    \alpha:L_{B/R}[-1]\longrightarrow M,
\]
there is a canonical equivalence
\[
\boxed{
    D^{\mathrm{sp}}_{B/R}(B_\alpha)
    \simeq
    \operatorname{Free}^{\pi,E_\infty}_{B/R}
    \left(
        M^\vee\xrightarrow{\alpha^\vee}
        \mathbb T^{\mathrm{pc}}_{B/R}[1]
    \right).}
\]
\end{proposition}

\begin{proof}
The square-zero object lies in the image of the right adjoint
$\operatorname{sqz}$.  Under a comonadic adjunction, applying the left
adjoint to an object in the image of the right adjoint gives the cofree
coalgebra for the induced comonad.  Continuous duality turns that cofree
coalgebra into the free algebra for the dual monad.
\end{proof}

\subsection{Filtered duality and spectral Artinian excision}

\begin{lemma}[Continuous dual of a complete almost-perfect filtration]
\label{lem:bnla-dual-complete-filtration}
Let
\[
    F^0M\supset F^1M\supset F^2M\supset\cdots
\]
be a complete descending filtration of an almost-perfect $B$-module whose
filtration pieces become arbitrarily connective.  Then
\[
\boxed{
    \operatorname*{colim}_n
    \left(F^0M/F^nM\right)^\vee
    \xrightarrow{\sim}
    \left(F^0M\right)^\vee}
\]
in $\operatorname{QC}^{\vee}_B$.
\end{lemma}

\begin{proof}
Continuous duals of coherent modules compactly generate
$\operatorname{QC}^{\vee}_B$.  Test the map against a compact generator
$N^\vee$, equivalently by $\operatorname{Map}_B(-,N)$ with $N$ coherent and
eventually coconnective.  Increasing connectivity of $F^nM$ implies
$\operatorname{Map}_B(F^nM,N)\simeq*$ for large $n$.  The fibre sequence
$F^nM\to F^0M\to F^0M/F^nM$ then gives the required equivalence on maps to
every compact generator.
\end{proof}

\begin{theorem}[Spectral square-zero excision for $D$]
\label{thm:bnla-spectral-D-excision}
Let $A\to B$ be spectral Artinian, let $M$ be a connective coherent
$B$-module, and let
\[
    \eta:A\longrightarrow B\oplus M[1]
\]
be a morphism over $B$.  Put
\[
    A_\eta:=A\times_{B\oplus M[1]}B.
\]
Then
\[
\boxed{
    D(A)
    \amalg_{D(B\oplus M[1])}
    D(B)
    \xrightarrow{\sim}
    D(A_\eta)}
\]
is an equivalence of spectral partition-Lie algebroids.
\end{theorem}

\begin{proof}
We use a proof predicate only during the induction.  For a complete
almost-finite arrow $C\to B$, let $P(C)$ denote the assertion that, for every
connective almost-perfect $B$-module $N$,
\[
    D(C)\amalg_{D(B)}D(B\oplus N)
    \longrightarrow
    D(C\times_B(B\oplus N))
\]
is an equivalence.  The symbol $P(C)$ is not a structure or admissibility
condition.

If $C=B_\alpha$ is square-zero, all terms are free algebroids by
Proposition~\ref{prop:bnla-square-zero-free-spectral-algebroid}.  The generator
for $C\times_B(B\oplus N)$ is the direct sum of the two anchored generators,
so $P(C)$ follows from preservation of pushouts by the free-algebroid left
adjoint.

Assume $P(A)$ and consider the displayed square-zero extension $A_\eta\to A$.
Take the canonical relative cobar resolution
\[
    A\longrightarrow A^\bullet\longrightarrow B\oplus M[1]
\]
from Corollary~\ref{cor:bnla-spectral-relative-square-zero-resolution}.  Its
positive terms are square-zero extensions of $B$, and
$\operatorname{cot}(A)\to\operatorname{cot}(A^\bullet)$ is split.

For a connective almost-perfect $B$-module $N$, let
\[
    B\oplus N\longrightarrow B\oplus M[1]
\]
be the composite of the augmentation $B\oplus N\to B$ with the zero section
$B\to B\oplus M[1]$, and define
\[
\boxed{
    A_{\eta,N}
    :=
    A\times_{B\oplus M[1]}(B\oplus N)
    \simeq
    A_\eta\times_B(B\oplus N).}
\]
Termwise define
\[
\boxed{
    A_{\eta,N}^n
    :=
    A^n\times_{B\oplus M[1]}(B\oplus N),
    \qquad
    A_{\eta,N}^\bullet
    :=
    A^\bullet\times_{B\oplus M[1]}(B\oplus N).}
\]
In particular
\[
    A_\eta^\bullet
    :=
    A^\bullet\times_{B\oplus M[1]}B
\]
is now defined, and pullback pasting gives
\[
    A_{\eta,N}^\bullet
    \simeq
    A_\eta^\bullet\times_B(B\oplus N).
\]
Thus there is a termwise Cartesian diagram
\[
\begin{tikzcd}
A_{\eta,N} \arrow[r] \arrow[d] &
A_{\eta,N}^\bullet \arrow[r] \arrow[d] &
B\oplus N \arrow[d]\\
A \arrow[r] & A^\bullet \arrow[r] & B\oplus M[1].
\end{tikzcd}
\]

After applying $D$, the augmented simplicial object
$D(A^\bullet)\to D(A)$ is a geometric-realization diagram: continuous dual of
the split cotangent cosimplicial object is split simplicial, and geometric
realizations of algebroids are created on underlying anchored pro-coherent
objects because the monad preserves sifted colimits.

Levelwise, the maps $A^n\to B\oplus M[1]$ arise from morphisms of relative
derivations.  Each positive $A^n$ is in the image of
$\operatorname{sqz}$, so Proposition~\ref{prop:bnla-square-zero-free-spectral-algebroid}
identifies its $D$-image with a free algebroid.  The termwise base changes
above are induced by the corresponding maps of anchored generators; hence the
levelwise $D$-squares are pushouts by the free-algebroid adjunction.

It remains to identify the realization of the base-changed cobar row.  Give
$B\oplus M[1]$ the two-step filtration
\[
    F^1=M[1],\qquad F^2=0,
\]
and put $A$, $A^\bullet$, and $B\oplus N$ in filtration weight zero.  Forming
the displayed homotopy pullbacks produces finite, hence complete, compatible
filtrations on $A_{\eta,N}$ and $A_{\eta,N}^\bullet$ and on their cotangent
fibres.  On associated graded the $M[1]$-valued classifying component has
positive filtration weight and therefore vanishes as a map out of the
weight-zero source.  The associated-graded square-zero geometry is the split
one: the weight-one fibre contributes $M$, while the independent split
coefficient $N$ remains in weight zero.  Thus the associated-graded
cotangent comparison is the direct-sum comparison governed on the augmented
term by $P(A)$ and termwise by the already proved square-zero predicates
$P(A^n)$.

The filtered square-zero finiteness theorem makes the cotangent objects
almost perfect; the filtrations are complete, and their positive pieces
become increasingly connective after passage through the canonical cobar
resolution.  Lemma~\ref{lem:bnla-dual-complete-filtration} therefore computes
continuous duality from the finite filtration quotients.  Since the
associated-graded augmented simplicial comparison is a colimit diagram, the
complete filtered comparison is a colimit diagram as well.  Equivalently,
\[
\boxed{
    \left|D(A_{\eta,N}^\bullet)\right|
    \xrightarrow{\sim}
    D(A_{\eta,N}).}
\]
Taking $N=0$ gives
\[
    D(A)
    \amalg_{D(B\oplus M[1])}
    D(B)
    \xrightarrow{\sim}
    D(A_\eta),
\]
while arbitrary $N$ proves $P(A_\eta)$.

By Proposition~\ref{prop:bnla-spectral-art-square-zero-generation}, every
spectral Artinian extension is obtained by finitely many such structured
square-zero stages.  Induction proves the theorem.
\end{proof}

\begin{corollary}[Finite spectral Artinian excision]
\label{cor:bnla-spectral-D-finite-excision}
A pullback square in $\operatorname{Art}^{\mathrm{sp}}_{R/B}$ with one
square-zero leg is carried by $D$ to a pushout square.  The same holds for a
finite composite of square-zero legs.
\end{corollary}

\begin{proof}
The one-stage assertion is
Theorem~\ref{thm:bnla-spectral-D-excision}; finite composites follow by
pasting pushouts.  A factorization may be chosen inside this proof but is not
retained by the conclusion.
\end{proof}

\subsection{Cellular deformation recognition}

We use the standard cellular restricted-Yoneda theorem for deformation
contexts in the following form.

\begin{theorem}[Cellular deformation theorem]
\label{thm:bnla-cellular-restricted-yoneda}
Let $\mathcal B$ be a presentable $\infty$-category with a set of adjunctions
\[
    L_\lambda:\operatorname{Sp}
    \rightleftarrows
    \mathcal B:R_\lambda
\]
such that every $R_\lambda$ preserves sifted colimits and the family is
jointly conservative.  Let $\mathcal B_0$ be the full finite cellular
subcategory generated from the initial object by pushouts of the negative
sphere cells $L_\lambda(\mathbb S[n])$, $n<0$.  Then restricted Yoneda
\[
    \mathcal B
    \longrightarrow
    \operatorname{Fun}(\mathcal B_0^{\mathrm{op}},\mathcal S)
\]
is fully faithful, with essential image the functors taking the initial object
to a point and the generating cellular pushouts to pullbacks.
\end{theorem}

\begin{proof}
This is the cellular deformation theorem used in
\cite{BrantnerMagidsonNuiten2025}.  Briefly, left Kan extension along
$\mathcal B_0\hookrightarrow\mathcal B$ is left adjoint to restricted Yoneda.
The right adjoints $R_\lambda$ are jointly conservative, so the unit is
detected on the sphere generators.  Preservation of sifted colimits permits
induction through finite cell attachments and then passage to the sifted
completion.  The same induction identifies the essential image with the
stated pullback conditions.
\end{proof}

\begin{lemma}[Anchored objects as negative-cell attachments]
\label{lem:bnla-anchored-cell-attachment}
Let $\mathcal C$ be a stable $\infty$-category and let $T\in\mathcal C$.
For $E\in\mathcal C$ write
\[
    Z(E[-1]):=(E[-1]\xrightarrow{0}T),
    \qquad
    \mathbf 0_T:=(0\to T)
\]
as objects of $\mathcal C_{/T}$.  There is a canonical equivalence
\[
\boxed{
    \operatorname{Map}_{\mathcal C_{/T}}
    \left(Z(E[-1]),\mathbf 0_T\right)
    \simeq
    \operatorname{Map}_{\mathcal C}(E,T).}
\]
If $\rho:E\to T$ corresponds to
$\partial_\rho:Z(E[-1])\to\mathbf 0_T$ and
$\partial_0$ denotes the map corresponding to the zero morphism, then the
canonical square
\[
\begin{tikzcd}
Z(E[-1]) \arrow[r,"\partial_\rho"] \arrow[d,"\partial_0"'] &
\mathbf 0_T \arrow[d]\\
\mathbf 0_T \arrow[r] &
(E\xrightarrow{\rho}T)
\end{tikzcd}
\]
is a pushout in $\mathcal C_{/T}$.
\end{lemma}

\begin{proof}
The mapping space in the slice is the homotopy fibre over the zero map of
\[
    \operatorname{Map}_{\mathcal C}(E[-1],0)
    \longrightarrow
    \operatorname{Map}_{\mathcal C}(E[-1],T).
\]
The source is contractible, so this fibre is the loop space at zero of
$\operatorname{Map}_{\mathcal C}(E[-1],T)$.  Stability gives
\[
    \Omega\operatorname{Map}_{\mathcal C}(E[-1],T)
    \simeq
    \operatorname{Map}_{\mathcal C}(E,T),
\]
which proves the first assertion.

Colimits in the slice are created in $\mathcal C$.  The underlying pushout of
$0\leftarrow E[-1]\to0$ is canonically $E$.  Its map to $T$ is obtained by
gluing the two specified nullhomotopies of the zero map $E[-1]\to T$; under
the preceding loop-space identification their difference is precisely
$\rho$.  Hence the induced object of the slice is
$(E\xrightarrow{\rho}T)$, proving the pushout assertion.
\end{proof}

\begin{lemma}[The spectral algebroid category satisfies the cellular hypotheses]
\label{lem:bnla-spectral-cellular-hypotheses}
For a connective coherent $B$-module $M$, define
\[
    R_M(\rho)
    :=
    \operatorname{map}_{\operatorname{QC}^{\vee}_B}
    \left(
        M^\vee,\operatorname{fib}U(\rho)
    \right)
\]
for
$\rho\in\operatorname{LieAlgd}^{\pi,E_\infty}_{B/R}$.  Then $R_M$ is a
right adjoint preserving sifted colimits.  As $M$ varies, the $R_M$ are jointly
conservative.  If $L_M$ denotes its left adjoint, then for every $n<0$
\[
\boxed{
    L_M(\mathbb S[n])
    \simeq
    \operatorname{Free}^{\pi,E_\infty}_{B/R}
    \left(M^\vee[n]\xrightarrow{0}
    \mathbb T^{\mathrm{pc}}_{B/R}[1]\right).}
\]
Equivalently, these negative sphere cells are the $D$-images of the split
square-zero extensions with the matching coherent shifts.
\end{lemma}

\begin{proof}
Anchor fibre is a right adjoint by
Theorem~\ref{thm:bnla-zero-anchor-spectral-algebroids}.  Mapping spectra out
of $M^\vee$ are right adjoints as well, so $R_M$ is a right adjoint.  The same
theorem proves that anchor fibre preserves sifted colimits, and $M^\vee$ is
compact in $\operatorname{QC}^{\vee}_B$; hence $R_M$ preserves sifted
colimits.

The objects $M^\vee$ for coherent $M$ compactly generate the pro-coherent
category.  Therefore the family detects equivalences of anchor fibres.  Since
the anchor target is fixed, an equivalence on fibres is an equivalence on the
source anchored objects, and the monadic forgetful functor detects
algebroid equivalences.  Thus the family is jointly conservative.

The left adjoint to $R_M$ first applies the free one-object
partition-Lie-algebra functor to the generator $M^\vee\otimes(-)$, includes
the result as a zero-anchor algebroid, and then uses the zero-anchor inclusion.
Consequently
\[
    L_M(\mathbb S[n])
    \simeq
    \operatorname{Free}^{\pi,E_\infty}_{B/R}
    \left(M^\vee[n]\xrightarrow{0}
    \mathbb T^{\mathrm{pc}}_{B/R}[1]\right).
\]
For $n<0$, put
\[
    N:=M[-n],
    \qquad
    M':=M[-n-1].
\]
Both modules are connective coherent, and $N\simeq M'[1]$.  Continuous
duality reverses suspension, so
\[
    N^\vee
    \simeq
    M^\vee[n].
\]
Proposition~\ref{prop:bnla-square-zero-free-spectral-algebroid} therefore
identifies the negative sphere cell literally as
\[
\boxed{
    L_M(\mathbb S[n])
    \simeq
    D^{\mathrm{sp}}_{B/R}
    \left(B\oplus M'[1]\longrightarrow B\right).}
\]
Thus every negative sphere cell is the $D$-image of one of the split
square-zero objects occurring in the stated Schlessinger generators, with no
undefined ``positive shift'' convention.

Finally, Lemma~\ref{lem:bnla-anchored-cell-attachment} explains how arbitrary
anchors enter the same cellular theory: the cell itself has zero anchor, while
the attaching morphism records the nonzero anchored class.  No additional
family of cellular generators is required.
\end{proof}

\begin{theorem}[Affine spectral formal-moduli/partition-Lie equivalence]
\label{thm:bnla-spectral-affine-recognition}
Let $R\to B$ be a morphism of connective $E_\infty$-rings with $B$ coherent.
There is a canonical equivalence
\[
\boxed{
    \operatorname{Diff}^{\pi,E_\infty}_{B/R}:
    \operatorname{FMP}^{\mathrm{sp}}_{B/R}
    \xrightarrow{\sim}
    \operatorname{LieAlgd}^{\pi,E_\infty}_{B/R}.}
\]
Its inverse is
\[
\boxed{
    \operatorname{Int}^{\pi,E_\infty}_{B/R}(\mathfrak a)(A)
    :=
    \operatorname{Map}_{\operatorname{LieAlgd}^{\pi,E_\infty}_{B/R}}
    \left(D^{\mathrm{sp}}_{B/R}(A),\mathfrak a\right).}
\]
For $F\in\operatorname{FMP}^{\mathrm{sp}}_{B/R}$, the equivalence also
canonically extends $F$ from Artinian arrows to the complete almost-finite
sector by
\[
\boxed{
    F^{\wedge\mathrm{aft}}(A)
    :=
    \operatorname{Map}_{\operatorname{LieAlgd}^{\pi,E_\infty}_{B/R}}
    \left(
        D^{\mathrm{sp}}_{B/R}(A),
        \operatorname{Diff}^{\pi,E_\infty}_{B/R}(F)
    \right),
    \qquad
    A\in\operatorname{CAlg}^{\mathrm{sp},\wedge\mathrm{aft}}_{R/B}.}
\]
Its restriction to $\operatorname{Art}^{\mathrm{sp}}_{R/B}$ is canonically
$F$.

Write
\[
\boxed{
    \mathbb T^{\mathrm{pc}}_{B/F}[1]
    \longrightarrow
    \mathbb T^{\mathrm{pc}}_{B/R}[1]}
\]
for the underlying anchored pro-coherent object of
$\operatorname{Diff}^{\pi,E_\infty}_{B/R}(F)$.  This tangent-fibre notation
is characterized intrinsically by structured square-zero tests: for every
\[
    \alpha:L_{B/R}[-1]\longrightarrow M
\]
with $M$ connective and almost perfect, there is a natural equivalence
\[
\boxed{
    F^{\wedge\mathrm{aft}}(B_\alpha)
    \simeq
    \operatorname{Map}_{\operatorname{Anch}^{\mathrm{pc}}_{B/R}}
    \left(
        M^\vee\xrightarrow{\alpha^\vee}
        \mathbb T^{\mathrm{pc}}_{B/R}[1],
        \mathbb T^{\mathrm{pc}}_{B/F}[1]
        \longrightarrow
        \mathbb T^{\mathrm{pc}}_{B/R}[1]
    \right).}
\]
When $B_\alpha\to B$ is Artinian, the canonical equivalence
$F^{\wedge\mathrm{aft}}(B_\alpha)\simeq F(B_\alpha)$ recovers the original
Artinian evaluation.
\end{theorem}

\begin{proof}
Apply Theorem~\ref{thm:bnla-cellular-restricted-yoneda} to
$\mathcal B=\operatorname{LieAlgd}^{\pi,E_\infty}_{B/R}$ and the adjunctions
of Lemma~\ref{lem:bnla-spectral-cellular-hypotheses}.  We first identify its
finite cellular subcategory.

The identity augmentation $B\to B$ satisfies
\[
    D(B)\simeq\mathbf 0_{B/R},
\]
the initial spectral partition-Lie algebroid.  Let $A\to B$ be Artinian, let
$M$ be a connective coherent $B$-module, and let
\[
    \eta:A\longrightarrow B\oplus M[1]
\]
be a morphism over $B$.  Put
\[
    A_\eta:=A\times_{B\oplus M[1]}B.
\]
By Corollary~\ref{cor:bnla-spectral-D-finite-excision},
\[
\boxed{
    D(A_\eta)
    \simeq
    D(A)\amalg_{D(B\oplus M[1])}D(B).}
\]
Lemma~\ref{lem:bnla-spectral-cellular-hypotheses}, with sphere degree
$n=-1$, identifies $D(B\oplus M[1])$ with the corresponding negative sphere
cell.  Thus a one-stage structured square-zero pullback is exactly a cellular
attachment on the algebroid side.  Proposition~\ref{prop:bnla-spectral-art-square-zero-generation}
therefore shows by induction that $D(A)$ is finite cellular for every
Artinian $A\to B$.

Conversely, suppose that $D(A)$ is already the $D$-image of an Artinian arrow
and attach a negative sphere cell $L_M(\mathbb S[n])$ with $n<0$.  Put
\[
    M':=M[-n-1].
\]
Then $M'$ is connective coherent and
Lemma~\ref{lem:bnla-spectral-cellular-hypotheses} gives
\[
    L_M(\mathbb S[n])
    \simeq
    D(B\oplus M'[1]).
\]
An attaching map
\[
    D(B\oplus M'[1])\longrightarrow D(A)
\]
corresponds, by the full faithfulness and contravariance of
$D$ from Theorem~\ref{thm:bnla-spectral-D-fully-faithful}, to a morphism
\[
    A\longrightarrow B\oplus M'[1]
\]
over $B$.  Forming
\[
    A_\eta=A\times_{B\oplus M'[1]}B
\]
gives another spectral Artinian arrow, and the excision equivalence above
identifies its $D$-image with the prescribed cellular pushout.  Induction on
the number of cells proves that every object of $\mathcal B_0$ is the
$D$-image of an Artinian arrow.  Full faithfulness of $D$ identifies the
morphisms as well.  Hence
\[
\boxed{
    \left(\operatorname{Art}^{\mathrm{sp}}_{R/B}\right)^{\mathrm{op}}
    \xrightarrow{\sim}
    \mathcal B_0.}
\]

The essential-image equations of cellular restricted Yoneda are exactly the
locality equations imposed in
Definition~\ref{def:bnla-schlessinger-generators}.  This proves the
equivalence and the formal-integration formula.

The formal-integration formula already makes sense on the entire complete
almost-finite sector because $D^{\mathrm{sp}}_{B/R}$ was constructed there.
Using the equivalence just proved to identify the controller of $F$ therefore
defines $F^{\wedge\mathrm{aft}}$, and on Artinian arrows the defining formula
is exactly $\operatorname{Int}(\operatorname{Diff}(F))\simeq F$.  This proves
the extension statement.

For the tangent statement, evaluate the extended formal integration on the
square-zero extension classified by
$\alpha:L_{B/R}[-1]\to M$.  By
Proposition~\ref{prop:bnla-square-zero-free-spectral-algebroid}, this is the
mapping space from the free algebroid on
\[
    M^\vee\longrightarrow\mathbb T^{\mathrm{pc}}_{B/R}[1].
\]
The free-forgetful adjunction identifies these values with the mapping space
in the anchored pro-coherent category from that anchored generator to the
underlying anchored object of
$\operatorname{Diff}^{\pi,E_\infty}_{B/R}(F)$.  This gives the displayed
characterization and defines the tangent-fibre notation without evaluating an
Artinian functor outside its domain.
\end{proof}

The coherent-base spectral recognition theorem above is therefore a result of
the constructed spectral cotangent--square-zero adjunction, its comonadic
cobar resolution, spectral adic completion, continuous pro-coherent duality,
the coherent one-object spectral partition-Lie monad, Artinian excision, and
cellular restricted Yoneda.  No recognition datum or admissibility package is
retained.
\section{Etale transport and global spectral formal moduli}
\label{sec:bnla-spectral-global-formal-moduli}

The affine recognition theorem is already canonical.  The remaining issue is
globalization.  The order of construction matters.  We do not independently
Schlessinger-reflect the quotient-lifting functor at every affine centre and
then ask those reflected functors to glue.  Instead, we first impose the
global spectral deformation equations---Postnikov convergence and nilpotent
excision, refining the Artinian equations---on the nonlinear quotient target
itself.  Affine formal leaves are then extracted from that single global target.  Their
etale cocartesianity will be proved from tangent base change and the affine
recognition equivalence.

\subsection{Etale transport of spectral Artinian extensions}

Fix first a coherent connective $E_\infty$-ring $R$ and let
$B\to B'$ be spectral etale between coherent connective $R$-algebras.

\begin{theorem}[Etale lifting of spectral Artinian extensions]
\label{thm:bnla-spectral-art-etale-lift}
For every $A\to B$ in $\operatorname{Art}^{\mathrm{sp}}_{R/B}$ there is a
contractibly unique commutative square
\[
\begin{tikzcd}
A \arrow[r,"\widetilde f"] \arrow[d] & A' \arrow[d]\\
B \arrow[r,"f"'] & B'
\end{tikzcd}
\]
with $\widetilde f$ spectral etale.  The induced arrow $A'\to B'$ is spectral
Artinian and the square is a pushout:
\[
\boxed{
    A'\otimes_A B\xrightarrow{\sim}B'.}
\]
These lifts are coherent for identity and composition and give cocartesian
transport
\[
    f_{\mathrm{Art},!}:
    \operatorname{Art}^{\mathrm{sp}}_{R/B}
    \longrightarrow
    \operatorname{Art}^{\mathrm{sp}}_{R/B'}.
\]
\end{theorem}

\begin{proof}
By Proposition~\ref{prop:bnla-spectral-art-square-zero-generation}, the
opposite affine map
$\operatorname{Spec}(B)\to\operatorname{Spec}(A)$ is a finite spectral
infinitesimal thickening in the sense of Chapter~9.  Finite-infinitesimal
etale-site rigidity identifies the etale categories of its source and target.
Therefore the etale object
$\operatorname{Spec}(B')\to\operatorname{Spec}(B)$ has a contractibly unique
etale lift
$\operatorname{Spec}(A')\to\operatorname{Spec}(A)$.  The affine form of that
lifting theorem gives the displayed pushout equivalence.

The intrinsic Artinian conditions are preserved.  Nilpotence and eventual
coconnectivity are preserved by flat etale base change, and almost finite
presentation is stable under base change.  Equivalently, one may base change
the finite coherent square-zero chain whose existence is proved in
Proposition~\ref{prop:bnla-spectral-art-square-zero-generation}; every stage
remains a structured square-zero extension with coherent coefficient.
Contractible uniqueness gives the cocartesian universal property and all
composition coherence.
\end{proof}

\begin{proposition}[Etale transport preserves Schlessinger generators]
\label{prop:bnla-etale-art-schlessinger}
The functor $f_{\mathrm{Art},!}$ preserves the identity Artinian object and
structured square-zero pullback diagrams.  Consequently restriction along
$f_{\mathrm{Art},!}$ preserves Schlessinger-local objects.
\end{proposition}

\begin{proof}
The identity lift is the identity.  For a structured square-zero pullback,
apply the unique etale lift to all vertices.  The pullback of the three lifted
input vertices is another etale lift of the original pullback vertex.
Contractible uniqueness identifies it with the selected lift.  The
Schlessinger assertion follows from Yoneda.
\end{proof}

\begin{proposition}[Affine formal-moduli centre change]
\label{prop:bnla-spectral-fmp-center-change}
Left Kan extension along $f_{\mathrm{Art},!}$ followed by Artinian
Schlessinger localization defines
\[
    f_{\mathrm{FMP},!}:
    \operatorname{FMP}^{\mathrm{sp}}_{B/R}
    \longrightarrow
    \operatorname{FMP}^{\mathrm{sp}}_{B'/R}.
\]
These functors are coherent for identities and composition.
\end{proposition}

\begin{proof}
The ambient functor categories are presentable, so left Kan extension exists.
By Proposition~\ref{prop:bnla-etale-art-schlessinger}, the right-adjoint
restriction preserves local objects.  Hence the induced left adjoint on local
objects is left Kan extension followed by
$L^{\mathrm{Sch,Art}}_{B'/R}$.  Identity and composition coherence are the
canonical coherences of Kan extension and localization.
\end{proof}

\subsection{Etale base change of the relative coLie and partition-Lie monads}

The base-change proof is carried out on relative derivations.  It does not
require an etale lift of an arbitrary complete almost-finite augmentation.

\begin{theorem}[Etale base change of the spectral algebroid monad]
\label{thm:bnla-spectral-algebroid-etale-base-change}
For spectral etale $f:B\to B'$ there is a canonical equivalence of monad-valued
functors
\[
\boxed{
    f^*\operatorname{Lie}^{\pi,E_\infty}_{B/R}
    \xrightarrow{\sim}
    \operatorname{Lie}^{\pi,E_\infty}_{B'/R}f^*.}
\]
Consequently extension of scalars lifts to
\[
    f^*:
    \operatorname{LieAlgd}^{\pi,E_\infty}_{B/R}
    \longrightarrow
    \operatorname{LieAlgd}^{\pi,E_\infty}_{B'/R},
\]
coherently in $f$.
\end{theorem}

\begin{proof}
Let
\[
    \alpha:L_{B/R}[-1]\longrightarrow M
\]
be a connective almost-perfect relative derivation.  Etale cotangent base
change gives
\[
    L_{B'/R}[-1]
    \simeq
    B'\otimes_BL_{B/R}[-1].
\]
Put
\[
    M':=B'\otimes_BM
\]
and let
\[
    \alpha':
    L_{B'/R}[-1]\longrightarrow M'
\]
be the transported derivation.

Form the structured square-zero extension
\[
    B_\alpha\longrightarrow B.
\]
Because $M$ is connective, the opposite affine morphism
$\operatorname{Spec}(B)\to\operatorname{Spec}(B_\alpha)$ is a one-stage
finite spectral infinitesimal thickening in the sense of Chapter~9.  Spectral
etale rigidity for finite infinitesimal thickenings therefore lifts
$\operatorname{Spec}(B')\to\operatorname{Spec}(B)$ through a contractible
space to a spectral etale morphism
\[
    \operatorname{Spec}(A')
    \longrightarrow
    \operatorname{Spec}(B_\alpha),
\]
with prescribed reduction $\operatorname{Spec}(B')$.  On coordinate rings
this gives a pushout square
\[
\begin{tikzcd}
B_\alpha \arrow[r] \arrow[d] &
A' \arrow[d]\\
B \arrow[r,"f"'] &
B'.
\end{tikzcd}
\]
The augmentation $A'\to B'$ is the affine base change of the structured
square-zero extension $B_\alpha\to B$.  Chapter~9's base-change theorem for
structured square-zero extensions therefore identifies it as a structured
square-zero extension with coefficient
\[
    M'=B'\otimes_BM.
\]
Naturality of the relative cotangent--derivation correspondence and the
etale cotangent equivalence
$L_{B'/R}\simeq B'\otimes_BL_{B/R}$ identify its classifying derivation with
\[
    \alpha':L_{B'/R}[-1]\longrightarrow M'.
\]
The structured square-zero classification now gives a canonical equivalence
\[
    A'\xrightarrow{\sim}B'_{\alpha'}
\]
over $B'$, through a contractible space of compatible identifications.
Substituting this equivalence in the preceding pushout gives the canonical
pushout square
\[
\begin{tikzcd}
B_\alpha \arrow[r] \arrow[d] &
B'_{\alpha'} \arrow[d]\\
B \arrow[r,"f"'] &
B'.
\end{tikzcd}
\]

Cotangent base change for this pushout gives
\[
\boxed{
    B'\otimes_B L_{B/B_\alpha}[-1]
    \xrightarrow{\sim}
    L_{B'/B'_{\alpha'}}[-1].}
\]
The anchor maps from $L_{B/R}[-1]$ and $L_{B'/R}[-1]$ are carried to one
another by naturality of cotangent transitivity.  Thus, on connective
almost-perfect derivations,
\[
\boxed{
    f^*\operatorname{coLie}^{\pi,E_\infty}_{B/R}
    \xrightarrow{\sim}
    \operatorname{coLie}^{\pi,E_\infty}_{B'/R}f^*.}
\]
The counit and comultiplication agree because all maps were obtained by
functoriality of the same
$\operatorname{cot}\dashv\operatorname{sqz}$ adjunction.

Continuous duality and
Lemma~\ref{lem:bnla-procoherent-finite-tor-bc} carry this equivalence to the
opposite monad comparison on
$\operatorname{dAPerfAnch}_{B/R}$.  Both sides preserve sifted colimits.
The uniqueness clause of
Theorem~\ref{thm:bnla-spectral-lie-algebroid-monad} therefore extends the
comparison uniquely to all anchored pro-coherent objects, including the monad
unit, multiplication, identity coherence, and composition coherence.
\end{proof}

\begin{theorem}[Etale naturality of affine spectral differentiation]
\label{thm:bnla-affine-spectral-recognition-etale}
The affine equivalences
\[
    \operatorname{Diff}^{\pi,E_\infty}_{B/R}:
    \operatorname{FMP}^{\mathrm{sp}}_{B/R}
    \xrightarrow{\sim}
    \operatorname{LieAlgd}^{\pi,E_\infty}_{B/R}
\]
commute with spectral etale centre change by canonical equivalences:
\[
\boxed{
    \operatorname{Diff}^{\pi,E_\infty}_{B'/R}
    f_{\mathrm{FMP},!}
    \simeq
    f^*
    \operatorname{Diff}^{\pi,E_\infty}_{B/R}.}
\]
\end{theorem}

\begin{proof}
It is enough to compare formal integration.  For an Artinian arrow
$A\to B$, let $A'\to B'$ be its canonical etale transport.  The construction
in the preceding proof, iterated through the canonical square-zero
resolution of $A$, gives
\[
\boxed{
    D^{\mathrm{sp}}_{B'/R}(A')
    \simeq
    f^*D^{\mathrm{sp}}_{B/R}(A).}
\]
Equivalently, this follows from cotangent base change on the canonical
comonadic cobar resolution together with
Lemma~\ref{lem:bnla-procoherent-finite-tor-bc}.  Insert the comparison into
the formal-integration formula
\[
    \operatorname{Int}(\mathfrak a)(A)
    =
    \operatorname{Map}
    \bigl(D(A),\mathfrak a\bigr).
\]
The resulting equivalence identifies left Kan centre change of formal
integration with extension of scalars of algebroids.  Differentiation is the
inverse equivalence, which gives the displayed naturality.
\end{proof}

\subsection{Two-variable affine-etale transport}

Let now $X\to S$ be a morphism of locally coherent qcqs spectral schemes.

\begin{definition}[Affine-etale pair category]
\label{def:bnla-affine-etale-pairs}
Let
\[
    \operatorname{AffEt}^{\mathrm{coh}}(X/S)
\]
be a small skeleton, in the universe enlargement fixed in the preceding
chapters, of the essentially small category whose objects are all pairs
\[
    V=\operatorname{Spec}(R)\longrightarrow S,
    \qquad
    U=\operatorname{Spec}(B)\longrightarrow X\times_SV
\]
with both displayed morphisms spectral etale and with $R$ and $B$ coherent.
A morphism $(U',V')\to(U,V)$ is a compatible pair of spectral etale geometric
maps.  In coordinate rings an arrow in the opposite category is a diagram
\[
\begin{tikzcd}
R \arrow[r] \arrow[d] & R' \arrow[d]\\
B \arrow[r] & B'
\end{tikzcd}
\]
with $R\to R'$ and $B\otimes_RR'\to B'$ spectral etale.
\end{definition}

\begin{theorem}[Two-variable etale base change]
\label{thm:bnla-two-variable-etale-base-change}
For every such ring diagram there are canonical coherent transports
\begin{enumerate}
    \item on spectral Artinian extensions from $R/B$ to $R'/B'$;
    \item on spectral formal moduli problems by left Kan extension followed by
    Schlessinger localization;
    \item on anchored pro-coherent modules and spectral partition-Lie
    algebroids by extension of scalars;
    \item and affine differentiation commutes with these transports.
\end{enumerate}
\end{theorem}

\begin{proof}
First base change a relative derivation along $R\to R'$ and then along the
spectral etale centre morphism
\[
    B\otimes_RR'\longrightarrow B'.
\]
Cotangent transitivity and etale vanishing give
\[
\boxed{
    L_{B'/R'}
    \simeq
    B'\otimes_BL_{B/R}.}
\]
Structured square-zero extensions commute with the same two-step transport.
The argument of
Theorem~\ref{thm:bnla-spectral-algebroid-etale-base-change} therefore gives
the relative coLie and partition-Lie comparisons directly on derivations.

For an Artinian arrow $A\to B$, base change first along $R\to R'$ and then
apply Theorem~\ref{thm:bnla-spectral-art-etale-lift} to the resulting centre
map.  The intrinsic Artinian conditions are preserved, giving the Artinian
transport.  Proposition~\ref{prop:bnla-etale-art-schlessinger} gives the
formal-moduli transport.  Finally
Theorem~\ref{thm:bnla-affine-spectral-recognition-etale} and the corresponding
base-change calculation for $D^{\mathrm{sp}}$ identify the two orders of
transport and differentiation.  Every comparison is built from the same
cotangent and square-zero base-change maps, so the identity, composition, and
higher two-variable pasting coherences agree.
\end{proof}

\subsection{Pro-coherent and algebroid etale descent}

\begin{construction}[Relative pro-coherent spectral etale coefficient family]
\label{con:bnla-relative-procoherent-etale-family}
Fix a coherent spectral base algebra $R$.  Let
\[
    \operatorname{CAlg}^{\mathrm{sp,coh,et}}_R
\]
be the category of coherent connective $R$-$E_\infty$-algebras and spectral
etale $R$-algebra morphisms.  Extension of scalars gives a cocartesian family
of stable module categories over this base.  Spectral etale pullback is
$t$-exact, preserves coherent and almost-perfect modules, finite
cosimplicial totalizations, and almost-eventually-constant Postnikov towers.
Consequently the opposite family on connective almost-perfect modules admits
the relative right-left completion obtained by imposing, fibrewise, the same
finite-totalization and almost-eventually-constant relations used in
Theorem~\ref{thm:bnla-spectral-lie-algebroid-monad}.

Denote the resulting cocartesian family by
\[
\boxed{
    \operatorname{QC}^{\vee,\mathrm{et}}_{/R}
    \longrightarrow
    \operatorname{CAlg}^{\mathrm{sp,coh,et}}_R.}
\]
Its fibre over $B$ is canonically $\operatorname{QC}^{\vee}_B$.  For a
spectral etale map $f:B\to B'$, its cocartesian transport is the functor
\[
    f^*:\operatorname{QC}^{\vee}_B
    \longrightarrow
    \operatorname{QC}^{\vee}_{B'}
\]
constructed in Lemma~\ref{lem:bnla-procoherent-finite-tor-bc}.  The
identification is canonical and coherent for identities and composition.
\end{construction}

\begin{proof}
The spectral module construction of Chapter~1 gives the cocartesian family
before imposing finiteness.  Since a spectral etale map is flat and of finite
presentation, extension of scalars is $t$-exact and preserves coherent and
almost-perfect modules.  Exactness preserves finite totalizations, while
$t$-exactness and Postnikov completeness preserve the
almost-eventually-constant tower relations.  Hence the relative right-left
completion construction of \cite{BrantnerCamposNuiten2025}, in the relative
form used in \cite{BrantnerMagidsonNuiten2025}, applies fibrewise and is
compatible with cocartesian transport.

On the fibre over $B$, the defining completion is exactly the completion used
in the proof of Theorem~\ref{thm:bnla-spectral-lie-algebroid-monad}; the
pro-coherent spectral coefficient theorem therefore identifies it with
$\operatorname{QC}^{\vee}_B$.  On continuous duals of connective
almost-perfect modules the transition is forced to be
\[
    M^\vee\longmapsto(B'\otimes_BM)^\vee.
\]
Lemma~\ref{lem:bnla-procoherent-finite-tor-bc} identifies this with the
already constructed $f^*$.  Both functors preserve the completion relations
and sifted colimits, so uniqueness of the relative right-left extension
identifies them on the entire fibre.  The relative completion universal
property supplies the identity, composition, and higher cocartesian
coherences.
\end{proof}

For every coherent affine-etale pair $(U,V)$ as above put
\[
    \mathscr{LieAlgd}^{\pi,E_\infty}(U,V)
    :=
    \operatorname{LieAlgd}^{\pi,E_\infty}_{B/R}
\]
and
\[
    \mathscr{FMP}^{\mathrm{sp}}(U,V)
    :=
    \operatorname{FMP}^{\mathrm{sp}}_{B/R}.
\]

\begin{proposition}[Etale descent for pro-coherent spectral modules]
\label{prop:bnla-procoherent-etale-descent}
Let $B\to B'$ be a faithfully spectral etale morphism between coherent
connective $E_\infty$-rings, and write
\[
    B^\bullet
    =
    \left(
        B'
        \rightrightarrows
        B'\otimes_BB'
        \rightrightarrows\cdots
    \right)
\]
for its Amitsur cosimplicial algebra.  Then extension of scalars induces a
canonical equivalence
\[
\boxed{
    \operatorname{QC}^{\vee}_B
    \xrightarrow{\sim}
    \operatorname*{Tot}_{[n]\in\Delta}
    \operatorname{QC}^{\vee}_{B^n}.}
\]
The same statement holds for a finite affine spectral etale covering family
after replacing the family by its finite product algebra.  The equivalence is
compatible with pullback, refinement, and Cech pasting.
\end{proposition}

\begin{proof}
A faithfully spectral etale morphism is spectral fppf and hence is a universal
descent morphism of connective commutative ring spectra.  Concretely, after
any base change the smallest stable retract-closed full subcategory of the
module category containing the restrictions of modules from the cover is the
whole module category.  Spectral etale maps have Tor-amplitude zero, so
Lemma~\ref{lem:bnla-procoherent-finite-tor-bc} supplies the
Beck--Chevalley equivalences for every square in the Amitsur diagram.

Let
\[
    \theta:
    \operatorname{QC}^{\vee}_B
    \longrightarrow
    \operatorname*{Tot}_{[n]\in\Delta}
    \operatorname{QC}^{\vee}_{B^n}
\]
be the canonical descent comparison.  The Barr--Beck descent criterion,
applied with the Beck--Chevalley equivalences just proved, reduces the
invertibility of $\theta$ to conservativity of
\[
    f^*:
    \operatorname{QC}^{\vee}_B
    \longrightarrow
    \operatorname{QC}^{\vee}_{B'}.
\]
We prove this conservativity directly.  Let
$F\in\operatorname{QC}^{\vee}_B$ satisfy $f^*F\simeq0$, and use the
canonical right-left exact extension of $F$ to the module category.  Let
\[
    \mathcal K_F
    :=
    \left\{
        N\in\operatorname{Mod}_B
        \mid
        F(N)\simeq0
    \right\}.
\]
Exactness makes $\mathcal K_F$ a stable full subcategory, and it is closed
under retracts.  For every $B'$-module $N'$, the pointwise description of
pro-coherent pullback gives
\[
    F\left(\operatorname{Res}^{B'}_BN'\right)
    \simeq
    (f^*F)(N')
    \simeq0.
\]
Thus $\mathcal K_F$ contains the restriction of every $B'$-module.  Universal
descent for $f$ therefore implies
\[
    \mathcal K_F=\operatorname{Mod}_B.
\]
In particular the right-left extension of $F$ vanishes on its entire defining
module sector, so $F\simeq0$.  Hence $f^*$ is conservative, and Barr--Beck
gives the desired equivalence
\[
    \operatorname{QC}^{\vee}_B
    \xrightarrow{\sim}
    \operatorname*{Tot}_{[n]\in\Delta}
    \operatorname{QC}^{\vee}_{B^n}.
\]

For a finite affine spectral etale covering family, replace the family by the
finite product of its members.  Modules, coherent modules, and therefore
pro-coherent modules over a finite product split componentwise, so the
Amitsur diagram of the product identifies with the componentwise Cech nerve.
The pointwise Beck--Chevalley equivalences of
Lemma~\ref{lem:bnla-procoherent-finite-tor-bc} are constructed from tensor
associativity and are therefore coherent under pullback, refinement, and
iterated Cech pasting.  This proves the final assertion.
\end{proof}

\begin{corollary}[Conservativity for spectral etale covers]
\label{cor:bnla-procoherent-etale-conservative}
If $B\to B'$ is faithfully spectral etale between coherent connective
$E_\infty$-rings, then
\[
    f^*:\operatorname{QC}^{\vee}_B
    \longrightarrow
    \operatorname{QC}^{\vee}_{B'}
\]
is conservative.
\end{corollary}

\begin{proof}
If $f^*E\simeq0$, then every term of the Amitsur descent datum of $E$ is zero.
Proposition~\ref{prop:bnla-procoherent-etale-descent} reconstructs $E$ as the
totalization of that zero diagram, so $E\simeq0$.
\end{proof}

\begin{proposition}[Etale descent for spectral partition-Lie algebroids]
\label{prop:bnla-spectral-algebroid-etale-descent}
The prestack
\[
    (U,V)
    \longmapsto
    \operatorname{LieAlgd}^{\pi,E_\infty}_{B/R}
\]
on $\operatorname{AffEt}^{\mathrm{coh}}(X/S)$ satisfies spectral etale
descent.  Denote the resulting etale stack by
\[
    \mathscr{LieAlgd}^{\pi,E_\infty}_{X/S}.
\]
\end{proposition}

\begin{proof}
Proposition~\ref{prop:bnla-procoherent-etale-descent} gives descent for the
pro-coherent coefficient categories.  Cotangent base change identifies the
anchor targets
\[
    \mathbb T^{\mathrm{pc}}_{B/R}[1]
\]
through every Cech degree.  Slicing over these compatible targets therefore
gives descent for
$\operatorname{Anch}^{\mathrm{pc}}_{B/R}$.

Theorem~\ref{thm:bnla-spectral-algebroid-etale-base-change} identifies the
partition-Lie monads under every Cech transition, including unit and
multiplication.  Algebras for a compatible diagram of accessible monads are
computed in the limit of the corresponding monadic diagrams.  Hence
\[
\boxed{
    \operatorname{LieAlgd}^{\pi,E_\infty}_{B/R}
    \xrightarrow{\sim}
    \operatorname*{Tot}
    \operatorname{LieAlgd}^{\pi,E_\infty}_{B^\bullet/R^\bullet}.}
\]
Thus the affine algebroid prestack is already an etale stack.
\end{proof}

\begin{theorem}[Global spectral formal-moduli stack]
\label{thm:bnla-global-fmp-localization}
The prestack
\[
    (U,V)\longmapsto
    \operatorname{FMP}^{\mathrm{sp}}_{B/R}
\]
with transition functors of
Theorem~\ref{thm:bnla-two-variable-etale-base-change} satisfies spectral
etale descent.  Denote it by
\[
    \mathscr{FMP}^{\mathrm{sp}}_{X/S}.
\]
There is a canonical equivalence of etale stacks
\[
\boxed{
    \operatorname{Diff}^{\pi,E_\infty}:
    \mathscr{FMP}^{\mathrm{sp}}_{X/S}
    \xrightarrow{\sim}
    \mathscr{LieAlgd}^{\pi,E_\infty}_{X/S}.}
\]
\end{theorem}

\begin{proof}
The affine differentiation equivalence is objectwise an equivalence by
Theorem~\ref{thm:bnla-spectral-affine-recognition}.  Its compatibility with
all one-variable and two-variable transition functors is
Theorems~\ref{thm:bnla-affine-spectral-recognition-etale} and
\ref{thm:bnla-two-variable-etale-base-change}.  Hence it is an equivalence of
presentable-category-valued prestacks.  The target prestack satisfies etale
descent by
Proposition~\ref{prop:bnla-spectral-algebroid-etale-descent}; therefore the
source prestack does as well.
\end{proof}

\begin{definition}[Global spectral formal moduli and partition-Lie algebroids]
\label{def:bnla-global-spectral-pair-sections}
Define
\[
\boxed{
    \operatorname{FMP}^{\mathrm{sp}}_{X/S}
    :=
    \Gamma_{\mathrm{et}}
    \left(
        \operatorname{AffEt}^{\mathrm{coh}}(X/S),
        \mathscr{FMP}^{\mathrm{sp}}_{X/S}
    \right)}
\]
and
\[
\boxed{
    \operatorname{LieAlgd}^{\pi,E_\infty}_{X/S}
    :=
    \Gamma_{\mathrm{et}}
    \left(
        \operatorname{AffEt}^{\mathrm{coh}}(X/S),
        \mathscr{LieAlgd}^{\pi,E_\infty}_{X/S}
    \right).}
\]
No affine presentation or separately retained descent datum is part of an
object.
\end{definition}

\begin{theorem}[Global spectral formal-moduli/partition-Lie equivalence]
\label{thm:bnla-global-spectral-recognition}
Taking etale global sections in
Theorem~\ref{thm:bnla-global-fmp-localization} gives a canonical equivalence
\[
\boxed{
    \operatorname{Diff}^{\pi,E_\infty}_{X/S}:
    \operatorname{FMP}^{\mathrm{sp}}_{X/S}
    \xrightarrow{\sim}
    \operatorname{LieAlgd}^{\pi,E_\infty}_{X/S}.}
\]
It is computed on every coherent affine-etale pair by
Theorem~\ref{thm:bnla-spectral-affine-recognition}.
\end{theorem}

\begin{proof}
Global sections preserve equivalences of category-valued etale stacks.
\end{proof}

\begin{theorem}[Affine restriction of global spectral formal moduli]
\label{prop:bnla-global-fmp-restriction}
Restriction from
$\operatorname{FMP}^{\mathrm{sp}}_{X/S}$ to a coherent affine-etale pair
$(U,V)=(\operatorname{Spec}B,\operatorname{Spec}R)$ is the ordinary affine
restriction functor to
$\operatorname{FMP}^{\mathrm{sp}}_{B/R}$.  A global object is reconstructed
from these restrictions by the canonical etale cocartesian transport and
Cech descent.
\end{theorem}

\begin{proof}
This is the defining global-section universal property of the etale stack
$\mathscr{FMP}^{\mathrm{sp}}_{X/S}$.
\end{proof}

\subsection{A global Artinian-deformation reflector on nonlinear targets}

We first construct the Artinian component of the target-level deformation
localization.  It is sufficient to extract a formal-moduli problem at one
fixed coherent centre.  The stronger convergence and nilpotent-excision
localization actually applied to the nonlinear quotient is constructed
immediately afterwards.

\begin{definition}[Global Artinian-deformation generators]
\label{def:bnla-global-artdef-generators}
Work in the native big spectral Nisnevich slice $\mathcal X_{/S}$.  For every
coherent affine diagram in the fixed coefficient universe
\[
    U=\operatorname{Spec}B
    \longrightarrow
    V=\operatorname{Spec}R
    \longrightarrow S,
\]
every spectral Artinian arrow $A\to B$, every connective coherent
$B$-module $M$, and every morphism over $B$
\[
    \eta:A\longrightarrow B\oplus M[1],
\]
put
\[
    A_\eta:=A\times_{B\oplus M[1]}B.
\]
Let $h_C$ denote the spectral Nisnevich representable associated to
$\operatorname{Spec}C$ over $S$.  Define
\[
\boxed{
    w_{A,\eta}:
    h_A\amalg_{h_{B\oplus M[1]}}h_B
    \longrightarrow
    h_{A_\eta}.}
\]
Let $\mathcal W^{\mathrm{ArtDef}}_S$ be the set of all such morphisms in the
fixed universe.  For $Y\in\mathcal X_{/S}$ write
\[
    Y(C):=\operatorname{Map}_S(\operatorname{Spec}C,Y).
\]
\end{definition}

\begin{theorem}[Global Artinian-deformation reflector]
\label{thm:bnla-global-artdef-reflector}
Accessible localization at
$\mathcal W^{\mathrm{ArtDef}}_S$ exists:
\[
\boxed{
    L^{\mathrm{ArtDef}}_S:
    \mathcal X_{/S}
    \longrightarrow
    \mathcal X^{\mathrm{ArtDef}}_{/S}.}
\]
An object $Y\to S$ is local precisely when, for every generator, the canonical
map
\[
\boxed{
    Y(A_\eta)
    \longrightarrow
    Y(A)\times_{Y(B\oplus M[1])}Y(B)}
\]
is an equivalence.

Every represented spectral scheme over $S$ is
$L^{\mathrm{ArtDef}}_S$-local.
\end{theorem}

\begin{proof}
The slice of the native big spectral Nisnevich topos is presentable, and the
chosen universe makes the generating family a set.  Accessible localization
therefore exists.  Yoneda identifies locality for $w_{A,\eta}$ with the
displayed mapping-space equation.

It remains to check represented objects.  Both algebra maps
\[
    A\longrightarrow B\oplus M[1],
    \qquad
    B\longrightarrow B\oplus M[1]
\]
are surjective on $\pi_0$: the first has the same degree-zero map as the
Artinian augmentation $A\to B$, while the second is an isomorphism on
$\pi_0$ because $M$ is connective.  The corresponding affine geometric maps
are therefore closed immersions.  Chapter~1's closed-gluing theorem identifies
\[
\boxed{
    \operatorname{Spec}(A_\eta)
    \simeq
    \operatorname{Spec}(A)
    \amalg_{\operatorname{Spec}(B\oplus M[1])}
    \operatorname{Spec}(B)}
\]
as a pushout of spectral schemes.  Mapping this pushout into any represented
spectral scheme $Y$ gives exactly the required pullback of mapping spaces.
Thus every represented spectral scheme is local.
\end{proof}

\begin{proposition}[The Artinian-deformation unit is reduced-classically invisible]
\label{prop:bnla-artdef-dred-equivalence}
For every $Y\in\mathcal X_{/S}$ the localization unit
\[
    Y\longrightarrow L^{\mathrm{ArtDef}}_S Y
\]
is a $D_{\mathrm{red}}$-equivalence:
\[
\boxed{
    D_{\mathrm{red}}Y
    \xrightarrow{\sim}
    D_{\mathrm{red}}L^{\mathrm{ArtDef}}_S Y.}
\]
\end{proposition}

\begin{proof}
For a spectral Artinian arrow $A\to B$, the nilpotent kernel on $\pi_0$ and
the finite square-zero generation theorem give
\[
    D_{\mathrm{red}}h_A
    \simeq
    D_{\mathrm{red}}h_B.
\]
The same holds for $B\oplus M[1]$ and for $A_\eta$.  Consequently
$D_{\mathrm{red}}$ carries every generator
$w_{A,\eta}$ to an equivalence.  By the universal property of localization,
$D_{\mathrm{red}}$ factors through $L^{\mathrm{ArtDef}}_S$.  Applying this
factorization to the localization unit gives the displayed equivalence.
\end{proof}


\subsection{A global spectral deformation reflector on nonlinear targets}
\label{sec:bnla-global-def-reflector}

The Artinian-deformation reflector is sufficient to extract a formal-moduli
problem at one fixed coherent centre.  The comparison of those centres requires
a stronger target-level property: the nonlinear target must itself have
spectral deformation theory, namely Postnikov convergence and excision for
nilpotent extensions.  We construct that property by a second accessible
localization.  It is applied globally once, before any affine leaf is
extracted.

\begin{definition}[Global spectral deformation generators]
\label{def:bnla-global-def-generators}
Work in the native big spectral Nisnevich slice $\mathcal X_{/S}$ and in the
fixed affine coefficient universe.

First, for every connective affine spectral test
\[
    \operatorname{Spec}A\longrightarrow S,
\]
let
\[
\boxed{
    c_A:
    \operatorname*{colim}_{n\geq0}h_{\tau_{\leq n}A}
    \longrightarrow
    h_A}
\]
be the morphism induced by the Postnikov truncation maps
$A\to\tau_{\leq n}A$.

Second, for every pullback square of connective commutative ring spectra over
an affine chart of $S$
\[
\begin{tikzcd}
C \arrow[r] \arrow[d] &
A_0 \arrow[d]\\
A_1 \arrow[r] &
A_{01},
\end{tikzcd}
\qquad
C\simeq A_0\times_{A_{01}}A_1,
\]
in which one of the two maps
\[
    A_i\longrightarrow A_{01}
\]
induces a surjection on $\pi_0$ with nilpotent kernel, let
\[
\boxed{
    w_{\square}:
    h_{A_0}\amalg_{h_{A_{01}}}h_{A_1}
    \longrightarrow
    h_C.}
\]
The connective pullback $C$ exists in connective commutative ring spectra:
the underlying fibre sequence
\[
    C\longrightarrow A_0\oplus A_1\longrightarrow A_{01}
\]
and surjectivity of one component on $\pi_0$ kill the only possible negative
homotopy group.

Let
\[
    \mathcal W^{\mathrm{Def}}_S
\]
be the set consisting of all $c_A$ and all $w_{\square}$ in the fixed
universe.
\end{definition}

\begin{lemma}[Spectral nilpotent affine gluing]
\label{lem:bnla-spectral-nilpotent-gluing}
Let
\[
\begin{tikzcd}
C \arrow[r] \arrow[d] &
A_0 \arrow[d]\\
A_1 \arrow[r] &
A_{01}
\end{tikzcd}
\]
be a pullback square of connective $E_\infty$-rings and suppose, say, that
$\pi_0(A_0)\to\pi_0(A_{01})$ is surjective with nilpotent kernel.  For every
represented spectral scheme $Y$ the canonical map
\[
\boxed{
    \operatorname{Map}(\operatorname{Spec}C,Y)
    \longrightarrow
    \operatorname{Map}(\operatorname{Spec}A_0,Y)
    \times_{\operatorname{Map}(\operatorname{Spec}A_{01},Y)}
    \operatorname{Map}(\operatorname{Spec}A_1,Y)}
\]
is an equivalence.
\end{lemma}

\begin{proof}
The assertion is immediate when $Y=\operatorname{Spec}D$ is affine, because
mapping out of $D$ carries the defining pullback of commutative ring spectra to
the displayed pullback of mapping spaces.

For general $Y$, the affine geometric morphism
\[
    \operatorname{Spec}A_{01}\longrightarrow\operatorname{Spec}A_0
\]
is a nilpotent closed immersion and hence a homeomorphism on underlying
topological spaces.  Its base change
\[
    \operatorname{Spec}A_1\longrightarrow\operatorname{Spec}C
\]
has the same property.  Thus a finite principal affine cover of
$\operatorname{Spec}C$, equivalently of $\operatorname{Spec}A_1$, may be
chosen subordinate to affine opens of $Y$ for a given compatible pair of
maps.  Compatibility on $\operatorname{Spec}A_{01}$ and the preceding
homeomorphism force the corresponding restrictions of
$\operatorname{Spec}A_0$ to land in the same affine opens.  On each member
and every finite intersection the affine calculation gives the desired
equivalence.  Spectral Zariski descent for maps into $Y$ then glues these
local equivalences.  The construction is natural in the square and in $Y$.
\end{proof}

\begin{theorem}[Global spectral deformation reflector]
\label{thm:bnla-global-def-reflector}
Accessible localization at $\mathcal W^{\mathrm{Def}}_S$ exists:
\[
\boxed{
    L^{\mathrm{Def}}_S:
    \mathcal X_{/S}
    \longrightarrow
    \mathcal X^{\mathrm{Def}}_{/S}.}
\]
An object $Y\to S$ is local if and only if the following two conditions hold:
\begin{enumerate}
    \item for every connective affine spectral test $A$ over $S$,
    \[
    \boxed{
        Y(A)
        \xrightarrow{\sim}
        \lim_nY(\tau_{\leq n}A);}
    \]
    \item for every pullback square of
    Definition~\ref{def:bnla-global-def-generators} with one nilpotent leg,
    \[
    \boxed{
        Y(C)
        \xrightarrow{\sim}
        Y(A_0)\times_{Y(A_{01})}Y(A_1).}
    \]
\end{enumerate}
Every represented spectral scheme over $S$ is
$L^{\mathrm{Def}}_S$-local.  Moreover every
$L^{\mathrm{Def}}_S$-local object is
$L^{\mathrm{ArtDef}}_S$-local.
\end{theorem}

\begin{proof}
The native slice is presentable and the fixed universe makes
$\mathcal W^{\mathrm{Def}}_S$ a set, so the accessible localization exists.
Yoneda identifies locality for the two classes of generators with the two
displayed conditions.

Let $Y$ be represented.  Nilpotent excision is
Lemma~\ref{lem:bnla-spectral-nilpotent-gluing}.  For convergence, first take
$Y=\operatorname{Spec}D$.  Operadic Postnikov convergence gives
\[
    A\simeq\lim_n\tau_{\leq n}A
\]
in connective commutative ring spectra, and therefore
\[
\begin{aligned}
    Y(A)
    &=
    \operatorname{Map}_{\operatorname{CAlg}(\operatorname{Sp})}(D,A)\\
    &\simeq
    \lim_n
    \operatorname{Map}_{\operatorname{CAlg}(\operatorname{Sp})}
    (D,\tau_{\leq n}A)
    =
    \lim_nY(\tau_{\leq n}A).
\end{aligned}
\]
For a general represented spectral scheme, the truncations
$\operatorname{Spec}\tau_{\leq n}A$ and $\operatorname{Spec}A$ have the same
underlying topological space.  Open spectral subschemes are detected on
classical truncation, so one may use the same finite principal cover in every
Postnikov degree.  The affine calculation on that cover and on all its finite
intersections glues by spectral Zariski descent to the required equivalence.

Finally every Artinian-deformation generator
\[
    A_\eta=A\times_{B\oplus M[1]}B
\]
is a special case of the nilpotent-excision generator: the map
\[
    B\longrightarrow B\oplus M[1]
\]
induces an isomorphism on $\pi_0$ for connective $M$.  Hence every
$L^{\mathrm{Def}}_S$-local object is
$L^{\mathrm{ArtDef}}_S$-local.
\end{proof}

\begin{proposition}[The deformation unit is reduced-classically invisible]
\label{prop:bnla-def-dred-equivalence}
For every $Y\in\mathcal X_{/S}$ the localization unit
\[
    Y\longrightarrow L^{\mathrm{Def}}_S Y
\]
is a $D_{\mathrm{red}}$-equivalence:
\[
\boxed{
    D_{\mathrm{red}}Y
    \xrightarrow{\sim}
    D_{\mathrm{red}}L^{\mathrm{Def}}_S Y.}
\]
\end{proposition}

\begin{proof}
Postnikov truncation does not change $\pi_0$, so $D_{\mathrm{red}}$ carries
every convergence generator $c_A$ to an equivalence.

For a nilpotent-excision generator, suppose
$A_0\to A_{01}$ has nilpotent degree-zero kernel.  Then
\[
    D_{\mathrm{red}}h_{A_0}
    \simeq
    D_{\mathrm{red}}h_{A_{01}},
\]
and its base change gives
\[
    D_{\mathrm{red}}h_C
    \simeq
    D_{\mathrm{red}}h_{A_1}.
\]
Thus $D_{\mathrm{red}}$ carries $w_\square$ to an equivalence as well.
The universal property of localization therefore factors
$D_{\mathrm{red}}$ through $L^{\mathrm{Def}}_S$, and the factorization applied
to the localization unit gives the result.
\end{proof}

\begin{theorem}[Spectral-etale base change of the deformation reflector]
\label{thm:bnla-def-etale-base-change}
Let $c:S'\to S$ be spectral etale.  Pullback induces a canonical equivalence
\[
\boxed{
    c^*L^{\mathrm{Def}}_S(Y)
    \xrightarrow{\sim}
    L^{\mathrm{Def}}_{S'}(c^*Y)}
\]
for every $Y\in\mathcal X_{/S}$, coherently in identities, composition, and
spectral-etale pullback pasting.
\end{theorem}

\begin{proof}
The assertion is local on the affine-etale basis of $S'$.  On an affine
diagram
\[
    R\longrightarrow R'
\]
with $R\to R'$ spectral etale, extension of scalars is $t$-exact.  Hence it
commutes with connective Postnikov truncation and carries every convergence
generator over $R$ to the corresponding convergence generator over $R'$.
It also carries a pullback square with a nilpotent leg to a pullback square
with a nilpotent leg.

Conversely, after restriction of scalars along $R\to R'$, every convergence
or nilpotent-excision generator over $R'$ is a generator of the same type over
$R$.  Thus the two strongly saturated localizing classes agree after
restriction to the affine-etale basis.  The universal property of accessible
localization gives the affine comparison, and spectral Nisnevich descent glues
it globally.  The coherence is inherited from scalar extension, restriction,
and the uniqueness of localization comparisons.
\end{proof}

\subsection{The globally formalized quotient target}

Let
\[
    \mathfrak G_\bullet\in\operatorname{FormLieGpd}(X/S)
\]
and write
\[
    q_{\mathfrak G}:X\twoheadrightarrow Q_{\mathfrak G}
\]
for its effective formal quotient.

\begin{definition}[Spectral deformation formalization of the quotient]
\label{def:bnla-artdef-formal-quotient}
Define
\[
\boxed{
    Q^{\mathrm{def}}_{\mathfrak G}
    :=
    L^{\mathrm{Def}}_S(Q_{\mathfrak G}).}
\]
The quotient map followed by the localization unit gives
\[
    X\longrightarrow Q^{\mathrm{def}}_{\mathfrak G}.
\]
\end{definition}

Since $q_{\mathfrak G}$ is a relative de Rham equivalence, the preceding
chapter and Chapter~9 give
\[
    D_{\mathrm{red}}X
    \xrightarrow{\sim}
    D_{\mathrm{red}}Q_{\mathfrak G}.
\]
Proposition~\ref{prop:bnla-def-dred-equivalence} therefore gives a canonical
equivalence
\[
    D_{\mathrm{red}}X
    \xrightarrow{\sim}
    D_{\mathrm{red}}Q^{\mathrm{def}}_{\mathfrak G}.
\]

\begin{definition}[Reduced-centred formal quotient target]
\label{def:bnla-reduced-centred-formal-quotient}
Define
\[
\boxed{
    Q^{\mathrm{fm}}_{\mathfrak G}
    :=
    Q^{\mathrm{def}}_{\mathfrak G}
    \times_{D_{\mathrm{red}}Q^{\mathrm{def}}_{\mathfrak G}}
    D_{\mathrm{red}}X.}
\]
The displayed reduced-classical equivalence makes the projection
\[
    Q^{\mathrm{fm}}_{\mathfrak G}
    \longrightarrow
    Q^{\mathrm{def}}_{\mathfrak G}
\]
an equivalence.  In particular $Q^{\mathrm{fm}}_{\mathfrak G}$ is
$L^{\mathrm{Def}}_S$-local.  We retain the pullback model because it carries a canonical
reduced centre and a canonical map
\[
\boxed{
    q^{\mathrm{fm}}_{\mathfrak G}:
    X\longrightarrow Q^{\mathrm{fm}}_{\mathfrak G}.}
\]
Moreover the localization unit
$Q_{\mathfrak G}\to Q^{\mathrm{def}}_{\mathfrak G}$ and the inverse of
$D_{\mathrm{red}}q_{\mathfrak G}$ determine a canonical morphism
\[
    Q_{\mathfrak G}\longrightarrow Q^{\mathrm{fm}}_{\mathfrak G}
\]
through a contractible space of choices.
\end{definition}

\begin{proposition}[Represented quotients are unchanged]
\label{prop:bnla-represented-quotient-artdef-local}
If $Q_{\mathfrak G}$ is represented by a spectral scheme, then
\[
\boxed{
    Q_{\mathfrak G}
    \xrightarrow{\sim}
    Q^{\mathrm{def}}_{\mathfrak G}
    \xleftarrow{\sim}
    Q^{\mathrm{fm}}_{\mathfrak G}.}
\]
\end{proposition}

\begin{proof}
The first equivalence is
Theorem~\ref{thm:bnla-global-def-reflector}, because represented spectral
schemes are deformation-local.  The second is the reduced-centred pullback
along an equivalence.
\end{proof}

\subsection{Affine leaves of the global formal quotient}

For a coherent affine-etale pair
\[
    V=\operatorname{Spec}R\longrightarrow S,
    \qquad
    U=\operatorname{Spec}B\longrightarrow X\times_SV,
\]
write
\[
    q^{\mathrm{fm}}_U:
    U\longrightarrow Q^{\mathrm{fm}}_{\mathfrak G}
\]
for the restriction of the formalized quotient map.

\begin{definition}[Affine formal leaf of the global quotient]
\label{def:bnla-affine-global-formal-leaf}
For $A\to B\in\operatorname{Art}^{\mathrm{sp}}_{R/B}$ and
$T=\operatorname{Spec}A$, define
\[
\boxed{
    \mathcal L^{\mathrm{fm,sp}}_{\mathfrak G,U/V}(A\to B)
    :=
    \operatorname{fib}_{q^{\mathrm{fm}}_U}
    \left(
        \operatorname{Map}_S(T,Q^{\mathrm{fm}}_{\mathfrak G})
        \longrightarrow
        \operatorname{Map}_S(U,Q^{\mathrm{fm}}_{\mathfrak G})
    \right).}
\]
\end{definition}

\begin{proposition}[Artinian leaf of an Artinian-deformation-local target]
\label{prop:bnla-affine-leaf-is-fmp}
Let
\[
    V=\operatorname{Spec}R\longrightarrow S,
    \qquad
    U=\operatorname{Spec}B\longrightarrow V
\]
be coherent affines, let $Y\to S$ be $L^{\mathrm{ArtDef}}_S$-local, and let
$u:U\to Y$ be a centre over $S$.  Define
\[
\boxed{
    F_{U/Y}(A\to B)
    :=
    \operatorname{fib}_{u}
    \left(
        \operatorname{Map}_S(\operatorname{Spec}A,Y)
        \longrightarrow
        \operatorname{Map}_S(U,Y)
    \right).}
\]
Then
\[
\boxed{
    F_{U/Y}\in\operatorname{FMP}^{\mathrm{sp}}_{B/R}.}
\]
In particular, for
$Y=Q^{\mathrm{fm}}_{\mathfrak G}$ and $u=q^{\mathrm{fm}}_U$, one has
\[
    F_{U/Y}
    =
    \mathcal L^{\mathrm{fm,sp}}_{\mathfrak G,U/V}.
\]
The morphism
$Q_{\mathfrak G}\to Q^{\mathrm{fm}}_{\mathfrak G}$ induces, on the
coherent Artinian subcategory of the finite infinitesimal tests, a natural
transformation from the raw quotient-lifting functor to this formal leaf.
\end{proposition}

\begin{proof}
For every $R$-algebra $C$ equipped with its structural map to $B$, pullback to
$V$ gives the canonical identification
\[
    \operatorname{Map}_V
    \left(\operatorname{Spec}C,Y\times_SV\right)
    \simeq
    \operatorname{Map}_S(\operatorname{Spec}C,Y),
\]
so the displayed formula is the centred Artinian leaf in either slice.  At
the identity object $B\to B$, the restriction map is the identity and hence
\[
    F_{U/Y}(B)\simeq *.
\]

Now let
\[
\begin{tikzcd}
A_\eta \arrow[r] \arrow[d] & B \arrow[d]\\
A \arrow[r,"\eta"'] & B\oplus M[1]
\end{tikzcd}
\]
be a Schlessinger square.  Since $Y$ is $L^{\mathrm{ArtDef}}_S$-local,
\[
    Y(A_\eta)
    \simeq
    Y(A)\times_{Y(B\oplus M[1])}Y(B).
\]
Taking the homotopy fibre of this equivalence over the compatible centre
$u\in Y(B)$ and using pullback pasting gives
\[
\begin{aligned}
    F_{U/Y}(A_\eta)
    &\simeq
    F_{U/Y}(A)
    \times_{F_{U/Y}(B\oplus M[1])}
    F_{U/Y}(B)\\
    &\simeq
    F_{U/Y}(A)
    \times_{F_{U/Y}(B\oplus M[1])}*.
\end{aligned}
\]
Thus $F_{U/Y}$ satisfies the Artinian Schlessinger equations.

For the stated specialization, the target
$Q^{\mathrm{fm}}_{\mathfrak G}$ is $L^{\mathrm{ArtDef}}_S$-local:
Definition~\ref{def:bnla-reduced-centred-formal-quotient} gives a canonical
equivalence
\[
    Q^{\mathrm{fm}}_{\mathfrak G}
    \xrightarrow{\sim}
    Q^{\mathrm{def}}_{\mathfrak G},
\]
and $Q^{\mathrm{def}}_{\mathfrak G}=L^{\mathrm{Def}}_S(Q_{\mathfrak G})$
is $L^{\mathrm{Def}}_S$-local by construction, hence
$L^{\mathrm{ArtDef}}_S$-local by
Theorem~\ref{thm:bnla-global-def-reflector}.  The final natural transformation
is induced by postcomposition with
$Q_{\mathfrak G}\to Q^{\mathrm{fm}}_{\mathfrak G}$.
\end{proof}

\subsection{Pre-tangent coefficients and almost-finite approximation}
\label{sec:bnla-pre-tangent-aft}

The affine formal leaf is defined only on Artinian tests, whereas etale
transport of its controller must be functorial before one knows that an
arbitrary coherent coefficient over the new centre descends from the old
one.  We therefore construct the square-zero tangent functor before applying
pro-coherent almost-finite approximation.

\begin{definition}[Accessible convergent pre-pro-coherent functors]
\label{def:bnla-pre-procoherent-functors}
For a coherent connective $E_\infty$-ring $B$, let
\[
    \mathcal D_B
\]
be the full stable subcategory of
$\operatorname{Fun}(\operatorname{Mod}_B,\operatorname{Sp})$ spanned by the
exact accessible convergent functors.  Here convergence means that for every
$B$-module $M$ the canonical map
\[
    F(M)
    \longrightarrow
    \lim_n F(\tau_{\leq n}M)
\]
is an equivalence.
\end{definition}

\begin{proposition}[Almost-finite approximation]
\label{prop:bnla-pre-tangent-aft-approximation}
For coherent $B$, continuous pro-coherent duality identifies
$\operatorname{QC}^{\vee}_B$ with the full subcategory of $\mathcal D_B$
consisting of the almost-finitely-presented functors.  For every
$F\in\mathcal D_B$, the slice
\[
    \left(\operatorname{QC}^{\vee}_B\right)_{/F}
\]
has a terminal object
\[
\boxed{
    F^{\mathrm{aft}}\longrightarrow F,}
\]
called the almost-finitely-presented approximation of $F$.  The construction
is functorial in $F$ inside the fixed coefficient fibre.  No assertion that
the fibrewise approximation functors commute with a change of coefficient
ring is made.
\end{proposition}

\begin{proof}
The spectral right-left description of pro-coherent modules in
\cite{BrantnerCamposNuiten2025} identifies
$\operatorname{QC}^{\vee}_B$ with exact accessible convergent functors which
are almost finitely presented.  The same right-left theorem proves the
universal approximation statement: after restricting to a sufficiently
large essentially small subcategory of compact Postnikov-bounded modules, the
pro-coherent subcategory is an accessible colimit-closed full subcategory of
the corresponding functor category, so its inclusion admits a right adjoint.
Applied to the restriction of $F$, this right adjoint gives the terminal
object $F^{\mathrm{aft}}\to F$, and convergence extends the universal property
from the bounded subcategory to all modules.

This construction is fibrewise in $B$.  In particular, nothing in the
argument identifies the right adjoints for two different coefficient rings.
The required coefficient-change functoriality is constructed below in the
anchored category, before any such identification is attempted.
\end{proof}

\begin{construction}[Coherent pre-tangent of a nonlinear target]
\label{con:bnla-coherent-pre-tangent}
Let
\[
    V=\operatorname{Spec}R\longrightarrow S,
    \qquad
    U=\operatorname{Spec}B\longrightarrow V
\]
be coherent affines and let $u:U\to Y$ be a morphism in
$\mathcal X_{/S}$, with $Y$ $L_S^{\mathrm{ArtDef}}$-local.  Write
\[
    F_{U/Y}\in\operatorname{FMP}^{\mathrm{sp}}_{B/R}
\]
for the centred Artinian leaf functor of $Y$ at $u$, whose formal-moduli
property was proved in Proposition~\ref{prop:bnla-affine-leaf-is-fmp}.

For a connective coherent $B$-module $M$ and $n\geq0$, put
\[
\begin{aligned}
    \Theta^n_{Y/V,u}(M)
    &:=F_{U/Y}(B\oplus M[n])\\
    &\simeq
    \operatorname{fib}_{u}
    \left(
        \operatorname{Map}_V
        \left(\operatorname{Spec}(B\oplus M[n]),Y\times_SV\right)
        \longrightarrow
        \operatorname{Map}_V(U,Y\times_SV)
    \right).
\end{aligned}
\]
The split square-zero identities
\[
    B\oplus M[n]
    \simeq
    B\times_{B\oplus M[n+1]}B
\]
and the Schlessinger equations give canonical equivalences
\[
    \Theta^n_{Y/V,u}(M)
    \xrightarrow{\sim}
    \Omega\Theta^{n+1}_{Y/V,u}(M).
\]
Thus $n\mapsto\Theta^n_{Y/V,u}(M)$ is canonically an $\Omega$-spectrum.
The square-zero excision equations for finite coherent coefficient diagrams
make the resulting spectrum-valued functor exact on the stable coherent
closure.  Its accessible Postnikov right-left extension defines a canonical
exact accessible convergent functor
\[
\boxed{
    \mathbb T^{\mathrm{pre}}_{Y/V,u}
    \in\mathcal D_B.}
\]

For the represented target $U$ itself, the relative cotangent universal
property gives
\[
\boxed{
    \mathbb T^{\mathrm{pre}}_{U/V}(M)
    \simeq
    \operatorname{map}_B(L_{B/R},M)}
\]
for all $M\in\operatorname{Mod}_B$.  Consequently
\[
\boxed{
    \left(\mathbb T^{\mathrm{pre}}_{U/V}\right)^{\mathrm{aft}}
    \simeq
    \mathbb T^{\mathrm{pc}}_{B/R}.}
\]
Postcomposition with $u$ induces a natural morphism
\[
    \mathbb T^{\mathrm{pre}}_{U/V}
    \longrightarrow
    \mathbb T^{\mathrm{pre}}_{Y/V,u}.
\]
Define
\[
    \mathbb T^{\mathrm{pre}}_{U/Y;V}
    :=
    \operatorname{fib}
    \left(
        \mathbb T^{\mathrm{pre}}_{U/V}
        \longrightarrow
        \mathbb T^{\mathrm{pre}}_{Y/V,u}
    \right)
    \in\mathcal D_B
\]
and
\[
\boxed{
    \mathbb T^{\mathrm{pc}}_{U/Y;V}
    :=
    \left(\mathbb T^{\mathrm{pre}}_{U/Y;V}\right)^{\mathrm{aft}}.}
\]
The composite
\[
    \mathbb T^{\mathrm{pc}}_{U/Y;V}
    \longrightarrow
    \mathbb T^{\mathrm{pre}}_{U/Y;V}
    \longrightarrow
    \mathbb T^{\mathrm{pre}}_{U/V}
\]
factors, by the terminal universal property of the almost-finite
approximation of $\mathbb T^{\mathrm{pre}}_{U/V}$, through a canonical anchor
\[
\boxed{
    \mathbb T^{\mathrm{pc}}_{U/Y;V}[1]
    \longrightarrow
    \mathbb T^{\mathrm{pc}}_{B/R}[1].}
\]
\end{construction}

\begin{proof}
For connective coherent $M$ and $n\geq0$, the split extension
$B\oplus M[n]\to B$ is spectral Artinian.  The displayed identity
\[
    B\oplus M[n]
    \simeq
    B\times_{B\oplus M[n+1]}B
\]
is the structured square-zero pullback associated to the zero derivation.
Since $F_{U/Y}(B)\simeq*$, the Schlessinger equation turns this pullback into
\[
    F_{U/Y}(B\oplus M[n])
    \simeq
    \Omega F_{U/Y}(B\oplus M[n+1]).
\]
The corresponding split square-zero pullbacks for finite direct sums and
cofiber sequences of coherent coefficients give the standard additive and
exactness diagrams for this spectrum-valued tangent functor.  Hence the
stable coherent closure carries an exact functor, and the same Postnikov
right-left extension used in the pro-coherent coefficient construction gives
the asserted object of $\mathcal D_B$.

For $Y=U$, a point of the square-zero lifting space is a section of
$B\oplus M\to B$ in connective commutative $R$-algebras.  The relative
cotangent universal property identifies its stabilized mapping spectrum with
$\operatorname{map}_B(L_{B/R},M)$.  The pro-coherent dual
$(L_{B/R})^\vee$ is, by the almost-finite approximation theorem, the terminal
almost-finitely-presented functor mapping to this corepresentable convergent
functor.  This gives
$\left(\mathbb T^{\mathrm{pre}}_{U/V}\right)^{\mathrm{aft}}
\simeq\mathbb T^{\mathrm{pc}}_{B/R}$.

Postcomposition of a split infinitesimal lift to $U$ with $u$ gives the map
of pre-tangent functors.  Exactness gives its fibre.  The final anchor is the
unique factorization supplied by the terminal almost-finite approximation of
the represented ambient pre-tangent.  No cotangent complex for $Y$ and no
commutation of continuous duality with coefficient base change is used.
\end{proof}


\begin{proposition}[Deformation-local evaluation of the pre-tangent]
\label{prop:bnla-pre-tangent-def-evaluation}
Let
\[
    V=\operatorname{Spec}R\longrightarrow S,
    \qquad
    U=\operatorname{Spec}B\longrightarrow V
\]
be coherent affines, let $Y\to S$ be $L^{\mathrm{Def}}_S$-local, and let
$u:U\to Y$ be a centre.  For every connective $B$-module $M$ and every
$n\geq0$, put
\[
\Theta^n_{Y/V,u}(M)
:=
\operatorname{fib}_{u}
\left(
    \operatorname{Map}_V
    \left(
        \operatorname{Spec}(B\oplus M[n]),
        Y\times_SV
    \right)
    \longrightarrow
    \operatorname{Map}_V(U,Y\times_SV)
\right).
\]
Then the split square-zero suspension maps make
$n\mapsto\Theta^n_{Y/V,u}(M)$ an $\Omega$-spectrum for every connective
$M$.  These spectra form an exact accessible Postnikov-convergent functor
\[
    \Theta^{\mathrm{sp}}_{Y/V,u}:
    \operatorname{Mod}_B
    \longrightarrow
    \operatorname{Sp}.
\]
The canonical comparison from the Postnikov right-left extension of
Construction~\ref{con:bnla-coherent-pre-tangent} is an equivalence
\[
\boxed{
    \mathbb T^{\mathrm{pre}}_{Y/V,u}
    \xrightarrow{\sim}
    \Theta^{\mathrm{sp}}_{Y/V,u}.}
\]
Thus, for a deformation-local target, the constructed pre-tangent is computed
by the actual split square-zero lifting spectrum on every connective
coefficient, not only on coherent coefficients.
\end{proposition}

\begin{proof}
The split identity
\[
    B\oplus M[n]
    \simeq
    B\times_{B\oplus M[n+1]}B
\]
is a nilpotent-excision square.  Hence
Theorem~\ref{thm:bnla-global-def-reflector}, followed by the homotopy fibre
over $u$, gives
\[
    \Theta^n_{Y/V,u}(M)
    \xrightarrow{\sim}
    \Omega\Theta^{n+1}_{Y/V,u}(M).
\]
This supplies the spectrum structure.

Let
\[
    M'\longrightarrow M\longrightarrow M''
\]
be a fibre sequence of connective $B$-modules.  The map
$\pi_0(M)\to\pi_0(M'')$ is surjective, and there is a pullback square
\[
\begin{tikzcd}
B\oplus M' \arrow[r] \arrow[d] &
B \arrow[d]\\
B\oplus M \arrow[r] &
B\oplus M''.
\end{tikzcd}
\]
The lower horizontal map is surjective on $\pi_0$ with square-zero kernel.
Nilpotent excision therefore identifies the corresponding square of
$Y$-mapping spaces as Cartesian.  Taking fibres over the centre and
stabilizing shows that
\[
    \Theta^{\mathrm{sp}}_{Y/V,u}(M')
    \longrightarrow
    \Theta^{\mathrm{sp}}_{Y/V,u}(M)
    \longrightarrow
    \Theta^{\mathrm{sp}}_{Y/V,u}(M'')
\]
is a fibre sequence.  Thus the connective construction is exact and extends
uniquely across shifts to the stable module category.

We next verify Postnikov convergence.  For connective $M$, deformation
locality gives
\[
\operatorname{Map}_S(\operatorname{Spec}(B\oplus M),Y)
\simeq
\lim_k
\operatorname{Map}_S
\left(
    \operatorname{Spec}\tau_{\leq k}(B\oplus M),Y
\right).
\]
There are canonical equivalences
\[
    \tau_{\leq k}(B\oplus M)
    \simeq
    \tau_{\leq k}B\oplus\tau_{\leq k}M.
\]
Applying the same formula to $B\oplus\tau_{\leq n}M$ and taking the iterated
limit in $k$ and $n$, the tower is eventually constant in the $M$-variable
for every fixed $k$.  Interchanging the two limits therefore gives
\[
    \Theta^{\mathrm{sp}}_{Y/V,u}(M)
    \xrightarrow{\sim}
    \lim_n
    \Theta^{\mathrm{sp}}_{Y/V,u}(\tau_{\leq n}M).
\]
Accessibility follows from the same bounded Postnikov presentation used in
Construction~\ref{con:bnla-coherent-pre-tangent}: in the fixed universe the
square-zero evaluation functors have a common accessibility rank, and finite
limits and stabilization preserve accessibility.

On connective coherent modules this construction is exactly the spectrum
constructed in
Construction~\ref{con:bnla-coherent-pre-tangent}.  Both objects are exact,
accessible, and Postnikov-convergent, and the latter construction is the
Postnikov right-left extension of that coherent restriction.  Its universal
property therefore gives the displayed equivalence.
\end{proof}


\begin{proposition}[Cartesian family of spectral pre-tangents]
\label{prop:bnla-pre-tangent-cartesian-family}
Let $Y\to S$ be $L^{\mathrm{Def}}_S$-local, let
\[
    V=\operatorname{Spec}R\longrightarrow S,
    \qquad
    U=\operatorname{Spec}B\longrightarrow V
\]
be a coherent affine centre $u:U\to Y$, and let
\[
\begin{tikzcd}
R \arrow[r] \arrow[d] & R' \arrow[d]\\
B \arrow[r,"f"'] & B'
\end{tikzcd}
\]
be a morphism of coherent affine-etale pairs.  Write
\[
    u':U'=\operatorname{Spec}B'\longrightarrow Y
\]
for the restricted centre.  Precomposition with restriction of scalars gives
\[
    f_{\mathcal D}^*:
    \mathcal D_B\longrightarrow\mathcal D_{B'}.
\]
There are canonical equivalences
\[
\boxed{
    f_{\mathcal D}^*
    \mathbb T^{\mathrm{pre}}_{Y/V,u}
    \xrightarrow{\sim}
    \mathbb T^{\mathrm{pre}}_{Y/V',u'}}
\]
and
\[
\boxed{
    f_{\mathcal D}^*
    \mathbb T^{\mathrm{pre}}_{U/Y;V}
    \xrightarrow{\sim}
    \mathbb T^{\mathrm{pre}}_{U'/Y;V'}.}
\]
They preserve identities, composition, and all higher affine-etale
pullback-pasting coherences.
\end{proposition}

\begin{proof}
Let $M'$ be a connective $B'$-module.  By
Proposition~\ref{prop:bnla-pre-tangent-def-evaluation}, both pre-tangent
functors are computed on this coefficient and its restriction by the actual
split square-zero lifting spectra.  There is a canonical pullback square of
connective commutative ring spectra
\[
\begin{tikzcd}
B\oplus\operatorname{Res}^{B'}_BM'
  \arrow[r] \arrow[d] &
B \arrow[d,"f"]\\
B'\oplus M' \arrow[r] &
B',
\end{tikzcd}
\]
where the lower horizontal arrow is the split square-zero augmentation.  The
latter is a nilpotent extension: it is surjective on $\pi_0$ with square-zero
kernel.  Since $Y$ is $L^{\mathrm{Def}}_S$-local,
Theorem~\ref{thm:bnla-global-def-reflector} gives a canonical pullback square
of mapping spaces
\[
\begin{tikzcd}
Y(B\oplus\operatorname{Res}^{B'}_BM')
  \arrow[r] \arrow[d] &
Y(B) \arrow[d]\\
Y(B'\oplus M') \arrow[r] &
Y(B').
\end{tikzcd}
\]
Taking the homotopy fibre over the compatible centres $u$ and $u'$ gives
\[
\boxed{
    \Theta_{Y/V,u}
    \left(\operatorname{Res}^{B'}_BM'\right)
    \xrightarrow{\sim}
    \Theta_{Y/V',u'}(M').}
\]
The same calculation applies to every nonnegative suspension of $M'$ and
is natural in morphisms of modules.  Proposition~\ref{prop:bnla-pre-tangent-def-evaluation}
identifies these stabilized lifting spectra with the values of the constructed
pre-tangents.  Hence
\[
    f_{\mathcal D}^*
    \mathbb T^{\mathrm{pre}}_{Y/V,u}
    \xrightarrow{\sim}
    \mathbb T^{\mathrm{pre}}_{Y/V',u'}
\]
in $\mathcal D_{B'}$.

For the represented ambient tangent, two-variable spectral-etale cotangent
base change gives
\[
    L_{B'/R'}
    \simeq
    B'\otimes_BL_{B/R}.
\]
Corepresenting this equivalence gives
\[
    f_{\mathcal D}^*
    \mathbb T^{\mathrm{pre}}_{U/V}
    \xrightarrow{\sim}
    \mathbb T^{\mathrm{pre}}_{U'/V'}.
\]
The two comparisons commute with postcomposition by the centre map to $Y$.
Since $f_{\mathcal D}^*$ is exact, taking their fibres gives the second boxed
equivalence.  Every comparison is induced by canonical pullback,
restriction-of-scalars, and cotangent base-change equivalences, so the
identity, composition, and higher-pasting coherences are inherited from those
constructions.
\end{proof}


\begin{lemma}[Anchored right-left pullback]
\label{lem:bnla-anchored-right-left-pullback}
Let
\[
    \iota_B:
    \operatorname{QC}^{\vee}_B
    \hookrightarrow
    \mathcal D_B
\]
be the fully faithful right-left inclusion and let
\[
    r_B:\mathcal D_B\longrightarrow\operatorname{QC}^{\vee}_B,
    \qquad
    F\longmapsto F^{\mathrm{aft}}
\]
be its right adjoint.  Let
\[
    T\in\mathcal D_B,
    \qquad
    T^{\mathrm{aft}}\longrightarrow T
\]
be the terminal almost-finitely-presented approximation.  Composition with
this counit induces an equivalence
\[
\boxed{
    \left(\operatorname{QC}^{\vee}_B\right)_{/T^{\mathrm{aft}}}
    \xrightarrow{\sim}
    \left\{
        (E\to T)\in(\mathcal D_B)_{/T}
        \mid E\in\operatorname{QC}^{\vee}_B
    \right\}.}
\]

Now let
\[
    P:\mathcal D_B\longrightarrow\mathcal D_{B'}
\]
be an exact functor carrying pro-coherent objects to pro-coherent objects, let
\[
    q:T'\longrightarrow P(T)
\]
be a morphism, and put
\[
    H:=(P(T))^{\mathrm{aft}}.
\]
Let $T'^{\mathrm{aft}}\to T'$ be the almost-finite approximation and assume
the two maps
\[
    P(T^{\mathrm{aft}})\longrightarrow P(T),
    \qquad
    T'^{\mathrm{aft}}\longrightarrow T'\xrightarrow{q}P(T)
\]
are factored through $H$ by its terminal universal property.  Then the
cartesian right-left transport of
\[
    E\longrightarrow T^{\mathrm{aft}}
\]
along $q$ is represented by
\[
\boxed{
    P(E)\times_H T'^{\mathrm{aft}}
    \longrightarrow
    T'^{\mathrm{aft}}.}
\]
This object is characterized as the terminal pro-coherent object mapping to
the ordinary pullback
\[
    P(E)\times_{P(T)}T'
\]
in $\mathcal D_{B'}$.  The construction is coherent under composition of such
squares.
\end{lemma}

\begin{proof}
For the first assertion, if $E$ is pro-coherent, every map
\[
    E\longrightarrow T
\]
factors through $T^{\mathrm{aft}}\to T$ through a contractible space, by the
terminal universal property of $T^{\mathrm{aft}}$ in
$(\operatorname{QC}^{\vee}_B)_{/T}$.  This proves full faithfulness and
essential surjectivity of the displayed functor.

For the second assertion put
\[
    K:=P(E)\times_{P(T)}T'.
\]
The object
\[
    K^{\mathrm{rl}}
    :=
    P(E)\times_H T'^{\mathrm{aft}}
\]
is pro-coherent because the pro-coherent category is stable and closed under
finite limits.  The two counit maps give a canonical morphism
\[
    K^{\mathrm{rl}}\longrightarrow K.
\]
Let $Z$ be pro-coherent and let $Z\to K$ be any morphism.  Its two components
give compatible maps
\[
    Z\longrightarrow P(E),
    \qquad
    Z\longrightarrow T'
\]
over $P(T)$.  The second map factors contractibly through
$T'^{\mathrm{aft}}$, while their common composite to $P(T)$ factors
contractibly through $H$.  By uniqueness of these two terminal
factorizations, the resulting maps to $P(E)$ and $T'^{\mathrm{aft}}$ agree
over $H$.  Hence $Z\to K$ factors through $K^{\mathrm{rl}}$ through a
contractible space.  Thus $K^{\mathrm{rl}}\to K$ is terminal among
pro-coherent objects over $K$.

The construction uses only terminal right-left factorizations and finite
pullbacks.  Their uniqueness and pullback pasting give the identity,
composition, and higher coherences.
\end{proof}


\begin{construction}[Anchored pre-tangent transport]
\label{con:bnla-anchored-pre-tangent-transport}
Let
\[
\begin{tikzcd}
R \arrow[r] \arrow[d] & R' \arrow[d]\\
B \arrow[r,"f"'] & B'
\end{tikzcd}
\]
be a morphism of coherent affine-etale pairs.  Restriction of scalars on modules
induces an exact coefficient-change functor
\[
    f_{\mathcal D}^*:\mathcal D_B\longrightarrow\mathcal D_{B'},
    \qquad
    (f_{\mathcal D}^*F)(M'):=F(\operatorname{Res}^{B'}_B M').
\]
It carries almost-finitely-presented functors to
almost-finitely-presented functors, and under
$\operatorname{QC}^{\vee}\simeq\mathcal D^{\mathrm{aft}}$ its restriction is
ordinary pro-coherent pullback
\[
    f^*:\operatorname{QC}^{\vee}_B
    \longrightarrow
    \operatorname{QC}^{\vee}_{B'}.
\]

Cotangent transitivity and the two-variable base-change map give a natural
morphism
\[
    \mathbb T^{\mathrm{pre}}_{U'/V'}
    \longrightarrow
    f_{\mathcal D}^*\mathbb T^{\mathrm{pre}}_{U/V}.
\]
Put
\[
    H_f
    :=
    \left(
        f_{\mathcal D}^*\mathbb T^{\mathrm{pre}}_{U/V}
    \right)^{\mathrm{aft}}
    \in\operatorname{QC}^{\vee}_{B'}.
\]
The terminal universal property gives canonical morphisms
\[
    f^*\mathbb T^{\mathrm{pc}}_{B/R}
    \longrightarrow H_f,
    \qquad
    \mathbb T^{\mathrm{pc}}_{B'/R'}
    \longrightarrow H_f.
\]
For an anchored pro-coherent object
\[
    \rho:E\longrightarrow\mathbb T^{\mathrm{pc}}_{B/R}[1]
\]
define
\[
\boxed{
    f_\sharp(\rho)
    :=
    \left(
        f^*E
        \times_{H_f[1]}
        \mathbb T^{\mathrm{pc}}_{B'/R'}[1]
        \longrightarrow
        \mathbb T^{\mathrm{pc}}_{B'/R'}[1]
    \right).}
\]
The construction is coherent for identity and composition.  If the morphism
of affine pairs is spectral etale, then
\[
    \mathbb T^{\mathrm{pre}}_{U'/V'}
    \xrightarrow{\sim}
    f_{\mathcal D}^*\mathbb T^{\mathrm{pre}}_{U/V}
\]
by spectral etale cotangent base change.  Hence
$\mathbb T^{\mathrm{pc}}_{B'/R'}\xrightarrow{\sim}H_f$ and
\[
\boxed{
    U(f_\sharp\rho)
    \simeq
    f^*E.}
\]
Thus the anchored pre-tangent transport reduces to ordinary pro-coherent
pullback on underlying sources in the spectral etale case.
\end{construction}

\begin{proof}
Precomposition with restriction of scalars preserves exactness,
accessibility, and Postnikov convergence.  If $F$ is almost finitely
presented, spectral-etale restriction of scalars preserves the boundedness
classes and filtered-colimit tests entering that condition, so
$f_{\mathcal D}^*F$ is again almost finitely presented.  On a continuous
dual one has
\[
\begin{aligned}
(f_{\mathcal D}^*M^\vee)(K')
&=\operatorname{map}_B(M,\operatorname{Res}^{B'}_BK')\\
&\simeq
\operatorname{map}_{B'}(B'\otimes_BM,K'),
\end{aligned}
\]
which is the pro-coherent pullback of
Lemma~\ref{lem:bnla-procoherent-finite-tor-bc}.  The relative right-left
construction extends this formula from continuous duals to the whole
pro-coherent fibre.

The cotangent morphism
\[
    B'\otimes_BL_{B/R}
    \longrightarrow
    L_{B'/R'}
\]
induces, after corepresenting, the displayed map of ambient pre-tangent
functors in the opposite direction.  Since both
$f^*\mathbb T^{\mathrm{pc}}_{B/R}$ and
$\mathbb T^{\mathrm{pc}}_{B'/R'}$ are almost finitely presented and map to
$f_{\mathcal D}^*\mathbb T^{\mathrm{pre}}_{U/V}$, terminality of $H_f$
produces the two comparison maps.

Apply Lemma~\ref{lem:bnla-anchored-right-left-pullback} with
\[
    P=f_{\mathcal D}^*,
    \qquad
    T=\mathbb T^{\mathrm{pre}}_{U/V}[1],
    \qquad
    T'=\mathbb T^{\mathrm{pre}}_{U'/V'}[1].
\]
The lemma identifies the cartesian right-left representative of the ordinary
pullback in the anchored overcategory with
\[
    f^*E
    \times_{H_f[1]}
    \mathbb T^{\mathrm{pc}}_{B'/R'}[1].
\]
This is exactly the formula defining $f_\sharp$.  In particular, the formula
does not assert that the fibrewise right adjoints
$(-)^{\mathrm{aft}}$ commute with coefficient change; it constructs the
cartesian transport in the anchored overcategory by the terminal
right-left factorization.

For a morphism of coherent affine-etale pairs,
Theorem~\ref{thm:bnla-two-variable-etale-base-change} gives
\[
    L_{B'/R'}\simeq B'\otimes_BL_{B/R}.
\]
Hence the ambient pre-tangent comparison is an equivalence.  Its terminal
almost-finite approximation is
$\mathbb T^{\mathrm{pc}}_{B'/R'}$, so the cartesian pullback above has
underlying source $f^*E$.  Cartesian lifts in an overcategory are coherent
for identities, composition, and higher pasting; these coherences give the
claimed coherence of $f_\sharp$.
\end{proof}

\begin{theorem}[Pre-tangent identification of the affine leaf controller]
\label{thm:bnla-pre-tangent-leaf-controller}
Let $Y$ be $L_S^{\mathrm{ArtDef}}$-local and let
$u:U=\operatorname{Spec}B\to Y$ be a coherent affine centre over
$V=\operatorname{Spec}R$.  Let
\[
    F_{U/Y}\in\operatorname{FMP}^{\mathrm{sp}}_{B/R}
\]
be its centred Artinian leaf functor.  Then the underlying anchored
pro-coherent object of
$\operatorname{Diff}^{\pi,E_\infty}_{B/R}(F_{U/Y})$ is canonically
\[
\boxed{
    \mathbb T^{\mathrm{pc}}_{U/Y;V}[1]
    \longrightarrow
    \mathbb T^{\mathrm{pc}}_{B/R}[1].}
\]
The identification is natural in maps of Artinian-local targets and coherent
centres.
\end{theorem}

\begin{proof}
It is enough to compare the two anchored objects on the dense dually
almost-perfect generators.  Let
\[
    \alpha:L_{B/R}[-1]\longrightarrow M
\]
with $M$ connective coherent and let $B_\alpha\to B$ be its structured
square-zero extension.  The defining pullback of $B_\alpha$ says exactly that
a lift of the centre through $\operatorname{Spec}B_\alpha$ is a lift of the
point $\alpha$ through the morphism of pre-tangent functors
\[
    \mathbb T^{\mathrm{pre}}_{U/V}[1]
    \longrightarrow
    \mathbb T^{\mathrm{pre}}_{Y/V,u}[1].
\]
By the terminal universal property of the almost-finite approximation, this
lifting space is the mapping space from the anchored generator
\[
    M^\vee\xrightarrow{\alpha^\vee}
    \mathbb T^{\mathrm{pc}}_{B/R}[1]
\]
to the anchored object
$\mathbb T^{\mathrm{pc}}_{U/Y;V}[1]\to
\mathbb T^{\mathrm{pc}}_{B/R}[1]$.

On the other hand, Theorem~\ref{thm:bnla-spectral-affine-recognition}, using
its canonical complete-almost-finite extension, identifies the same
square-zero lifting space with the mapping space from that anchored generator
to the underlying anchored object of
$\operatorname{Diff}^{\pi,E_\infty}_{B/R}(F_{U/Y})$.  Continuous duals of
connective coherent modules generate the pro-coherent category, so the two
anchored objects are canonically equivalent.  All maps used in the
comparison are functorial, which gives naturality.
\end{proof}

\begin{lemma}[Formal-etale inheritance of Artinian-deformation locality]
\label{lem:bnla-artdef-formal-etale-inheritance}
Let
\[
    h:Y'\longrightarrow Y
\]
be formally etale on finite spectral infinitesimal tests.  If $Y$ is
$L^{\mathrm{ArtDef}}_S$-local, then $Y'$ is
$L^{\mathrm{ArtDef}}_S$-local.
\end{lemma}

\begin{proof}
Fix an Artinian-deformation generator
\[
\begin{tikzcd}
A_\eta \arrow[r] \arrow[d] & B \arrow[d]\\
A \arrow[r,"\eta"'] & B\oplus M[1].
\end{tikzcd}
\]
By Proposition~\ref{prop:bnla-spectral-art-square-zero-generation}, the
augmentation $A\to B$ is a finite spectral infinitesimal extension.  The
split square-zero augmentation $B\oplus M[1]\to B$ is one-stage finite
infinitesimal, and the pullback $A_\eta\to B$ is again finite infinitesimal.
Formal etaleness therefore gives canonical equivalences
\[
    Y'(C)
    \simeq
    Y(C)\times_{Y(B)}Y'(B),
    \qquad
    C\in\{A,\,B\oplus M[1],\,A_\eta\}.
\]
Pullback pasting then gives
\[
\begin{aligned}
    Y'(A)\times_{Y'(B\oplus M[1])}Y'(B)
    &\simeq
    \left(
        Y(A)\times_{Y(B\oplus M[1])}Y(B)
    \right)
    \times_{Y(B)}Y'(B)\\
    &\simeq
    Y(A_\eta)\times_{Y(B)}Y'(B)\\
    &\simeq
    Y'(A_\eta).
\end{aligned}
\]
The middle equivalence is the Artinian-deformation locality of $Y$.  Thus
$Y'$ is local for every generator.
\end{proof}

\begin{theorem}[Etale centre transport and formally-etale target invariance]
\label{thm:bnla-pre-tangent-etale-base-change}
Let
\[
    f:(U',V')\longrightarrow(U,V)
\]
be a morphism of coherent affine-etale pairs, and let $Y$ be an
$L_S^{\mathrm{Def}}$-local target equipped with a centre $u:U\to Y$.
Restriction of $u$ gives a centre $u':U'\to Y$.  Then relative tangent
formation is cartesian for the anchored transport of
Construction~\ref{con:bnla-anchored-pre-tangent-transport}:
\[
\boxed{
    f_\sharp
    \left(
        \mathbb T^{\mathrm{pc}}_{U/Y;V}[1]
        \longrightarrow
        \mathbb T^{\mathrm{pc}}_{B/R}[1]
    \right)
    \xrightarrow{\sim}
    \left(
        \mathbb T^{\mathrm{pc}}_{U'/Y;V'}[1]
        \longrightarrow
        \mathbb T^{\mathrm{pc}}_{B'/R'}[1]
    \right).}
\]
In particular, because $f$ is spectral etale, the underlying source
comparison is the pro-coherent pullback equivalence
\[
\boxed{
    f^*\mathbb T^{\mathrm{pc}}_{U/Y;V}[1]
    \xrightarrow{\sim}
    \mathbb T^{\mathrm{pc}}_{U'/Y;V'}[1].}
\]

Moreover, if $h:Y'\to Y$ is formally etale on finite spectral infinitesimal
tests and $u':U'\to Y'$ is a lift of a centre $U'\to Y$, then
Lemma~\ref{lem:bnla-artdef-formal-etale-inheritance} makes $Y'$ automatically
$L^{\mathrm{ArtDef}}_S$-local, so both relative pre-tangent constructions are
defined.  Postcomposition with $h$ induces a canonical equivalence
\[
\boxed{
    \mathbb T^{\mathrm{pc}}_{U'/Y';V'}
    \xrightarrow{\sim}
    \mathbb T^{\mathrm{pc}}_{U'/Y;V'}.}
\]
All comparisons are coherent for identities, composition, and Cartesian
pasting.
\end{theorem}

\begin{proof}
We first vary the centre while keeping the nonlinear target fixed.
Proposition~\ref{prop:bnla-pre-tangent-cartesian-family} gives the canonical
equivalence of relative pre-tangent objects
\[
    f_{\mathcal D}^*
    \mathbb T^{\mathrm{pre}}_{U/Y;V}
    \xrightarrow{\sim}
    \mathbb T^{\mathrm{pre}}_{U'/Y;V'}.
\]
Construction~\ref{con:bnla-anchored-pre-tangent-transport} identifies the
cartesian lift in the fibrewise right-left anchored overcategory with
$f_\sharp$.  Applying that cartesian lift to the almost-finitely-presented
relative tangent object over $(U,V)$ therefore gives the relative tangent
object over $(U',V')$.  This proves the first displayed equivalence.  Notice
that this argument forms the cartesian object in the anchored overcategory;
it does not commute the fibrewise right adjoints
$(-)^{\mathrm{aft}}$ with coefficient change.

For a morphism of coherent affine-etale pairs, the ambient pre-tangent
comparison is an equivalence by two-variable spectral-etale cotangent base
change.  Construction~\ref{con:bnla-anchored-pre-tangent-transport} therefore
identifies $f_\sharp$ with ordinary $f^*$ on the underlying pro-coherent
source, giving the second displayed equivalence.

Now let $h:Y'\to Y$ be formally etale.  By
Lemma~\ref{lem:bnla-artdef-formal-etale-inheritance}, $Y'$ is
$L^{\mathrm{ArtDef}}_S$-local.  For every finite spectral infinitesimal
thickening of the centre, formal etaleness identifies the space of dotted
lifts to $Y'$ with the space of dotted lifts to $Y$ over the specified
centre.  Hence the square-zero lifting functors $\Theta$ are canonically
equivalent.  Their spectrum objects, Postnikov right-left
extensions, relative pre-tangent fibres, and terminal almost-finite
approximations are therefore canonically equivalent.  This proves target
invariance.  Since every comparison is produced by restriction, universal
lifting, and terminal approximation, the identity, composition, and
Cartesian-pasting coherences are inherited from those constructions.
\end{proof}

\subsection{Centred completion and etale tangent base change}

For a map $u:U\to Y$ in $\mathcal X_{/S}$ define its reduced-centred
completion by
\[
\boxed{
    C_U(Y)
    :=
    Y\times_{D_{\mathrm{red}}Y}D_{\mathrm{red}}U.}
\]

\begin{lemma}[Artinian-deformation locality of reduced-centred completion]
\label{lem:bnla-artdef-centred-locality}
Let $u:U\to Y$ be a morphism in $\mathcal X_{/S}$.  If $Y$ is
$L^{\mathrm{ArtDef}}_S$-local, then
\[
    C_U(Y)=Y\times_{D_{\mathrm{red}}Y}D_{\mathrm{red}}U
\]
is $L^{\mathrm{ArtDef}}_S$-local.
\end{lemma}

\begin{proof}
It is enough to verify locality on the generators of
Definition~\ref{def:bnla-global-artdef-generators}.  Fix a Schlessinger
square
\[
\begin{tikzcd}
A_\eta \arrow[r] \arrow[d] & B \arrow[d]\\
A \arrow[r,"\eta"'] & B\oplus M[1].
\end{tikzcd}
\]
The four connective rings
$A_\eta$, $A$, $B\oplus M[1]$, and $B$ have canonically equivalent reduced
classical truncations: the Artinian augmentation has nilpotent degree-zero
kernel and $B\oplus M[1]$ has the same degree-zero ring as $B$.  For the
pullback $A_\eta$, the underlying fibre sequence
\[
    A_\eta\longrightarrow A\oplus B\longrightarrow B\oplus M[1]
\]
has surjective degree-zero right-hand map because its $B$-component is the
identity.  Hence $A_\eta$ is connective and the degree-zero kernel identifies
canonically with $\pi_0(A)$ via the graph of
$\pi_0(A)\to\pi_0(B)$.  Thus all four reduced degree-zero rings are
canonically $\pi_0(B)_{\mathrm{red}}$.  Put
\[
    C_0:=H\!\left(\pi_0(B)_{\mathrm{red}}\right)
\]
for the common reduced classical coefficient ring.  Every
morphism in the square induces the identity of $C_0$ under these canonical
identifications.

By the definition of $D_{\mathrm{red}}$, evaluation of the reduced-centred
pullback on any one of the four rings $C$ is
\[
    C_U(Y)(C)
    \simeq
    Y(C)\times_{Y(C_0)}U(C_0).
\]
Therefore pullback pasting gives
\[
\begin{aligned}
&C_U(Y)(A)
\times_{C_U(Y)(B\oplus M[1])}
C_U(Y)(B)
\\
&\qquad\simeq
\left(
    Y(A)\times_{Y(B\oplus M[1])}Y(B)
\right)
\times_{Y(C_0)}U(C_0).
\end{aligned}
\]
Since $Y$ is $L^{\mathrm{ArtDef}}_S$-local, its Artinian-deformation equation
identifies the parenthesized factor with $Y(A_\eta)$.  Hence
\[
\begin{aligned}
&C_U(Y)(A)
\times_{C_U(Y)(B\oplus M[1])}
C_U(Y)(B)
\\
&\qquad\simeq
Y(A_\eta)\times_{Y(C_0)}U(C_0)
\simeq
C_U(Y)(A_\eta).
\end{aligned}
\]
Thus $C_U(Y)$ is local for every Artinian-deformation generator, which proves
the claim.
\end{proof}

\begin{lemma}[Transitivity of reduced-centred completion]
\label{lem:bnla-artdef-centred-transitivity}
For composable maps
\[
    U'\longrightarrow U\longrightarrow Y
\]
there is a canonical equivalence
\[
\boxed{
    C_{U'}(C_U(Y))
    \xrightarrow{\sim}
    C_{U'}(Y).}
\]
If $Y$ is $L^{\mathrm{ArtDef}}_S$-local, then both sides are
$L^{\mathrm{ArtDef}}_S$-local.
\end{lemma}

\begin{proof}
Left exactness and idempotence of $D_{\mathrm{red}}$ give
\[
\begin{aligned}
C_{U'}(C_U(Y))
&\simeq
\left(
Y\times_{D_{\mathrm{red}}Y}D_{\mathrm{red}}U
\right)
\times_{D_{\mathrm{red}}U}
D_{\mathrm{red}}U'
\\
&\simeq
Y\times_{D_{\mathrm{red}}Y}D_{\mathrm{red}}U'
\\
&=
C_{U'}(Y).
\end{aligned}
\]
If $Y$ is $L^{\mathrm{ArtDef}}_S$-local, then
Lemma~\ref{lem:bnla-artdef-centred-locality} makes $C_{U'}(Y)$ local.  The
right-hand side of the displayed equivalence is therefore local, and the
left-hand side is local because it is equivalent to that object.  Equivalently,
one may first apply the same lemma to $C_U(Y)$, which is itself local by
Lemma~\ref{lem:bnla-artdef-centred-locality}.
\end{proof}

\begin{lemma}[Centred completion does not change the Artinian leaf]
\label{lem:bnla-centred-completion-leaf-invariance}
Let $u:U=\operatorname{Spec}B\to Y$ and let $A\to B$ be spectral Artinian.
Then the canonical map
\[
\begin{aligned}
&\operatorname{fib}_{u}
\left(
    \operatorname{Map}_S(\operatorname{Spec}A,C_U(Y))
    \to
    \operatorname{Map}_S(U,C_U(Y))
\right)
\\
&\qquad\longrightarrow
\operatorname{fib}_{u}
\left(
    \operatorname{Map}_S(\operatorname{Spec}A,Y)
    \to
    \operatorname{Map}_S(U,Y)
\right)
\end{aligned}
\]
is an equivalence.
\end{lemma}

\begin{proof}
A spectral Artinian thickening has the same reduced classical affine as its
centre:
\[
    D_{\mathrm{red}}\operatorname{Spec}A
    \simeq
    D_{\mathrm{red}}U.
\]
Mapping into the pullback
$C_U(Y)=Y\times_{D_{\mathrm{red}}Y}D_{\mathrm{red}}U$
therefore adds a reduced component which is uniquely forced by the prescribed
map on the centre.  The two homotopy fibres are consequently equivalent.
\end{proof}

\begin{lemma}[Reduced-classical Cartesianity of spectral etale maps]
\label{lem:bnla-dred-etale-cartesian}
Let
\[
    f:U'=\operatorname{Spec}(B')
    \longrightarrow
    U=\operatorname{Spec}(B)
\]
be a spectral etale morphism of affine spectral schemes.  Then the naturality
square of the reduced-classical localization unit
\[
\begin{tikzcd}
U' \arrow[r,"f"] \arrow[d] & U \arrow[d]\\
D_{\mathrm{red}}U' \arrow[r,"D_{\mathrm{red}}f"'] &
D_{\mathrm{red}}U
\end{tikzcd}
\]
is Cartesian in the native spectral Nisnevich $\infty$-topos.  The
comparison is natural in spectral etale morphisms and compatible with
composition and base change.
\end{lemma}

\begin{proof}
Evaluate the square on an arbitrary affine spectral test
$T=\operatorname{Spec}(C)$ and put
\[
    T_{\mathrm{red}}
    :=
    \operatorname{Spec}H(\pi_0(C)_{\mathrm{red}}).
\]
By the definition of $D_{\mathrm{red}}$, Cartesianity on this test is the
assertion that
\[
\boxed{
    \operatorname{Map}(T,U')
    \xrightarrow{\sim}
    \operatorname{Map}(T,U)
    \times_{\operatorname{Map}(T_{\mathrm{red}},U)}
    \operatorname{Map}(T_{\mathrm{red}},U').}
\]
Fix a point $x:T\to U$ and form the spectral etale base change
\[
    V:=T\times_UU'\longrightarrow T.
\]
The homotopy fibre of the left-hand restriction map over $x$ is the space of
sections of $V\to T$, while the corresponding fibre on the right is the
space of sections of
\[
    V\times_TT_{\mathrm{red}}
    \longrightarrow
    T_{\mathrm{red}}.
\]

Chapter~3 identifies spectral etale geometry over $T$ with ordinary etale
geometry over $\operatorname{Spec}\pi_0(C)$.  Ordinary etale geometry is
unchanged by reduction.  For completeness, write
$R=\pi_0(C)$ and $N=\sqrt{0}\subset R$.  Since an etale algebra and a
morphism between etale algebras are finitely presented, each descends along
\[
    R\longrightarrow R/N
    =\operatorname*{colim}_{J\subset N,\,J\text{ finitely generated}}R/J
\]
to some quotient $R/J$; after enlarging $J$ inside $N$, the descended
algebra may be taken etale because etaleness is a finitely presented
condition.  A finitely generated nil ideal $J$ is nilpotent; choose $m$ with
$J^m=0$.  The quotient $R\to R/J$ factors through the finite tower
\[
    R=R/J^m\longrightarrow R/J^{m-1}\longrightarrow\cdots\longrightarrow
    R/J,
\]
whose successive kernels $J^{k-1}/J^k$ are square-zero.  Viewed as discrete
spectral thickenings, Chapter~9's finite-infinitesimal etale-site rigidity
identifies the etale categories across every stage.  Thus etale objects and
their morphisms lift through the tower with contractible uniqueness.  It
follows that restriction induces an equivalence
\[
    \operatorname{Et}(T)
    \xrightarrow{\sim}
    \operatorname{Et}(T_{\mathrm{red}}).
\]
Under this equivalence the terminal etale object and the etale object
$V\to T$ are carried to their reductions.  Therefore restriction induces an
equivalence on section spaces
\[
    \operatorname{Map}_{\operatorname{Et}(T)}(T,V)
    \xrightarrow{\sim}
    \operatorname{Map}_{\operatorname{Et}(T_{\mathrm{red}})}
    \left(
        T_{\mathrm{red}},
        V\times_TT_{\mathrm{red}}
    \right).
\]
Thus the two fibres above are equivalent for every $x$, proving the boxed
mapping-space equivalence.  Yoneda gives the Cartesian square.  Naturality,
composition, and base-change compatibility follow from the corresponding
functoriality of etale base change and reduced classical truncation.
\end{proof}

\begin{proposition}[Etale base change of centred Artinian tangents]
\label{prop:bnla-centred-artinian-tangent-base-change}
Let
\[
\begin{tikzcd}
R \arrow[r] \arrow[d] & R' \arrow[d]\\
B \arrow[r,"f"'] & B'
\end{tikzcd}
\]
be a morphism of coherent affine-etale pairs, and let
$U'=\operatorname{Spec}B'\to U=\operatorname{Spec}B$ be the induced centre
map.  Let $Y$ be $L^{\mathrm{Def}}_S$-local.  Write
\[
    F_{U/Y}\in\operatorname{FMP}^{\mathrm{sp}}_{B/R},
    \qquad
    F_{U'/Y}\in\operatorname{FMP}^{\mathrm{sp}}_{B'/R'}
\]
for the centred Artinian leaf functors of $Y$.  If
\[
    E_{U/Y}[1]\longrightarrow
    \mathbb T^{\mathrm{pc}}_{B/R}[1]
\]
and
\[
    E_{U'/Y}[1]\longrightarrow
    \mathbb T^{\mathrm{pc}}_{B'/R'}[1]
\]
are the underlying anchored objects of their affine spectral
differentiations, then there is a canonical equivalence in the anchored
pro-coherent category
\[
\boxed{
    f^*E_{U/Y}[1]
    \xrightarrow{\sim}
    E_{U'/Y}[1].}
\]
The comparison is coherent for identity and composition of affine-etale
centre changes.
\end{proposition}

\begin{proof}
Theorem~\ref{thm:bnla-pre-tangent-leaf-controller} identifies the two
underlying anchored objects with the relative tangent coefficients
\[
    \mathbb T^{\mathrm{pc}}_{U/Y;V}[1]
    \longrightarrow
    \mathbb T^{\mathrm{pc}}_{B/R}[1]
\]
and
\[
    \mathbb T^{\mathrm{pc}}_{U'/Y;V'}[1]
    \longrightarrow
    \mathbb T^{\mathrm{pc}}_{B'/R'}[1].
\]
Because $Y$ is $L^{\mathrm{Def}}_S$-local,
Theorem~\ref{thm:bnla-pre-tangent-etale-base-change} identifies the second
with the $f_\sharp$-transport of the first.  For a morphism of coherent
affine-etale pairs,
Construction~\ref{con:bnla-anchored-pre-tangent-transport} identifies the
underlying source of $f_\sharp$ with ordinary pro-coherent pullback.  Hence
\[
    f^*E_{U/Y}[1]
    \xrightarrow{\sim}
    E_{U'/Y}[1].
\]
The identity, composition, and higher-pasting coherences are those of the
cartesian anchored pre-tangent transport.
\end{proof}

\subsection{The missing cocartesian comparison}

Let
\[
    f:(U',V')\longrightarrow(U,V)
\]
be a morphism in
$\operatorname{AffEt}^{\mathrm{coh}}(X/S)$, written in coordinate rings as
above.  Put
\[
    F:=\mathcal L^{\mathrm{fm,sp}}_{\mathfrak G,U/V},
    \qquad
    F':=\mathcal L^{\mathrm{fm,sp}}_{\mathfrak G,U'/V'}.
\]

Composition of a lift on an Artinian test with the canonical etale lift of
that test gives a natural transformation
\[
    F\longrightarrow f_{\mathrm{Art},!}^{\,*}F'.
\]
Adjunction and Schlessinger localization give the canonical centre-change
morphism
\[
\boxed{
    \lambda_f:
    f_{\mathrm{FMP},!}F
    \longrightarrow
    F'.}
\]

\begin{theorem}[Etale cocartesianity of the formal leaves]
\label{thm:bnla-formal-leaf-etale-cocartesian}
For every morphism of coherent affine-etale pairs, the canonical map
\[
\boxed{
    \lambda_f:
    f_{\mathrm{FMP},!}
    \mathcal L^{\mathrm{fm,sp}}_{\mathfrak G,U/V}
    \xrightarrow{\sim}
    \mathcal L^{\mathrm{fm,sp}}_{\mathfrak G,U'/V'}}
\]
is an equivalence.  These equivalences are coherent for identities,
composition, and higher etale pasting.
\end{theorem}

\begin{proof}
Apply affine spectral differentiation.  By
Theorem~\ref{thm:bnla-affine-spectral-recognition-etale},
\[
\operatorname{Diff}_{B'/R'}^{\pi,E_\infty}
\left(
    f_{\mathrm{FMP},!}F
\right)
\simeq
f^*
\operatorname{Diff}_{B/R}^{\pi,E_\infty}(F).
\]
Proposition~\ref{prop:bnla-centred-artinian-tangent-base-change} identifies
the underlying anchored pro-coherent object of the right-hand side with the
underlying anchored object of
$\operatorname{Diff}_{B'/R'}^{\pi,E_\infty}(F')$.  Under this
identification, the underlying map of
$\operatorname{Diff}(\lambda_f)$ is the canonical base-change equivalence
constructed in that proposition.

The forgetful functor from spectral partition-Lie algebroids to anchored
pro-coherent modules is conservative because the algebroid category is
monadic.  Therefore
$\operatorname{Diff}(\lambda_f)$ is an equivalence.  Affine differentiation
itself is an equivalence, so $\lambda_f$ is an equivalence.

The proof identifies $\lambda_f$ by the canonical derivation,
square-zero, cotangent, and finite-Tor Beck--Chevalley comparisons.  The
coherence already proved for those comparisons supplies the identity,
composition, and higher-pasting coherence.
\end{proof}

\subsection{The global spectral formal-leaf object}

\begin{definition}[Global spectral formal-leaf envelope]
\label{def:bnla-global-spectral-leaf-envelope}
The family
\[
    (U,V)\longmapsto
    \mathcal L^{\mathrm{fm,sp}}_{\mathfrak G,U/V},
\]
together with the cocartesian equivalences of
Theorem~\ref{thm:bnla-formal-leaf-etale-cocartesian}, is an etale global
section of $\mathscr{FMP}^{\mathrm{sp}}_{X/S}$.  Define
\[
\boxed{
    \mathcal L^{\mathrm{fm,sp}}_{\mathfrak G}
    \in
    \operatorname{FMP}^{\mathrm{sp}}_{X/S}}
\]
to be this global section.
\end{definition}

Thus the formal leaf is not obtained by independently reflecting affine
quotient-lifting functors and then adding descent.  It is extracted from the
single globally spectral-deformation-local target
$Q^{\mathrm{fm}}_{\mathfrak G}$, and its cocartesian transition maps are
theorem-level equivalences.

\begin{theorem}[Functoriality and base change of the global leaf envelope]
\label{thm:bnla-global-leaf-functoriality}
For fixed $X/S$, the assignment
\[
    \mathfrak G
    \longmapsto
    \mathcal L^{\mathrm{fm,sp}}_{\mathfrak G}
\]
is functorial on $\operatorname{FormLieGpd}(X/S)$.

More generally, let
\[
    F:(\mathfrak G,X)\longrightarrow(\mathfrak H,Y)
\]
be a functor of formal Lie groupoids over $S$ whose object map
$f:X\to Y$ is spectral etale, and suppose both sides remain in the locally
coherent coefficient sector.  Restriction of coherent affine-etale pairs
along $f$ defines
\[
    f^*:
    \operatorname{FMP}^{\mathrm{sp}}_{Y/S}
    \longrightarrow
    \operatorname{FMP}^{\mathrm{sp}}_{X/S},
\]
and $F$ induces canonically
\[
\boxed{
    \mathcal L^{\mathrm{fm,sp}}_{\mathfrak G}
    \longrightarrow
    f^*\mathcal L^{\mathrm{fm,sp}}_{\mathfrak H}.}
\]
These maps are coherently compatible with identities and composition of
spectral-etale object maps.

If $c:S'\to S$ is spectral etale and the base-changed geometry remains in the
locally coherent coefficient sector, then
\[
\boxed{
    c^*
    \mathcal L^{\mathrm{fm,sp}}_{\mathfrak G}
    \xrightarrow{\sim}
    \mathcal L^{\mathrm{fm,sp}}_{\mathfrak G\times_SS'}.}
\]
For a more general change of spectral base within the coefficient sector there
is the canonical mate comparison.  No functoriality statement for an
arbitrary varying object map $X\to Y$ is asserted by the present global
coefficient construction.
\end{theorem}

\begin{proof}
First fix $X/S$.  A morphism
$\mathfrak G\to\mathfrak H$ in
$\operatorname{FormLieGpd}(X/S)$ induces, by the quotient functoriality of the
preceding chapter, a morphism
\[
    Q_{\mathfrak G}\longrightarrow Q_{\mathfrak H}
\]
under $X$.  Functoriality of $L^{\mathrm{Def}}_S$, of $D_{\mathrm{red}}$,
and of finite pullbacks therefore gives a morphism
\[
    Q^{\mathrm{fm}}_{\mathfrak G}
    \longrightarrow
    Q^{\mathrm{fm}}_{\mathfrak H}
\]
under the common centre $X$.  Postcomposition gives a morphism on every
affine Artinian leaf.  Naturality of
Theorem~\ref{thm:bnla-formal-leaf-etale-cocartesian} makes these maps commute
with the cocartesian transition equivalences, so they assemble to a map of
global sections.  Identities and composition are inherited from the quotient,
localization, and mapping-space functors.

Now let $F:(\mathfrak G,X)\to(\mathfrak H,Y)$ cover a spectral etale
$f:X\to Y$.  Every coherent affine-etale pair
\[
    V=\operatorname{Spec}R\to S,
    \qquad
    U=\operatorname{Spec}B\to X\times_SV
\]
for $X/S$ is, by composition with $f\times_SV$, also a coherent affine-etale
pair for $Y/S$.  Hence restriction along this functor of affine-etale pair
categories defines the displayed $f^*$ on global formal-moduli sections.

The formal-groupoid functor induces a morphism of effective quotients
\[
    Q_{\mathfrak G}\longrightarrow Q_{\mathfrak H}
\]
whose restriction to the object spaces is $f$.  Applying
$L^{\mathrm{Def}}_S$, $D_{\mathrm{red}}$, and the defining reduced-centred
pullbacks gives a compatible morphism
\[
    Q^{\mathrm{fm}}_{\mathfrak G}
    \longrightarrow
    Q^{\mathrm{fm}}_{\mathfrak H}.
\]
For every coherent affine-etale pair $(U,V)$ of $X/S$, postcomposition
therefore gives
\[
    \mathcal L^{\mathrm{fm,sp}}_{\mathfrak G,U/V}
    \longrightarrow
    \mathcal L^{\mathrm{fm,sp}}_{\mathfrak H,U/V},
\]
where the target is the restriction of the $Y/S$ leaf to the same pair.
Naturality of the etale cocartesian comparison with respect to maps of
nonlinear targets makes these affine maps compatible with all pair
transitions.  They hence glue to
\[
    \mathcal L^{\mathrm{fm,sp}}_{\mathfrak G}
    \longrightarrow
    f^*\mathcal L^{\mathrm{fm,sp}}_{\mathfrak H}.
\]
The quotient functoriality and functoriality of restriction supply the
identity, composition, and higher coherence.

Finally let $c:S'\to S$ be spectral etale.  Pullback between the native
slices preserves colimits and finite limits, and
Theorem~\ref{thm:bnla-def-etale-base-change} gives the canonical equivalence
\[
\boxed{
    c^*L^{\mathrm{Def}}_S(Y)
    \simeq
    L^{\mathrm{Def}}_{S'}(c^*Y).}
\]
Chapter~9 gives the corresponding equivalence for $D_{\mathrm{red}}$, and the
preceding chapter gives
\[
    c^*Q_{\mathfrak G}
    \simeq
    Q_{\mathfrak G\times_SS'}.
\]
Finite-limit pasting therefore identifies the formalized quotient targets.
The affine leaf functors are then identified by the mapping-space
pullback--restriction adjunction, and the cocartesian theorem glues these
equivalences globally.

For an arbitrary change of spectral base, the same universal properties
produce the canonical mate comparison, but the generator argument need not
make it invertible.  For an arbitrary varying object map $X\to Y$, a coherent
affine-etale chart over $X$ need not remain etale over $Y$; the present
etale-global coefficient categories therefore do not supply a general pullback
functor in that variance.
\end{proof}

\section{The spectral partition-Lie algebroid of a formal groupoid}
\label{sec:bnla-spectral-groupoid-algebroid}

\begin{definition}[Brave-new spectral partition-Lie algebroid]
\label{def:bnla-groupoid-spectral-algebroid}
In the locally coherent spectral coefficient sector define
\[
\boxed{
    \mathbb A^{\pi,E_\infty}_{\mathfrak G}
    :=
    \operatorname{Diff}^{\pi,E_\infty}_{X/S}
    \left(\mathcal L^{\mathrm{fm,sp}}_{\mathfrak G}\right)
    \in
    \operatorname{LieAlgd}^{\pi,E_\infty}_{X/S}.}
\]
\end{definition}

\begin{theorem}[Spectral differentiation of brave-new formal Lie groupoids]
\label{thm:bnla-spectral-groupoid-differentiation}
For fixed locally coherent $X/S$, the assignment
\[
    \mathfrak G\longmapsto
    \mathbb A^{\pi,E_\infty}_{\mathfrak G}
\]
is functorial on $\operatorname{FormLieGpd}(X/S)$.  Its underlying anchored
pro-coherent object is the shifted tangent-fibre morphism
\[
\boxed{
    \mathbb T^{\mathrm{pc}}_{X/
      \mathcal L^{\mathrm{fm,sp}}_{\mathfrak G}}[1]
    \longrightarrow
    \mathbb T^{\mathrm{pc}}_{X/S}[1].}
\]
All spectral partition-Lie operations and their coherences are the algebra
structure over the constructed monad
$\operatorname{Lie}^{\pi,E_\infty}_{X/S}$.

If
\[
    F:(\mathfrak G,X)\longrightarrow(\mathfrak H,Y)
\]
covers a spectral etale morphism $f:X\to Y$ over $S$, and both geometries lie
in the locally coherent coefficient sector, then the restriction functor on
the global etale coefficient stacks gives a canonical morphism
\[
\boxed{
    \mathbb A^{\pi,E_\infty}_{\mathfrak G}
    \longrightarrow
    f^*\mathbb A^{\pi,E_\infty}_{\mathfrak H}.}
\]
These morphisms are coherent for identities and composition.  No arbitrary
varying-object functoriality is included without a separate spectral families
construction for partition-Lie algebroids.
\end{theorem}

\begin{proof}
For fixed $X/S$, the formalized quotient target and its affine leaf functors
are functorial on $\operatorname{FormLieGpd}(X/S)$ by
Theorem~\ref{thm:bnla-global-leaf-functoriality}.  Apply the global recognition
equivalence Theorem~\ref{thm:bnla-global-spectral-recognition}.  The underlying
anchor is the affine tangent-fibre object of
Theorem~\ref{thm:bnla-spectral-affine-recognition}; the cocartesian
equivalences of Theorem~\ref{thm:bnla-formal-leaf-etale-cocartesian} identify
those affine anchors on overlaps, so they descend to the displayed global
morphism.

Now let $F$ cover spectral etale $f:X\to Y$.  The second part of
Theorem~\ref{thm:bnla-global-leaf-functoriality} supplies
\[
    \mathcal L^{\mathrm{fm,sp}}_{\mathfrak G}
    \longrightarrow
    f^*\mathcal L^{\mathrm{fm,sp}}_{\mathfrak H}.
\]
On every coherent affine-etale pair $(U,V)$ of $X/S$, this is a morphism of
formal-moduli problems over the same affine coefficient map $R\to B$.
Applying affine differentiation gives a morphism of
$\operatorname{Lie}^{\pi,E_\infty}_{B/R}$-algebras.  The etale naturality and
descent theorems identify these affine morphisms with the restriction
$f^*\mathbb A^{\pi,E_\infty}_{\mathfrak H}$ and glue them to the displayed
global algebroid morphism.  Coherence follows from the coherent restriction,
leaf, and affine-differentiation constructions.
\end{proof}

\subsection{Compatibility with represented first order}

\begin{theorem}[Represented quotient first-order comparison]
\label{thm:bnla-spectral-represented-first-order-comparison}
Assume that $X$, $S$, and $Q_{\mathfrak G}$ are represented spectral schemes
in the locally coherent coefficient sector.  Then the underlying anchored
pro-coherent object of
$\mathbb A^{\pi,E_\infty}_{\mathfrak G}$ is canonically
\[
\boxed{
    \left(L_{X/Q_{\mathfrak G}}\right)^\vee[1]
    \longrightarrow
    \left(L_{X/S}\right)^\vee[1],}
\]
the continuous pro-coherent dual of cotangent transitivity.  Consequently
\[
    \left(L_{X/Q_{\mathfrak G}}\right)^\vee
    \simeq
    \left(\mathfrak a^\vee_{\mathfrak G}\right)^\vee
\]
by Theorem~\ref{thm:bnla-relative-cotangent-identification}.  No dualizability
is required.
\end{theorem}

\begin{proof}
By Proposition~\ref{prop:bnla-represented-quotient-artdef-local}, the global
spectral deformation formalization does not change a represented quotient:
\[
    Q^{\mathrm{fm}}_{\mathfrak G}
    \simeq
    Q_{\mathfrak G}.
\]
Hence the affine formal leaf of
Definition~\ref{def:bnla-affine-global-formal-leaf} is the ordinary Artinian
lifting functor of the represented morphism
\[
    q_{\mathfrak G}:X\longrightarrow Q_{\mathfrak G}.
\]

The assertion is spectral-etale local on $X$ and $Q_{\mathfrak G}$.  Choose
affine charts
\[
    U=\operatorname{Spec}B\subset X,
    \qquad
    W=\operatorname{Spec}C\subset Q_{\mathfrak G},
\]
with $q_{\mathfrak G}(U)\subset W$, and choose a coherent affine-etale base
chart $V=\operatorname{Spec}R\to S$.  Thus we have
\[
    R\longrightarrow C\longrightarrow B.
\]
For a structured square-zero test classified by
\[
    \alpha:L_{B/R}[-1]\longrightarrow M,
\]
the relative cotangent lifting theorem identifies the space of lifts of
$U\to W$ across the corresponding Artinian thickening with the mapping space
in the derivation undercategory from
\[
    L_{B/R}[-1]\longrightarrow L_{B/C}[-1]
\]
to $\alpha$.  The square-zero characterization in
Theorem~\ref{thm:bnla-spectral-affine-recognition} therefore identifies the
underlying anchored pro-coherent controller with
\[
\boxed{
    (L_{B/C})^\vee[1]
    \longrightarrow
    (L_{B/R})^\vee[1].}
\]
These local identifications are compatible with spectral etale restriction by
cotangent base change and
Lemma~\ref{lem:bnla-procoherent-finite-tor-bc}.  They therefore descend to
\[
    (L_{X/Q_{\mathfrak G}})^\vee[1]
    \longrightarrow
    (L_{X/S})^\vee[1].
\]
Finally
Theorem~\ref{thm:bnla-relative-cotangent-identification} identifies
$L_{X/Q_{\mathfrak G}}$ with
$\mathfrak a^\vee_{\mathfrak G}$.
\end{proof}

\begin{corollary}[Identity spectral algebroid]
\label{cor:bnla-identity-spectral-algebroid}
The identity formal groupoid differentiates to the zero spectral
partition-Lie algebroid
\[
\boxed{
    \mathbb A^{\pi,E_\infty}_{0_X}
    \simeq
    \left(
        0\longrightarrow\mathbb T^{\mathrm{pc}}_{X/S}[1]
    \right).}
\]
\end{corollary}

\begin{proof}
Its effective quotient is $X\xrightarrow{\operatorname{id}}X$.  On every
represented affine chart Theorem~\ref{thm:bnla-spectral-represented-first-order-comparison}
gives zero underlying source because $L_{X/X}=0$.  The zero source admits the
unique zero-anchor partition-Lie algebra structure of
Theorem~\ref{thm:bnla-zero-anchor-spectral-algebroids}; etale descent gives the
global statement.
\end{proof}

\begin{corollary}[Maximal spectral algebroid under image and deformation locality]
\label{cor:bnla-maximal-spectral-algebroid}
Assume that the effective de Rham image is itself relative de Rham-local,
equivalently that the canonical monomorphism
\[
    Q_{X/S}\longrightarrow X_{\mathrm{dR}/S}
\]
is an equivalence, and assume in addition that $Q_{X/S}$ is
$L^{\mathrm{Def}}_S$-local.  Then the maximal formal relation differentiates
to
\[
\boxed{
    \mathbb A^{\pi,E_\infty}_{R_{X/S}}
    \simeq
    \left(
        \mathbb T^{\mathrm{pc}}_{X/S}[1]
        \xrightarrow{\operatorname{id}}
        \mathbb T^{\mathrm{pc}}_{X/S}[1]
    \right).}
\]
The second hypothesis is automatic whenever $Q_{X/S}$ is
$D_{\mathrm{red}}$-local.  Thus, in the represented smooth range, the
statement applies whenever smooth effectivity
$Q_{X/S}\simeq X_{\mathrm{dR}/S}$ and the Chapter~9 convergence criterion for
reduced-classical locality both hold.
\end{corollary}

\begin{proof}
The deformation-locality hypothesis and
Definition~\ref{def:bnla-artdef-formal-quotient} give
\[
    Q^{\mathrm{def}}_{R_{X/S}}
    \simeq
    Q_{X/S},
\]
and the reduced-centred pullback is therefore also unchanged:
\[
    Q^{\mathrm{fm}}_{R_{X/S}}
    \simeq
    Q_{X/S}.
\]
Relative de Rham-locality gives, for every finite spectral infinitesimal
thickening $U\hookrightarrow T$ over $S$,
\[
    \operatorname{Map}_S(T,Q_{X/S})
    \xrightarrow{\sim}
    \operatorname{Map}_S(U,Q_{X/S}).
\]
Thus every affine formal leaf of the maximal relation is the terminal
relative formal-moduli problem.  Under affine recognition, the terminal
formal-moduli problem corresponds to the tautological identity anchor
\[
    \mathbb T^{\mathrm{pc}}_{B/R}[1]
    \xrightarrow{\operatorname{id}}
    \mathbb T^{\mathrm{pc}}_{B/R}[1].
\]
The etale cocartesianity theorem identifies these affine objects on overlaps,
so they descend to the displayed global identity-anchor algebroid.

If $Q_{X/S}$ is $D_{\mathrm{red}}$-local, then it is Postnikov convergent and
preserves pullbacks along nilpotent extensions because its affine evaluation
depends only on the reduced degree-zero algebra.  Hence it is
$L^{\mathrm{Def}}_S$-local by
Theorem~\ref{thm:bnla-global-def-reflector}.
\end{proof}

\begin{remark}[Why both locality hypotheses are stated]
The preceding corollary does not infer relative de Rham locality merely from
the fact that $X\to Q_{X/S}$ is a de Rham equivalence.  Chapter~9 treats
locality of the effective image as an additional conclusion obtained from its
image-locality criterion.  The stronger nonlinear deformation formalization
used in this chapter also imposes Postnikov convergence and nilpotent
excision, so relative de Rham locality alone does not identify
$Q_{X/S}$ with its $L^{\mathrm{Def}}_S$-localization.  The Chapter~9
convergence criterion supplies a sufficient additional hypothesis.  The
unconditional intrinsic stable calculation for $R_{X/S,\bullet}$ remains
valid without either additional hypothesis.
\end{remark}

\begin{remark}[Two spectral Lie objects]
The integral isotropy primitive Lie algebra
\[
    \mathfrak l^{\mathrm{sp}}_{\mathfrak G}
\]
and the many-object spectral partition-Lie algebroid
\[
    \mathbb A^{\pi,E_\infty}_{\mathfrak G}
\]
are different constructions.  The former is obtained from derived primitives
of the unreduced suspension Hopf object of actual formal isotropy.  The latter
is obtained from formal-moduli differentiation of the effective formal
quotient.  Their direct operation-level comparison belongs to the subsequent
Atiyah--Spencer calculus.
\end{remark}
\section{Integral animated comparison interface}
\label{sec:bnla-animated-shadow}

The spectral partition-Lie algebroid constructed above is primary.  An
integral comparison with the animated partition-Lie algebroids of
\cite{BrantnerMagidsonNuiten2025} would require a functor from animated
Artinian tests to spectral Artinian tests which is compatible with structured
square-zero extensions, Schlessinger pullbacks, cotangent fibres, and
\`etale change of centre.

There is a canonical Eilenberg--Mac Lane realization from animated
commutative rings to connective $E_\infty$-rings, but the additive
Dold--Kan equivalence does not by itself imply the finite-limit and
commutative-algebra compatibility required for such an Artinian comparison.
In particular, this chapter does not identify an animated classified
square-zero pullback
\[
    C_\eta
    \simeq
    C\times_{C\oplus M[1]}C
\]
with a corresponding pullback of connective commutative ring spectra merely
by applying Dold--Kan to underlying additive objects.

\begin{remark}[Integral animated comparison is not an input]
\label{rem:bnla-animated-comparison-deferred}
No functor
\[
    \operatorname{FMP}^{\mathrm{sp}}_{A/K}
    \longrightarrow
    \operatorname{FMP}^{\Delta}_{\pi_0(A)/\pi_0(K)}
\]
and no comparison
\[
    \mathbb A^{\pi,E_\infty}_{\mathfrak G}
    \longrightarrow
    \mathbb A^{\pi,\Delta}_{\mathfrak G}
\]
is postulated in this chapter.  Such a comparison may be introduced only
after a separate theorem constructs the required comparison of the two
Artinian deformation contexts and proves its compatibility with the
corresponding formal-integration equivalences.  No animated compatibility
condition is therefore hidden in the definition of a brave-new spectral
partition-Lie algebroid.
\end{remark}

The characteristic-zero comparison below does not depend on such an integral
animated comparison.  It is constructed directly from rational symmetric
monoidal rectification of the spectral coefficient theory.
\section{Characteristic-zero dg Lie algebroid comparison}
\label{sec:bnla-dg-lie-rectification}

In characteristic zero the spectral partition-Lie operations rectify to the
pro-coherent shifted dg Lie operations.  This comparison is constructed directly
from the spectral cotangent--square-zero and pro-coherent monad constructions;
it does not pass through an integral animated comparison.

\subsection{Rational rectification of coefficients and the relative monad}

\paragraph{Rational and grading convention.}
Throughout this section, a \emph{rational connective $E_\infty$-ring} means a
connective commutative $H\mathbb Q$-algebra.  Differential graded objects are
homologically graded: a dg module $M$ is connective when
\[
    H_i(M)=0
    \qquad(i<0),
\]
and $M[1]$ denotes homological suspension.  Thus the spectral convention
``connective with shift $[1]$'' and the dg convention ``nonnegative homology
with shift $[1]$'' agree under rectification.

Let
\[
    \Phi_{\mathbb Q}
\]
denote the rational symmetric monoidal rectification equivalence.  For a
rational coherent morphism of connective $E_\infty$-rings
\[
    R\longrightarrow B
\]
write
\[
    R_{\mathrm{dg}}\longrightarrow B_{\mathrm{dg}}
\]
for its image.  We write
\[
    \mathbb T^{\mathrm{pc,dg}}_{B_{\mathrm{dg}}/R_{\mathrm{dg}}}
    :=
    \left(L_{B_{\mathrm{dg}}/R_{\mathrm{dg}}}\right)^\vee
\]
for the continuous pro-coherent dg tangent object.  After desuspending a
shifted anchor we also write
\[
    T_{B_{\mathrm{dg}}/R_{\mathrm{dg}}}
    :=
    \mathbb T^{\mathrm{pc,dg}}_{B_{\mathrm{dg}}/R_{\mathrm{dg}}}
\]
when using the ordinary unshifted dg Lie-algebroid convention.

\begin{proposition}[Rectification of cotangent, square-zero, and pro-coherent coefficients]
\label{prop:bnla-rational-coefficient-rectification}
Rational rectification identifies:
\begin{enumerate}
    \item relative cotangent complexes and derivation mapping spaces;
    \item structured square-zero extensions and the relative
    $\operatorname{cot}\dashv\operatorname{sqz}$ adjunction;
    \item coherent and almost-perfect modules, and therefore the
    pro-coherent categories and continuous duality;
    \item the complete almost-finite sectors and their canonical comonadic
    cobar resolutions;
    \item the anchored coefficient categories, through the canonical
    equivalence
    \[
        \Phi_{\mathbb Q}
        \left(\mathbb T^{\mathrm{pc}}_{B/R}[1]\right)
        \simeq
        \mathbb T^{\mathrm{pc,dg}}_{B_{\mathrm{dg}}/R_{\mathrm{dg}}}[1].
    \]
\end{enumerate}
\end{proposition}

\begin{proof}
The rectification equivalence is symmetric monoidal, exact, and compatible
with the connective $t$-structures.  The cotangent complex is characterized
by the universal derivation mapping property, so it is preserved.  A
structured square-zero extension is the pullback of a classifying derivation
against the zero derivation, so it is preserved as well.  Coherence and
almost-perfectness are detected by finite Postnikov approximations in the
corresponding module categories and are invariant under the exact module
equivalence.  Continuous duality is defined by mapping into coherent modules,
so the pro-coherent categories identify.  The complete almost-finite sector
is cut out by surjectivity on degree zero, almost-perfectness of the
augmentation module, and derived completion, all of which are preserved and
reflected.  Naturality identifies the comonadic cobar resolutions.  Finally
continuous duality carries
\[
    L_{B/R}[-1]
\]
to the dg object
\[
    L_{B_{\mathrm{dg}}/R_{\mathrm{dg}}}[-1]
\]
and therefore identifies the corresponding shifted pro-coherent tangent
anchors.
\end{proof}

\begin{definition}[Relative dg derivations and square-zero formation]
\label{def:bnla-relative-dg-cot-sqz}
Define
\[
    \operatorname{Der}^{\mathrm{dg}}_{B_{\mathrm{dg}}/R_{\mathrm{dg}}}
    :=
    \left(
        \operatorname{Mod}_{B_{\mathrm{dg}}}
    \right)_{L_{B_{\mathrm{dg}}/R_{\mathrm{dg}}}[-1]/}.
\]
Thus an object is a morphism
\[
    \alpha:
    L_{B_{\mathrm{dg}}/R_{\mathrm{dg}}}[-1]
    \longrightarrow M.
\]
For an augmented commutative dg algebra
\[
    A_{\mathrm{dg}}\longrightarrow B_{\mathrm{dg}}
\]
define
\[
    \operatorname{cot}^{\mathrm{dg}}_{B_{\mathrm{dg}}/R_{\mathrm{dg}}}
    (A_{\mathrm{dg}})
    :=
    \left(
    L_{B_{\mathrm{dg}}/R_{\mathrm{dg}}}[-1]
    \longrightarrow
    L_{B_{\mathrm{dg}}/A_{\mathrm{dg}}}[-1]
    \right).
\]
For $\alpha$ as above, let
\[
    d_\alpha,d_0:
    B_{\mathrm{dg}}
    \rightrightarrows
    B_{\mathrm{dg}}\oplus M[1]
\]
be the classified and zero sections and put
\[
\boxed{
    (B_{\mathrm{dg}})_\alpha
    :=
    B_{\mathrm{dg}}
    \times_{B_{\mathrm{dg}}\oplus M[1]}
    B_{\mathrm{dg}}.}
\]
Define
\[
    \operatorname{sqz}^{\mathrm{dg}}_{B_{\mathrm{dg}}/R_{\mathrm{dg}}}
    (M,\alpha)
    :=
    \left((B_{\mathrm{dg}})_\alpha\to B_{\mathrm{dg}}\right).
\]
\end{definition}

\begin{proposition}[Relative dg cotangent--square-zero comonad]
\label{prop:bnla-relative-dg-colie}
The preceding constructions form an adjunction
\[
    \operatorname{cot}^{\mathrm{dg}}_{B_{\mathrm{dg}}/R_{\mathrm{dg}}}
    \dashv
    \operatorname{sqz}^{\mathrm{dg}}_{B_{\mathrm{dg}}/R_{\mathrm{dg}}}.
\]
It restricts to the dg analogue of the complete almost-finite augmented sector
and the connective almost-perfect derivation sector, and the restricted left
adjoint is comonadic.  Define
\[
\boxed{
    \operatorname{coLie}^{\mathrm{dg}}_{B_{\mathrm{dg}}/R_{\mathrm{dg}}}
    :=
    \operatorname{cot}^{\mathrm{dg}}_{B_{\mathrm{dg}}/R_{\mathrm{dg}}}
    \operatorname{sqz}^{\mathrm{dg}}_{B_{\mathrm{dg}}/R_{\mathrm{dg}}}.}
\]
Rational rectification identifies this comonad, including its counit and
comultiplication, with the rectification of
$\operatorname{coLie}^{\pi,E_\infty}_{B/R}$.
\end{proposition}

\begin{proof}
The dg cotangent universal property identifies a lift through
$(B_{\mathrm{dg}})_\alpha\to B_{\mathrm{dg}}$ with a morphism
\[
    L_{B_{\mathrm{dg}}/A_{\mathrm{dg}}}[-1]
    \longrightarrow M
\]
whose restriction along
$L_{B_{\mathrm{dg}}/R_{\mathrm{dg}}}[-1]$ is $\alpha$; this is exactly the
adjunction.  Proposition~\ref{prop:bnla-rational-coefficient-rectification}
identifies the dg adjunction with the rectification of the spectral
cotangent--square-zero adjunction, including the connective almost-perfect and
complete almost-finite subcategories.  It therefore transports
Theorem~\ref{thm:bnla-spectral-cot-sqz-comonadic} to the dg categories.
Comonadicity and the comonad structure are invariant under equivalence of
adjunctions, giving the final assertion.
\end{proof}

\begin{definition}[Anchored pro-coherent dg coefficients]
\label{def:bnla-dg-procoherent-anchors}
Put
\[
    \operatorname{QC}^{\vee,\mathrm{dg}}_{B_{\mathrm{dg}}}
    :=
    \operatorname{Ind}
    \left(
        \operatorname{Coh}(B_{\mathrm{dg}})^{\mathrm{op}}
    \right)
\]
and
\[
    \operatorname{Anch}^{\mathrm{pc,dg}}_{B_{\mathrm{dg}}/R_{\mathrm{dg}}}
    :=
    \left(
        \operatorname{QC}^{\vee,\mathrm{dg}}_{B_{\mathrm{dg}}}
    \right)_{/
      \mathbb T^{\mathrm{pc,dg}}_{B_{\mathrm{dg}}/R_{\mathrm{dg}}}[1]}.
\]
Let
\[
    \operatorname{dAPerfAnch}^{\mathrm{dg}}_{B_{\mathrm{dg}}/R_{\mathrm{dg}}}
\]
be the full subcategory whose source is the continuous dual of a connective
almost-perfect dg $B_{\mathrm{dg}}$-module.
\end{definition}

\begin{proposition}[Construction of the relative pro-coherent dg monad]
\label{prop:bnla-relative-dg-monad-construction}
On
$\operatorname{dAPerfAnch}^{\mathrm{dg}}_{B_{\mathrm{dg}}/R_{\mathrm{dg}}}$,
continuous duality transports the opposite comonad
\[
    \left(
    \operatorname{coLie}^{\mathrm{dg}}_{B_{\mathrm{dg}}/R_{\mathrm{dg}}}
    \right)^{\mathrm{op}}
\]
to a monad.  This monad extends uniquely to a sifted-colimit-preserving monad
\[
\boxed{
    \operatorname{Lie}^{\mathrm{pc}}_{\mathrm{dg},
      B_{\mathrm{dg}}/R_{\mathrm{dg}}}:
    \operatorname{Anch}^{\mathrm{pc,dg}}_{B_{\mathrm{dg}}/R_{\mathrm{dg}}}
    \longrightarrow
    \operatorname{Anch}^{\mathrm{pc,dg}}_{B_{\mathrm{dg}}/R_{\mathrm{dg}}}.}
\]
Define
\[
\boxed{
    \operatorname{LieAlgd}^{\mathrm{pc}}_{\mathrm{dg}}
    (B_{\mathrm{dg}}/R_{\mathrm{dg}})
    :=
    \operatorname{Alg}_{
      \operatorname{Lie}^{\mathrm{pc}}_{\mathrm{dg},
      B_{\mathrm{dg}}/R_{\mathrm{dg}}}}
    \left(
      \operatorname{Anch}^{\mathrm{pc,dg}}_{B_{\mathrm{dg}}/R_{\mathrm{dg}}}
    \right).}
\]
This construction is intrinsic to the rectified dg cotangent--square-zero
adjunction; no spectral algebroid is part of the definition of an object on
the right-hand side.
\end{proposition}

\begin{proof}
On connective almost-perfect derivations, the dg coLie comonad is identified
by Proposition~\ref{prop:bnla-relative-dg-colie} with the rational
rectification of the spectral coLie comonad.  Hence the four continuity
properties of Lemma~\ref{lem:bnla-spectral-colie-continuity} transport to the
dg category: preservation of the connective almost-perfect sector, sifted
colimits within that sector, finite cosimplicial totalizations, and inverse
limits of almost eventually constant towers.

The rectified pro-coherent category is obtained from continuous duals of
connective almost-perfect dg modules by the same right-left completion
relations as on the spectral side.  Dualizing the dg coLie comonad therefore
preserves precisely the relations required for this completion.  The proof of
Theorem~\ref{thm:bnla-spectral-lie-algebroid-monad} applies verbatim: the
opposite comonad extends uniquely to a sifted-colimit-preserving endofunctor,
and its counit and comultiplication extend uniquely to the unit and
multiplication of a monad.  The monad identities hold on the dense dually
almost-perfect subcategory and hence on the extension.
\end{proof}

\begin{definition}[Rational dg Artinian extensions]
\label{def:bnla-rational-dg-artinian}
Let
\[
    R_{\mathrm{dg}}\longrightarrow B_{\mathrm{dg}}
\]
be the connective commutative dg $\mathbb Q$-algebras corresponding to
$R\to B$.  Define
\[
    \operatorname{Art}^{\mathrm{dg}}_{R_{\mathrm{dg}}/B_{\mathrm{dg}}}
\]
to be the full subcategory of connective augmented commutative dg
$R_{\mathrm{dg}}$-algebras
\[
    A_{\mathrm{dg}}\longrightarrow B_{\mathrm{dg}}
\]
for which:
\begin{enumerate}
    \item $H_0(A_{\mathrm{dg}})\to H_0(B_{\mathrm{dg}})$ is surjective with
    nilpotent kernel;
    \item the augmentation is almost finitely presented in the Postnikov
    sense;
    \item its fibre is eventually coconnective.
\end{enumerate}
Define
$\operatorname{FMP}^{\mathrm{dg}}_{B_{\mathrm{dg}}/R_{\mathrm{dg}}}$
by the dg square-zero Schlessinger equations obtained from structured dg
square-zero extensions.
\end{definition}

\begin{theorem}[Rational Artinian rectification]
\label{thm:bnla-rational-artinian-rectification}
Rational symmetric monoidal rectification restricts to a canonical
equivalence
\[
\boxed{
    \operatorname{Art}^{\mathrm{sp}}_{R/B}
    \xrightarrow{\sim}
    \operatorname{Art}^{\mathrm{dg}}_{R_{\mathrm{dg}}/B_{\mathrm{dg}}}.}
\]
Under this equivalence:
\begin{enumerate}
    \item structured spectral square-zero extensions correspond to structured
    dg square-zero extensions with the rectified coefficient and derivation;
    \item every spectral Schlessinger pullback
    \[
    A_\eta=A\times_{B\oplus M[1]}B
    \]
    corresponds to the dg pullback
    \[
    (A_{\mathrm{dg}})_\eta
    =
    A_{\mathrm{dg}}
    \times_{B_{\mathrm{dg}}\oplus M_{\mathrm{dg}}[1]}
    B_{\mathrm{dg}};
    \]
    \item the identity Artinian object corresponds to the identity dg object.
\end{enumerate}
Consequently precomposition with the Artinian equivalence gives a canonical
equivalence
\[
\boxed{
    \operatorname{FMP}^{\mathrm{sp}}_{B/R}
    \xrightarrow{\sim}
    \operatorname{FMP}^{\mathrm{dg}}_{B_{\mathrm{dg}}/R_{\mathrm{dg}}}.}
\]
The comparison is natural under rational spectral etale change of centre.
\end{theorem}

\begin{proof}
The rational rectification equivalence is exact, symmetric monoidal, and
t-exact on the connective categories.  It therefore identifies $\pi_0$ with
$H_0$, preserves fibres, and preserves nilpotence of the degree-zero
augmentation ideal.  Almost finite presentation is the preservation of
filtered colimits by mapping spaces on each truncated connective algebra
category; an equivalence of the corresponding truncated algebra categories
preserves and reflects this property.  Eventual coconnectivity of the fibre
is likewise preserved and reflected by the t-exact equivalence.  Hence
rectification restricts to the displayed equivalence of intrinsic Artinian
categories.

By
Proposition~\ref{prop:bnla-rational-coefficient-rectification}, relative
cotangent complexes and derivation spaces correspond.  A structured
square-zero extension is the pullback of its classifying derivation against
the zero derivation.  Since rectification preserves finite limits, it carries
that pullback to the dg structured square-zero pullback.  The same
finite-limit argument identifies every Schlessinger square, and the identity
object is preserved.

Thus the Artinian equivalence carries the generating localization morphisms
of Definition~\ref{def:bnla-schlessinger-generators} bijectively, up to
canonical equivalence, to the dg Schlessinger generators.  Precomposition is
therefore an equivalence of the ambient functor categories which identifies
their local subcategories.  This gives the formal-moduli equivalence.
Etale naturality follows from rational symmetric monoidal base change and the
uniqueness of the etale Artinian transports.
\end{proof}


\begin{lemma}[Rational disappearance of divided-power corrections]
\label{lem:bnla-rational-tate-vanishing}
For a finite group $G$ and a rational module spectrum $M$,
\[
    M^{tG}\simeq0.
\]
Consequently the norm
\[
    M_{hG}\xrightarrow{\sim}M^{hG}
\]
is an equivalence.  Under rational rectification the pro-coherent
PD-composition product governing the spectral partition-Lie monad therefore
identifies with ordinary symmetric composition.
\end{lemma}

\begin{proof}
Restriction followed by transfer acts on rational homotopy as multiplication
by $|G|$, which is invertible.  Hence the norm map is an equivalence and its
Tate cofiber vanishes.  The refined pro-coherent composition formulas of
\cite{BrantnerCamposNuiten2025} express the divided-power correction through
these finite-group homotopy-fixed-point/norm terms, which therefore reduce to
the characteristic-zero symmetric formula.
\end{proof}

\begin{theorem}[Rational rectification of the spectral algebroid monad]
\label{thm:bnla-rational-monad-rectification}
Under rational symmetric monoidal rectification there is a canonical
equivalence
\[
\boxed{
    \operatorname{LieAlgd}^{\pi,E_\infty}_{B/R}
    \xrightarrow{\sim}
    \operatorname{LieAlgd}^{\mathrm{pc}}_{\mathrm{dg}}
    (B_{\mathrm{dg}}/R_{\mathrm{dg}}),}
\]
where the right-hand side is the dg category constructed in
Proposition~\ref{prop:bnla-relative-dg-monad-construction}.  Its zero-anchor
restriction is canonically the ordinary pro-coherent shifted dg Lie monad.
Accordingly we call its algebras \emph{pro-coherent shifted dg Lie
algebroids}.  In coefficient ranges for which a tame chain model presents
this pro-coherent $\infty$-category, it may equivalently be represented by the
corresponding tame dg Lie-algebroid model.  No tame point-set presentation is
required here, and no identification with the naive quasi-isomorphism
localization of arbitrary dg modules is asserted.

A shifted anchor
\[
    E\longrightarrow
    \mathbb T^{\mathrm{pc,dg}}_{B_{\mathrm{dg}}/R_{\mathrm{dg}}}[1]
\]
corresponds, after desuspending the source, to the ordinary dg anchor
\[
\boxed{
    E[-1]
    \longrightarrow
    T_{B_{\mathrm{dg}}/R_{\mathrm{dg}}}.}
\]
The comparison is natural under rational spectral etale base change.
\end{theorem}

\begin{proof}
Proposition~\ref{prop:bnla-rational-coefficient-rectification} identifies
continuous duality and the anchored coefficient categories, while
Proposition~\ref{prop:bnla-relative-dg-colie} identifies the relative spectral
and dg coLie comonads on connective almost-perfect derivations.  Opposite
continuous duality therefore identifies their monads on the dually
almost-perfect anchor categories.  The spectral monad and the dg monad of
Proposition~\ref{prop:bnla-relative-dg-monad-construction} are both the unique
sifted-colimit-preserving extensions of those restrictions.  Hence the
comparison extends uniquely to an equivalence of monads on all anchored
pro-coherent objects, including unit and multiplication.  Passing to algebras
gives the displayed equivalence.

It remains to identify the zero-anchor dg monad.  By
Proposition~\ref{prop:bnla-one-object-spectral-partition-monad}, the spectral
zero-anchor restriction is the pro-coherent PD Koszul dual of nonunital
$E_\infty$-algebras.  Lemma~\ref{lem:bnla-rational-tate-vanishing} identifies
homotopy orbits and homotopy fixed points for every finite symmetric group,
so the divided-power correction in the pro-coherent composition product
disappears.  The characteristic-zero Koszul dual of the commutative operad is
the shifted Lie operad.  Thus the zero-anchor restriction of the constructed
dg monad is the ordinary pro-coherent shifted dg Lie monad.

Finally, continuous duality gives
\[
    \left(L_{B_{\mathrm{dg}}/R_{\mathrm{dg}}}[-1]\right)^\vee
    \simeq
    \mathbb T^{\mathrm{pc,dg}}_{B_{\mathrm{dg}}/R_{\mathrm{dg}}}[1].
\]
Hence an algebra whose shifted underlying anchor is
$E\to\mathbb T^{\mathrm{pc,dg}}[1]$ has ordinary dg source $E[-1]$ and anchor
$E[-1]\to T_{B_{\mathrm{dg}}/R_{\mathrm{dg}}}$.  Every comparison above is
built from symmetric monoidal rectification, cotangent base change, and the
same pro-coherent completion universal property, so it is natural under
rational spectral etale base change.
\end{proof}

The affine equivalences descend through the global etale coefficient diagram.
Write
\[
    \mathcal R_{\mathbb Q}
\]
for the resulting global rectification/desuspension equivalence.

\begin{definition}[Rectified characteristic-zero formal leaf]
\label{def:bnla-dg-formal-leaf}
For a rational locally coherent spectral formal groupoid, define
\[
    \mathcal L^{\mathrm{dg}}_{\mathfrak G}
\]
by applying the formal-moduli equivalence of
Theorem~\ref{thm:bnla-rational-artinian-rectification} to every coherent
affine-etale restriction of
\[
    \mathcal L^{\mathrm{fm,sp}}_{\mathfrak G}
\]
and gluing the resulting dg formal-moduli objects by the transported etale
cocartesian equivalences.  Thus the dg formal leaf is defined on the actual
dg Artinian category, not merely by a coefficient-level analogy.

Its tangent-fibre object is, by construction, the desuspension of the
rectified spectral tangent-fibre object
\[
    \mathbb T^{\mathrm{pc}}_{X/
      \mathcal L^{\mathrm{fm,sp}}_{\mathfrak G}}[1].
\]
\end{definition}

\begin{definition}[Characteristic-zero dg controller]
\label{def:bnla-dg-controller}
For a rational locally coherent spectral formal groupoid define
\[
\boxed{
    \mathbb A^{\mathrm{Lie}}_{\mathfrak G}
    :=
    \mathcal R_{\mathbb Q}
    \left(\mathbb A^{\pi,E_\infty}_{\mathfrak G}\right).}
\]
\end{definition}

\begin{theorem}[Underlying characteristic-zero anchor]
\label{thm:bnla-dg-lie-rectification}
The underlying anchored dg coefficient of
$\mathbb A^{\mathrm{Lie}}_{\mathfrak G}$ is the desuspension of the rectified
spectral tangent-fibre morphism:
\[
\boxed{
    T_{X/\mathcal L^{\mathrm{dg}}_{\mathfrak G}}
    \longrightarrow
    T_{X/S}.}
\]
On every affine chart satisfying the hypotheses of the formal-moduli/dg
Lie-algebroid equivalence of \cite{Nuiten2019}, this dg controller is
canonically identified with the controller obtained by applying that theorem
to the same rectified formal-moduli problem.
\end{theorem}

\begin{proof}
The spectral controller has underlying shifted anchor
\[
    \mathbb T^{\mathrm{pc}}_{X/
      \mathcal L^{\mathrm{fm,sp}}_{\mathfrak G}}[1]
    \longrightarrow
    \mathbb T^{\mathrm{pc}}_{X/S}[1].
\]
Theorem~\ref{thm:bnla-rational-monad-rectification} rectifies and desuspends
it to the displayed dg tangent morphism.

On a chart in the scope of \cite{Nuiten2019}, the rectified spectral
controller formally integrates, by construction, to the rectified
formal-moduli object.  Nuiten's differentiation equivalence sends that same
formal-moduli object to its dg controller.  Applying the inverse equivalences
and their units and counits gives the canonical identification.  No extension
of Nuiten's stated coefficient range is used.
\end{proof}

\begin{remark}[Bracket comparison is deferred, not bracket existence]
\label{rem:bnla-dg-bracket-comparison-deferred}
The dg Lie bracket and the spectral partition-Lie operations have now both
been constructed.  What remains for the Atiyah--Spencer chapter is their
direct comparison with the nonlinear isotropy commutator and the normalized
jet--cobar operation calculus.
\end{remark}
\section{Actual isotropy versus homotopy anchor fibres in algebraic differentiation}
\label{sec:bnla-formal-isotropy-recognition}

The intrinsic stable construction supplied
\[
    \mathbb I^{\mathrm{st}}_{\mathfrak G}
    \xrightarrow{\gamma^{\mathrm{st}}_{\mathfrak G}}
    \mathbb F^{\mathrm{st}}_{\mathfrak G}.
\]
We now construct the corresponding comparison in spectral partition-Lie
algebroids and transport it through rational rectification.

\subsection{The formal isotropy groupoid}

Let
\[
    I^{\mathrm{for}}_{\mathfrak G}\longrightarrow X
\]
be formal isotropy.  Its group-object structure determines the one-object
fixed-object groupoid
\[
    B_\bullet I^{\mathrm{for}}_{\mathfrak G}
    \rightrightarrows X.
\]

\begin{proposition}[The formal isotropy groupoid is formal]
\label{prop:bnla-formal-isotropy-groupoid}
The groupoid
\[
    \mathfrak I_{\mathfrak G,\bullet}
    :=
    B_\bullet I^{\mathrm{for}}_{\mathfrak G}
\]
is a brave-new formal Lie groupoid on $X/S$.  The isotropy inclusion defines
a fixed-object formal-groupoid morphism
\[
    \iota_{\mathfrak G}:
    \mathfrak I_{\mathfrak G,\bullet}
    \longrightarrow
    \mathfrak G_\bullet.
\]
Its canonical nonlinear anchor factors, through a contractible space of
compatible factorizations, as
\[
    \mathfrak I_{\mathfrak G,\bullet}
    \longrightarrow
    0_{X,\bullet}
    \longrightarrow
    R_{X/S,\bullet}.
\]
\end{proposition}

\begin{proof}
The arrow group $I^{\mathrm{for}}_{\mathfrak G}$ is complete along its
identity section.  Every higher bar degree is a finite fibre product over
$X$ of copies of the arrow object.  Centred completion preserves those
pullbacks and fixes $X$, so every bar degree is complete along its fully
degenerate unit.  Hence the bar groupoid is formal.

The inclusion of formal isotropy into $\mathfrak G_1$ is a homomorphism of
group objects over $X$ and therefore extends simplicially.  The
anchor-kernel theorem of the preceding chapter gives the specified
factorization
\[
    I^{\mathrm{for}}_{\mathfrak G}
    \longrightarrow
    X
    \xrightarrow{u_R}
    R_{X/S}
\]
in degree one.  Compatibility with multiplication and degeneracies extends
it through the bar construction.  Contractibility is inherited from the
defining anchor-kernel pullback and from the terminality of the identity
arrows in $0_{X,\bullet}$.
\end{proof}

\subsection{Spectral partition-Lie differentiated isotropy}

Assume that $X/S$ lies in the locally coherent spectral coefficient sector.

\begin{definition}[Spectral differentiated actual isotropy]
\label{def:bnla-spectral-differentiated-isotropy}
Define
\[
\boxed{
    \mathfrak i^{\pi,E_\infty}_{\mathfrak G}
    :=
    \mathbb A^{\pi,E_\infty}_{\mathfrak I_{\mathfrak G}}.}
\]
\end{definition}

\begin{proposition}[Actual isotropy has zero spectral anchor]
\label{prop:bnla-spectral-isotropy-zero-anchor}
The spectral algebroid
$\mathfrak i^{\pi,E_\infty}_{\mathfrak G}$ has canonically zero anchor.
Consequently the zero-anchor equivalence of
Theorem~\ref{thm:bnla-zero-anchor-spectral-algebroids} regards it canonically
as a spectral partition-Lie algebra on $X$.
\end{proposition}

\begin{proof}
Apply spectral differentiation to the factorization
\[
    \mathfrak I_{\mathfrak G}
    \longrightarrow
    0_X
    \longrightarrow
    R_{X/S}
\]
of Proposition~\ref{prop:bnla-formal-isotropy-groupoid}.  By
Corollary~\ref{cor:bnla-identity-spectral-algebroid}, the middle groupoid
differentiates to
\[
    0\longrightarrow\mathbb T^{\mathrm{pc}}_{X/S}[1].
\]
Hence the anchor of
$\mathfrak i^{\pi,E_\infty}_{\mathfrak G}$ factors canonically through the
zero source.  The zero-anchor theorem then supplies the partition-Lie
algebra interpretation; no zero-anchor structure is added as data.
\end{proof}

Write
\[
    \rho^{\pi,E_\infty}_{\mathfrak G}:
    U\mathbb A^{\pi,E_\infty}_{\mathfrak G}
    \longrightarrow
    \mathbb T^{\mathrm{pc}}_{X/S}[1]
\]
for the underlying anchor.

\begin{construction}[Canonical zero and tautological anchor algebroids]
\label{con:bnla-zero-tautological-anchor-algebroids}
Let
\[
    \mathbf 0_{X/S}
    :=
    \left(
        0\longrightarrow\mathbb T^{\mathrm{pc}}_{X/S}[1]
    \right)
\]
with the canonical spectral partition-Lie algebroid structure supplied
etale-locally by
Theorem~\ref{thm:bnla-zero-anchor-spectral-algebroids} and then descended.
Let $\ast_{X/S}$ denote the terminal object of
$\operatorname{FMP}^{\mathrm{sp}}_{X/S}$ and define
\[
    \mathbf T_{X/S}
    :=
    \operatorname{Diff}^{\pi,E_\infty}_{X/S}(\ast_{X/S}).
\]
The square-zero characterization in
Theorem~\ref{thm:bnla-spectral-affine-recognition} identifies its underlying
anchor with
\[
    \mathbb T^{\mathrm{pc}}_{X/S}[1]
    \xrightarrow{\operatorname{id}}
    \mathbb T^{\mathrm{pc}}_{X/S}[1].
\]
For every spectral partition-Lie algebroid $\mathfrak a$, the unique map from
its formal integration to $\ast_{X/S}$ differentiates to a canonical
algebroid morphism
\[
    \mathfrak a\longrightarrow\mathbf T_{X/S}
\]
whose underlying morphism of anchored pro-coherent objects is its anchor.
The unique morphism
$\mathbf 0_{X/S}\to\mathbf T_{X/S}$ has underlying source map
$0\to\mathbb T^{\mathrm{pc}}_{X/S}[1]$.
\end{construction}

\begin{definition}[Spectral partition anchor fibre]
\label{def:bnla-spectral-partition-anchor-fibre}
Define
\[
\boxed{
    \mathbb F^{\pi,E_\infty}_{\mathfrak G}}
\]
to be the zero-anchor algebroid obtained by applying the right-adjoint
anchor-fibre functor of
Theorem~\ref{thm:bnla-zero-anchor-spectral-algebroids} to
$\mathbb A^{\pi,E_\infty}_{\mathfrak G}$.  Equivalently, using
Construction~\ref{con:bnla-zero-tautological-anchor-algebroids}, it is the
pullback in the category of spectral partition-Lie algebroids
\[
\begin{tikzcd}
\mathbb F^{\pi,E_\infty}_{\mathfrak G}
  \arrow[r] \arrow[d] &
\mathbb A^{\pi,E_\infty}_{\mathfrak G}
  \arrow[d]\\
\mathbf 0_{X/S} \arrow[r] &
\mathbf T_{X/S}.
\end{tikzcd}
\]
Its underlying pro-coherent object is canonically
\[
    U\mathbb F^{\pi,E_\infty}_{\mathfrak G}
    \simeq
    \operatorname{fib}
    \left(
        \rho^{\pi,E_\infty}_{\mathfrak G}
    \right).
\]
\end{definition}

\begin{theorem}[Spectral partition isotropy comparison]
\label{thm:bnla-spectral-partition-isotropy-comparison}
Differentiation of the nonlinear isotropy inclusion and its constructed
anchor nullhomotopy gives a canonical morphism of zero-anchor spectral
partition-Lie algebroids
\[
\boxed{
    \gamma^{\pi,E_\infty}_{\mathfrak G}:
    \mathfrak i^{\pi,E_\infty}_{\mathfrak G}
    \longrightarrow
    \mathbb F^{\pi,E_\infty}_{\mathfrak G}.}
\]
Define the spectral partition isotropy discrepancy in the underlying stable
pro-coherent category by
\[
\boxed{
    \mathbb C^{\mathrm{iso},\pi,E_\infty}_{\mathfrak G}
    :=
    \operatorname{cofib}
    \left(
        U\gamma^{\pi,E_\infty}_{\mathfrak G}
    \right).}
\]
The comparison and discrepancy are functorial and compatible with spectral
étale base change.
\end{theorem}

\begin{proof}
The formal-groupoid morphism
\[
    \mathfrak I_{\mathfrak G}\longrightarrow\mathfrak G
\]
gives, after the functorial global formal-moduli localization and then
spectral differentiation, a morphism
\[
    \mathfrak i^{\pi,E_\infty}_{\mathfrak G}
    \longrightarrow
    \mathbb A^{\pi,E_\infty}_{\mathfrak G}.
\]
The factorization through the identity groupoid in
Proposition~\ref{prop:bnla-formal-isotropy-groupoid} supplies a specified
nullhomotopy of the composite with the anchor.  By the pullback universal
property in
Definition~\ref{def:bnla-spectral-partition-anchor-fibre}, this determines
the displayed morphism $\gamma$.

This argument does not require a formal-moduli reflector to preserve the
nonlinear isotropy pullback.  It transports a constructed morphism and its
constructed nullhomotopy through functorial differentiation.  Étale
compatibility follows from that of the spectral differentiation equivalence.
The discrepancy is then formed functorially in the underlying stable
pro-coherent category.
\end{proof}

\subsection{Anchor transitivity in spectral coefficients}

Write the underlying tangent-fibre anchor as
\[
    \mathbb T^{\mathrm{pc}}_{X/\mathcal L^{\mathrm{fm,sp}}_{\mathfrak G}}[1]
    \longrightarrow
    \mathbb T^{\mathrm{pc}}_{X/S}[1].
\]

\begin{definition}[Formal-leaf ambient tangent coefficient]
\label{def:bnla-formal-leaf-ambient-tangent}
Define
\[
\boxed{
    \mathbb T^{\mathrm{pc}}_{\mathcal L_{\mathfrak G}/S}\big|_X
    :=
    \operatorname{cofib}
    \left(
        \mathbb T^{\mathrm{pc}}_{X/\mathcal L^{\mathrm{fm,sp}}_{\mathfrak G}}
        \longrightarrow
        \mathbb T^{\mathrm{pc}}_{X/S}
    \right).}
\]
When the formal-moduli object is represented by a formal leaf stack
$X\to\mathfrak X_{\mathfrak G}$, ordinary tangent transitivity identifies
this object with
$q^*\mathbb T^{\mathrm{pc}}_{\mathfrak X_{\mathfrak G}/S}$.
\end{definition}

\begin{proposition}[Spectral anchor-fibre transitivity]
\label{prop:bnla-spectral-anchor-fibre-transitivity}
There is a canonical equivalence of underlying pro-coherent objects
\[
\boxed{
    U\mathbb F^{\pi,E_\infty}_{\mathfrak G}
    \simeq
    \mathbb T^{\mathrm{pc}}_{\mathcal L_{\mathfrak G}/S}\big|_X.}
\]
Thus the spectral partition-Lie anchor has a canonical fibre sequence on
underlying pro-coherent objects
\[
\boxed{
    U\mathbb F^{\pi,E_\infty}_{\mathfrak G}
    \longrightarrow
    U\mathbb A^{\pi,E_\infty}_{\mathfrak G}
    \longrightarrow
    \mathbb T^{\mathrm{pc}}_{X/S}[1].}
\]
\end{proposition}

\begin{proof}
Let
\[
    A\longrightarrow T\longrightarrow C
\]
be the cofiber sequence in the stable pro-coherent category underlying the
unshifted tangent morphism
\[
    \mathbb T^{\mathrm{pc}}_{X/\mathcal L^{\mathrm{fm,sp}}_{\mathfrak G}}
    \longrightarrow
    \mathbb T^{\mathrm{pc}}_{X/S}.
\]
After shifting the first two terms by $[1]$, stability gives
\[
    \operatorname{fib}(A[1]\to T[1])\simeq C.
\]
By Definition~\ref{def:bnla-formal-leaf-ambient-tangent}, this $C$ is the
formal-leaf ambient tangent coefficient.  The zero-anchor fibre theorem
identifies it with the underlying object of
$\mathbb F^{\pi,E_\infty}_{\mathfrak G}$.
\end{proof}

\subsection{Characteristic-zero image}

In the rational coefficient sector define
\[
    \mathfrak i^{\mathrm{Lie}}_{\mathfrak G}
    :=
    \mathcal R_{\mathbb Q}
    \left(
        \mathfrak i^{\pi,E_\infty}_{\mathfrak G}
    \right)
\]
and let
\[
    \mathbb F^{\mathrm{Lie}}_{\mathfrak G}
\]
be the zero-anchor dg Lie algebra obtained by applying the dg anchor-fibre
right adjoint to $\mathbb A^{\mathrm{Lie}}_{\mathfrak G}$.

\begin{corollary}[Characteristic-zero isotropy comparison]
\label{cor:bnla-dg-isotropy-comparison}
There is a canonical dg Lie algebra morphism
\[
\boxed{
    \gamma^{\mathrm{Lie}}_{\mathfrak G}:
    \mathfrak i^{\mathrm{Lie}}_{\mathfrak G}
    \longrightarrow
    \mathbb F^{\mathrm{Lie}}_{\mathfrak G}.}
\]
Its underlying cofiber
\[
    \mathbb C^{\mathrm{iso},\mathrm{Lie}}_{\mathfrak G}
    :=
    \operatorname{cofib}
    \left(
        U\gamma^{\mathrm{Lie}}_{\mathfrak G}
    \right)
\]
is the rational isotropy discrepancy.
\end{corollary}

\begin{proof}
Apply the rational rectification equivalence to
Theorem~\ref{thm:bnla-spectral-partition-isotropy-comparison}.  Since this is
an equivalence of algebroid categories, it preserves the zero-anchor
pullback and its universal factorization.
\end{proof}

\begin{proposition}[Characteristic-zero anchor transitivity]
\label{prop:bnla-dg-anchor-transitivity}
There is a canonical equivalence
\[
\boxed{
    U\mathbb F^{\mathrm{Lie}}_{\mathfrak G}
    \simeq
    \left(
        \mathbb T^{\mathrm{pc}}_{\mathcal L_{\mathfrak G}/S}\big|_X
    \right)_{\mathrm{dg}}[-1].}
\]
In particular there is a fibre sequence of underlying dg modules
\[
\boxed{
    U\mathbb F^{\mathrm{Lie}}_{\mathfrak G}
    \longrightarrow
    U\mathbb A^{\mathrm{Lie}}_{\mathfrak G}
    \longrightarrow
    T_{X/S}.}
\]
\end{proposition}

\begin{proof}
Rational rectification sends the shifted spectral anchor to the unshifted dg
anchor.  If $A\to T\to C$ is the corresponding unshifted tangent cofiber
sequence, then
\[
    \operatorname{fib}(A\to T)\simeq\Omega C=C[-1].
\]
Apply Proposition~\ref{prop:bnla-spectral-anchor-fibre-transitivity}.
\end{proof}

\paragraph{Atiyah terminology.}
Whenever, for a specified class of formal groupoids, one proves that one of
the canonical comparison maps
\[
    \gamma^{\mathrm{st}}_{\mathfrak G},
    \qquad
    \gamma^{\pi,E_\infty}_{\mathfrak G},
    \qquad
    \gamma^{\mathrm{Lie}}_{\mathfrak G}
\]
is an equivalence, we call the resulting anchor-fibre sequence the
\emph{Atiyah sequence} of that groupoid in the corresponding coefficient
theory.  No ``Atiyah identification'' is retained as additional data.
\section{Gauge groupoids and effective-image correction}
\label{sec:bnla-gauge-specialization}

Let
\[
    P_\chi=X\times_{BG}S\longrightarrow X
\]
be the principal right $G$-object classified by
\[
    \chi:X\longrightarrow BG.
\]
The preceding chapter defined
\[
    \operatorname{Gauge}(P_\chi)
    :=
    \check C_\bullet(\chi)
\]
and
\[
    \operatorname{Gauge}^{\mathrm{for}}(P_\chi)
    :=
    \widehat{\operatorname{Gauge}(P_\chi)}^{\,\mathrm{dR}}.
\]

Factor the classifying map by its effective image in the ambient $\infty$-topos:
\[
    X\xrightarrow{e_\chi}I_\chi
    \xrightarrow{m_\chi}BG,
\]
where $e_\chi$ is an effective epimorphism and $m_\chi$ is a monomorphism.

\begin{theorem}[Effective quotient of the gauge groupoid]
\label{thm:bnla-gauge-effective-image}
There is a canonical equivalence of simplicial objects
\[
    \boxed{
    \check C_\bullet(\chi)
    \simeq
    \check C_\bullet(e_\chi).}
\]
Consequently
\[
    \left|\operatorname{Gauge}(P_\chi)\right|
    \simeq
    I_\chi.
\]
The effective quotient of the formal gauge groupoid is the quotient
formalization of $e_\chi$:
\[
    \boxed{
    Q_{\operatorname{Gauge}^{\mathrm{for}}(P_\chi)}
    \simeq
    \widehat I_\chi,}
\]
where
\[
    X\twoheadrightarrow\widehat I_\chi
    \longrightarrow I_\chi
\]
is the canonical quotient formalization constructed in the preceding chapter.
\end{theorem}

\begin{proof}
Because $m_\chi:I_\chi\to BG$ is a monomorphism, its diagonal is an
equivalence.  Hence
\[
    X\times_{BG}X
    \simeq
    X\times_{I_\chi}X,
\]
and the same argument gives, for every $n$,
\[
    \underbrace{X\times_{BG}\cdots\times_{BG}X}_{n+1}
    \simeq
    \underbrace{X\times_{I_\chi}\cdots\times_{I_\chi}X}_{n+1}.
\]
These equivalences respect all Cech operators, proving the simplicial
comparison.  Since $e_\chi$ is an effective epimorphism, its Cech nerve
realizes to $I_\chi$.

The preceding chapter proves that unitwise formalization of a Cech groupoid is
the Cech groupoid of the quotient formalization.  Apply that theorem to
$e_\chi$.
\end{proof}

\begin{remark}[The classifying stack is not the quotient]
The classifying map $X\to BG$ is used to construct the gauge groupoid, but the
effective quotient of its Cech nerve is $I_\chi$, not $BG$ unless the
classifying map is itself an effective epimorphism.  All differentiated
quotient constructions therefore use $\widehat I_\chi$.
\end{remark}

\begin{definition}[Gauge formal-moduli envelope]
\label{def:bnla-gauge-formal-envelope}
Define the gauge spectral formal-moduli envelope by applying
Definition~\ref{def:bnla-global-spectral-leaf-envelope} to
\[
    \operatorname{Gauge}^{\mathrm{for}}(P_\chi)
\]
and hence, by Theorem~\ref{thm:bnla-gauge-effective-image}, to the formal
quotient
\[
    X\twoheadrightarrow\widehat I_\chi.
\]
Its spectral partition-Lie controller and, in the rational coefficient
sector, its dg controller are obtained by the differentiation and
rectification functors constructed above.  No map to $BG$ is used in place
of the actual effective quotient.
\end{definition}

\begin{proposition}[Gauge anchor-fibre sequences]
\label{prop:bnla-gauge-anchor-fibre}
In the coherent spectral coefficient sector the formal gauge groupoid has the
canonical spectral partition-Lie fibre sequence on underlying pro-coherent
objects
\[
\boxed{
    U\mathbb F^{\pi,E_\infty}_{P_\chi}
    \longrightarrow
    U\mathbb A^{\pi,E_\infty}_{P_\chi}
    \longrightarrow
    \mathbb T^{\mathrm{pc}}_{X/S}[1],}
\]
together with the canonical actual-isotropy comparison
\[
    \mathfrak i^{\pi,E_\infty}_{P_\chi}
    \longrightarrow
    \mathbb F^{\pi,E_\infty}_{P_\chi}.
\]
In the rational coefficient sector its dg image is
\[
    U\mathbb F^{\mathrm{Lie}}_{P_\chi}
    \longrightarrow
    U\mathbb A^{\mathrm{Lie}}_{P_\chi}
    \longrightarrow
    T_{X/S},
\]
with differentiated actual gauge isotropy mapping canonically to the first
term.  No equivalence with the anchor fibre is asserted without a separate
proof.
\end{proposition}

\begin{proof}
Apply
Theorem~\ref{thm:bnla-spectral-partition-isotropy-comparison} and
Proposition~\ref{prop:bnla-spectral-anchor-fibre-transitivity} to the formal
gauge groupoid, using
Theorem~\ref{thm:bnla-gauge-effective-image} for its actual effective
quotient.  Apply rational rectification for the dg sequence.
\end{proof}

\section{Fundamental differentiation calculations}
\label{sec:bnla-basic-calculations}

\subsection{The identity formal groupoid}

\begin{proposition}[Identity groupoid]
\label{prop:bnla-zero-groupoid}
For
\[
    0_{X,\bullet}
    =
    \check C_\bullet(\operatorname{id}_X),
\]
there are canonical equivalences
\[
    \mathbb A^{\mathrm{st}}_{0_X}\simeq0,
\]
\[
    \boxed{
    \mathbb F^{\mathrm{st}}_{0_X}
    \simeq
    \Omega\mathbb T^{\mathrm{st}}_{X/S},}
\]
and
\[
    \boxed{
    \mathbb N^{\mathrm{st}}_{0_X}
    \simeq
    \mathbb T^{\mathrm{st}}_{X/S}.}
\]
The actual formal isotropy remains the formal inertia kernel constructed in the
preceding chapter, and there is a canonical comparison
\[
    \mathbb I^{\mathrm{st}}_{0_X}
    \longrightarrow
    \Omega\mathbb T^{\mathrm{st}}_{X/S}.
\]
If $X\to S$ is $0$-truncated---in particular, for a morphism of ordinary
schemes---actual formal isotropy is trivial while the displayed anchor fibre
need not vanish.
\end{proposition}

\begin{proof}
The arrow object is the identity-pointed object, whose reduced stabilization is
zero by Lemma~\ref{lem:bnla-reduced-identity-zero}.  Thus the stable anchor is
$0\to\mathbb T^{\mathrm{st}}_{X/S}$.  In a stable category its fibre is the
loop of the target and its cofiber is the target.  The isotropy comparison is
Definition~\ref{def:bnla-stable-isotropy-comparison}.

If $X\to S$ is $0$-truncated, its diagonal
$X\to X\times_SX$ is $(-1)$-truncated.  Hence the relative inertia
\[
    I_{X/S}
    =
    X\times_{X\times_SX}X
\]
is the trivial group object $X$.  The preceding chapter identifies the raw
isotropy of the identity groupoid with this relative inertia and its formal
isotropy with the centred formal inertia kernel.  Therefore
$I^{\mathrm{for}}_{0_X}\simeq X$ as a group object over $X$, and reduced
stabilization gives $\mathbb I^{\mathrm{st}}_{0_X}\simeq0$.  This does not
force $\Omega\mathbb T^{\mathrm{st}}_{X/S}$ to vanish.
\end{proof}

\begin{remark}[Why this example fixes the terminology]
The identity groupoid shows that the homotopy anchor fibre cannot be called
isotropy without further argument.  The two objects have different meanings
already in the simplest case.
\end{remark}

\subsection{The maximal relative de Rham groupoid}

\begin{proposition}[Maximal formal relation]
\label{prop:bnla-maximal-relation}
For
\[
    \mathfrak G_\bullet=R_{X/S,\bullet},
\]
the stable anchor is the identity
\[
    \mathbb T^{\mathrm{st}}_{X/S}
    \xrightarrow{\operatorname{id}}
    \mathbb T^{\mathrm{st}}_{X/S}.
\]
Consequently
\[
    \mathbb F^{\mathrm{st}}_R\simeq0,
    \qquad
    \mathbb N^{\mathrm{st}}_R\simeq0.
\]
Its formal isotropy is the trivial group object $X$, so
\[
    \mathbb I^{\mathrm{st}}_R\simeq0,
    \qquad
    \mathbb C^{\mathrm{iso,st}}_R\simeq0.
\]
\end{proposition}

\begin{proof}
The canonical anchor of the terminal formal relation is its identity.  Its
formal isotropy is the fibre of the identity relation over its unit, hence the
identity group object on $X$.  Apply the definitions.
\end{proof}

\subsection{One-object formal groups}

Let $G/S$ be a group object and put
\[
    \mu:=\mu^{\mathrm{mod}}_{G/S}
    =
    G\times_{G_{\mathrm{dR}/S}}S.
\]
The preceding chapter gives
\[
    \widehat{B_\bullet G}^{\,\mathrm{dR}}
    \simeq
    B_\bullet\mu.
\]

\begin{proposition}[One-object differentiation]
\label{prop:bnla-one-object-differentiation}
For the completed one-object groupoid,
\[
    \mathbb A^{\mathrm{st}}_{B\mu}
    \simeq
    \widetilde\Sigma^\infty_S\mu,
\]
\[
    \mathbb T^{\mathrm{st}}_{S/S}\simeq0,
\]
so
\[
    \boxed{
    \mathbb F^{\mathrm{st}}_{B\mu}
    \simeq
    \widetilde\Sigma^\infty_S\mu,}
\]
and
\[
    \boxed{
    \mathbb N^{\mathrm{st}}_{B\mu}
    \simeq
    \Sigma\widetilde\Sigma^\infty_S\mu.}
\]
Actual formal isotropy is the whole arrow group $\mu$, and
\[
    \gamma^{\mathrm{st}}_{B\mu}:
    \mathbb I^{\mathrm{st}}_{B\mu}
    \xrightarrow{\sim}
    \mathbb F^{\mathrm{st}}_{B\mu}
\]
is an equivalence.  The integral spectral Lie algebra is
\[
    \boxed{
    \mathfrak l^{\mathrm{sp}}_{G/S}
    =
    \operatorname{Prim}^{\mathrm{der}}
    \left(
        \Sigma^\infty_{+,S}\mu
    \right).}
\]
\end{proposition}

\begin{proof}
The maximal relative relation for $S/S$ is the identity relation, so its
reduced stable coefficient vanishes.  Hence the anchor is the zero map from
$\widetilde\Sigma^\infty_S\mu$ to zero.  The fibre and cofiber formulas follow
from stability.  In a one-object groupoid every arrow is isotropy, so the
canonical isotropy comparison is the identity under these identifications.
The spectral Lie statement is the definition of derived primitives.
\end{proof}

\subsection{Formal action groupoids}

Let a group object $G/S$ act on $X$, and again write
$\mu=\mu^{\mathrm{mod}}_{G/S}$.  The preceding chapter gives
\[
    \widehat{G\ltimes X}^{\,\mathrm{dR}}
    \simeq
    \mu\ltimes X.
\]
Let $p:X\to S$ be the structural morphism.

\begin{proposition}[Stable differential of a formal action groupoid]
\label{prop:bnla-action-stable}
There is a canonical equivalence
\[
    \mathbb A^{\mathrm{st}}_{\mu\ltimes X}
    \simeq
    p^*\widetilde\Sigma^\infty_S\mu.
\]
Let
\[
    \kappa_R:R_{X/S}\longrightarrow X\times_SX
\]
be the centred-completion counit.  Under the displayed identification, the
stable anchor is the reduced stabilization of the canonical formal lift
\[
    \phi_{\mu\ltimes X,1}:
    \mu\times_SX\longrightarrow R_{X/S}
\]
of the source--action map
\[
    (\operatorname{pr}_X,\operatorname{act}):
    \mu\times_SX\longrightarrow X\times_SX,
    \qquad
    \kappa_R\phi_{\mu\ltimes X,1}
    \simeq
    (\operatorname{pr}_X,\operatorname{act}).
\]
Thus the stable infinitesimal-action morphism is
\[
    p^*\widetilde\Sigma^\infty_S\mu
    \longrightarrow
    \mathbb T^{\mathrm{st}}_{X/S}.
\]
Actual formal isotropy at a generalized point is the completed stabilizer
subgroup, and its stable coefficient maps canonically to the fibre of this
action anchor.
\end{proposition}

\begin{proof}
The arrow object of the action groupoid is $\mu\times_SX$, over $X$ by the
source projection and pointed by $(e,\operatorname{id}_X)$.  The base-change
equivalence for reduced stabilization identifies its coefficient with
$p^*\widetilde\Sigma^\infty_S\mu$.  The source--target morphism of the action
groupoid is
\[
    (\operatorname{pr}_X,\operatorname{act}):
    \mu\times_SX\longrightarrow X\times_SX.
\]
Because $\mu\ltimes X$ is formal, the canonical nonlinear anchor is the
contractibly unique centred formal lift of this pointed morphism to
$R_{X/S}$.  Reduced stabilization of that lift is therefore exactly the
stable anchor displayed in the proposition.  The isotropy statement is the
general comparison theorem applied to the stabilizer pullback.
\end{proof}

\begin{remark}[No semidirect Lie presentation is inserted integrally]
The stable action anchor is constructed without choosing a Lie algebra model
for $\mu$.  In coefficient theories where formal group/Lie recognition is
proved, its algebraic image may be identified with the corresponding action
algebroid by a separate comparison theorem.
\end{remark}

\subsection{Spectral partition-Lie images of the basic calculations}

\begin{proposition}[Identity spectral algebroid]
\label{prop:bnla-extremal-spectral-algebroids}
In the coherent spectral coefficient sector the identity formal groupoid
satisfies
\[
    \mathbb A^{\pi,E_\infty}_{0_X}
    \simeq
    \left(
        0\longrightarrow
        \mathbb T^{\mathrm{pc}}_{X/S}[1]
    \right),
\]
and therefore
\[
    U\mathbb F^{\pi,E_\infty}_{0_X}
    \simeq
    \Omega\mathbb T^{\mathrm{pc}}_{X/S}[1]
    \simeq
    \mathbb T^{\mathrm{pc}}_{X/S}.
\]
If, in addition, the effective de Rham image is relative de Rham-local, then
for the maximal relation
\[
    \mathbb A^{\pi,E_\infty}_{R_{X/S}}
    \simeq
    \left(
        \mathbb T^{\mathrm{pc}}_{X/S}[1]
        \xrightarrow{\operatorname{id}}
        \mathbb T^{\mathrm{pc}}_{X/S}[1]
    \right),
    \qquad
    U\mathbb F^{\pi,E_\infty}_{R_{X/S}}
    \simeq0.
\]
The additional hypotheses hold in the represented smooth range whenever
Theorem~10.2.1, using Theorem~9.16.2, gives smooth effectivity and the
Chapter~9 convergence criterion gives reduced-classical locality.
\end{proposition}

\begin{proof}
The identity calculation is
Corollary~\ref{cor:bnla-identity-spectral-algebroid}; apply the anchor-fibre
right adjoint of Theorem~\ref{thm:bnla-zero-anchor-spectral-algebroids} and
use stability.  The maximal calculation is
Corollary~\ref{cor:bnla-maximal-spectral-algebroid}, whose image-locality
hypothesis is repeated here explicitly.
\end{proof}

\begin{proposition}[One-object spectral partition Lie algebra]
\label{prop:bnla-one-object-spectral-partition}
For the completed one-object groupoid $B_\bullet\mu$ over $S/S$, the spectral
partition-Lie algebroid has zero anchor and is therefore a spectral partition
Lie algebra
\[
\boxed{
    \mathbb A^{\pi,E_\infty}_{B\mu}
    \in
    \operatorname{Alg}_{\operatorname{Lie}^{\pi}_{S,E_\infty}}
    (\operatorname{QC}^{\vee}_S).}
\]
Its formal integration is the spectral Artinian formal-moduli object of the
formal identity group $\mu$.  There is no unconditional identification
\[
    \mathbb A^{\pi,E_\infty}_{B\mu}
    \simeq
    \mathfrak l^{\mathrm{sp}}_{G/S};
\]
the first object is the partition-Lie formal-moduli controller and the second
is obtained from Hopf derived primitives.
\end{proposition}

\begin{proof}
The ambient tangent target for $S/S$ is zero.  Apply the zero-anchor theorem
\ref{thm:bnla-zero-anchor-spectral-algebroids} and one-object recovery
Corollary~\ref{cor:bnla-spectral-one-object-recovery}.  Formal integration is
inverse to the differentiation equivalence used to define the controller.
The final non-identification records the two distinct constructions.
\end{proof}

\begin{proposition}[Spectral action algebroid]
\label{prop:bnla-spectral-action-algebroid}
For a formal action groupoid $\mu\ltimes X$, spectral differentiation gives a
spectral partition-Lie algebroid
\[
    \mathbb A^{\pi,E_\infty}_{\mu\ltimes X}
    \longrightarrow
    \mathbb T^{\mathrm{pc}}_{X/S}[1]
\]
whose underlying anchor is the tangent-fibre differentiation of the nonlinear
action quotient.  Differentiated completed stabilizer isotropy maps
canonically to its anchor fibre.
\end{proposition}

\begin{proof}
Apply Theorem~\ref{thm:bnla-spectral-groupoid-differentiation} and the actual
isotropy comparison
Theorem~\ref{thm:bnla-spectral-partition-isotropy-comparison}.
\end{proof}

\section{Classical smooth groupoids: direct anchor and bracket construction}
\label{sec:bnla-classical-groupoids}

Let
\[
    \mathcal H_1
    \mathop{\rightrightarrows}^{s}_{t}
    X
\]
be an ordinary smooth Lie groupoid over an ordinary smooth characteristic-zero
base $S$, viewed in the manuscript through discrete spectral enhancement.  We
do not assume that this ordinary groupoid is already unitwise de Rham-complete;
first form its canonical unitwise formalization
\[
    \widehat{\mathcal H}^{\,\mathrm{dR}}_\bullet.
\]
By
Theorem~\ref{thm:bnla-first-neighbourhood-formalization}, its first-order unit
coefficient is still the coefficient of the original smooth arrow geometry.

\begin{theorem}[Classical anchor from the formalized groupoid]
\label{thm:bnla-classical-anchor}
There is a canonical identification
\[
    \mathfrak a_{\mathcal H}
    \simeq
    u^*T_{\mathcal H_1/X,s}
    =
    \ker(ds)|_{u(X)},
\]
and the first-order anchor is
\[
    \boxed{
    \rho_{\mathcal H}
    =
    u^*(dt)|_{\ker(ds)}:
    \mathfrak a_{\mathcal H}
    \longrightarrow
    T_{X/S}.}
\]
The identification is the first-order differential of the canonical nonlinear
anchor of the unitwise formalization.
\end{theorem}

\begin{proof}
Smoothness makes $L_{\mathcal H_1/X,s}$ finite locally free, and its dual is
the source-vertical tangent bundle.  Pulling back along the unit gives the
first formula.  The equivalent target formula for the cotangent anchor proved
in the preceding chapter dualizes to $dt$ restricted to source-vertical
tangent vectors.  Theorem~\ref{thm:bnla-first-neighbourhood-formalization}
identifies this first-order coefficient with the first neighbourhood carried
into the unitwise completion.
\end{proof}

\subsection{Right-invariant vector fields and the bracket}

For a local section
\[
    \alpha\in\Gamma(U,\mathfrak a_{\mathcal H}),
\]
right multiplication in the groupoid uniquely extends $\alpha$ to a
right-invariant $s$-vertical vector field
\[
    \widetilde\alpha
\]
on $t^{-1}(U)$.  The exact right/left convention is the multiplication
convention already fixed for the groupoid in the preceding chapter.

\begin{definition}[Classical groupoid bracket]
\label{def:bnla-classical-bracket}
Define
\[
    [\alpha,\beta]_{\mathcal H}
    :=
    u^*
    [\widetilde\alpha,\widetilde\beta]_{\mathrm{vf}},
\]
where the bracket on the right is the ordinary commutator of vector fields.
\end{definition}

\begin{theorem}[Classical Lie-algebroid identities]
\label{thm:bnla-classical-bracket}
The operation of
Definition~\ref{def:bnla-classical-bracket} is a Lie bracket on local sections
of $\mathfrak a_{\mathcal H}$ and, together with
$\rho_{\mathcal H}$, satisfies
\[
    \boxed{
    [\alpha,f\beta]_{\mathcal H}
    =
    \rho_{\mathcal H}(\alpha)(f)\,\beta
    +
    f[\alpha,\beta]_{\mathcal H}.}
\]
Thus
\[
    \left(
        \mathfrak a_{\mathcal H},
        [\ ,\ ]_{\mathcal H},
        \rho_{\mathcal H}
    \right)
\]
is the ordinary Lie algebroid differentiated from the smooth groupoid.
\end{theorem}

\begin{proof}
Right multiplication identifies source fibres and carries the unit value of a
source-vertical vector field uniquely to a right-invariant source-vertical
field.  The commutator of two right-invariant source-vertical vector fields is
again right-invariant and source-vertical.  Hence restriction along the unit
defines the displayed bracket.  Antisymmetry and Jacobi follow from the
corresponding identities for vector fields.

For $f\in\mathcal O_X$, right invariance gives
\[
    \widetilde{f\beta}
    =
    (t^*f)\widetilde\beta.
\]
The ordinary vector-field Leibniz rule yields
\[
    [\widetilde\alpha,(t^*f)\widetilde\beta]
    =
    \widetilde\alpha(t^*f)\widetilde\beta
    +(t^*f)
    [\widetilde\alpha,\widetilde\beta].
\]
Restrict to the unit.  Since $tu=\operatorname{id}_X$ and the target
differential of $\widetilde\alpha$ at the unit is
$\rho_{\mathcal H}(\alpha)$, the first term becomes
$\rho_{\mathcal H}(\alpha)(f)\beta$ and the second becomes
$f[\alpha,\beta]_{\mathcal H}$.
\end{proof}

\begin{remark}[Comparison with a recognized dg bracket]
The theorem constructs the classical bracket directly from nonlinear groupoid
multiplication.  A statement identifying this bracket with a dg bracket
obtained independently from a formal-moduli recognition theorem requires a
comparison between the two formal integrations.  That comparison is not
replaced by uniqueness rhetoric here; it is part of the later
Atiyah--Spencer/integration comparison.
\end{remark}

\subsection{Classical Lie groups}

\begin{corollary}[Classical Lie algebra]
\label{cor:bnla-classical-lie-algebra}
For an ordinary smooth Lie group $G/S$, the object space is $S$, the anchor
target vanishes, and
\[
    \mathfrak g
    =
    e^*T_{G/S}.
\]
The bracket of
Theorem~\ref{thm:bnla-classical-bracket} is the right-invariant Lie-algebra
convention.  Differential of inversion at the identity,
\[
    d(\operatorname{inv})_e=-\operatorname{id}_{\mathfrak g},
\]
canonically identifies this convention with the usual left-invariant Lie
algebra convention.
\end{corollary}

\begin{proof}
Apply the theorem to the one-object groupoid.  The source and target maps both
land in $S$, so the anchor is zero and the invariant groupoid vector fields are
ordinary right-invariant vector fields on $G$.

For a matrix group, if $X^R(g)=Xg$ and $Y^R(g)=Yg$, then
\[
    [X^R,Y^R](g)
    =(YX-XY)g
    =-[X,Y]_{\mathrm{left}}^R(g).
\]
Thus the right-invariant commutator convention is the opposite of the usual
left-invariant matrix convention.  Inversion exchanges left and right
translations and has derivative $-\operatorname{id}$ at the identity, so its
differential is a Lie algebra isomorphism from the right-invariant convention
to the usual left-invariant convention.  This fixes the sign canonically.
\end{proof}

\section{Classical principal bundles and the Atiyah sequence}
\label{sec:bnla-classical-atiyah}

Let
\[
    P\longrightarrow X
\]
be an ordinary smooth principal right $G$-bundle over a smooth
characteristic-zero base $S$, where $G\to S$ is an ordinary smooth group
scheme.

\begin{definition}[Classical Atiyah and adjoint bundles]
\label{def:bnla-classical-atiyah-bundles}
Define
\[
    \operatorname{At}(P)
    :=
    T_{P/S}/G
\]
and
\[
    \operatorname{ad}(P)
    :=
    P\times^G\mathfrak g,
\]
where $G$ acts on $TP$ through the differential of the principal action and on
$\mathfrak g$ by the adjoint action.
\end{definition}

The differential of the bundle projection is $G$-equivariant and therefore
descends to
\[
    \rho_P:
    \operatorname{At}(P)
    \longrightarrow
    T_{X/S}.
\]

\begin{theorem}[Classical gauge algebroid and Atiyah sequence]
\label{thm:bnla-classical-atiyah}
There is a canonical identification of Lie algebroids
\[
\boxed{
    A\bigl(\operatorname{Gauge}(P)\bigr)
    \simeq
    T_{P/S}/G
    =
    \operatorname{At}(P).}
\]
Under this identification the anchor is induced by $d\pi:T_{P/S}\to
\pi^*T_{X/S}$ and there is a canonical short exact sequence of locally free
sheaves
\[
\boxed{
    0
    \longrightarrow
    \operatorname{ad}(P)
    \longrightarrow
    \operatorname{At}(P)
    \xrightarrow{\rho_P}
    T_{X/S}
    \longrightarrow
    0.}
\]
The first arrow identifies differentiated actual gauge isotropy with the
ordinary kernel of the Atiyah anchor.  No connection or splitting is used in
the construction.
\end{theorem}

\begin{proof}
Write the ordinary gauge groupoid as
\[
    \operatorname{Gauge}(P)_1
    =
    (P\times P)/G
    \rightrightarrows X,
\]
where $G$ acts diagonally on the right.  Pull this groupoid and its Lie
algebroid back along the effective principal cover
\[
    \pi:P\twoheadrightarrow X.
\]
Over $P$ the pulled-back gauge groupoid has its canonical principal
trivialization: choosing the first point of the principal fibre identifies an
arrow with the second point.  Along the unit, its source-vertical tangent is
therefore canonically $T_{P/S}$.  The descent action on this tangent bundle is
the differential of the principal right $G$-action.  Effective descent for
vector bundles gives
\[
    A\bigl(\operatorname{Gauge}(P)\bigr)
    \simeq
    T_{P/S}/G.
\]
The target differential descends under the same identification to the map
induced by $d\pi$, which is $\rho_P$.

Because $\pi:P\to X$ is a smooth principal bundle, its relative tangent
sequence is exact:
\[
    0
    \longrightarrow
    T_{P/X}
    \longrightarrow
    T_{P/S}
    \xrightarrow{d\pi}
    \pi^*T_{X/S}
    \longrightarrow
    0.
\]
To match the right-invariant groupoid bracket of
Theorem~\ref{thm:bnla-classical-bracket} with the usual left-invariant
convention on $\mathfrak g$, use only the algebraic differential of the
principal action.  For $p\in P$, let
\[
    a_p:G\longrightarrow P,
    \qquad
    h\longmapsto ph
\]
be the orbit morphism and define
\[
\boxed{
    \xi^\#_p
    :=
    d(a_p)_e\!\left(d(\operatorname{inv})_e(\xi)\right)
    =
    -\,d(a_p)_e(\xi).}
\]
No exponential map is used.

With the standard associated-bundle right action
\[
    (p,\xi)\cdot g
    =
    (pg,\operatorname{Ad}_{g^{-1}}\xi),
\]
this gives a canonical $G$-equivariant identification
\[
    P\times\mathfrak g
    \xrightarrow{\sim}
    T_{P/X}.
\]
Indeed, the identity of orbit morphisms
\[
    a_{pg}(h)
    =
    R_g\!\left(a_p(ghg^{-1})\right)
\]
gives after differentiation
\[
    \left(\operatorname{Ad}_{g^{-1}}\xi\right)^\#_{pg}
    =
    dR_g(\xi^\#_p).
\]
Moreover the normalization is Lie-compatible.  In a local trivialization, an
equivariant section of $P\times\mathfrak g$ determined by
$\xi\in\mathfrak g$ gives the negative of the right-invariant vector field
$X^R(g)=\xi g$.  By
Corollary~\ref{cor:bnla-classical-lie-algebra}, the commutator of
right-invariant fields is the opposite of the usual left-invariant bracket;
the additional minus sign therefore gives
\[
    [\xi^\#,\eta^\#]
    =
    [\xi,\eta]^\#.
\]
Thus the induced $G$-action on $\mathfrak g$ under change of principal point
is the adjoint action with the standard Lie-algebra convention.
Quotienting the entire exact sequence by effective principal descent therefore gives
\[
    0
    \longrightarrow
    P\times^G\mathfrak g
    \longrightarrow
    T_{P/S}/G
    \longrightarrow
    T_{X/S}
    \longrightarrow
    0.
\]
This is the displayed Atiyah sequence.

Finally, the actual isotropy group of the gauge groupoid is the conjugation
bundle
\[
    P\times^G G.
\]
Differentiating its identity section gives
$P\times^G\mathfrak g$.  In the principal trivialization, the isotropy
inclusion differentiates by the same normalized fundamental-field map
$\xi\mapsto\xi^\#$ used above; equivalently, this is the inversion
normalization of
Corollary~\ref{cor:bnla-classical-lie-algebra}.  Under the preceding
identification of the gauge Lie algebroid, its image is exactly the kernel of
$\rho_P$.  Hence the directly differentiated classical isotropy morphism is
the inclusion of the ordinary kernel and is an equivalence onto that kernel
in this classical vector-bundle setting.
\end{proof}

\begin{corollary}[Classical gauge Atiyah identification]
\label{cor:bnla-classical-gauge-atiyah-identification}
For an ordinary smooth principal bundle, direct differentiation of the
ordinary gauge groupoid identifies actual differentiated isotropy with
$\operatorname{ad}(P)$ and identifies the ordinary anchor kernel with the same
vector bundle.  Consequently the directly differentiated classical sequence
is the classical Atiyah sequence.
\end{corollary}

\begin{proof}
The preceding proof identifies both objects with
$\operatorname{ad}(P)$ on every local trivialization.  The identifications are
equivariant under transition functions and therefore glue.
\end{proof}

\begin{remark}[Comparison with the rational formal-moduli controller]
The preceding theorem is a direct classical groupoid calculation.  An
identification of the independently constructed rational formal-moduli
comparison
\[
    \gamma^{\mathrm{Lie}}_{\mathfrak G}
\]
with this classical inclusion requires a comparison between the corresponding
formal integrations.  No such identification is inserted as an axiom here;
the operation-level and integration comparison is deferred to the
Atiyah--Spencer chapter.
\end{remark}

\section{Audit of constructions and coefficient theories}
\label{sec:bnla-recognition-audit}

The chapter contains several differentiation functors.  Their separation
records different constructions and theorem scopes; it does not define
several competing axiomatic packages.

\subsection{Unconditional intrinsic constructions}

For every brave-new formal Lie groupoid in the native big spectral Nisnevich
$\infty$-topos, without representability, the chapter constructs
\[
    \mathbb A^{\mathrm{st}}_{\mathfrak G},
    \qquad
    \mathbb T^{\mathrm{st}}_{X/S},
    \qquad
    a^{\mathrm{st}}_{\mathfrak G},
\]
the stable fibre/cofiber triangle, the actual-isotropy comparison
\[
    \mathbb I^{\mathrm{st}}_{\mathfrak G}
    \longrightarrow
    \mathbb F^{\mathrm{st}}_{\mathfrak G},
\]
its discrepancy, the isotropy suspension Hopf object, and the derived-primitive
spectral Lie algebra
\[
    \mathfrak l^{\mathrm{sp}}_{\mathfrak G}.
\]
All are outputs of stabilization, stable finite limits and colimits, group
object suspension, or the spectral Lie--Hopf adjunction.

\subsection{Represented first-order constructions}

When the arrow geometry is represented, the split first-order deformation
functor is corepresented by
\[
    \mathfrak a^\vee_{\mathfrak G}
    =u^*L_{\mathfrak G_1/X,s}.
\]
If the effective quotient is represented, this coefficient is canonically
\[
    \mathfrak a^\vee_{\mathfrak G}
    \simeq L_{X/Q_{\mathfrak G}},
\]
and the cotangent anchor is transitivity.  Unitwise formalization does not
need the completed arrow object to remain represented: split square-zero
probes of the completion are still corepresented by the original coefficient.

In the represented locally coherent quotient range,
Theorem~\ref{thm:bnla-spectral-represented-first-order-comparison} identifies
the underlying shifted spectral partition-Lie anchor with the continuous dual
of the same cotangent transitivity morphism.  The relation between first order
and formal-moduli differentiation is therefore a theorem.

\subsection{Spectral coherent-base recognition}

For a coherent spectral coefficient $B$, the chapter constructs rather than
postulates:
\begin{enumerate}
    \item the adic filtration as the left adjoint to augmentation, its
    filtered/graded cotangent compatibility, and its intrinsic associated-graded model;
    \item the spectral associated-graded formula
    \[
        \operatorname{Gr}\widehat{\operatorname{adic}}(A\to B)
        \simeq
        \operatorname{Sym}^{\mathrm{gr}}_B
        (L_{B/A}[-1](1));
    \]
    \item the positive spectral relation-cell construction and the
    surjective criterion identifying almost finite presentation with
    almost-perfectness of the augmentation module;
    \item the intrinsic complete almost-finite subcategory cut out by
    surjectivity on $\pi_0$, almost-perfectness of $B$ over $A$, and derived
    completion;
    \item the unrestricted spectral cotangent--square-zero adjunction, its
    connective restriction, and the spectral cotangent--square-zero
    comonadicity theorem with its canonical cobar resolution;
    \item the relative coLie comonad and its anchor-weight filtration;
    \item the relative spectral partition-Lie monad obtained by continuous
    duality and the pro-coherent sifted-completion theorem;
    \item the zero-anchor partition-Lie algebra theorem and the anchor-fibre
    right adjoint;
    \item the intrinsic spectral Artinian category defined by nilpotence,
    almost finite presentation, and eventual coconnectivity;
    \item the theorem that every Artinian extension admits a finite coherent
    square-zero factorization;
    \item the Artinian Schlessinger localization;
    \item the Koszul-dual functor $D$, its full faithfulness, free
    square-zero calculation, and Artinian excision;
    \item the anchored negative-cell attachment calculation identifying
    arbitrary anchored square-zero stages with attachments of zero-anchor
    cellular generators;
    \item the affine recognition equivalence
    \[
        \operatorname{FMP}^{\mathrm{sp}}_{B/R}
        \simeq
        \operatorname{LieAlgd}^{\pi,E_\infty}_{B/R},
    \]
    together with the canonical complete almost-finite extension
    $F^{\wedge\mathrm{aft}}$ used for square-zero tangent tests outside the
    Artinian subcategory;
\end{enumerate}

The spectral relative recognition theorem is therefore not imported from the
coherent one-object theory of \cite{BrantnerCamposNuiten2025}; that source
supplies the one-object pro-coherent Koszul-dual monad, while the relative
spectral cotangent, comonadicity, Artinian, and recognition statements are
proved in this chapter.

The anchor-weight filtration used in the relative monad construction is a
filtration of the actual coLie object.  Filtered spectral square-zero
finiteness proves its cotangent object almost perfect, and filtered
almost-perfect completeness proves that it is already complete; no completion
is inserted into the definition.  Likewise the zero-anchor theorem is induced
by the canonical subcomonad equivalence
\[
    \operatorname{coLie}^{\pi,E_\infty}_{B/R}z
    \simeq
    z\operatorname{coLie}^{\pi,E_\infty}_{B/B},
\]
not by an associated-graded identification.

\subsection{Globalization}

The affine spectral formal-moduli categories and affine spectral
partition-Lie algebroid categories carry canonical two-variable etale
transport.  The coherent spectral module family is passed through the
relative pro-coherent right-left completion, producing a pro-coherent etale
coefficient family whose transition maps are identified by the finite-Tor
Beck--Chevalley theorem.  Spectral fppf universal descent, together with that
pointwise Beck--Chevalley theorem and Barr--Beck, proves etale descent for
this coefficient family; the compatible relative partition-Lie monads then
give descent for
partition-Lie algebroids.  Etale naturality of affine differentiation
transports that descent theorem to the formal-moduli categories.  Therefore
\[
    \operatorname{FMP}^{\mathrm{sp}}_{X/S}
    \quad\text{and}\quad
    \operatorname{LieAlgd}^{\pi,E_\infty}_{X/S}
\]
are constructed as etale global-section categories and
\[
    \operatorname{FMP}^{\mathrm{sp}}_{X/S}
    \simeq
    \operatorname{LieAlgd}^{\pi,E_\infty}_{X/S}.
\]

The formal leaf of a groupoid is not obtained by separately reflecting raw
affine leaf functors.  The nonlinear quotient is first sent to
\[
    Q^{\mathrm{def}}_{\mathfrak G}
    =
    L^{\mathrm{Def}}_S(Q_{\mathfrak G})
\]
and then to its reduced-centred model
\[
    Q^{\mathrm{fm}}_{\mathfrak G}
    =
    Q^{\mathrm{def}}_{\mathfrak G}
    \times_{D_{\mathrm{red}}Q^{\mathrm{def}}_{\mathfrak G}}
    D_{\mathrm{red}}X.
\]
Every affine Artinian leaf of this one global target is already a spectral
formal moduli problem.  Before comparing centres, its coherent square-zero
lifting functor is stabilized to a convergent pre-tangent functor and then
sent to its universal almost-finitely-presented pro-coherent approximation.
The global deformation equations make these pre-tangent functors cartesian
under spectral-etale changes of centre.  The relative right-left anchored
overcategory then produces the cartesian transport $f_\sharp$ without
asserting that the fibrewise almost-finite approximation right adjoints commute
with coefficient change.  Etale naturality of affine differentiation and
conservativity of the monadic forgetful functor then make the centre-change
morphism an equivalence.  Consequently these affine leaves form a genuine
etale cocartesian global
section
\[
    \mathcal L^{\mathrm{fm,sp}}_{\mathfrak G}
    \in
    \operatorname{FMP}^{\mathrm{sp}}_{X/S}.
\]
No chosen affine presentation, chosen affine-etale basis, independently
retained descent datum, or cocartesian-section reflector is part of the
construction.

For a fixed object space $X/S$, these global leaf and algebroid constructions
are functorial for arbitrary morphisms in $\operatorname{FormLieGpd}(X/S)$.
For a functor of formal groupoids whose object map $X\to Y$ is spectral etale,
restriction of coherent affine-etale pairs constructs the varying-object maps
\[
    \mathcal L^{\mathrm{fm,sp}}_{\mathfrak G}
    \longrightarrow
    f^*\mathcal L^{\mathrm{fm,sp}}_{\mathfrak H},
    \qquad
    \mathbb A^{\pi,E_\infty}_{\mathfrak G}
    \longrightarrow
    f^*\mathbb A^{\pi,E_\infty}_{\mathfrak H}.
\]
An arbitrary varying-object map is outside the variance of the etale-global
coefficient categories constructed here and would require a separate spectral
families formalism; no such functoriality is used below.

\subsection{Rational comparison theory}

For rational coherent spectral coefficients, where rational means connective
commutative $H\mathbb Q$-algebra and dg grading is homological, symmetric
monoidal rectification identifies cotangent complexes, structured square-zero
extensions, complete almost-finite objects, pro-coherent coefficients, and the
relative coLie construction with their dg analogues.  The dg
cotangent--square-zero adjunction, its comonad, and its pro-coherent dual monad
are constructed explicitly before comparison with the spectral monad.  Rational
rectification also restricts to
\[
    \operatorname{Art}^{\mathrm{sp}}_{R/B}
    \simeq
    \operatorname{Art}^{\mathrm{dg}}_{R_{\mathrm{dg}}/B_{\mathrm{dg}}},
\]
carries the spectral Schlessinger generators to the dg generators, and hence
identifies the corresponding formal-moduli categories.  Finite-group Tate
terms vanish rationally, so the PD spectral partition-Lie monad rectifies to
the pro-coherent shifted dg Lie-algebroid monad.  This constructs
\[
    \mathbb A^{\mathrm{Lie}}_{\mathfrak G}.
\]
Comparison with \cite{Nuiten2019} is then invoked only in that theorem's
coefficient and boundedness range.

\subsection{Deferred integral animated comparison}

No integral spectral-to-animated Artinian comparison is used.  An animated
comparison may be added only after a separate theorem constructs a comparison
of the spectral and animated Artinian deformation contexts compatible with
square-zero extensions, cotangent fibres, Schlessinger localization, and
formal integration.  Such compatibility is neither assumed nor encoded in
the spectral algebroid.

\subsection{No axiom disguised as a package, datum, presentation, or admissibility condition}

None of the following is input data of a brave-new formal Lie groupoid or of
its differentiated object:
\begin{enumerate}
    \item a finite infinitesimal factorization;
    \item an Artinian square-zero chain;
    \item a coefficient module or derivation presenting an already
    constructed infinitesimal arrow;
    \item an adic filtration, ideal generator, topology, or completion
    presentation;
    \item a formal-moduli recognition structure;
    \item an admissibility structure;
    \item a chosen nonlinear or linear anchor;
    \item a transitivity or surjectivity condition;
    \item a chosen spectral partition-Lie presentation;
    \item a zero-anchor presentation;
    \item an identification of actual isotropy with a homotopy anchor fibre;
    \item a primitive base-change equivalence;
    \item a connection or splitting;
    \item a dg rectification presentation;
    \item an etale descent datum retained independently of the descended
    object;
    \item a distinguished affine presentation or chosen affine-etale basis of
    the global spectral base.
\end{enumerate}

Finite factorizations and coefficient filtrations occur inside proofs because
the intrinsic hypotheses imply that such constructions exist.  They are never
remembered by the resulting objects.  The temporary predicate used in the
proof of Artinian excision is likewise only an induction device.

\subsection{Theorem-scope hypotheses}

The following are genuine domains of proved constructions rather than fields
of structure:
\begin{enumerate}
    \item representability when the cotangent complex of a groupoid arrow or
    quotient is used as a geometric quasi-coherent object;
    \item coherence of spectral coefficient rings for the pro-coherent
    spectral partition-Lie theory;
    \item local coherence and qcqs hypotheses for the global coherent etale
    coefficient system;
    \item rationality for dg rectification;
    \item the coefficient and boundedness scope of \cite{Nuiten2019} for the
    comparison with that formal-moduli theorem;
    \item fixed-object functoriality for arbitrary formal-groupoid morphisms,
    and varying-object spectral partition-Lie functoriality for spectral etale
    object maps; arbitrary varying-object maps require a separate spectral
    families construction and are not used in this chapter;
    \item relative de Rham image-locality together with
    $L^{\mathrm{Def}}_S$-locality for the simplified identity-anchor
    calculation of the maximal spectral controller.  The deformation-locality
    condition is automatic under reduced-classical locality; in the
    represented smooth range it therefore suffices to combine smooth
    effectivity from Theorem~10.2.1, using Theorem~9.16.2, with the Chapter~9
    convergence criterion.
\end{enumerate}
None is retained as a datum on $\mathfrak G$.

\section{The brave-new differentiation package}
\label{sec:bnla-synthesis}

\begin{theorem}[Brave-New Differentiation Package]
\label{thm:bnla-differentiation-package}
Let
\[
    \mathfrak G_\bullet
    \in
    \operatorname{FormLieGpd}(X/S).
\]
The constructions of this chapter canonically provide the following outputs.
\begin{enumerate}
    \item The intrinsic stable coefficients
    \[
        \mathbb A^{\mathrm{st}}_{\mathfrak G}
        =\widetilde\Sigma^\infty_X(\mathfrak G_1,u),
        \qquad
        \mathbb T^{\mathrm{st}}_{X/S}
        =\widetilde\Sigma^\infty_X(R_{X/S},u_R),
    \]
    and the stable anchor between them.

    \item The stable anchor fibre/cofiber triangle with
    \[
        \mathbb N^{\mathrm{st}}_{\mathfrak G}
        \simeq
        \Sigma\mathbb F^{\mathrm{st}}_{\mathfrak G}.
    \]

    \item Actual stable formal isotropy, the canonical comparison
    \[
        \gamma^{\mathrm{st}}_{\mathfrak G}:
        \mathbb I^{\mathrm{st}}_{\mathfrak G}
        \longrightarrow
        \mathbb F^{\mathrm{st}}_{\mathfrak G},
    \]
    and its discrepancy cofiber.

    \item The cocommutative isotropy suspension Hopf object, its augmentation
    ideal, and the integral derived-primitive spectral Lie algebra
    $\mathfrak l^{\mathrm{sp}}_{\mathfrak G}$ with enveloping-Hopf counit.

    \item Formal adjoint transport on isotropy Hopf objects and the
    always-defined primitive base-change comparison
    \[
        f^*\operatorname{Prim}(H)
        \longrightarrow
        \operatorname{Prim}(f^*H).
    \]

    \item In the represented range, the first-order vertical deformation
    coefficient
    \[
        \mathfrak a^\vee_{\mathfrak G}
        =u^*L_{\mathfrak G_1/X,s},
    \]
    and, for represented $Q_{\mathfrak G}$,
    \[
        \mathfrak a^\vee_{\mathfrak G}
        \simeq L_{X/Q_{\mathfrak G}},
    \]
    with cotangent anchor given by transitivity.

    \item The first-order stable comparison, the retractive centred-coreflection
    theorem, and the theorem that split square-zero probes of a unitwise
    formalized represented arrow remain corepresented by the original
    first-order coefficient, with the first structured unit neighbourhood
    mapping compatibly to the canonical completed nonlinear anchor.

    \item The raw finite spectral quotient-lifting functor on the Chapter~9
    finite infinitesimal test category.

    \item In coherent spectral coefficients, the unrestricted
    cotangent--square-zero adjunction, its connective restriction, the
    constructed adic filtration, its cotangent-detection and associated-graded
    descriptions, the intrinsic complete almost-finite augmentation category,
    and spectral cotangent--square-zero comonadicity with its canonical cobar
    resolution.

    \item The relative coLie comonad, its functorial anchor-weight filtration,
    the filtered spectral square-zero finiteness theorem proving that this
    filtration is already complete, and the spectral partition-Lie algebroid
    monad
    \[
        \operatorname{Lie}^{\pi,E_\infty}_{B/R}
    \]
    produced by continuous pro-coherent duality and the proved completion
    relations.

    \item The zero-derivation subcomonad equivalence, the resulting zero-anchor
    equivalence with spectral partition-Lie algebras, and the constructively
    lifted canonical anchor-fibre right adjoint.

    \item The intrinsic spectral Artinian category and the theorem that every
    Artinian arrow has a finite factorization by structured square-zero
    extensions with coherent coefficients, without retaining such a
    factorization.

    \item The Artinian Schlessinger localization, the fully faithful spectral
    Koszul-dual functor $D^{\mathrm{sp}}$, the square-zero/free theorem,
    Artinian excision, the anchored negative-cell attachment theorem, and the
    resulting identification of the Artinian opposite category with the
    finite cellular algebroid subcategory.  These give the affine equivalence
    \[
        \operatorname{FMP}^{\mathrm{sp}}_{B/R}
        \simeq
        \operatorname{LieAlgd}^{\pi,E_\infty}_{B/R},
    \]
    together with the canonical complete almost-finite extension
    $F^{\wedge\mathrm{aft}}$ used by the square-zero tangent formula.

    \item Canonical two-variable spectral etale transport, the finite-Tor
    pro-coherent Beck--Chevalley theorem, the relative pro-coherent spectral
    etale coefficient family and its etale descent, etale descent for the
    affine partition-Lie coefficient system, and the global
    equivalence
    \[
        \operatorname{FMP}^{\mathrm{sp}}_{X/S}
        \simeq
        \operatorname{LieAlgd}^{\pi,E_\infty}_{X/S}.
    \]

    \item The global spectral deformation reflector
    \[
        L^{\mathrm{Def}}_S,
    \]
    refining the Artinian-deformation reflector by Postnikov convergence and
    nilpotent excision, and the formalized nonlinear quotient targets
    \[
        Q^{\mathrm{def}}_{\mathfrak G},
        \qquad
        Q^{\mathrm{fm}}_{\mathfrak G},
    \]
    the theorem that every affine formal leaf is a spectral formal moduli
    problem, the coherent pre-tangent functor and universal
    almost-finitely-presented tangent approximation, the cartesian
    affine-etale pre-tangent family and anchored $f_\sharp$ transport, the
    formally-etale pre-tangent base-change theorem, the etale cocartesian centre-change
    equivalence, and the resulting global formal leaf
    \[
        \mathcal L^{\mathrm{fm,sp}}_{\mathfrak G}
        \in
        \operatorname{FMP}^{\mathrm{sp}}_{X/S}.
    \]
    Its brave-new spectral partition-Lie algebroid is
    \[
        \mathbb A^{\pi,E_\infty}_{\mathfrak G}
        =
        \operatorname{Diff}^{\pi,E_\infty}_{X/S}
        (\mathcal L^{\mathrm{fm,sp}}_{\mathfrak G}).
    \]
    For fixed $X/S$ this construction is functorial on
    $\operatorname{FormLieGpd}(X/S)$; a formal-groupoid functor covering a
    spectral etale object map $f:X\to Y$ induces canonically
    \[
        \mathbb A^{\pi,E_\infty}_{\mathfrak G}
        \longrightarrow
        f^*\mathbb A^{\pi,E_\infty}_{\mathfrak H}.
    \]

    \item In the represented locally coherent quotient range, the underlying
    anchor
    \[
        (L_{X/Q_{\mathfrak G}})^\vee[1]
        \longrightarrow
        (L_{X/S})^\vee[1].
    \]

    \item The rational Artinian equivalence
    \[
        \operatorname{Art}^{\mathrm{sp}}_{R/B}
        \simeq
        \operatorname{Art}^{\mathrm{dg}}_{R_{\mathrm{dg}}/B_{\mathrm{dg}}},
        \qquad
        \operatorname{FMP}^{\mathrm{sp}}_{B/R}
        \simeq
        \operatorname{FMP}^{\mathrm{dg}}_{B_{\mathrm{dg}}/R_{\mathrm{dg}}},
    \]
    together with the algebroid rectification
    \[
        \mathcal R_{\mathbb Q}
        (\mathbb A^{\pi,E_\infty}_{\mathfrak G})
        \simeq
        \mathbb A^{\mathrm{Lie}}_{\mathfrak G},
    \]
    and its separately scoped comparison with dg formal-moduli
    differentiation in the range of \cite{Nuiten2019}.

    \item Spectral differentiated actual isotropy
    $\mathfrak i^{\pi,E_\infty}_{\mathfrak G}$, its zero anchor, the canonical
    map to the spectral anchor fibre, the rational dg image, and the
    corresponding discrepancy objects.

    \item The spectral anchor-fibre transitivity equivalence on underlying
    pro-coherent objects and its deshifted rational formula.

    \item The effective-image gauge formula
    \[
        Q_{\operatorname{Gauge}(P_\chi)}\simeq I_\chi,
        \qquad
        Q_{\operatorname{Gauge}^{\mathrm{for}}(P_\chi)}
        \simeq\widehat I_\chi,
    \]
    together with the stable, spectral partition-Lie, and rational dg gauge
    anchor-fibre sequences.

    \item The identity, maximal, one-object, and action calculations in the
    intrinsic stable theory; the identity, one-object, and action spectral
    calculations; and the maximal spectral identity-anchor calculation under
    the explicit image-locality hypothesis of
    Corollary~\ref{cor:bnla-maximal-spectral-algebroid}.

    \item Direct recovery of the classical smooth Lie-groupoid anchor and
    right-invariant bracket, with the inversion comparison to the ordinary
    left-invariant convention.

    \item Direct recovery of the classical gauge Lie algebroid
    \[
        A(\operatorname{Gauge}(P))
        \simeq T_{P/S}/G
    \]
    and the classical Atiyah sequence
    \[
        0\longrightarrow P\times^G\mathfrak g
        \longrightarrow T_{P/S}/G
        \longrightarrow T_{X/S}
        \longrightarrow0
    \]
    without a connection or splitting.
\end{enumerate}

Every item is produced by a functor, adjunction, localization, monadic or
comonadic construction, canonical cobar resolution, continuous duality,
etale global-section construction, or universal property proved in this chapter.  No binary
bracket is postulated on the intrinsic stable arrow coefficient; no spectral
algebroid structure is attached to a groupoid as extra data; no
isotropy/fibre identification is made without a constructed comparison; and
no connection, splitting, finite factorization, adic presentation,
recognition datum, admissibility structure, or chosen descent presentation is
part of the output.

No integral spectral-to-animated comparison is included.  Such a comparison
requires a separate theorem comparing the two Artinian deformation contexts.
Likewise no arbitrary varying-object spectral algebroid pullback is asserted
beyond the spectral-etale variance constructed by the global coefficient
system.
\end{theorem}

\begin{proof}
Items concerning intrinsic stabilization, stable isotropy, Hopf objects, and
primitives are
Theorems~\ref{thm:bnla-intrinsic-stable-differentiation},
\ref{thm:bnla-stable-anchor-triangle},
\ref{thm:bnla-isotropy-augmentation-ideal},
\ref{thm:bnla-integral-spectral-isotropy}, and
\ref{thm:bnla-differentiated-adjoint}.

The represented and formalization statements are
Theorems~\ref{thm:bnla-source-coefficient-corepresents},
\ref{thm:bnla-relative-cotangent-identification},
\ref{thm:bnla-first-order-stable-comparison},
\ref{thm:bnla-first-probe-centred-completion}, and
\ref{thm:bnla-first-neighbourhood-formalization}.

The spectral algebraic construction is
Proposition~\ref{prop:bnla-cot-sqz-adjunction},
Lemma~\ref{lem:bnla-sqz-connective-restriction},
Lemma~\ref{lem:bnla-adic-cotangent-compatibility},
Lemma~\ref{lem:bnla-graded-cotangent-detection},
Theorem~\ref{thm:bnla-adic-associated-graded},
Proposition~\ref{prop:bnla-adic-pi-zero},
Lemma~\ref{lem:bnla-relative-spectral-hurewicz},
Lemma~\ref{lem:bnla-positive-relation-cell},
Lemma~\ref{lem:bnla-spectral-cell-aft},
Proposition~\ref{prop:bnla-spectral-aft-criterion},
Lemma~\ref{lem:bnla-coherent-nilpotent-ring},
Theorems~\ref{thm:bnla-spectral-cot-sqz-comonadic},
\ref{thm:bnla-spectral-lie-algebroid-monad},
\ref{thm:bnla-zero-anchor-spectral-algebroids},
and Proposition~\ref{prop:bnla-spectral-art-square-zero-generation}.
Artinian formal-moduli recognition is
Theorems~\ref{thm:bnla-spectral-D-fully-faithful},
\ref{thm:bnla-spectral-D-excision},
Lemma~\ref{lem:bnla-anchored-cell-attachment}, and
Theorem~\ref{thm:bnla-spectral-affine-recognition}.

Globalization is
Lemma~\ref{lem:bnla-procoherent-finite-tor-bc},
Construction~\ref{con:bnla-relative-procoherent-etale-family},
Theorems~\ref{thm:bnla-two-variable-etale-base-change},
\ref{thm:bnla-global-fmp-localization}, and
\ref{thm:bnla-global-spectral-recognition}, together with
Propositions~\ref{prop:bnla-procoherent-etale-descent} and
\ref{prop:bnla-spectral-algebroid-etale-descent}.  The nonlinear target and
formal-leaf construction is
Theorem~\ref{thm:bnla-global-artdef-reflector},
Theorem~\ref{thm:bnla-global-def-reflector},
Propositions~\ref{prop:bnla-artdef-dred-equivalence} and
\ref{prop:bnla-def-dred-equivalence},
Proposition~\ref{prop:bnla-affine-leaf-is-fmp},
Proposition~\ref{prop:bnla-pre-tangent-aft-approximation},
Construction~\ref{con:bnla-coherent-pre-tangent},
Proposition~\ref{prop:bnla-pre-tangent-def-evaluation},
Proposition~\ref{prop:bnla-pre-tangent-cartesian-family},
Lemma~\ref{lem:bnla-anchored-right-left-pullback},
Construction~\ref{con:bnla-anchored-pre-tangent-transport},
Theorems~\ref{thm:bnla-pre-tangent-leaf-controller} and
\ref{thm:bnla-pre-tangent-etale-base-change},
Proposition~\ref{prop:bnla-centred-artinian-tangent-base-change},
Theorem~\ref{thm:bnla-formal-leaf-etale-cocartesian}, and
Theorem~\ref{thm:bnla-global-leaf-functoriality}.  Groupoid differentiation and the
represented comparison are
Theorems~\ref{thm:bnla-spectral-groupoid-differentiation} and
\ref{thm:bnla-spectral-represented-first-order-comparison}.

Rational rectification is
Proposition~\ref{prop:bnla-rational-coefficient-rectification},
Proposition~\ref{prop:bnla-relative-dg-colie},
Proposition~\ref{prop:bnla-relative-dg-monad-construction},
Theorem~\ref{thm:bnla-rational-artinian-rectification},
Theorem~\ref{thm:bnla-rational-monad-rectification}, and
Theorem~\ref{thm:bnla-dg-lie-rectification}.  The algebraic isotropy statements are
Theorem~\ref{thm:bnla-spectral-partition-isotropy-comparison},
Corollary~\ref{cor:bnla-dg-isotropy-comparison}, and
Propositions~\ref{prop:bnla-spectral-anchor-fibre-transitivity} and
\ref{prop:bnla-dg-anchor-transitivity}.

The remaining examples and classical calculations are
Theorem~\ref{thm:bnla-gauge-effective-image},
Propositions~\ref{prop:bnla-zero-groupoid},
\ref{prop:bnla-maximal-relation},
\ref{prop:bnla-one-object-differentiation},
\ref{prop:bnla-action-stable},
\ref{prop:bnla-extremal-spectral-algebroids},
\ref{prop:bnla-one-object-spectral-partition}, and
\ref{prop:bnla-spectral-action-algebroid}, together with
Theorems~\ref{thm:bnla-classical-anchor},
\ref{thm:bnla-classical-bracket}, and
\ref{thm:bnla-classical-atiyah}.
\end{proof}

\section{Handoff to Atiyah--Spencer calculus and higher Lie operations}
\label{sec:bnla-handoff-spencer}

The nonlinear and differentiated structures have now been constructed at the
level required for the operation calculus.  The formal groupoid supplies
\[
    \mathfrak G_\bullet,
    \qquad
    J_{\mathfrak G},
    \qquad
    I^{\mathrm{for}}_{\mathfrak G},
    \qquad
    \operatorname{Ad}^{\mathrm{for}}(\mathfrak G),
\]
while intrinsic stabilization supplies
\[
    \mathbb A^{\mathrm{st}}_{\mathfrak G},
    \qquad
    \mathbb I^{\mathrm{st}}_{\mathfrak G}
    \longrightarrow
    \mathbb F^{\mathrm{st}}_{\mathfrak G},
    \qquad
    \mathfrak l^{\mathrm{sp}}_{\mathfrak G}.
\]
The global spectral deformation formalization of the quotient and the
etale cocartesian leaf theorem supply
\[
    \mathcal L^{\mathrm{fm,sp}}_{\mathfrak G},
\]
and spectral differentiation supplies
\[
    \mathbb A^{\pi,E_\infty}_{\mathfrak G}.
\]
In the rational coefficient sector its dg image is
\[
    \mathbb A^{\mathrm{Lie}}_{\mathfrak G}.
\]

The next chapter begins with
\[
    \mathfrak G_\bullet
    \simeq
    \check C_\bullet
    \left(X\to Q_{\mathfrak G}\right),
    \qquad
    J_{\mathfrak G}=q_{\mathfrak G}^*\Pi_{q_{\mathfrak G}},
\]
forms the intrinsic jet--cobar and normalized Spencer objects, and compares
their operation algebra with the algebraic operations constructed here.  The
progression is
\[
\boxed{
\begin{aligned}
\text{groupoid descent}
&\longrightarrow \text{jet--cobar object}\\
&\longrightarrow \text{Spencer normalization}\\
&\longrightarrow \text{Atiyah representations}\\
&\longrightarrow \text{higher operations}\\
&\longrightarrow \text{bracket and Maurer--Cartan comparison}.
\end{aligned}}
\]

The preceding chapter constructed the nonlinear commutator on formal isotropy.
The present chapter has constructed its suspension Hopf object and derived
primitives, the spectral partition-Lie differentiation of actual isotropy,
the many-object spectral algebroid of the localized formal leaf, the canonical
isotropy-to-anchor-fibre comparison, and the rational dg comparison object.
The next chapter can therefore compare nonlinear commutators, differentiated
conjugation, and normalized jet--cobar operations with the arity-two and
higher spectral partition-Lie operations without introducing a new algebraic
controller.

The integral spectral-to-animated comparison remains a separate comparison
problem and is not needed for the Atiyah--Spencer construction.
The Atiyah--Spencer constructions below use the fixed-object functoriality and
the spectral-etale descent variance proved here; no arbitrary changing-object
spectral partition-Lie pullback is required.

The four-stage organization is
\[
\boxed{
\text{formal groupoid first,}
\qquad
\text{differentiation second,}
\qquad
\text{operations third,}
\qquad
\text{integration/comparison last.}}
\]
